%=====================================================================
% PES University --- M.Tech (Data Science and Machine Learning)
% Project Phase-2 Report / Thesis
% Tinker RL Lab --- Post-Training of Large Language Models
% Author: Arvind C R (SRN )   Guide: Ramesh Prakash Guledgudd
% Built with tectonic:  cd <dir> && tectonic thesis_master.tex
%=====================================================================
\documentclass[12pt,a4paper]{report}

% ---- binding margin: extra room on the left (spine) side --------------
\usepackage[a4paper,left=1.5in,right=1in,top=1in,bottom=1in]{geometry}
\usepackage{setspace}
\onehalfspacing

\usepackage{amsmath,amssymb}
\usepackage[T1]{fontenc}
\usepackage{newtxtext,newtxmath}   % Times, per PES front-page sample
\usepackage{graphicx}
\usepackage{array}
\usepackage{booktabs}
\usepackage{longtable}
\usepackage{calc}          % required by pandoc's table column-width arithmetic
\newcolumntype{L}[1]{>{\raggedright\arraybackslash\hspace{0pt}}p{#1}}
\usepackage{xcolor}
\definecolor{pesmaroon}{RGB}{150,30,40}
\usepackage{enumitem}
\usepackage{needspace}
\usepackage[numbers,sort&compress]{natbib}
\usepackage[hidelinks]{hyperref}
\hypersetup{
  pdftitle={Derived Public Revision - Tinker RL Lab: Post-Training of Large Language Models --- A multi-framework training harness and studies of GRPO signal starvation},
  pdfauthor={Arvind C R},
  pdfsubject={Derived public C1 coverage revision dated 2026-10-03 of M.Tech Project Phase-2 Report, PES University},
  pdfkeywords={GRPO, RLVR, zero-variance fraction, reproducibility, telemetry, benchmark evaluation}
}
\usepackage{tikz}
\usepackage{pgfplots}
\pgfplotsset{compat=1.18}
\usetikzlibrary{arrows.meta,positioning,shapes.geometric,fit,backgrounds,calc,
                decorations.pathreplacing,patterns,pgfplots.groupplots}

% ---- shared tikz styles (identical to the ESA / Phase-1 figure set) ----
\tikzset{
  blk/.style={rectangle, rounded corners=2pt, draw=black!70, thick, fill=blue!6,
              minimum height=8mm, minimum width=20mm, align=center, font=\small},
  io/.style={rectangle, rounded corners=6pt, draw=black!70, thick, fill=green!8,
             minimum height=8mm, align=center, font=\small},
  proc/.style={rectangle, draw=black!70, thick, fill=orange!10, minimum height=8mm,
               align=center, font=\small},
  dec/.style={diamond, aspect=2, draw=black!70, thick, fill=yellow!15, align=center,
              font=\footnotesize},
  warnbox/.style={rectangle, rounded corners=2pt, draw=red!60!black, thick, fill=red!7,
                  minimum height=8mm, align=center, font=\small},
  oklbox/.style={rectangle, rounded corners=2pt, draw=green!45!black, thick, fill=green!8,
                 minimum height=8mm, align=center, font=\small},
  ar/.style={-{Latex[length=2mm]}, thick},
  arl/.style={-{Latex[length=2mm]}, thick, dashed, draw=black!55},
  lbl/.style={font=\footnotesize\itshape, text=black!60},
}
\definecolor{good}{RGB}{34,139,34}
\definecolor{bad}{RGB}{178,34,34}
\definecolor{warn}{RGB}{218,165,32}
\definecolor{pesblue}{RGB}{0,45,98}
\definecolor{complete}{RGB}{30,125,79}
\definecolor{partial}{RGB}{176,106,0}
\definecolor{blocked}{RGB}{166,43,43}

\newcommand{\zvf}{\mathrm{ZVF}}
\renewcommand{\bibname}{References}

\setlength{\parskip}{0.4em}
\setlength{\parindent}{0pt}
\setlength{\emergencystretch}{4em}

%=====================================================================
% ---- LONG UNBREAKABLE STRINGS: allow line breaks inside \texttt ----
%=====================================================================
% The corpus contains file paths, JSON fragments, model ids, W&B keys and
% SHA-256 digests far wider than the text column and containing no spaces.
% Measured on the 2026-09-25 build with pdftotext -bbox: 72 of 217 pages
% (33%) carried ink past the right text margin at 523.3pt, worst case
% 602pt on a 595pt-wide A4 page -- i.e. text physically off the paper.
% Every offending token was an unbreakable \texttt path/hash/URI.
%
% Fix, verified in isolation (max ink 523.3pt vs 523.3pt limit, 0 overflow):
%   * xurl     -- breaks \url{} and path-like strings at their delimiters
%   * seqsplit -- breaks arbitrary long unbroken tokens (hashes, URIs) at
%                 any character, via the \longtok wrapper
% The \longtok/sub-token rewriting is applied automatically by a source-level
% pass in build_thesis.py; short identifiers are left alone so ordinary words
% such as logging-group-14 are NOT hyphenated mid-word (the [htt]hyphenat
% option was tried and removed for exactly that reason).
\usepackage{xurl}
\usepackage{seqsplit}
\DeclareRobustCommand{\longtok}[1]{\texttt{\seqsplit{#1}}}
\tolerance=6000
\hbadness=10000
\vbadness=10000
\sloppy

% ---- pandoc compatibility -------------------------------------------
% pandoc's LaTeX fragments reference macros that live in pandoc's own
% default template, not in the fragment. Define them here.
\providecommand{\real}[1]{#1}          % table column-width arithmetic
\providecommand{\tightlist}{%          % compact lists
  \setlength{\itemsep}{0pt}\setlength{\parskip}{0pt}}
\providecommand{\pandocbounded}[1]{#1} % image bounding wrapper
\providecommand{\passthrough}[1]{#1}
\newenvironment{Shaded}{}{}            % syntax-highlighted code block
\newenvironment{Highlighting}{}{}

% pandoc emits \XxxTok{...} for syntax-highlighted spans. Render each as
% monospace with no colour so highlighted code compiles under any engine.
\providecommand{\NormalTok}[1]{\texttt{#1}}
\providecommand{\KeywordTok}[1]{\texttt{#1}}
\providecommand{\DataTypeTok}[1]{\texttt{#1}}
\providecommand{\DecValTok}[1]{\texttt{#1}}
\providecommand{\BaseNTok}[1]{\texttt{#1}}
\providecommand{\FloatTok}[1]{\texttt{#1}}
\providecommand{\ConstantTok}[1]{\texttt{#1}}
\providecommand{\CharTok}[1]{\texttt{#1}}
\providecommand{\SpecialCharTok}[1]{\texttt{#1}}
\providecommand{\StringTok}[1]{\texttt{#1}}
\providecommand{\VerbatimStringTok}[1]{\texttt{#1}}
\providecommand{\SpecialStringTok}[1]{\texttt{#1}}
\providecommand{\ImportTok}[1]{\texttt{#1}}
\providecommand{\CommentTok}[1]{\texttt{#1}}
\providecommand{\DocumentationTok}[1]{\texttt{#1}}
\providecommand{\AnnotationTok}[1]{\texttt{#1}}
\providecommand{\CommentVarTok}[1]{\texttt{#1}}
\providecommand{\OtherTok}[1]{\texttt{#1}}
\providecommand{\FunctionTok}[1]{\texttt{#1}}
\providecommand{\VariableTok}[1]{\texttt{#1}}
\providecommand{\ControlFlowTok}[1]{\texttt{#1}}
\providecommand{\OperatorTok}[1]{\texttt{#1}}
\providecommand{\BuiltInTok}[1]{\texttt{#1}}
\providecommand{\ExtensionTok}[1]{\texttt{#1}}
\providecommand{\PreprocessorTok}[1]{\texttt{#1}}
\providecommand{\AttributeTok}[1]{\texttt{#1}}
\providecommand{\RegionMarkerTok}[1]{\texttt{#1}}
\providecommand{\InformationTok}[1]{\texttt{#1}}
\providecommand{\WarningTok}[1]{\texttt{#1}}
\providecommand{\AlertTok}[1]{\texttt{#1}}
\providecommand{\ErrorTok}[1]{\texttt{#1}}

% ---- unicode glyphs absent from Latin Modern ----
\usepackage{newunicodechar}
\newunicodechar{§}{\S}
\newunicodechar{−}{\ensuremath{-}}
\newunicodechar{…}{\ldots{}}
\newunicodechar{‖}{\ensuremath{\Vert}}
\newunicodechar{†}{\ensuremath{\dagger}}
\newunicodechar{—}{\textnormal{\textemdash}}
\newunicodechar{∈}{\ensuremath{\in}}
\newunicodechar{∉}{\ensuremath{\notin}}
\newunicodechar{≤}{\ensuremath{\leq}}
\newunicodechar{≥}{\ensuremath{\geq}}
\newunicodechar{≠}{\ensuremath{\neq}}
\newunicodechar{≈}{\ensuremath{\approx}}
\newunicodechar{±}{\ensuremath{\pm}}
\newunicodechar{×}{\ensuremath{\times}}
\newunicodechar{→}{\ensuremath{\rightarrow}}
\newunicodechar{←}{\ensuremath{\leftarrow}}
\newunicodechar{⁻}{\ensuremath{^{-}}}
\newunicodechar{⁰}{\ensuremath{^{0}}}
\newunicodechar{¹}{\ensuremath{^{1}}}
\newunicodechar{²}{\ensuremath{^{2}}}
\newunicodechar{³}{\ensuremath{^{3}}}
\newunicodechar{⁴}{\ensuremath{^{4}}}
\newunicodechar{⁵}{\ensuremath{^{5}}}
\newunicodechar{⁶}{\ensuremath{^{6}}}
\newunicodechar{⁷}{\ensuremath{^{7}}}
\newunicodechar{⁸}{\ensuremath{^{8}}}
\newunicodechar{⁹}{\ensuremath{^{9}}}
\newunicodechar{₀}{\ensuremath{_{0}}}
\newunicodechar{₁}{\ensuremath{_{1}}}
\newunicodechar{₂}{\ensuremath{_{2}}}
\newunicodechar{₃}{\ensuremath{_{3}}}
\newunicodechar{σ}{\ensuremath{\sigma}}
\newunicodechar{π}{\ensuremath{\pi}}
\newunicodechar{θ}{\ensuremath{\theta}}
\newunicodechar{Δ}{\ensuremath{\Delta}}
\newunicodechar{μ}{\ensuremath{\mu}}
\newunicodechar{α}{\ensuremath{\alpha}}
\newunicodechar{β}{\ensuremath{\beta}}
\newunicodechar{₈}{\ensuremath{_{8}}}
\newunicodechar{₄}{\ensuremath{_{4}}}
\newunicodechar{₅}{\ensuremath{_{5}}}
\newunicodechar{₆}{\ensuremath{_{6}}}
\newunicodechar{₇}{\ensuremath{_{7}}}
\newunicodechar{₉}{\ensuremath{_{9}}}
\newunicodechar{↔}{\ensuremath{\leftrightarrow}}
\newunicodechar{∇}{\ensuremath{\nabla}}
\newunicodechar{≳}{\ensuremath{\gtrsim}}
\newunicodechar{≲}{\ensuremath{\lesssim}}
\newunicodechar{⌊}{\ensuremath{\lfloor}}
\newunicodechar{⌋}{\ensuremath{\rfloor}}
\newunicodechar{⌈}{\ensuremath{\lceil}}
\newunicodechar{⌉}{\ensuremath{\rceil}}
\newunicodechar{⇒}{\ensuremath{\Rightarrow}}
\newunicodechar{∝}{\ensuremath{\propto}}
\newunicodechar{∂}{\ensuremath{\partial}}
\newunicodechar{∅}{\ensuremath{\emptyset}}
\newunicodechar{γ}{\ensuremath{\gamma}}
\newunicodechar{δ}{\ensuremath{\delta}}
\newunicodechar{ε}{\ensuremath{\varepsilon}}
\newunicodechar{η}{\ensuremath{\eta}}
\newunicodechar{κ}{\ensuremath{\kappa}}
\newunicodechar{λ}{\ensuremath{\lambda}}
\newunicodechar{ρ}{\ensuremath{\rho}}
\newunicodechar{τ}{\ensuremath{\tau}}
\newunicodechar{φ}{\ensuremath{\phi}}
\newunicodechar{ω}{\ensuremath{\omega}}
\newunicodechar{ν}{\ensuremath{\nu}}
\newunicodechar{χ}{\ensuremath{\chi}}
\newunicodechar{ψ}{\ensuremath{\psi}}
\newunicodechar{ξ}{\ensuremath{\xi}}
\newunicodechar{ζ}{\ensuremath{\zeta}}
\newunicodechar{ι}{\ensuremath{\iota}}
\newunicodechar{Σ}{\ensuremath{\Sigma}}
\newunicodechar{Ω}{\ensuremath{\Omega}}
\newunicodechar{Λ}{\ensuremath{\Lambda}}
\newunicodechar{Φ}{\ensuremath{\Phi}}
\newunicodechar{Γ}{\ensuremath{\Gamma}}
\newunicodechar{√}{\ensuremath{\surd}}
\newunicodechar{∞}{\ensuremath{\infty}}
\newunicodechar{∑}{\ensuremath{\sum}}
\newunicodechar{≪}{\ensuremath{\ll}}
\newunicodechar{≫}{\ensuremath{\gg}}

% ---- figure inclusion helper: \fig{name}{caption}{label} ----
\newcommand{\fig}[4]{%
  \begin{figure}[htbp]\centering
  \includegraphics[width=0.94\linewidth]{figures/#1.pdf}%
  \caption[#2]{#3}\label{#4}%
  \end{figure}}

\begin{document}
\pagenumbering{roman}
\hypersetup{pageanchor=false}

%=====================================================================
% Front matter: title page, certificate, declaration, acknowledgement,
% abstract, contents, list of figures, list of tables.
%=====================================================================
\begin{titlepage}
\begin{tikzpicture}[remember picture,overlay]
  \draw[line width=2.5pt] ([xshift=0.6in,yshift=-0.6in]current page.north west)
    rectangle ([xshift=-0.6in,yshift=0.6in]current page.south east);
\end{tikzpicture}
\centering
\thispagestyle{empty}
\begin{spacing}{1.0}
\vspace*{0.02in}

\includegraphics[width=1.6in]{figures/pes_logo.png}\par
\vspace{0.2in}
{\large\itshape Dissertation on\par}
\vspace{0.5em}
{\Large\bfseries ``Tinker RL Lab:\\[0.35em]
Post-Training of Large Language Models''\par}
\vspace{0.4em}
{\normalsize A multi-framework training harness and studies of GRPO signal starvation\par}
\vspace{0.4em}
{\small\bfseries DERIVED PUBLIC C1 COVERAGE REVISION\quad 3 October 2026\par}

\vfill

{\normalsize\itshape Submitted in partial fulfilment of the requirements for the award of the degree of\par}
\vspace{0.6em}
{\large\bfseries Master of Technology\par}
\vspace{0.2em}
{\normalsize\bfseries in\par}
\vspace{0.2em}
{\large\bfseries Data Science and Machine Learning\par}
\vspace{0.2em}
{\large\bfseries UE20CS972 -- Project Phase -- 2\par}

\vfill

{\normalsize\itshape Submitted by:\par}
\vspace{0.3em}
{\large\bfseries Arvind C R\par}

\vspace{0.5em}
{\normalsize\itshape Under the guidance of\par}
\vspace{0.2em}
{\large\bfseries Ramesh Prakash Guledgudd\par}

\vspace{0.6em}
{\normalsize\bfseries\color{pesmaroon} May 2026 -- October 2026\par}

\vfill

{\normalsize\bfseries DEPARTMENT OF COMPUTER SCIENCE AND ENGINEERING\par}
{\normalsize FACULTY OF ENGINEERING\par}
{\large\bfseries PES UNIVERSITY\par}
{\footnotesize (Established under Karnataka Act No.\ 16 of 2013)\\
Electronic City, Hosur Road, Bengaluru -- 560 100, Karnataka, India\par}
\end{spacing}
\end{titlepage}

%---------------------------------------------------------------------


%---------------------------------------------------------------------
\chapter*{Public Derivative Revision Record}
\addcontentsline{toc}{chapter}{Public Derivative Revision Record}
This derived public revision is based on the 311-page coverage-recovery manuscript dated 3 October 2026. It preserves the scientific findings and their stated limitations while omitting private operational identifiers, account-wide inventories, access correspondence, monetary account details, signature images and student registration data. Private provider exports and original recovery archives are not included. This derivative is not a new institutional certification and does not carry the historical certificate. Historical manuscripts and evidence remain unchanged.

Six original files resolved five named primary-file gaps in the source audit. The public text retains the scientific recovery conclusions and the E14 denominator correction without disclosing the private source inventory. Appendix I explains the publication boundary; the public availability manifest describes what is included and withheld.

The frozen C1 endpoint is 11/64 (17.19\%; Wilson 95\% interval 9.88--28.21\%). Reference errors, material ambiguities and unit mismatches limit its mathematical interpretation. A separate conservative review corroborates three numerical recovery witnesses on the unchanged 64-question denominator. The two counts must be read together; neither is evidence of training improvement.

New material appears in Sections 1.5, 4.6, 7.5, 9.3 and B.13 and Appendix G. Historical V2 results in Section 7.4 remain separate and unpooled. Appendix H is the expanded evidence map. The included case ledger retains classifications, hashes and disagreements. Source descriptions distinguish bundled editable content from evidence retained elsewhere. No frozen source, reference label, selection or primary score was changed. Appendix I reports bounded scientific coverage methods, public-source verification and the remaining limits without private provider inventories. Appendix A corrects historical registry scope and the P11 evidence-tier wording; the reported experimental values are unchanged.

\chapter*{Declaration}
\addcontentsline{toc}{chapter}{Declaration}
\thispagestyle{plain}

I declare that this project report is my own work. Where the work of others has
been used, it is cited and acknowledged. The matter embodied in this report has
not been submitted, in part or in full, to this or any other university or
institution for the award of any other degree or diploma.

The multi-framework benchmark infrastructure, the baseline framework
integrations, and the initial GRPO/PPO/DPO experimentation were produced by the
Semester-3 project group, of which I was a member, and are inherited here as
the shared foundation. The research direction, the experimental expansion, the
statistical analyses, the P1--P8 internal reports, the reporting and registry
artefacts, the adaptive group-size controller, the applied fraud study, and the
E1--E14 evaluation campaign reported in this document are my
individual contribution, carried out under the guidance of my project guide.

Quantitative results are tied to named evidence. Appendix B distinguishes
recovered primary campaign files, inspectable artefacts and the remaining
limits on underlying payload verification. Incomplete, blocked and negative
results are reported rather than omitted. The boundary between inherited and
individual work is stated precisely in \texttt{platform\_hybrid/sem 4
work/PROVENANCE.md}.

\vspace{0.3in}
\noindent\begin{tabular}{@{}p{2.6in}@{}}

\dotfill \\
\textbf{Arvind C R} \\

\end{tabular}

%---------------------------------------------------------------------
% NEEDS AUTHOR CONFIRMATION: drafted from the tools recorded in the repository
% (git Co-Authored-By trailers, frontier_synthesis_*.tex, curriculum_opening/
% FINDINGS.md, zvf-program/jev_lab/, thesis_review_2026-09-27/). Confirm the
% list is complete and matches PES policy wording before submission.
\chapter*{Declaration on the Use of AI Tools}
\addcontentsline{toc}{chapter}{Declaration on the Use of AI Tools}
\thispagestyle{plain}

The following AI tools were used in this work, as recorded in the project
repository.

\begin{itemize}
\item \textbf{Coding and writing assistants.} Anthropic Claude models and
OpenAI Codex were used to write and revise code, analysis scripts and
documents, including drafts of this report. Commits they contributed carry a
\texttt{Co-Authored-By} trailer in the git history.
\item \textbf{Frontier-model synthesis.} ChatGPT Pro Extended and Gemini Deep
Think produced proposed derivations and reasoning recorded in
\longtok{platform\_hybrid/paper/sections/frontier\_synthesis\_*.tex}, which those
documents mark as not being experiments. Where this report uses such a result,
it states its own derivation (Chapter 6 \S6.3.3).
\item \textbf{Independent verification.} An LLM verification pass (Kimi,
6 July 2026) reviewed two single-seed results and flagged both as
\textsc{suspect}; both were then re-tested and reported as nulls (Chapter 6
\S6.4.3).
\item \textbf{Evaluation triage.} The Jev (TypeSafe) judgment harness was used
in the E1--E14 ledger only for classification, faithfulness and value
ordering, never for arithmetic (Chapter 4 \S4.5).
\item \textbf{Review.} The pre-defence self-review of 27 September 2026 and the
revision of this report that followed it were AI-assisted.
\item \textbf{Submission preparation.} An AI coding assistant helped revise
reporting language, harden build and test scripts, check stored arithmetic,
and inspect the final PDF. Gemini 3.1 Pro reviewed the build and verification
code through the \texttt{agy} command-line tool. This preparation ran no new
training or benchmark evaluations and changed no stored experimental results.
\end{itemize}

No AI assistant generated experimental data: every reported measurement comes
from an executed run or evaluation recorded in the repository, and simulated
values are labelled as simulations. The research direction,
the experimental design, the choice of claims and their final wording are the
author's responsibility.

%---------------------------------------------------------------------
\chapter*{Acknowledgement}
\addcontentsline{toc}{chapter}{Acknowledgement}
\thispagestyle{plain}

I thank my project guide, \textbf{Ramesh Prakash Guledgudd}, for direction and
sustained feedback throughout this work, and for the freedom to pursue the
conservative and negative results that form a significant part of it.

I thank the Semester-3 project group for the shared multi-framework benchmark
and the baseline integrations on which the Semester-4 studies build. I thank
the maintainers of the benchmark suites used in the E1--E14 campaign, whose
native evaluators made an independently checkable evaluation possible, and the
platforms --- Tinker, Modal, and Google Cloud --- on which the experiments ran.

Finally, I thank the reviewers whose questions at earlier reviews shaped the
reporting discipline this report depends on: the insistence on stating what a
number does \emph{not} establish has improved the work materially.

\vspace{0.3in}
\noindent\textbf{Arvind C R}

%---------------------------------------------------------------------
\chapter*{Abstract}
\addcontentsline{toc}{chapter}{Abstract}
\thispagestyle{plain}

This report studies group-relative reinforcement-learning post-training of
large language models under one methodological commitment: a comparative
claim is made only when the whole training stack is held fixed except the
factor under test. It makes four contributions.

\textbf{(1) A recomputable signal-starvation diagnostic and a loss-form
audit.} The Zero-Variance Fraction (ZVF) is logged per step and recomputed
from stored reward tensors, with a closed-form Bernoulli null. Reading the
runners' code distinguishes custom REINFORCE group-mean loops from clipped
surrogates, nonzero-KL configuration records and unknown managed settings.

\textbf{(2) Stack-conditioned reporting.} A minimum reporting standard and a
machine-readable stack registry with a \texttt{stackdiff} tool that flags
label-flip risk between two runs from their records alone.

\textbf{(3) Calibrated negative results.} No reliable scaling law is recovered
from reward curves across twelve fitted anchors. A same-stack PPO/GRPO
contrast (0.992 against 0.990, $p = 0.374$, five seeds) and a group-size sweep
find no detectable difference on saturated, in-distribution arithmetic draws;
neither establishes equivalence or an optimum. On an unsaturated GSM8K rerun, GRPO exceeds an
untuned-critic PPO by 5.0 points (paired-$t$ 95\% CI $[1.5, 8.5]$;
$p=0.016$, exact sign-flip $p=0.0625$ at five seeds), a suggestive
configuration-specific result; $G = 2$ and $G = 8$ show no detectable
difference, without establishing equivalence. ZVF has no detectable
association with held-out accuracy in this sample
($\rho = 0.27$, 95\% CI $[-0.37, 0.88]$). An adaptive group-size controller
derived from ZVF reports zero rescue under a plug-in counterfactual model;
92\% of its escalations target observed all-correct groups, without proving
latent certainty or impossibility of recovery with fresh rollouts.

\textbf{(4) An evaluation-governance protocol.} One lightly trained actor
(Qwen3.6-35B-A3B with a LoRA adapter, 40 training steps) is evaluated on
fourteen suites outside its training data under native graders at pinned revisions, with
receipts, a no-pooling rule and a per-lane serving audit that traced a
zero score to a serving defect. Complete scopes include VerilogEval
($129/312 = 41.35\%$), LAB-Bench ($0.2288$) and AgentDojo benign utility
($88/97$); six replacement scopes completed on 27 September 2026
(Chapter~8, Table~8.C); seven original contracts remain externally blocked. A
paired trained-versus-base rerun on 5--151 identical items per lane finds no
significant difference in any lane (inconclusive at those sizes). No
claim of campaign-level improvement over the base model is made.


\textbf{C1 diagnostic extension.} Benchmark reference errors, ambiguous wording and equivalent units limit the interpretation of a new frozen-policy study on public GSM8K train. The first 64 of 87 eligible questions from 4,096 screened questions received 24 fresh responses after eight valid, uncapped parser-zero responses. The frozen endpoint is 11/64 (17.19\%; Wilson 95\% interval 9.88--28.21\%). Assistant-assisted reconciliation corroborates three numerical witnesses on the same denominator, while preserving convention-dependent and ambiguous cases and one wrong-reference match. This is conditional diagnostic evidence, not a population estimate, reasoning validation or training gain. A separate, undeployed parser candidate preserves 2,025 previously valid parse objects and admits four unambiguous extractions in a predefined future subset of 2,048 outputs. It does not alter C1's frozen parser-v2 score.

\vspace{0.25in}
\noindent\textbf{Keywords:} GRPO, RLVR, zero-variance fraction, signal
starvation, stack-conditioned results, benchmark evaluation, reproducibility,
evaluation governance, attribution.

%---------------------------------------------------------------------
\tableofcontents
\listoffigures
\listoftables


\cleardoublepage
\pagenumbering{arabic}
\hypersetup{pageanchor=true}


\part{Foundations}


\chapter{Introduction}\label{introduction}

\section{Background: critic-free RL post-training and the group-relative
estimator}\label{background-critic-free-rl-post-training-and-the-group-relative-estimator}

Post-training is now the stage at which a base language model is turned
into something that reliably follows instructions and solves problems.
For several years the dominant recipe was reinforcement learning from
human feedback (RLHF), in which a learned reward model is trained on
human preference comparisons and then used as the reward signal for a
policy-optimisation loop. RLHF works, but it introduces a second learned
model whose errors the policy can exploit, and it makes reward
construction expensive and hard to audit. The alternative that has taken
hold more recently is reinforcement learning with verifiable rewards
(RLVR): instead of a learned reward model, the reward comes from a
programmatic checker --- exact-match against a reference answer for
grade-school mathematics, unit tests for code, a formal or semi-formal
validator for structured output. The checker is deterministic,
transparent, and free to run. What it gives up in generality it gains in
measurability: every reward is a fact about a generated string rather
than a prediction from a model (Figure~\ref{fig:grpoloop}).


\fig{fig_grpo_loop}
{The GRPO training loop}
{The GRPO training loop. A prompt is sampled $G$ times, each completion is scored by a reward function, advantages are formed relative to the group, and the policy is updated. When all $G$ rewards coincide the group-relative advantage collapses to zero and no reward-relative gradient flows, although the rollouts are still paid for. Number line: the Phase-1 P2 batch ZVF of 0.72--0.77 (four methods, 40 steps, 16 prompts $\times$ 8 samples, GSM8K, seed 0, Qwen3.5-4B) and the group-size sweep on Qwen2.5-0.5B two-digit addition over three seeds, where mean ZVF falls from 0.838 at $G=2$ to 0.631 at $G=16$ (a fall of 0.207).}
{fig:grpoloop}



The workhorse optimiser for both settings has been proximal policy
optimisation (PPO), which optimises a clipped surrogate objective
against a value function critic. The critic is a second network, usually
of comparable size to the policy, trained to predict the expected return
of the current state; its predictions supply the baseline that the
advantage estimator subtracts. This works, but it is expensive and
fragile. The critic consumes memory that could hold policy parameters,
it has its own learning rate and its own failure modes, and its errors
propagate directly into the policy gradient as biased advantages.

Group Relative Policy Optimization (GRPO) removes the critic entirely.
Rather than learn a value function, GRPO samples a group of \(G\)
completions for the same prompt, scores all of them with the verifier,
and derives the advantage of each completion from the group's own reward
statistics. Concretely, for a group with rewards \(r_1, \dots, r_G\),
the advantage assigned to completion \(i\) is a within-group
standardisation,

\[A_i = \frac{r_i - \mu_g}{\sigma_g + \varepsilon}, \qquad \mu_g = \frac{1}{G}\sum_j r_j, \qquad \sigma_g^2 = \frac{1}{G}\sum_j (r_j - \mu_g)^2,\]

which in the original GRPO formulation is then used in a clipped,
KL-regularised policy-gradient update. Several inspected custom runners
omit KL regularisation; the custom-loss Tinker reference loop is
REINFORCE with a group-mean baseline and has no ratio or clipping. This
does not describe every run: TRL configuration records include a
reference policy and nonzero KL, while managed or unrecoverable settings
remain unknown. Table 4.2 gives row-wise loss evidence; the campaign
actor's original execution source is not recoverable (Appendix D). The
empirical baseline is the group mean; there is no value network to
train, no value loss to weight, and no critic to keep in sync with the
policy. The trade is a much simpler optimiser for a much smaller
effective batch, since \(G\) samples of the same prompt now carry one
learning signal rather than \(G\) independent ones.

That trade has a sharp consequence, and it is the observation this
project is organised around. Because the advantage is a
\emph{within-group} standardisation, its numerator is a deviation from
the group mean and its denominator is the group's own spread. When every
completion in a group receives the same verifier outcome --- all
correct, or all incorrect --- the group contributes nothing:
\(\sigma_g = 0\), every numerator is zero, and every advantage \(A_i\)
is exactly zero. The group has been fully sampled and fully scored, at
full token cost, and the reward-relative part of the update is
identically zero. In an objective with a KL penalty, that penalty would
still pull the policy toward the reference; for verified zero-KL
group-relative policy surrogates there is no such contribution, although
auxiliary/value losses and optimiser state can still update parameters.
Exact ties imply zero reward-relative advantage; the tolerance-based
diagnostic in §4.3 also counts near-ties and does not assert exactly
zero gradient for them. We call the fraction of a training batch that
falls into this state the \textbf{Zero-Variance Fraction (ZVF)}, and it
is the single measurement that most of this report's diagnostics are
built around.

ZVF is not a rare edge case, but its size depends on the regime. The
Phase-1 main configuration (Qwen3.5-4B, GSM8K, 40 steps, \(G = 8\))
records a ZVF of approximately 0.72--0.77, meaning that roughly three
quarters of that rollout batch contributes no reward-relative gradient.
Other measured runs span 0.16 (Qwen3-8B on held-out GSM8K at \(G = 8\))
to 0.84 (Qwen2.5-0.5B arithmetic at \(G = 2\)), so no single value is
typical; the measured values are collected in the canonical-numbers
table (Table 1.1). The fraction responds to the one knob that most
directly controls it: on the Qwen2.5-0.5B arithmetic sweep, mean ZVF
falls from 0.838 at \(G = 2\) to 0.631 at \(G = 16\), a decline of 0.207
across that range. Larger groups buy more contrast, with diminishing
returns that saturate well before the group is small enough to be cheap.

There is a second and less comfortable observation in the background,
and it is about attribution rather than mechanism. ``GRPO'' is not a
single algorithm. It is a label that covers implementations in Tinker,
TRL, SkyRL, verl, OpenRLHF and a growing number of research forks, each
of which chooses its own answer parser, sampling temperature, top-\(p\),
KL coefficient, clip range, LoRA target set, precision and evaluation
harness. Two runs that both say ``GRPO'' in their titles may differ in a
dozen ways that plausibly move the outcome, and the reported headline
number does not distinguish between them. The project's own
stack-sensitivity probe makes this concrete: a nominally matched
configuration (Qwen3-8B, seed 42, \(G = 8\), learning rate
\(1\times10^{-5}\), GSM8K-500, 30 steps, visible hyperparameters held
equal, although the managed Tinker run used the Qwen3-8B-Base checkpoint
and TRL the instruction-tuned Qwen3-8B) was run across four framework
paths, two of which completed and two of which are dry-run placeholders,
with results recorded in
\longtok{experiments/results/framework\_comparison.json}. Four separate
marginal one-way ANOVAs over 42 experiments give \(\eta^2 = 0.546\) for
library or framework, \(\eta^2 = 0.558\) for training algorithm,
\(\eta^2 = 0.471\) for model family and \(\eta^2 = 0.322\) for model
scale --- with the important caveat, stated by the source itself, that
these factors are confounded with one another and with back-end, reward
implementation, hardware and task. Because each value is a marginal
per-factor \(\eta^2\) on a collinear corpus, the four sum to 1.897 and
are not a variance decomposition; none estimates a factor's independent
contribution (Chapter 6 §6.1; source:
platform\_hybrid/paper/sections/\_shared\_methods.tex). The framework
factor is nonetheless not a nuisance parameter to be wished away:
marginally, outcomes vary with it about as strongly as with the
algorithm label.

\section{A Motivating Example}\label{a-motivating-example}

Consider a single training step with a group size of \(G = 8\) on an
easy GSM8K prompt that the current policy already solves. All eight
sampled completions produce the correct final answer, so each receives
reward \(r_i = 1\). The group mean is \(\mu_g = 1\) and the group
standard deviation is \(\sigma_g = 0\), so every advantage \(A_i = 0\).
The compute spent on that group is real --- eight full forward-decode
passes at whatever temperature the sampler is configured with --- and
the group supplies no task-learning signal. It contributes zero
reward-relative surrogate gradient in a verified group-centred
implementation; KL and auxiliary losses, and optimiser state, remain
separate. Now imagine a batch of 16 such prompts in which 12 are already
solved and 4 are not: \(12/16 = 0.75\) of the batch is ``dark'' compute.
What the Semester 4 work adds to this picture is that 0.75 is not a
worst case: it is close to the 0.72--0.77 measured on the Phase-1 main
configuration. It is not a universal value either, since the measured
runs range from 0.16 to 0.84 depending on model, task and \(G\) (Table
1.1), so any claim that GRPO is learning efficiently has to be read
against the ZVF of its own regime (Figure~\ref{fig:confounding}).


\fig{fig_confounding}
{The confounding problem this work addresses}
{The confounding problem this work addresses. A single reported reward figure sits on top of an unstated stack; without fixing and reporting that stack the number is not attributable to the method under test. Top right: of the four frameworks in the nominal configuration only Tinker and TRL were run; veRL and OpenRLHF were not installed in the sandbox and have no measurement. The $17\times$ gap between the two measured stacks is an under-specification exhibit, not a clean back-end-only effect, because the managed stack also pinned a different base checkpoint (Qwen3-8B-Base vs.\ Qwen3-8B). Bottom left: one-way $\eta^2$ shares over 42 reported experiments; they are confounded with one another and with back-end, reward implementation, hardware and task, so they do not sum to one. Bottom right: the PPO-versus-GRPO rows of the Phase-1 table, one run per arm, shown without inferential statistics; the Qwen arms did not share a back-end (GRPO on the managed API, PPO on a Modal H100 cluster), so no directional Qwen algorithm claim is made.}
{fig:confounding}



Now suppose, having noticed this, that one wants to answer the obvious
follow-up question: is the group-relative estimator actually a better
choice than a learned critic? The natural experiment is a matched
comparison --- same model, same training prompts, same verifier, same
step budget, same evaluation, with PPO and GRPO as the only difference.
The project's first attempt at that contrast is recorded in a historical
internal report under the label ``Qwen3-8B'', with a Welch \(t\)-test of
\(p = 0.7605\) computed over the final ten steps of a single run per
arm. Those ten steps are autocorrelated observations of one run, not ten
independent runs, and the project's later methods note re-examined the
row and marks it ``provenance conflict; not estimable''. No effect size,
confidence interval or power figure is therefore quoted for it here, and
the Welch value is carried only as a provenance record (Chapter 6
§6.2.3). The project's own rule is that single-seed Tinker experiments
have ``zero statistical power'' and are ``treated as descriptive only
and are not subject to significance testing''.

Two further observations from the same historical source make the
example sharper rather than cleaner. First, the same table reports the
PPO-versus-GRPO contrast on Llama-3.1-8B-Instruct in the opposite
direction; it is also one run per arm, so again no inferential statistic
is quoted, but the direction alone shows that the apparent answer is not
stable across model families. Second, the Qwen arms were not executed on
a common back-end --- the GRPO arm ran through Tinker-managed
infrastructure and the PPO arm on a Modal H100 cluster --- and the
source states this explicitly and declines to draw the algorithmic
conclusion it invites: ``We therefore treat the PPO/GRPO rows as
stack-conditioned evidence, not as an algorithm diagnostic baseline'',
and ``We therefore withhold a directional Qwen algorithm claim''. The
same section records an aggregation gap that illustrates why such care
is warranted: the Qwen PPO last-10 GSM8K figure is 22.5\% in the primary
results ledger but 35.0\% in the statistical summary, and the source
names the ledger value as primary and treats the summary value as
evidence of an aggregation defect rather than as a second measurement.

The genuinely same-stack evidence in the repository is also a null, and
its scope has to be stated with it. The same-stack PPO/GRPO run
(Qwen2.5-0.5B, 40 steps, five seeds, paired by seed but not by
evaluation item) is a two-digit addition task on which both arms are at
ceiling: evaluation accuracy is 0.990 for GRPO and 0.992 for PPO, the
paired difference is −0.002 with p = 0.374, and the 200 evaluation
prompts are fresh random draws from the same generator used in training
rather than a disjoint split; the two arms of a seed share only 3--7 of
their 200 evaluation items, and PPO had drawn about 40\% of its
evaluation items during training against about 7\% for GRPO (§6.2.3;
sources: platform\_hybrid/experiments/results/samestack\_\allowbreak{}ppo\_\allowbreak{}grpo.json;
platform\_hybrid/experiments/modal/modal\_\allowbreak{}samestack\_\allowbreak{}ppo\_\allowbreak{}grpo.py). At
that ceiling a null cannot separate ``no effect'' from ``no headroom'',
so the result is that no difference was detectable on a saturated task,
not that the estimators are equivalent; a paired TOST at a ±0.005 margin
gives p = 0.10 (Table 4.2 gives the exact loss form of both arms). Rerun
on the unsaturated Qwen2.5-1.5B GSM8K-CoT configuration with a disjoint
test split, the contrast is suggestive and baseline-specific: GRPO gains
+4.5 points while a minimal untuned-critic PPO shows no detectable gain,
a post-training gap of +0.050 {[}+0.015, +0.085{]} over five paired
seeds (paired-t p = 0.016; exact sign-flip p = 0.0625) (§6.5.5). The
variant audit (not to be confused with campaign lane E1) takes four
GRPO-family variants (DAPO, GSPO, Dr.~GRPO, AERO) plus the GRPO
reference, runs five arms across eight paired seeds for 40 units at a
frozen stack fingerprint, and evaluates on 500 held-out GSM8K examples
per unit. Every one of the four contrasts is inconclusive. The figures
below are paired-\(t\) 95\% intervals and paired-\(t\) \(p\)-values over
the eight seeds, recomputed from
\longtok{zvf-program/audit/results/full/*.json}: DAPO \(+0.00100\), CI
\([-0.00632, +0.00832]\), \(p = 0.756\); GSPO \(+0.00500\), CI
\([-0.00357, +0.01357]\), \(p = 0.210\); Dr.~GRPO \(-0.00200\), CI
\([-0.01272, +0.00872]\), \(p = 0.673\); AERO \(-0.00075\), CI
\([-0.01029, +0.00879]\), \(p = 0.858\). These paired-\(t\) intervals
replace the audit's own intervals (corrected; see Table E.1). After
Benjamini--Hochberg correction every adjusted \(p\) is at least 0.84.
The pre-audit aggregator's verdict that the differences ``disappear'' is
itself declared invalid, because the locked exact noncentral
paired-\(t\) minimum detectable effect of \(0.010115893\) exceeds the
\(0.010\) margin the study was designed around. The audit's permitted
wording is deliberately narrow: across eight paired seeds no arm
produced a BH-significant held-out difference from GRPO, and none met
the locked exact-power equivalence rule --- which is not the same as
saying the methods are the same, nor that they fail to improve. That
audit runs on Qwen3-8B GSM8K with 500 evaluation examples per unit
(Appendix A, P11).

Taken together, these two halves of the example are the whole motivation
for what follows. The mechanism half shows that a large fraction of the
compute in a GRPO run is not doing what a casual reading of the loss
curve suggests it is doing. The attribution half shows that the most
natural comparative question one would ask about that mechanism --- is
this estimator better than the alternative? --- cannot be answered from
the label ``GRPO'' and ``PPO'' alone, and that the answer on the
project's own evidence is a non-significant contrast that is either
back-end-confounded or measured at a performance ceiling, rather than a
ranking. Neither observation is a defect of the GRPO estimator; both are
defects of what gets measured and what gets reported.

\section{Problem Statement}\label{problem-statement}

The project is organised around four problems, in decreasing order of
how often the literature acknowledges them.

\textbf{Attribution is missing.} Comparative claims about RL
post-training are routinely made across frameworks that differ in
parser, sampler, KL coefficient, clip range, LoRA targets, precision and
evaluator. A reported gain from ``GRPO'' over ``PPO'' is, on most
published evidence, a gain from one complete stack over another complete
stack, with no way to say which component produced it. The project's own
marginal per-factor \(\eta^2\) values put the framework factor at a
magnitude comparable to the algorithm factor, although they are not a
partition of the variance, and its same-stack PPO/GRPO contrast, run on
a saturated 0.5B addition task (evaluation accuracy 0.990 GRPO, 0.992
PPO), returns a paired difference of −0.002 with \(p = 0.374\). The
problem is not that these numbers are surprising; it is that no paper in
this area is obliged to report the information that would let a reader
separate them.

\textbf{Signal starvation is invisible.} Standard training logs report
reward, loss, gradient norm, entropy and completion length. None of
these surfaces the fraction of a batch whose group advantage is exactly
zero. A run can show a rising reward curve and a falling ZVF, or a
rising reward curve with a ZVF pinned near 0.8, and the two look
identical in a loss plot. The fraction of sampling budget buying no
reward-relative gradient is, in the project's own measurements, the
largest measured source of zero-advantage sampling on this fixed setup,
and it is not part of any standard telemetry.

\textbf{There is no reporting or provenance standard.} Even where
authors are willing to describe their stack, there is no
machine-readable form in which to do it, no validator that checks
whether a claim's recorded configuration is internally consistent, and
no accepted minimum set of fields that a stack-conditioned result must
carry. Consequently, results cannot be pooled, cannot be compared across
papers, and cannot be audited after the fact.

\textbf{Evaluation results are rarely reported with their terminal
states.} Benchmark suites in the agentic and reasoning evaluation space
are frequently multi-scope (a public split versus a private original
contract, a benign-utility scope versus an adversarial scope),
quota-limited, or gated behind private payloads. Reports tend either to
omit such suites or to collapse them into a single number, which
conceals how much of the suite was actually graded. The project treats
this as a first-class concern rather than a presentational one: a result
that is genuinely partial should be recorded as partial, with its
coverage denominator and its reopen condition, and a lane that cannot be
run at all should be recorded as externally blocked rather than quietly
dropped.

\section{Objectives}\label{objectives}

The project pursues eight objectives. They are stated as they were
pursued, not retrofitted to the results.

\textbf{O1 --- Build a single harness that can dispatch the same
experiment to multiple back-ends.} Provide one control surface over
Tinker, TRL, SkyRL, verl and OpenRLHF paths so that a configuration can
be held fixed while the framework varies. The harness is the
deliverable; it is not a completed multi-framework benchmark. Only
Tinker and TRL produced completed runs, and those two used different
base checkpoints (Qwen3-8B-Base and Qwen3-8B Instruct), so no framework
comparison is claimed (Chapter 5 §5.2).

\textbf{O2 --- Instrument the training loop so that signal starvation is
directly observable.} Emit per-step telemetry carrying ZVF, gradient
utilisation, policy entropy, completion length and KL divergence, and
persist the raw per-group reward tensors so that any derived statistic
can be deterministically recomputed after the fact rather than trusted
from a log line.

\textbf{O3 --- Measure the zero-variance regime on real reward tensors
and test the levers that might move it.} Characterise ZVF empirically on
genuine GRPO reward tensors, and investigate re-baselining and
curriculum interventions as ways of converting dark compute into
contrast.

\textbf{O4 --- Establish and hold a scientific protocol.} Multi-seed
runs, matched baselines, held-out evaluation with explicit denominators,
deterministic recomputation, declared (not pre-registered) effect-size
and power accounting, and reporting of negatives and nulls --- including
the cases where the answer is that the study was underpowered to answer.

\textbf{O5 --- Scope the most original direction in the field.} A
literature scan for white space, which in this project identified stack
provenance and reporting standards as the unclaimed ground and produced
the minimum reporting standard (P5) and the machine-readable stack
registry (P6).

\textbf{O6 --- Turn the diagnostic into an intervention.} Derive an
adaptive group-size controller from the ZVF measurement and evaluate it
as a training-time policy rather than a post-hoc observation (P7).

\textbf{O7 --- Build and exercise an evaluation-governance protocol on
one trained actor.} Run a single trained checkpoint against fourteen
agentic and reasoning suites that were not part of its training data,
each with its own native evaluator at a pinned revision, under receipt,
budget and decontamination controls, and record every lane's terminal
state. The contribution is the governance protocol and the receipt-bound
record it produces. The campaign is not evidence about GRPO: the actor's
training run was short (40 steps, Chapter 8 §8.1.1), and the only
trained-versus-base comparison is the small-sample Table 8.B.

\textbf{O8 --- Report under a deliberate asymmetry.} Where the evidence
supports a structural or descriptive claim, state it; where it does not,
withhold it. No cross-suite aggregate, no pooling of original-contract
and replacement-scope numbers, and no claim of improvement over the base
model in the E1--E14 campaign. Its paired trained-versus-base comparison
(Chapter 8, Table 8.B) is inconclusive in every lane. The controlled
studies in Chapter 6 separately report pre-to-post gains and the M1b
GRPO--PPO contrast, limited to the tested stack and an untuned critic;
five seeds give an exact sign-flip p-value of 0.0625 despite the paired
t-test p-value of 0.016.

The objectives are what the project set out to build. The results
chapters answer four research questions that the objectives make
testable. Each is mapped to the table that answers it, and §9.3 gives a
one-line answer to each.

\begin{itemize}
\tightlist
\item
  \textbf{RQ1.} Is the Zero-Variance Fraction recomputable from stored
  reward tensors, and does it carry information about training outcome?
  Answered by the ZVF rows of Table 1.1 and §6.3.
\item
  \textbf{RQ2.} With the training stack held fixed, do algorithm labels
  (PPO vs GRPO, GRPO vs Dr.~GRPO) or group size separate on evaluation
  accuracy? Addressed by the same-stack rows of Table 1.1, the loss-form
  table in §4.3, and §6.2.3, §6.4 and §6.5. The experiments reported
  here can bound this question but not answer it: the PPO/GRPO,
  group-size and arithmetic GRPO/Dr.~GRPO contrasts were all run on a
  saturated task at evaluation accuracy near 0.99, and the two
  unsaturated GRPO/Dr.~GRPO comparisons (§1.2, §6.5.4) are inconclusive
  and establish no equivalence (§9.3). A same-stack rerun on an
  unsaturated task (§6.5.5) suggests a GRPO advantage over a minimal
  untuned-critic PPO, sensitive to the test used, while the G = 2 versus
  G = 8 contrast is inconclusive.
\item
  \textbf{RQ3.} Can a machine-readable stack manifest flag label-flip
  risk between two runs from their records alone? Answered by §7.2
  (registry health and \texttt{stackdiff}).
\item
  \textbf{RQ4.} What does a governance-first evaluation of one trained
  actor yield, and what can it not establish? Answered by Table C.3 and
  Tables 8.A and 8.B.
\end{itemize}

\textbf{Contributions.} Measured against these objectives, the thesis
makes four contributions. Everything else in it is supporting evidence
for one of them.

\begin{enumerate}
\def\labelenumi{\arabic{enumi}.}
\tightlist
\item
  \textbf{A recomputable signal-starvation diagnostic and a loss-form
  audit.} ZVF is logged per step and recomputed from stored reward
  tensors against a closed-form Bernoulli null (§4.3, §6.3). Reading the
  runner code distinguishes custom REINFORCE group-mean loops from
  clipped surrogates, nonzero-KL configuration records and unknown
  managed settings (Table 4.2).
\item
  \textbf{Stack-conditioned reporting.} A minimum reporting standard
  (P5) and a machine-readable stack registry with \texttt{stackdiff},
  which flags label-flip risk between two runs from their records alone
  (Chapter 7).
\item
  \textbf{Calibrated negative results.} No reliable scaling law is
  recovered from reward curves; the same-stack PPO/GRPO and group-size
  contrasts are nulls on a saturated task, and on an unsaturated rerun
  only the PPO/GRPO gap appears, from an untuned critic; ZVF has no
  detectable association with held-out accuracy in the measured sample;
  and the ZVF-derived controller's zero-rescue counterfactual is imposed
  by a plug-in model fitted to observed all-equal groups, not a test of
  fresh-rollout recoverability. Each is reported with the power
  limitation that bounds it (Chapters 6--7).
\item
  \textbf{An evaluation-governance protocol.} Native graders at pinned
  revisions, per-lane receipts, a no-pooling rule, a per-lane serving
  audit and a paired same-engine trained-versus-base design, exercised
  on one lightly trained actor across fourteen suites outside its
  training data (Chapter 8). The contribution is the protocol and the
  record it produces, not the actor's scores.
\end{enumerate}

\section{Scope and Non-Goals}\label{scope-and-non-goals}

The C1 extension dated 3 October 2026 adds a frozen-policy diagnostic
whose chief limitation is the reference itself. Exact-match zeros can
reflect an incorrect key, ambiguous wording or equivalent units as well
as a wrong response. Among the first 64 eligible questions from a fixed
4,096-question screen, the frozen score records 11 fresh-response
events; assistant-assisted review corroborates three numerical witnesses
on the same denominator. This extension separates scoring from semantic
support and does not imply reasoning improvement, useful RL gradients or
training gain. Section 7.5 and Appendix G retain the full accounting,
including disagreements. Historical V2 remains separate.

The scope of this report is measurement and diagnostics, not leaderboard
performance. The controlled studies and the campaign actor use
small-to-mid open models --- the Qwen2.5, Qwen3, Qwen3.5 and Qwen3.6
families, from roughly 0.5B parameters up to the 35B sparse
mixture-of-experts configuration used for the campaign actor. The P1
cross-scale roster is wider, 0.6B to 1T total parameters, but its
frontier-scale anchors are single-seed runs through a managed API and
are read descriptively (Chapter 6 §6.1; source:
platform\_hybrid/paper/sections/p1\_abstract.tex). Verifiable rewards
come from mathematics and code checkers, principally GSM8K exact-match
and HumanEval-style unit tests. Held-out evaluation uses suites' own
native evaluators at pinned revisions.

Several things are explicitly out of scope. This is not an attempt to
set a state-of-the-art result on any benchmark, and no such claim is
made anywhere in the report. It is not a study of frontier reasoning
models. The evaluation sets used in the controlled comparisons are
small, and their sizes differ by study: 200 fresh prompts for the 0.5B
arithmetic runs, 200 per seed for the 1.5B GSM8K chain-of-thought runs,
500 for the variant audit, and only 8 to 20 in several Phase-1 openings
(the layer-freeze result rests on n = 8). Table 1.1 gives the n behind
each headline number. Claims about novel algorithms are not made: the
length-weighting intervention (P4) and the adaptive controller (P7) were
designed during the work and are reported as designed and simulated
systems with their evaluation limits attached, not as demonstrated
improvements.

The report states negative, null and withdrawn results alongside
positive ones; this is its one general statement about reporting
practice, and later chapters report the results without repeating it.
Three evidentiary rules apply throughout and are worth stating here
because they constrain what the later chapters are allowed to say.
First, every number in this report is traceable to a repository file
that was actually read; where a figure could not be sourced, the finding
is described qualitatively and the figure is attributed to its source
rather than reconstructed. Second, original-contract and
replacement-scope benchmark results are never pooled and never averaged
across suites --- the two are different measurements of different things
under different licences, and treating them as one column would
manufacture a number that no suite reported. Third, no improvement over
a baseline or over a base model is claimed unless the source states it
with a test; where the source reports a difference without a test
supporting it, or reports an inconclusive test, that is what is written
here.

\section{Contribution of the
Candidate}\label{contribution-of-the-candidate}

This report covers work that spans two semesters and two very different
modes of authorship, and the distinction matters for how the
contribution should be assessed. Semester 3 was a group project (Group
6). The multi-framework repository, the baseline integrations, the
initial Tinker, SkyRL and GRPO experiments, the literature foundation,
and the early benchmark, evaluation and reproducibility scaffolding were
produced by that group, of which the candidate was a member. Semester 4
is the candidate's individual continuation, and the new research
direction, the experiment expansion, the analysis, the report writing
and the P1--P8 internal reports are the candidate's solo work (Figure~\ref{fig:contribution}).


\fig{fig_contribution_map}
{Boundary between the inherited Semester-3 group infrastructure and the individual Semester-4 contribution reported in this document}
{Boundary between the inherited Semester-3 group infrastructure and the individual Semester-4 contribution reported in this document. Both columns are drawn from the repository's own provenance record (the Semester-3 tag and boundary commit, the post-boundary commit and file counts, the solo-authorship commit, and the Semester-4 README with its paper-to-evidence map); where the record does not settle a claim, the claim is left out rather than inferred.}
{fig:contribution}



The boundary is recorded precisely rather than loosely. The inherited
state is frozen at the Git tag \texttt{capstone-final-2026-04-25},
corresponding to boundary commit
\longtok{21a99ef766543a1ca86bfbad188445d3596ea73c} of 23 April 2026. At
the time the separation was documented, 584 commits had been made after
that boundary, touching 5,050 changed files --- a figure the provenance
record explicitly declines to present as a measure of academic
contribution, since it counts mechanical and generated changes alongside
substantive ones. The transition to explicit solo authorship is marked
by commit \longtok{282e6b4f0f158de11b3f5bb4656e2d871f68f5c0} of 7 July
2026.

What was added after that boundary is the substance of this report. It
includes the eight unpublished internal research reports: P1, a
cross-library and cross-scale study of GRPO behaviour; P2, the
Zero-Variance Fraction as a descriptive diagnostic; P3, group size,
contrast density and the relationship to DPO; P4, length bias and
held-out generalisation; P5, the minimum reporting standard for
stack-conditioned results; P6, a machine-readable GRPO stack registry;
P7, an adaptive group-size controller derived from the signal-starvation
diagnostic; and P8, a side study comparing a language model against
gradient-boosted trees in credit-card fraud detection. It includes the
post-capstone experimental, audit, synthesis and report iterations ---
206 recorded analysis iterations --- the Berkeley-course-derived
research audits and prototypes, the expanded statistical analyses, and
the new GRPO diagnostics and controllers. And it includes the E1--E14
benchmark campaign, in which a single Tinker-trained actor (base
\texttt{Qwen/Qwen3.6-35B-A3B} at commit \texttt{995ad96e…}, with the
\texttt{campaign-adapter} adapter at commit
\texttt{withheld-adapter-revision…}, in bfloat16) was evaluated against
fourteen suites under receipt, budget and decontamination governance.
Its results are reported per suite in Chapter 8 and never pooled; Table
C.3 is the single source for each lane's scope and terminal state.

Two further contributions deserve naming because they are contributions
of restraint rather than of result. The first is that the campaign's
negative findings were kept rather than discarded. A later
BankerToolBench report records 100 base-model trials with mean reward
0.0. The underlying trajectories remain unavailable, and the earlier
recovered finish ledger predates that rerun. The bridge's missing stop
sequences provide a plausible failure mechanism, but the available
comparison does not identify its contribution to the full aggregate
zero. The finding is a source-limited, confounded diagnostic rather than
a controlled model-capability estimate (Chapter 8 §8.5). The second is
the discipline applied to the project's own claims: the same-stack audit
that invalidated the pre-audit aggregator's ``disappears'' verdict did
so on the strength of the study's own locked minimum detectable effect,
and the permitted wording it left in place is narrower than the finding
it refuted.

The report does not reallocate the shared runtime code. The inherited
repository remains where it was; the Semester 4 internal reports cite
their canonical LaTeX sources and evidence directories rather than
duplicating them, and the distinction recorded here is a distinction
about academic contribution, not about file ownership.

\section{Organisation of the Report}\label{organisation-of-the-report}

The report is organised in four parts.

\textbf{Part I --- Foundations} contains the present chapter, followed
by Chapter 2, the literature survey, which situates the work against the
RLHF and RLVR literature, the critic-free line of estimators, the
group-size and advantage-estimation literature, and the reproducibility
and reporting literature on which the provenance argument depends; and
Chapter 3, the system requirements specification, which states the
functional and non-functional requirements the harness, the telemetry
layer and the campaign infrastructure were built to satisfy.

\textbf{Part II --- Design and Implementation} contains Chapter 4, the
proposed methodology, which sets out the measurement design, the
statistical protocol, the evidence rules and the study-to-evidence
mapping; and Chapter 5, the implementation details, covering the
multi-framework harness, the per-step telemetry and raw tensor
persistence, the campaign launcher and its receipt and budget
governance, and the registry and validator components.

\textbf{Part III --- Results} contains three chapters. Chapter 6 reports
the core GRPO studies --- the cross-library and cross-scale scaling
work, the ZVF diagnostic, the group-size and contrast-density sweep, and
the length-bias and held-out generalisation study, corresponding to P1
through P4. Chapter 7 reports the reporting, registry and controller
work, corresponding to P5 through P7, including the minimum reporting
standard, the machine-readable stack registry, and the simulated
behaviour of the adaptive group-size controller. Chapter 8 reports the
E1--E14 benchmark campaign: the actor, the governance model, the
strictly completed scopes, the original-contract scores, the incomplete
and blocked lanes, the verification results and the campaign's terminal
states.

\textbf{Part IV --- Synthesis} contains Chapter 9, which draws the
cross-study synthesis, states the threats to validity, gives the
conclusions the evidence will bear, and sets out future work and the
reopen conditions under which the blocked lanes could be completed.

Nine appendices follow. Appendix A gives the historical run registry and
its identity reconciliation. Appendix B documents reproducibility and
provenance, including the inherited-versus-individual boundary and the
environment records needed to reproduce the reported results. Appendix C
gives the notation, symbols and abbreviations used throughout, and its
Table C.3 is the single record of each E1--E14 lane's scope, status and
receipt. Appendix D holds the lane-by-lane campaign detail moved out of
Chapter 8. Appendix E lists every correction applied to an internal
source document. Appendix F reports the applied fraud-detection side
study (P8), which uses synthetic data and supervised fine-tuning and so
sits outside the RL contribution. Appendix G records the C1 case ledger,
reference classifications and source bindings. Appendix H is the
evidence map for the preceding scientific chapters and appendices,
arranged section by section. Appendix I records the cross-service
experiment and checkpoint coverage audit, including exact inventory
boundaries and unresolved evidence gaps.

\section{Canonical Numbers}\label{canonical-numbers}

Several quantities in this report were measured more than once, in
different regimes, and appear in the sources with different values.
Table 1.1 fixes one row per measurement, with the model, task and group
size it belongs to, the number of seeds and evaluation items behind it,
and the file it comes from. Where the text quotes one of these
quantities, this table is the reference, and a value that differs from
its row is a different measurement, not a correction. Per-lane E1--E14
figures are kept in Table C.3.

\textbf{Canonical numbers.}

\begingroup\scriptsize\setlength{\tabcolsep}{3pt}
\begin{longtable}[]{@{}
  >{\raggedright\arraybackslash\hspace{0pt}}p{(\linewidth - 10\tabcolsep) * \real{0.2000}}
  >{\raggedright\arraybackslash\hspace{0pt}}p{(\linewidth - 10\tabcolsep) * \real{0.2000}}
  >{\raggedright\arraybackslash\hspace{0pt}}p{(\linewidth - 10\tabcolsep) * \real{0.2000}}
  >{\raggedright\arraybackslash\hspace{0pt}}p{(\linewidth - 10\tabcolsep) * \real{0.2000}}
  >{\raggedright\arraybackslash\hspace{0pt}}p{(\linewidth - 10\tabcolsep) * \real{0.2000}}@{}}
\caption{Canonical numbers. ``results/'' is
\protect\longtok{platform\_hybrid/experiments/results/}; ``sections/'',
``archive/'' and ``neurips\_\allowbreak{}2026\_\allowbreak{}variants/'' are under
\protect\longtok{platform\_hybrid/paper/}. The Qwen3-8B GSM8K GRPO
training-reward rows use different summaries (last-10 mean, trace mean,
pooled) and the sources do not establish that they come from the same
run, so they are not interchangeable.}\tabularnewline
\toprule\noalign{}
\begin{minipage}[b]{\linewidth}\raggedright
Quantity
\end{minipage} & \begin{minipage}[b]{\linewidth}\raggedright
Value
\end{minipage} & \begin{minipage}[b]{\linewidth}\raggedright
Model / task / G
\end{minipage} & \begin{minipage}[b]{\linewidth}\raggedright
n
\end{minipage} & \begin{minipage}[b]{\linewidth}\raggedright
Source
\end{minipage} \\
\midrule\noalign{}
\endfirsthead
\toprule\noalign{}
\begin{minipage}[b]{\linewidth}\raggedright
Quantity
\end{minipage} & \begin{minipage}[b]{\linewidth}\raggedright
Value
\end{minipage} & \begin{minipage}[b]{\linewidth}\raggedright
Model / task / G
\end{minipage} & \begin{minipage}[b]{\linewidth}\raggedright
n
\end{minipage} & \begin{minipage}[b]{\linewidth}\raggedright
Source
\end{minipage} \\
\midrule\noalign{}
\endhead
\bottomrule\noalign{}
\endlastfoot
Phase-1 batch ZVF & 0.72--0.77 (all-correct 0.65--0.71) & Qwen3.5-4B,
GSM8K, G = 8; four methods × 40 steps, 16 prompts per step & 1 seed &
\longtok{Phase1\_\allowbreak{}Project\_\allowbreak{}Report\_\allowbreak{}ZVF.tex} \\
ZVF, sampling only (P2 recomputation) & 0.130 / 0.190 / 0.155; pooled
0.158 & Qwen3-8B, GSM8K test problems, G = 8, T = 1.0, no training & 3
seeds × 200 problems &
\longtok{results/tinker\_\allowbreak{}gsm8k\_\allowbreak{}zvf\_\allowbreak{}s\{42,123,456\}.json} \\
ZVF by group size & 0.838 / 0.764 / 0.691 / 0.631 at G = 2 / 4 / 8 / 16
(fall of 0.207) & Qwen2.5-0.5B, two-digit addition & 3 seeds × 40 steps
per G & \longtok{results/groupsize\_\allowbreak{}zvf\_\allowbreak{}sweep.tsv} \\
ZVF, nine-row mitigation panel (``vanilla GRPO'' 0.481) & declared
simulation projection, not a measurement & synthetic, K = 8 & --- &
\longtok{sections/zvf.tex} \\
ZVF vs final held-out outcome & ρ ≈ 0.27, 95\% CI {[}−0.37, 0.88{]} &
pooled local analysis cells, GSM8K & n = 23 cells &
\longtok{sections/p2\_abstract.tex}; Table A.5 (P2-C2) \\
ZVF vs performance, pooled across task types & r = −0.769 (confounded by
task type) & 5 model families, tool-use and GSM8K & N = 15 experiments &
\longtok{neurips\_\allowbreak{}2026\_\allowbreak{}variants/paper\_\allowbreak{}P8\_\allowbreak{}workshop.tex} \\
ZVF vs final accuracy & r ≈ 0.09 & instrumented runs & n = 12 &
\longtok{sections/appendix\_\allowbreak{}zvf\_\allowbreak{}formalization.tex} \\
Same-stack PPO vs GRPO & GRPO 0.990, PPO 0.992; paired Δ = −0.002, p =
0.374 & Qwen2.5-0.5B, two-digit addition, G = 8 (PPO: 128 prompts per
step) & 5 seeds (paired by seed, not by item) × 200 fresh draws per arm
from the training generator &
\longtok{results/samestack\_\allowbreak{}ppo\_\allowbreak{}grpo.json} \\
Same-stack GRPO vs PPO, unsaturated & post-training accuracy GRPO (G =
8) 0.245, GRPO (G = 2) 0.250, PPO 0.195; GRPO G8 − PPO +0.050 {[}+0.015,
+0.085{]}, paired-t p = 0.016, exact p = 0.0625; G8 − G2 −0.005
{[}−0.038, +0.028{]} & Qwen2.5-1.5B-Instruct, GSM8K-CoT, 200-token cap,
64 generations per step, 30 steps & 5 paired seeds × 200 GSM8K test
items & \longtok{results/samestack\_\allowbreak{}gsm8k\_\allowbreak{}cot.json} \\
Group-size sweep, evaluation accuracy & 0.982 / 0.988 / 0.990 / 0.978 at
G = 2 / 4 / 8 / 16 & Qwen2.5-0.5B, two-digit addition & 3 seeds × 200
fresh draws & \longtok{results/groupsize\_\allowbreak{}zvf\_\allowbreak{}sweep.tsv} \\
GRPO vs Dr.~GRPO, arithmetic & 0.987 vs 0.992, p = 0.35; completions 4.7
tokens & Qwen2.5-0.5B, G = 8, 40 steps & 5 seeds × 200 &
\longtok{results/drgrpo\_\allowbreak{}vs\_\allowbreak{}grpo.json} \\
GRPO / Dr.~GRPO pre → post & 0.202 → 0.263 / 0.205 → 0.255 &
Qwen2.5-1.5B-Instruct, GSM8K-CoT, 200-token cap, G = 8 & 3 seeds × 200 &
\longtok{results/drgrpo\_\allowbreak{}gsm8k\_\allowbreak{}cot.json} \\
Qwen3-8B pre → post (GRPO) & 82.0\% → 83.3\%, p = 0.26 &
Qwen3-8B-Instruct, GSM8K; 82.0\% is the same checkpoint before RL & 5
seeds & \longtok{archive/absorbed/U01\_\allowbreak{}main\_\allowbreak{}compendium/main.tex} \\
Variant audit, DAPO − GRPO & Δ = +0.001, paired-t 95\% CI {[}−0.00632,
+0.00832{]}; no BH rejection & Qwen3-8B, GSM8K, five arms & 8 seeds ×
500 & \longtok{zvf-program/audit/results/full/};
\longtok{statistical\_refutation.md} \\
Tinker vs TRL, Qwen3-8B GSM8K & last-10 training reward 0.856 (Tinker,
Qwen3-8B-Base) vs 0.050 (TRL, Qwen3-8B Instruct) & GSM8K-500, G = 8, 30
steps & 1 seed each & \longtok{results/framework\_comparison.json} \\
Qwen3-8B GSM8K GRPO, reproduction headline & last-10 training reward
34.4\%, peak 62.5\% & Qwen3-8B, LoRA 32, G = 8, 30 steps & seed 42 &
\longtok{REPRODUCE.md} \\
Qwen3-8B GSM8K GRPO, P1 anchor fit & mean training reward 0.285, peak
0.625 & Qwen3-8B, GSM8K & 1 seed &
\longtok{sections/scaling\_laws.tex} \\
Pooled Tinker vs TRL training reward & 99.9\% vs 73.4\% & Tinker runs
pooled over 4B--235B models vs TRL on Qwen2.5-0.5B & pooled &
\longtok{LIMITATIONS\_\allowbreak{}AND\_\allowbreak{}IMPACT.md} §6.2 \\
Layer-freeze held-out gain & +0.0625 (full) vs +0.000 (frozen) &
Qwen2.5-1.5B-Instruct & 2 seeds × 8 items &
\longtok{results/p1\_layerfreeze/FREEZE\_FINDINGS.md} \\
Curriculum opening held-out gain & +0.05 in both arms & Qwen3.5-4B,
GSM8K, G = 4 & 1 seed × 20 items &
\longtok{results/curriculum\_opening/FINDINGS.md} \\
E1--E14 lane figures & see Table C.3 & Qwen3.6-35B-A3B + seed809 adapter
& per-lane native denominators &
\longtok{outputs/E1\_\allowbreak{}E14\_\allowbreak{}FINAL\_\allowbreak{}RESULTS\_\allowbreak{}2026-09-19.md} \\
\end{longtable}
\endgroup



\chapter{Literature Survey}\label{literature-survey}

This chapter surveys the prior work on which the project is built and
the gap it occupies. It revises the survey of the Phase-1 report,
retaining that survey's seven-part structure and updating each part with
the material the Semester 4 work added: the P1--P8 internal reports and
the E1--E14 evaluation campaign. Material on RLHF, PPO, DPO, GRPO,
DeepSeek-R1, parameter-efficient fine-tuning and evaluation benchmarks
also draws on the Semester-3 Group 6 capstone literature survey
\citep{group6capstone2026survey}, a group work of six authors of whom
the present author was one; where this chapter reuses that survey's
findings they are attributed to it here. Two rules of construction carry
over from Phase 1. The first concerns proportion. The contribution of
this project is overwhelmingly a contribution of measurement and
reporting practice rather than of a new training algorithm, so the
literature that establishes measurement practice receives more room here
than the literature that establishes methods. The second concerns
evidence. A figure is quoted only where a repository source actually
states it; where a source states a finding without giving a number, the
finding is described qualitatively and the chapter says so. Where the
project's own documents state the same quantity differently, the
divergence is reported rather than silently reconciled, and Section 2.5
records one such divergence in full.

\section{From RLHF and PPO to Group-Relative
Methods}\label{from-rlhf-and-ppo-to-group-relative-methods}

Policy-gradient post-training of language models begins with the RLHF
pipeline. Christiano et al. \citep{christiano2017deeprlhf} established
the pattern of reward-model training on human pairwise preferences
followed by reinforcement learning against that reward model, on Atari
and MuJoCo tasks starting from a randomly initialised policy; the
supervised fine-tuning stage that precedes both in the language-model
pipeline was added in later work
\citep{stiennon2020summarize, ouyang2022instructgpt}. Ziegler et al.
\citep{ziegler2019finetuning} adapted the pattern to language,
collecting approximately 60,000 comparisons and, following Jaques et al.
\citep{jaques2017sequencetutor, jaques2019wayoffpolicy}, adopting the
per-token KL penalty against a frozen reference policy that has been
standard ever since. Stiennon et al. \citep{stiennon2020summarize}
scaled the recipe to a 6.7B summarisation model with 64,832 comparisons,
demonstrated transfer to an out-of-distribution summarisation set
(CNN/DailyMail), and produced early empirical evidence of reward
over-optimisation: continued optimisation against a frozen reward model
eventually degrades the true objective. Ouyang et al.
\citep{ouyang2022instructgpt} then built InstructGPT, whose three-stage
pipeline used roughly 40 contractors to produce about 13,000 supervised
fine-tuning prompts, about 33,000 reward-model prompts, and about 31,000
PPO prompts. Two details of that work matter for this project. First,
the reward model was 6B rather than 175B because the larger reward model
proved unstable during training --- a reminder that the reward model is
itself a trained component with its own failure modes. Second, the
preference data were collected as rankings of K ∈ \{4, \ldots, 9\}
responses, which yields C(K, 2) pairwise comparisons from a single
annotation batch; this is precisely the within-prompt comparative
structure that group-relative methods would later exploit directly.
InstructGPT's headline results were that a 1.3B InstructGPT was
preferred to the 175B GPT-3 base model, that the 175B InstructGPT was
preferred to GPT-3 in 85\% of evaluations, that hallucination on a
closed-domain task fell from 41\% to 21\%, that toxicity was reduced by
roughly 25\%, and that TruthfulQA performance roughly doubled;
inter-annotator agreement was approximately 73\%. The same source
records four limitations --- annotation cost, reward-model
over-optimisation, the difficulty of aligning to user intent rather than
to annotator preference, and the residual capability costs of alignment (Figure~\ref{fig:lineage}).


\fig{fig_lineage}
{Algorithmic lineage from RLHF through PPO to GRPO, with DPO as an orthogonal branch}
{Algorithmic lineage from RLHF through PPO to GRPO, with DPO as an orthogonal branch. Each transition removes a component or substitutes a cheaper estimator; GRPO removes the critic and replaces it with a group-relative baseline. DPO (dashed frame) removes the on-policy loop and the reward model rather than simplifying the critic, trading online exploration for stability. Dates are first publication or first use at scale. GRPO variants (DAPO, Dr.~GRPO and others) are discussed in the text and not drawn.}
{fig:lineage}



The optimiser for that pipeline was PPO \citep{schulman2017ppo}. PPO's
clipping mechanism, with ε ≈ 0.2, restricts each update's policy ratio,
and generalised advantage estimation \citep{schulman2016gae} with its
standard defaults (λ = 0.95, γ = 0.99) provides the advantage estimate.
Recasting language generation as an RL problem makes the state the
prompt plus generated prefix, the action the next token, the policy the
language model, and the reward sparse and terminal. The practical
consequence is a four-model architecture --- policy, reference, reward,
and value --- whose mixed-precision memory footprint is roughly four
times that of supervised fine-tuning alone; PPO-ptx mitigates capability
regression by mixing pretraining gradients into the update. Zheng et
al.'s ``Secrets of RLHF'' study \citep{zheng2023secretsrlhf} found that
policy constraints, not algorithmic choice, are the single most critical
factor in whether RLHF succeeds: without them, training exhibits reward
hacking, length explosion, and catastrophic pattern collapse. That
finding is the intellectual ancestor of the zero-variance analysis in
this project, because it locates the failure mode in the constraint
structure rather than in the objective.

Group Relative Policy Optimization (GRPO), introduced with DeepSeekMath
\citep{shao2024deepseekmath}, removes the learned value network
altogether. For each prompt it samples a group of G completions, scores
them with a reward function, and standardises each completion's reward
against the group mean and standard deviation; the resulting z-score
serves as the advantage for every token of that completion. In the
DeepSeekMath configuration G = 64 and β = 0.04. Eliminating the value
model reduces the number of trained and served models from four to
three; the paper motivates this as reducing PPO's memory usage but does
not quantify the saving. DeepSeekMath-RL 7B reached 51.7\% top-1
accuracy on MATH against a 46.8\% supervised fine-tuning baseline,
60.9\% with self-consistency over 64 samples, GSM8K improved from 82.9\%
to 88.2\%, and CMATH improved from 84.6\% to 88.8\% despite no Chinese
data in the RL training set. The paper's iterative variant retains 10\%
of historical data in a replay buffer to prevent catastrophic forgetting
in the co-evolving reward model. RLOO \citep{ahmadian2024backtobasics}
is the closely related leave-one-out estimator that baselines each
sample against the mean of the \emph{other} samples in the group; it
already performs a within-group baseline, and is therefore directly
relevant to the signal-starvation analysis this project conducts.

DeepSeek-R1 \citep{guo2025deepseekr1} then demonstrated that pure
outcome-reward RL can elicit chain-of-thought behaviour from a base
model without supervised cold-start data. R1-Zero was trained on
DeepSeek-V3-Base, a 671B mixture-of-experts model, using rule-based
rewards only and no neural reward model. AIME 2024 pass@1 rose from
15.6\% to 71.0\% over training in the original preprint --- the
\emph{Nature} version cited here reports 77.9\% after extended training
--- with consensus decoding over 64 samples (cons@64) reaching 86.7\%
and MATH-500 reaching 95.9\%. The source states the training
configuration as 16 sampled outputs per question (the group size), a
training batch of 512 outputs (32 questions × 16 samples), a learning
rate of 3 × 10⁻⁶, a KL coefficient of 0.001, and approximately 10,400
policy update steps. The paper documents spontaneous chain-of-thought,
self-verification, and the widely reported ``aha moment'' as emergent
rather than taught behaviours. The full R1 model, trained with a
four-stage pipeline that includes a language-consistency reward, reports
79.8 on AIME, 97.3 on MATH-500, and 71.5 on GPQA, against 79.2 / 96.4 /
75.7 for o1-1217. Distillation into six students from 1.5B to 70B
produced Distill-Qwen-1.5B at 28.9\% AIME and 83.9\% MATH-500, and
Distill-Qwen-32B at 72.6\% and 94.3\%, against 47.0\% and 91.6\% for
Qwen2.5-32B-Zero trained from scratch; the 7B and 16B base models failed
to show meaningful reasoning gains from direct RL.

For this project, the significance of that lineage is narrower than the
headline numbers suggest. GRPO's defining operation is a within-group
standardisation, and a standardisation is undefined when the group's
rewards are all identical. That degenerate case is not an edge case: it
is the modal case early in training on a hard task and late in training
on an easy one. The rest of this chapter and the empirical chapters that
follow are concerned with what that degeneracy does to the training
signal, and with how honestly the field currently reports it.

\section{Direct Preference Optimisation as a
Contrast}\label{direct-preference-optimisation-as-a-contrast}

An orthogonal branch of post-training removes on-policy sampling
entirely. Direct Preference Optimisation (DPO) \citep{rafailov2023dpo}
exploits the fact that the RLHF objective with a KL constraint admits a
closed-form optimal policy, which permits the reward to be
reparameterised in terms of the policy itself. The result is a simple
classification loss over a fixed dataset of preference pairs, with no
reward model, no value model, and no sampling loop. This removes
components of the PPO training stack; no fixed memory or wall-clock
saving is asserted here, because it depends on the implementation and
workload. Four limitations are documented in the same source: the
training data are static and offline, so the implicit reward goes stale
as the policy drifts; there is no online exploration beyond the
dataset's support; the loss can lower the absolute likelihood of
preferred completions while widening their margin over dispreferred ones
(likelihood displacement); and the method has weaker credit assignment
for long-horizon reasoning, where correctness of later steps depends on
earlier ones.

The contrast is structurally instructive for this project. DPO cannot
suffer zero-variance collapse, because it never standardises within a
sampled group --- its contrast is supplied by the dataset rather than
constructed at training time. But it equally cannot exploit a verifier
that is available in the loop, which is the property that makes RLVR
attractive on mathematics. GRPO's zero-variance pathology is therefore
the price of its on-policy, verifier-in-the-loop design, and measuring
that price is exactly what makes the pathology worth quantifying. The
DPO family has since diversified into inverse-preference and non-paired
variants --- IPO \citep{azar2023ipo}, KTO \citep{ethayarajh2024kto},
ORPO \citep{hong2024orpo}, SimPO \citep{meng2024simpo}, and SLiC
\citep{zhao2023slichf} among them --- which the project's cross-library
protocol treats as no-rollout baselines.

The boundary between the two families is also less sharp than it first
appears. Wu et al. \citep{wu2025ittakestwo} argue that GRPO is
``secretly DPO'' --- that its efficacy stems from an implicit
contrastive objective, which makes it structurally related to DPO ---
and report that 2-GRPO retains 97.6\% of the performance of 16-GRPO at
12.5\% of the rollouts (arXiv v3; the earlier v2 reported 98.1\%). That
reduction is the theoretical counterpart of the PCD statistic this
project develops: if a group is exactly as informative as the
correct/incorrect pairs it contains, then the pairwise-contrast density
PCD = ((G − 1)/G) · E{[}p(1 − p){]} is the natural magnitude-aware
summary of a group's usable signal, and the group-size question becomes
a question about how much contrast a given sampling budget buys. The
present project is careful about what it inherits from that argument:
the reduction is cited as prior theory, while the empirical claims
resting on it are scoped to this project's own measurements, reported in
Chapter 6.

\section{Fixing GRPO's Biases and
Collapse}\label{fixing-grpos-biases-and-collapse}

A dense and fast-growing literature targets specific GRPO pathologies.
Dr.~GRPO \citep{liu2025understanding} removes two normalisation terms
from the standard objective --- the per-response 1/L factor and the
question-difficulty amplification --- restoring the sequence-level
gradient up to a global scale, and is motivated as a fix for an
optimisation bias in GRPO that artificially inflates response length,
especially of incorrect responses. DAPO \citep{yu2025dapo}, from
ByteDance, changes four things in one release: asymmetric Clip-Higher
clipping, dynamic sampling that resamples until the batch carries
signal, token-level loss normalisation, and over-long filtering; the
reported result is 50 points on AIME 2024 with a 32B model. GSPO
\citep{zheng2025gspo} (Qwen team, 2025) moves the importance ratio from
the token level to the sequence level and was credited with stabilising
the Qwen3 series, particularly its mixture-of-experts variants.
Production-side variants include VAPO \citep{yue2025vapo}, a value-based
PPO framework for long chain-of-thought reasoning.

A third thread attacks the zero-variance failure mode directly, and this
is the thread the project's own diagnostic sits in. GRESO
\citep{zheng2025greso} pre-screens prompts before sampling so that
uninformative prompts are never paid for. CPPO \citep{lin2025cppo}
prunes completions after rollout, reporting up to 7.98× speedup on
GSM8K. NGRPO \citep{nan2025ngrpo} adds advantage calibration and
asymmetric clipping so that all-incorrect groups still carry gradient.
Scaf-GRPO \citep{zhang2025scafgrpo} injects tiered scaffolding hints
when learning stagnates, reporting a 44.3\% relative improvement on AIME
2024. MO-GRPO \citep{ichihara2025mogrpo} addresses vanilla GRPO's
collapse to a single reward component under multiple objectives through
variance-based reward re-weighting. SSPO \citep{yang2025sspo} operates
at subsentence granularity to balance GSPO's stability against
token-level GRPO's expressiveness, raising GRPO's five-benchmark
mathematics average from 43.01 to 46.72 on Qwen2.5-Math-1.5B. RL-ZVP
\citep{le2025rlzvp} takes the opposite approach to filtering: it
extracts signal \emph{from} zero-variance prompts through entropy-guided
advantage shaping. AERO \citep{zhang2026aero} instead adapts the rollout
allocation and prunes rollouts to avoid zero-advantage dead zones. GVPO
\citep{zhang2025gvpo} re-weights by within-group variance. The tail
continues with λ-GRPO \citep{wang2025lambdagrpo}, MC-GRPO
\citep{kim2026mcgrpo} (median-centred advantage), G²RPO-A
\citep{guo2025g2rpoa}, ConSPO \citep{zhang2026conspo}, and Goldilocks RL
\citep{mahrooghi2026goldilocks}.

One act of bookkeeping in this literature deserves emphasis, because it
is a claim about what may \emph{not} be inferred. AVSPO
\citep{he2026avspo} and other interventions are described in the
project's own related-work source as prior or concurrent interventions,
not as evidence for this project's runner, and the internal report
explicitly claims neither priority nor an intervention effect over them.
The present thesis inherits that disclaimer without qualification.

Against this background, the project's own contribution to the area is
deliberately meagre in algorithmic ambition and deliberately demanding
in evidential standard. The Zero-Variance Fraction (ZVF) --- the
fraction of groups whose completions are all correct or all incorrect
--- was measured across libraries, model scales, and training phases in
the companion ZVF report (P2), and that report stopped at description:
ZVF tracks collapse, drifts upward, and is mechanically entangled with
accuracy and group size, so it declined to promote ZVF to a predictive
or causal statistic. A minimal binomial model of group degeneracy was
then proposed, in which the expected ZVF indicator is
\(h_G(p) = p^G + (1-p)^G\) and the usable signal is
\(S = p(1-p)(1-h_G(p))\); three regularities follow in closed form, each
paired with an empirical check drawn from a 368-run Weights \& Biases
audit that counts all logged runs and is therefore a superset of the
benchmark's 70+ core training runs. Because \(h_G(p)\) is symmetric
under \(p \leftrightarrow 1-p\), ZVF assigns the same value to a
mastered prompt and to a hopeless one, producing the U-shape measured
across six model labels in a 368-run audit; and exact-tie ZVF is
sensitive to reward rounding: jitter of order 10⁻⁴ eliminates exact ties
while shifting measured within-group contrast by less than 10⁻⁶. This
does not hold for the declared sample-variance threshold of 10⁻⁶, which
retains 95/600 groups before and after jitter (§6.3.1).
Pairwise-contrast density is proposed as a magnitude-aware candidate
control signal, without this test establishing its superiority over
tolerance-based ZVF, and the toy-scale validation reports a
gradient-norm correlation of corr = +0.71, marked directional only. The
word \emph{controller} in that work names the implemented policy, not an
established improvement: the decisive static-G = 16 matched-token
bakeoff remains unrun.

The one genuinely measured GRPO-variant comparison in the repository is
a Dr.~GRPO versus GRPO control. On Qwen2.5-0.5B two-digit arithmetic at
group size 8, 40 steps, and 5 seeds, both arms were at ceiling and
indistinguishable: GRPO 0.987 ± 0.003 against Dr.~GRPO 0.992 ± 0.003
(mean ± SE over 5 seeds; SDs 0.006 and 0.008), paired difference +0.5
percentage points, t = 1.05, df = 4, p = 0.35; mean completion length
was 4.74 against 4.71 tokens. At that ceiling, and with completions
under five tokens, the task cannot show either an accuracy or a length
difference. In a second regime --- Qwen2.5-1.5B-Instruct on GSM8K with a
200-token chain-of-thought cap and 3 seeds --- GRPO moved from 20.2\% to
26.3\% and Dr.~GRPO from 20.5\% to 25.5\%; per-seed McNemar tests were
significant in 2 of 3 seeds for GRPO (p = 0.071, 0.036, 0.016) and 1 of
3 for Dr.~GRPO (p = 0.122, 0.017, 0.122), and the pooled p \textless{}
10⁻³ is expressly rejected as overstating significance; with the seed as
the unit, one-sample t-tests on the three per-seed gains give p = 0.005
(GRPO) and p = 0.010 (Dr.~GRPO) (recomputed from
drgrpo\_\allowbreak{}gsm8k\_\allowbreak{}cot.json). Neither arm inflated length over 30 steps
(GRPO 195 → 180 tokens; Dr.~GRPO 194 → 187). The same source contains
head-to-head and rank-correlation tables whose values are
\emph{projections}, not measurements: the GRPO baseline row is a
synthetic baseline produced by a dry run of
\longtok{platform\_hybrid/experiments/variance\_\allowbreak{}mitigation\_\allowbreak{}integration.py},
the file states that it makes no measurement claim for any row of that
table, the correlation coefficients (0.78 / 0.66 / 0.63 / 0.31) carry
the same caveat, and the configuration proxies under
\texttt{platform\_local/unified/} are explicitly not faithful
reimplementations. Those numbers are reported here only to mark them as
unusable as evidence. The nearest thing to a headline cross-library
effect in the project is also the one that must be read most carefully:
holding the visible group size, learning rate, data, seed, and label
fixed while changing stacks moved final GSM8K reward from 85.6\% to
5.0\%, a 17× swing; the audit that followed found the managed stack had
also pinned a different base checkpoint (Qwen3-8B-Base on Tinker against
the instruction-tuned Qwen3-8B on TRL), so the source describes this as
an under-specification exhibit rather than a backend-only causal effect.
The complementary telemetry fact is that running ``DAPO'' on an open
Colab trainer whose dynamic sampling demonstrably filters zero-variance
groups produced mean ZVF 0.00, while running ``DAPO'' on a closed
managed stack where only an asymmetric-clip surrogate is user-settable
produced mean ZVF 0.58 --- the same label, opposite telemetry.

\section{Verifiable Rewards, Frameworks and Reporting
Practice}\label{verifiable-rewards-frameworks-and-reporting-practice}

RLVR --- reinforcement learning from verifiable rewards --- was
formalised for open post-training in TÜLU 3 \citep{lambert2024tulu3},
which is the source this project follows in treating a programmatic
checker rather than a learned preference model as the reward. The
implementation landscape the project unifies is crowded and
heterogeneous. On the training side it includes TRL, OpenRLHF
\citep{hu2024openrlhf}, veRL (HybridFlow) \citep{sheng2025hybridflow},
NeMo-RL, DeepSpeed-Chat \citep{yao2023deepspeedchat}, trlX, Axolotl, and
Tinker as a managed back-end; each makes different default choices about
advantage normalisation, KL estimators, tokenisation loss masking,
rollout batching, and importance-sampling granularity, and the project
identifies those differences as the principal source of the
cross-library gap it quantifies. On the inference side, vLLM
\citep{kwon2023vllm} and SGLang \citep{zheng2024sglang} supply the
rollout backends on which modern frameworks depend, with FlashAttention
\citep{dao2022flashattention} and LoRA \citep{hu2021lora} as
load-bearing kernels. Tinker is adopted here through its
\longtok{message\_\allowbreak{}with\_\allowbreak{}tool} interface as the neutral substrate for
tool-use extensions, and the same source explicitly cites tool-use and
agentic RL settings as concurrent environments in which this project's
ZVF measurements will need to be re-evaluated rather than assumed to
transfer.

The evaluation benchmarks are equally well established and equally
awkward. GSM8K \citep{cobbe2021gsm8k} comprises 8,500 problems split
7,473 / 1,319, on which a fine-tuned GPT-3 verifier reached 35\% at
release while current systems exceed 95\%; GSM8K-Platinum
\citep{vendrow2025platinum} has since documented label noise in the
original set, which is why the source describes the benchmark as
necessary but not sufficient. MATH \citep{hendrycks2021math} comprises
12,500 problems across seven subjects and five difficulty levels, on
which the original results reached 6.9\%; MATH-500 is the 500-problem
subset introduced with the Lightman process-supervision work
\citep{lightman2023lets}, on which current systems sit at roughly 99\%.
AIME comprises 15 problems per year with integer answers in 0--999;
reported AIME 2024 accuracies include GPT-5 at 95.7\% and Grok 4 at
94.3\%. Because these benchmarks are close to saturation, the source
frames chain-of-thought as the action space that makes them interesting
at all, and identifies process reward models as the bridge from terminal
outcome signals to step-level credit.

Reporting practice is where the project's own contribution is anchored,
and its lineage is explicit. Model Cards \citep{mitchell2019modelcards}
introduced documentation of intended use, out-of-scope use, and
per-slice evaluation results; Datasheets for Datasets
\citep{gebru2021datasheets} extended the idea to datasets, covering
composition, preprocessing, annotation process, and
discriminatory-impact mitigation; Data Statements
\citep{bender2018datastatements} added the speaker-demographic axis; and
Data Cards \citep{pushkarna2022datacards} bridged research and industry
practice with schema, provenance, sampling-distribution, and
audit-transparency fields. A cross-coupling audit at iteration 109 maps
these four works onto the project's legacy 18-field experimental
manifest, showing that they cover 12 distinct fields, that the remaining
six are either RL-specific (per-step ZVF, group-size schedule) or
stack-coverage items with no direct analogue in the prior literature,
and that items 9--12 --- annotation process and annotator demographics
--- have the lowest evidence density in the live manifests. The same
iteration verified all four entries against CrossRef metadata (4/4 pass
on year, five-gram title, and author-family overlap ≥ 0.6) and
retrofitted two entries that had been present in the bibliography but
never cited. The resulting standard, MIN-REPORT-RL, has eight items:
seven run-manifest fields and one evaluation item (held-out pass@k),
which belongs to the evaluation report rather than to run-start stack
identity. The count was aligned at eight on 2026-07-11, and this thesis
uses ``the eight-item standard'' and ``the seven manifest fields'' in
that sense throughout.

The reproducibility literature to which this standard is a deliberate
echo is older and better established. Henderson et al.
\citep{henderson2018deep} showed that deep-RL results depend on
codebase, hyperparameters, and even the choice of random seeds, to the
point that the same algorithm from two repositories behaves as two
algorithms. Colas et al. \citep{colas2019hitchhiker} and Agarwal et al.
\citep{agarwal2021deep} (\texttt{rliable}) supplied the statistical
repair kit --- rigorous comparison protocols and small-sample-aware
aggregate metrics with interquartile-mean reporting and performance
profiles. Jordan et al. \citep{jordan2024benchmarking} argued that
rigorous benchmarking in RL is computationally prohibitive at the scale
it would require, and called for an additional experimentation paradigm
to complement it. Dodge et al. \citep{dodge2019show} proposed reporting
expected validation performance as a function of compute budget; the
project's run manifest is the RL-for-LLM analogue, replacing
\emph{budget} with \emph{stack}. Pineau et al.
\citep{pineau2021reproducibility} formalised the NeurIPS reproducibility
criteria, and the project draws the explicit lesson that self-attested
checklists have limits, which motivates an artifact-computed badge
rather than a declared one. Statistical treatments of evaluation noise
(Miller \citep{miller2024errorbars}) quantify the uncertainty floor
against which the project's own within-noise findings should be read,
and the evaluation-stack sensitivity documented by the LM Evaluation
Harness effort \citep{biderman2024lmeval}, by audits of reasoning
evaluations \citep{hochlehnert2025sober}, and by contamination studies
of the data side of the same coin, together establish that RL
\emph{training} inherits all of this --- the verifier is an evaluator in
the loop --- plus the additional vectors unique to on-policy
optimisation.

Three existing resources come closest to the project's registry and are
the reason it exists in its present form. HELM \citep{liang2022helm}
aggregates evaluation across a large model and scenario matrix, but its
unit is the (model, scenario) cell, not the training stack. The LM
Evaluation Harness \citep{biderman2024lmeval} standardises evaluation by
\emph{executing} implementations, whereas the registry here standardises
by \emph{describing} stacks it cannot execute --- many of them closed
--- which is precisely why the project required a null-versus-absent
convention and an R5 rung that records un-auditability rather than
converting it into a false value. Papers with Code catalogued
paper--code--leaderboard linkage at scale, stopped at the repository
boundary, and was discontinued in July 2025; the source treats its fate
as a governance cautionary tale. The project's own claim to novelty is
narrow and stated negatively: no existing resource provides the executed
loss form at GRPO-variant granularity, per-component implementation
status, or an unauditability marker.

\section{Scaling and Layer-Adaptation
Context}\label{scaling-and-layer-adaptation-context}

Compute-optimal pretraining is governed by well-characterised scaling
laws (Kaplan et al. \citep{kaplan2020scaling}; Hoffmann et al.,
Chinchilla \citep{hoffmann2022chinchilla}), but RL post-training is much
less well characterised, and the sources the project builds on separate
cleanly into those that predict and those that describe. Hilton et al.
\citep{hilton2023rlscaling} derived compute-efficiency frontiers for
single-agent RL across Procgen games, Dota 2 and a toy MNIST-based task.
For RLHF, Gao et al. \citep{gao2023scaling} characterised the
over-optimisation curve of reward models --- the point at which
additional optimisation against a frozen reward model buys a decreasing
true objective --- and Rafailov et al. \citep{rafailov2024rlhfscaling}
extended over-optimisation scaling laws to direct alignment algorithms
(DPO, IPO and SLiC), which have no separate proxy reward model.
Nimmaturi et al. \citep{nimmaturi2025predictive} derived an empirical
scaling law specifically for GRPO, identifying three training phases
(slow start, rapid improvement, plateau) and reporting that most of the
benefit is exhausted by roughly 80\% of a single epoch. Tan et al.
\citep{tan2025scalingrl} treat group size as a budget knob in an
empirical scaling study of GRPO post-training, in an appendix analysis
showing that the optimal rollout size shifts with the training budget.
Snell et al. \citep{snell2024testtime} treat inference-time compute as a
scaling axis in its own right, and the project's own related-work source
cites these as the works against which its measurements should be read.

The project's scaling study covers 70+ runs across seven RL libraries
and five model families. This thesis uses one model range throughout:
0.6B--1T, with the anchors above 235B single-seed and descriptive. The
rule is that the range spans every model with at least one completed run
in the P1 roster. The source documents state narrower ranges because
each covers a subset: 0.6B--\textasciitilde671B in the P2, P3 and P4
introductions, and 0.6B--235B for the multi-family data in the
related-work source. The principal empirical result of the scaling work
is a null: across that range, no reliable scaling law is recovered from
reward curves. The companion group-size study adds a second bounded
negative: a three-seed sweep does not establish an optimum, since the
measured arms are statistically indistinguishable, and the measured
signal-to-noise ratio grows at only about 52\% of the √G ideal, with G ∈
\{2, 4, 8, 16\} measured and G ∈ \{4, 8, 16, 32, 64\} explicitly
reconstructed rather than run.

Layer adaptation supplies the second half of this section's context, and
connects the project's white-box investigation to the
parameter-efficient fine-tuning literature. LoRA \citep{hu2021lora}
freezes the pretrained weights and injects trainable low-rank matrices
into the attention projections, so that a d × d projection at d = 4096
with rank r = 16 trains 2dr rather than d² parameters, a 128× reduction
(this chapter's arithmetic; the paper itself reports up to 10,000× fewer
trainable parameters for GPT-3 175B). QLoRA \citep{dettmers2023qlora}
extends this by quantising the frozen base to NF4 and adding double
quantisation, which saves about 0.37 bits per parameter ---
approximately 3 GB for a 65B model --- and permits 65B fine-tuning on a
single 48 GB GPU. The motivating arithmetic is that Adam at mixed
precision costs roughly 16N bytes, exceeding 100 GB for a 7B model.
Prefix tuning, prompt tuning, and adapter layers are the earlier members
of the same family, and LoRA's advantages for RL specifically are that
the frozen base doubles as the reference policy and that the trainable
parameter count is small enough for modest hardware. The project's
layer-adaptation pillar asks a narrower and orthogonal question: where
inside the network GRPO-style adaptation concentrates, and whether the
layers that will matter can be identified early enough to freeze the
rest. The reported finding is moderate concentration of per-layer
adaptation, with the important negative that step-1 predictability does
\emph{not} survive scaling. This connects directly to Mukherjee et al.
\citep{mukherjee2025rlsubnetworks}, who found that RL fine-tuning
updates only 5--30\% of model parameters across seven algorithms and ten
model families, suggesting that RL activates a sparse subnetwork. The
length-bias pillar completes the scaling picture with its own bounded
null: under a 200-token GSM8K generation cap, GRPO and Dr.~GRPO produce
held-out gains that are indistinguishable at the project's seed budget,
and neither arm shows length inflation; the source is careful to state
that the 200-token cap is part of the estimand rather than a nuisance
detail, so the null bounds the capped, short-horizon regime and cannot
adjudicate uncapped reasoning or long agent trajectories.

\section{Process-, Tree- and Stability-Aware
Variants}\label{process--tree--and-stability-aware-variants}

Beyond the group-baseline fixes, a parallel line of work changes
\emph{what} signal drives the update. Process reward models supply
denser, step-level feedback instead of a single terminal reward, which
reduces the chance that a whole group collapses to identical terminal
rewards. Let's Verify Step by Step \citep{lightman2023lets} showed that
step-level PRMs substantially outperform outcome reward models on MATH;
Math-Shepherd \citep{wang2024mathshepherd} derives process rewards from
automatically generated process annotations, avoiding costly human
labels; PRM800K \citep{lightman2023lets} and PRMBench
\citep{song2025prmbench} supply step-level evaluation data; OmegaPRM
\citep{luo2024omegaprm} uses Monte-Carlo tree search to supervise
process rewards without human-in-the-loop; ReasonEval
\citep{xia2025reasoneval} scores reasoning steps for validity and
redundancy with trained LLM evaluators; and Implicit PRM
\citep{yuan2024implicitprm} learns process rewards implicitly from
outcome labels, with VersaPRM \citep{zeng2025versaprm} extending process
rewards to multi-domain tasks. Tree-of-thought search
\citep{yao2023tot}, self-consistency \citep{wang2023selfconsistency},
and critic-style post-hoc verifiers all either benefit from or
substitute for process supervision. The same source documents the
counter-current that matters most here: ``A Sober Look''
\citep{hochlehnert2025sober} warns that reward hacking and evaluation
artefacts can inflate perceived gains on reasoning benchmarks, which is
the project's stated motivation for insisting on binary verifiable
rewards in its own mathematics work.

Stability-aware variants constrain the update itself. GSPO's
sequence-level importance ratio, discussed above, is the clearest
instance, credited with stabilising the Qwen3 mixture-of-experts series.
SSPO's subsentence-granularity objective is explicitly framed as
balancing GSPO's stability against token-level GRPO's expressiveness.
Entropy-guided advantage shaping in RL-ZVP regulates the update when the
group carries no reward contrast, and GRESO's pre-screening is a
sampling-side stability control. On the reasoning-RL side, GRPO-Lead
\citep{zhang2025grpolead} and LCPO \citep{aggarwal2025l1} tune reasoning
length through length-aware rewards, LIMR \citep{li2025limr} selects a
small high-impact subset of RL training data, Search-R1
\citep{jin2025searchr1} trains reasoning agents to interleave search
with chain-of-thought, and TTRL \citep{zuo2025ttrl} performs test-time
RL on unlabelled data using majority-vote rewards.

The Phase-1 survey's judgement on this body of work is retained without
amendment, because the Semester 4 evidence has not changed it. These
methods propose \emph{levers}; this project first builds the
\emph{instrumentation} needed to tell whether any such lever converts
wasted rollouts into held-out accuracy --- ZVF, group utilisation,
pairwise-contrast density, and the collapse decomposition. The nulls
reported in the results chapters are exactly the kind of evidence that
instrumentation is designed to produce, and they constitute a standing
caution against adopting a variant on the strength of a single-seed
improvement. The project applies that caution to its own proposals: the
adaptive group-size controller is described as an implemented policy
rather than an established improvement, and its decisive matched-token
comparison remains unrun.

\section{Positioning the Portfolio}\label{positioning-the-portfolio}

The portfolio's positioning against prior art is a per-pillar statement
of what is claimed and, more importantly, what is not. Table 2.1
reproduces the Phase-1 positioning and updates it with the Semester 4
evidence (Figure~\ref{fig:questions}).


\fig{fig_research_questions}
{The eight studies P1--P8 grouped by theme, with the question each one addresses and the state of its evidence}
{The eight studies P1--P8 grouped by theme, with the question each one addresses and the state of its evidence. Every question is answered from the repository's own record; two of the eight answer with a declared null, and a study whose scope narrowed between the group semester and this one is marked partial rather than complete.}
{fig:questions}



\textbf{Positioning of the P1--P8 studies against closest prior art.}

\begin{longtable}[]{@{}
  >{\raggedright\arraybackslash\hspace{0pt}}p{(\linewidth - 6\tabcolsep) * \real{0.3333}}
  >{\raggedright\arraybackslash\hspace{0pt}}p{(\linewidth - 6\tabcolsep) * \real{0.3333}}
  >{\raggedright\arraybackslash\hspace{0pt}}p{(\linewidth - 6\tabcolsep) * \real{0.3333}}@{}}
\caption{Positioning of the eight portfolio studies (P1--P8) against
closest prior art and the narrowed axis each claims.}\tabularnewline
\toprule\noalign{}
\begin{minipage}[b]{\linewidth}\raggedright
Pillar
\end{minipage} & \begin{minipage}[b]{\linewidth}\raggedright
Closest prior art
\end{minipage} & \begin{minipage}[b]{\linewidth}\raggedright
Narrowed axis claimed
\end{minipage} \\
\midrule\noalign{}
\endfirsthead
\toprule\noalign{}
\begin{minipage}[b]{\linewidth}\raggedright
Pillar
\end{minipage} & \begin{minipage}[b]{\linewidth}\raggedright
Closest prior art
\end{minipage} & \begin{minipage}[b]{\linewidth}\raggedright
Narrowed axis claimed
\end{minipage} \\
\midrule\noalign{}
\endhead
\bottomrule\noalign{}
\endlastfoot
P1 (scaling) & whole-model scaling and saturation laws; Mukherjee et al.
\citep{mukherjee2025rlsubnetworks} on sparse RL subnetworks & Per-layer
adaptation is moderately concentrated; step-1 predictability does not
survive scaling; no reliable reward-curve scaling law recovered \\
P2 (ZVF) & DAPO dynamic sampling, GVPO, AERO, RL-ZVP, GRESO, CPPO, NGRPO
& \emph{Measure} ZVF as an existence statistic; show it is entangled
with accuracy and group size; propose PCD as the magnitude-aware
replacement \\
P3 (group size) & integer-G sweeps; MC-GRPO; Tan et al.
\citep{tan2025scalingrl} budget-knob treatment & Trainability is
non-monotone in G; a three-seed sweep finds its arms statistically
indistinguishable and does not establish an optimum; measured SNR grows
at \textasciitilde52\% of the √G ideal \\
P4 (length bias) & Dr.~GRPO and λ-GRPO length fixes & A bounded null:
under a 200-token cap, GRPO and Dr.~GRPO held-out gains are
indistinguishable with no length inflation \\
P5 (reporting standard) & Datasheets, Model Cards, Data Statements, Data
Cards & An eight-item RL-run reporting standard (seven manifest fields
and one evaluation item) targeting the stack rather than the label \\
P6 (registry) & static run-config registries; HELM; LM Evaluation
Harness; Papers with Code & Machine-readable stack records with
per-component implementation status and an unauditability marker \\
P7 (controller) & adaptive training-signal control (G²RPO-A adaptive
guidance \citep{guo2025g2rpoa}); DAPO dynamic sampling & Closed-form
\(h_G(p)\) prediction of ZVF, with a controller described as
implemented, not established \\
P8 (fraud study) & tuned gradient-boosted trees as the tabular
fraud-detection baseline & A measurement-discipline side study comparing
a fine-tuned LLM with a gradient-boosted tree on synthetic credit-card
fraud, with no RL claim (Appendix F); P8 was archived and absorbed into
the portfolio \\
\end{longtable}

The two crowded pillars are P2 and P4, where the literature is dense and
the project accordingly narrows its claim to measurement rather than to
a new algorithm. The widest open ground remains the systems and
reporting side, where no GRPO-specific machine-readable reporting
standard exists --- the gap Section 2.4 establishes from three
directions: HELM's unit is the (model, scenario) cell rather than the
training stack, the LM Evaluation Harness standardises evaluation by
executing code while the registry must describe stacks it cannot
execute, and Papers with Code stopped at the repository boundary and was
discontinued in July 2025.

The Semester 4 evaluation campaign extends this positioning from method
comparison to benchmark practice. Rather than adding one more
leaderboard number, it evaluates a single trained actor against fourteen
agentic and reasoning suites outside its training data, each through its
own native grader at a pinned revision, under receipt, budget and
decontamination controls, and it reports every lane's terminal state
rather than a cross-suite aggregate. Its results are in Chapter 8; Table
C.3 gives each lane's scope and status. One finding bears directly on
the benchmark-practice literature: E4's mean reward of 0.0 across 100
base-model trials was caused not by the model but by a serving bridge
that passed no stop sequences (Section 8.5), which is exactly the kind
of unreported stack detail this positioning argues must be recorded.

Two gaps in this survey should be stated plainly, since the chapter's
purpose is to locate the work rather than to flatter it. First, the
primary papers for the fourteen evaluation suites are not surveyed here
at the depth the algorithmic literature receives; the suites are
discussed as instruments, sourced to the campaign's own results
document, and their native protocols are documented in Chapter 8 and its
receipt trail rather than in this chapter. Second, the survey's
treatment of tool-use and agentic RL is inherited from the project's own
related-work source and is explicitly flagged there as a concurrent
setting in which the ZVF measurements will need to be re-evaluated
rather than assumed to transfer. What the chapter does establish is that
the project's contribution is a measurement and reporting contribution,
that the literature it sits in is dense enough that claim-narrowing
rather than claim-extension is the appropriate move, and that the
specific gap the portfolio occupies --- machine-readable, per-component,
executed-loss-form disclosure for group-relative RL runs --- is not
filled by any of the resources surveyed above.



\chapter{System Requirements
Specification}\label{system-requirements-specification}

This chapter states what the Tinker RL Lab system must do, what
qualities it must hold while doing it, and how each requirement is
verified. It inherits the specification of the Phase-1 report and
extends it across the Semester-4 additions: the multi-framework
benchmark rosters, the Zero-Variance Fraction (ZVF) diagnostic and the
adaptive group-size controller derived from it, the minimum reporting
standard and the machine-readable GRPO stack registry, and the E1--E14
evaluation campaign with its fail-closed completion gate. It is written
as a specification rather than as a retrospective description, but every
requirement below corresponds to a component that exists in the
repository, and every functional requirement is paired with the concrete
test that checks it. Where a requirement is not satisfied --- the cloud
capacity for two evaluation lanes, the white-box gradient path at the
largest model scale --- the shortfall is named in §3.6 and §3.8 rather
than written around.

\section{Overview and Actors}\label{overview-and-actors}

The central object of the Phase-1 specification was the \textbf{run}:
one configuration, one seed, one framework, one training budget,
producing per-step telemetry and a persisted set of raw group reward
tensors. Semester 4 introduces two further first-class objects. The
first is the \textbf{lane}: a single benchmark suite evaluated under a
named scope --- an original contract with the suite's owner, or a
declared replacement scope --- carrying its own provider, payload,
licence, quota and terminal state. The second is the \textbf{receipt}:
an immutable, hash-bound record that a declared step actually happened,
against declared inputs, at a declared revision. Requirements that
concern runs, lanes and receipts are distinguished below where the
distinction matters.

Four actors interact with the system, and the fourth is new in Semester
4.

\textbf{The experimenter} declares work. Under the original design this
meant declaring a run through a single configuration object and nothing
else: the model, the dataset, the reward function, the decoding
parameters (temperature, top-\(p\), maximum tokens), the group size
\(G\), the learning rate, the KL coefficient \(\beta\) and the seed all
live in one file, and the framework is a switch rather than a rewrite.
Under the campaign design the same actor declares a lane through a
sealed resource request: the lane specification fixes an immutable task
bundle, a pinned native runtime, a native grader, a budget and an
applicable licence in advance, and those declared inputs are then
checked for internal consistency before any launch is permitted.

\textbf{The back-end} executes the training loop and, under the managed
path, is the Tinker service; under the white-box path it is a local or
rented GPU. For evaluation, the corresponding actor is the
\textbf{provider grader}: the benchmark owner's own native evaluator at
a pinned revision. The system is permitted to record a suite score only
when that evaluator produced it, or when a declared replacement harness
reproduces it under a scope that is labelled as a replacement.

\textbf{The auditor} is a downstream script or an independent reader who
recomputes published diagnostics from stored raw artefacts. In Semester
4 this role acquires force: the audit suite of §3.4 is executed as a
gate with a non-zero exit status rather than as advisory tooling, and
the completion gate for the campaign lanes is fail-closed by
construction, in that a missing or malformed declaration stops the
procedure instead of degrading it.

\textbf{The campaign operator} holds launch authority for paid work.
Spend on any lane requires a bound authorization receipt issued by the
operator; the system will not mint one for itself, and lanes are
recorded as blocked when the receipt or the technical chain behind it is
absent.

A fifth party, the \textbf{external access holder}, is an actor only in
the negative sense: several lanes depend on a private payload, a hosted
evaluation service or a quota grant that only the provider can issue.
The specification requires that such a lane be recorded with an explicit
terminal state and a named reopen condition rather than silently
dropped, which is why the campaign's terminal-state vocabulary exists as
a first-class output.

\section{Functional Requirements}\label{functional-requirements}

Table 3.1 lists the functional requirements. Each row states the
requirement and names the source that defines it; §3.5 maps each to the
artefact that implements it and the test that verifies it.

\textbf{Functional requirements.}

\begin{longtable}[]{@{}
  >{\raggedright\arraybackslash\hspace{0pt}}p{(\linewidth - 6\tabcolsep) * \real{0.3333}}
  >{\raggedright\arraybackslash\hspace{0pt}}p{(\linewidth - 6\tabcolsep) * \real{0.3333}}
  >{\raggedright\arraybackslash\hspace{0pt}}p{(\linewidth - 6\tabcolsep) * \real{0.3333}}@{}}
\caption{Functional requirements. Sources: the Phase-1 report is
\protect\longtok{platform\_tinker/reports/esa\_phase1/Phase1\_\allowbreak{}Project\_\allowbreak{}Report\_\allowbreak{}ZVF.tex};
other files are under \protect\longtok{platform\_hybrid/},
\protect\longtok{platform\_local/} or
\protect\longtok{zvf-program/flagship/}.}\tabularnewline
\toprule\noalign{}
\begin{minipage}[b]{\linewidth}\raggedright
ID
\end{minipage} & \begin{minipage}[b]{\linewidth}\raggedright
Requirement
\end{minipage} & \begin{minipage}[b]{\linewidth}\raggedright
Source
\end{minipage} \\
\midrule\noalign{}
\endfirsthead
\toprule\noalign{}
\begin{minipage}[b]{\linewidth}\raggedright
ID
\end{minipage} & \begin{minipage}[b]{\linewidth}\raggedright
Requirement
\end{minipage} & \begin{minipage}[b]{\linewidth}\raggedright
Source
\end{minipage} \\
\midrule\noalign{}
\endhead
\bottomrule\noalign{}
\endlastfoot
FR-1 & One declarative file fully determines a run (model, dataset,
reward, decoding, G, learning rate, KL coefficient, seed, output); no
framework-specific parameters are needed on any supported back-end &
Phase-1 report \\
FR-2 & A configuration is executable on multiple frameworks through thin
adapters that share the reward grader and decoding configuration. Two
rosters are in scope and are kept apart: the cross-RL-library roster
(TRL, Stable-Baselines3, CleanRL, Tianshou, PufferLib, rl\_games,
d3rlpy) and the cross-launcher LLM-RL roster (Tinker, TRL, SkyRL, veRL,
OpenRLHF, Atropos, HF reference launchers). Completed-run parity is
\textbf{not} met: only Tinker and TRL produced completed runs; the veRL
and OpenRLHF entries are dry-run placeholders &
\longtok{\_shared\_methods.tex}; \longtok{framework\_comparison.json} \\
FR-3 & Every step computes ZVF and its all-correct / all-wrong
decomposition from the reward tensors & Phase-1 report \\
FR-4 & Every step emits a structured telemetry record (reward, ZVF,
gradient utilisation, collapse counts, entropy, completion length, KL)
to a log and to the tracker; the controller work adds a per-step
counterfactual against a matched fixed-G control & Phase-1 report;
\longtok{controller\_\allowbreak{}cf\_\allowbreak{}per\_\allowbreak{}step.tsv} \\
FR-5 & Raw per-group reward tensors are persisted for every run, so
diagnostics can be recomputed & Phase-1 report \\
FR-6 & Sweeps over G, seed and baseline-versus-intervention arms run
under a fixed wall-clock or token budget & Phase-1 report \\
FR-7 & For single-GPU models, per-layer LoRA gradient norms are
recorded; not required at managed scale, where they are not exposed &
Phase-1 report \\
FR-8 & Every run emits a provenance record binding configuration, grader
version and rollout hashes; a campaign lane additionally declares its
provider package and grant documents against published schema versions,
each with a 64-character hash & Phase-1 report;
\longtok{e1\_\allowbreak{}e14\_\allowbreak{}completion\_\allowbreak{}gate.py} \\
FR-9 & A stack-conditioned result carries the seven manifest fields of
the eight-item standard (loss form, reference KL, sampler/backend with
base-checkpoint revision and hash, telemetry, group-size schedule,
held-out split, decontamination) plus its eighth, evaluation item
(held-out pass@k); a 0--100 badge is emitted, and an unreported field is
distinguished from one reported as absent &
\longtok{paper\_\allowbreak{}P6\_\allowbreak{}registry.tex};
\longtok{minreport\_\allowbreak{}audit\_\allowbreak{}summary.json} \\
FR-10 & The stack registry is queryable by stack, field or status &
\longtok{registry/schema.json}; \longtok{registry/query.py} \\
FR-11 & Each named variant stores an explicit delta against the GRPO
reference, and \longtok{stackdiff} returns an R0--R5 flip-risk verdict
for same-label entries & \longtok{paper\_\allowbreak{}P6\_\allowbreak{}registry.tex};
\longtok{registry/provenance/} \\
FR-12 & An adaptive group-size controller with escalation asymmetry and
hysteresis records its counterfactual against the best static recipe at
matched and unequal rollout counts; it measures the trade and does not
claim a win & \longtok{paper\_\allowbreak{}P7\_\allowbreak{}zvf\_\allowbreak{}controller.tex};
\longtok{controller\_\allowbreak{}cf\_\allowbreak{}summary.json} \\
FR-13 & Each lane records results against a declared evidence class
(exact complete, partial exact, partial recovery, externally blocked,
local-setup-ready-provider-input-required) tied to the source status &
\longtok{e1\_\allowbreak{}e14\_\allowbreak{}completion\_\allowbreak{}gate.py} \\
FR-14 & The completion gate validates, and never executes, a lane
package; a valid package proves only that its inputs are consistent and
available; signatures are checked against provider trust roots, and the
gate refuses rather than warns &
\longtok{e1\_\allowbreak{}e14\_\allowbreak{}completion\_\allowbreak{}gate.py} \\
FR-15 & Original-contract and replacement-scope results are never
pooled, no cross-suite aggregate is computed, and every accuracy carries
its denominator & \longtok{E1\_\allowbreak{}E14\_\allowbreak{}FINAL\_\allowbreak{}RESULTS\_\allowbreak{}2026-09-19.md} \\
FR-16 & The export path builds an anonymised review package whose build
runs the audit suite, and refuses a package built with audits skipped &
\longtok{run\_\allowbreak{}all\_\allowbreak{}audits.py}; \longtok{export\_\allowbreak{}guard\_\allowbreak{}audit.py} \\
\end{longtable}

\section{Non-Functional Requirements}\label{non-functional-requirements}

\textbf{Non-functional requirements.}

\begin{longtable}[]{@{}
  >{\raggedright\arraybackslash\hspace{0pt}}p{(\linewidth - 6\tabcolsep) * \real{0.3333}}
  >{\raggedright\arraybackslash\hspace{0pt}}p{(\linewidth - 6\tabcolsep) * \real{0.3333}}
  >{\raggedright\arraybackslash\hspace{0pt}}p{(\linewidth - 6\tabcolsep) * \real{0.3333}}@{}}
\caption{Non-functional requirements.}\tabularnewline
\toprule\noalign{}
\begin{minipage}[b]{\linewidth}\raggedright
ID
\end{minipage} & \begin{minipage}[b]{\linewidth}\raggedright
Requirement
\end{minipage} & \begin{minipage}[b]{\linewidth}\raggedright
Source
\end{minipage} \\
\midrule\noalign{}
\endfirsthead
\toprule\noalign{}
\begin{minipage}[b]{\linewidth}\raggedright
ID
\end{minipage} & \begin{minipage}[b]{\linewidth}\raggedright
Requirement
\end{minipage} & \begin{minipage}[b]{\linewidth}\raggedright
Source
\end{minipage} \\
\midrule\noalign{}
\endhead
\bottomrule\noalign{}
\endlastfoot
NFR-1 Attribution & Comparative claims hold the stack fixed except the
factor under test; confounded factors are reported as confounded &
\longtok{\_shared\_methods.tex} \\
NFR-2 Reproducibility & Multi-arm claims rest on at least three seeds;
single-seed results are descriptive only & Phase-1 report \\
NFR-3 Recomputability & Every derived diagnostic is reproducible from
stored raw artefacts, independent of logged summaries & Phase-1
report \\
NFR-4 Denominator discipline & Held-out figures carry their n; negative,
null and underpowered results are reported as such &
\longtok{main\_\allowbreak{}eai\_\allowbreak{}body.tex} \\
NFR-5 Isolation & Concurrent experiments run in separate processes, not
threads (after the tracker thread-safety defect) & Phase-1 report \\
NFR-6 Budget & Spend is recorded per run against a declared envelope;
per-run bounds and lead-issued authorisation receipts apply even under
an unlimited total & \longtok{Pending\_Experiments.md} \\
NFR-7 Fail-closed & Validation defaults to refusal; an unverifiable
package is not scored & \longtok{e1\_\allowbreak{}e14\_\allowbreak{}completion\_\allowbreak{}gate.py} \\
NFR-8 Least privilege & Evidence records hold no key material;
credential checks record only liveness and scope &
\longtok{Pending\_Experiments.md} \\
NFR-9 Lost-source auditability & Execution sources for sealed artefacts
live in tracked paths; sealed requests and receipts suffice to pin a
faithful re-implementation &
\longtok{E1\_\allowbreak{}E14\_\allowbreak{}FINAL\_\allowbreak{}RESULTS\_\allowbreak{}2026-09-19.md} \\
NFR-10 External traceability & External payloads are bound by hash and
validated against trust roots and schema versions; an unreported field
is null, not zero or false & \longtok{e1\_\allowbreak{}e14\_\allowbreak{}completion\_\allowbreak{}gate.py};
\longtok{paper\_\allowbreak{}P6\_\allowbreak{}registry.tex} \\
\end{longtable}

\section{Test-Case Specification}\label{test-case-specification}

The test cases below are the concrete checks attached to the
requirements. TC-1 through TC-5 are inherited from the Phase-1
specification and are restated here because later chapters depend on
them; TC-6 onward are the Semester-4 additions, and TC-6 is the suite
that carries the largest share of the functional requirements.

\textbf{TC-1 --- Recompute ZVF and the collapse decomposition from
stored tensors.} Recompute the published ZVF and collapse figures from
the persisted per-group reward tensors of the main configuration, using
the documented analysis script. The published values are a ZVF of
approximately 0.72--0.77 and an all-correct fraction of 0.65--0.71;
recomputation must match. Verifies FR-3 and FR-5.

\textbf{TC-2 --- Recompute the detector metric from raw tensors.}
Recompute the Phase-3 detector's area under the ROC curve from raw
tensors under five-fold stratified cross-validation. The published
figure is an AUROC of \(0.84 \pm 0.01\) against \(0.43\) for the
reward-only comparator. Verifies FR-5 and NFR-3.

\textbf{TC-3 --- Re-run a matched-budget arm comparison.} Re-run the
baseline and the curriculum arm at a matched token budget over three
seeds; the Phase-1 result is that both arms average \(+0.028\) on the
held-out set with no curriculum advantage. Verifies FR-1, FR-2 and FR-6.

\textbf{TC-4 --- Recompute headline diagnostics.} Recompute each
headline diagnostic quoted in a results chapter directly from stored
artefacts, with the recomputation recorded next to the published value.
Verifies NFR-3.

\textbf{TC-5 --- Provenance verification.} Verify a run's provenance
record end to end, from configuration to rollout hashes to the grader
version. Verifies FR-8.

\textbf{TC-6 --- Execute the submission audit suite as a gate.} Run
\texttt{python\ platform\_\allowbreak{}local/\allowbreak{}run\_\allowbreak{}\allowbreak{}all\_\allowbreak{}\allowbreak{}audits.py}. The suite executes
nine audits --- \texttt{claim\_issues}, \texttt{sync\_issues},
\texttt{anon\_issues}, \texttt{strength\_issues},
\texttt{package\_issues}, \texttt{workflow\_issues},
\longtok{export\_\allowbreak{}guard\_\allowbreak{}issues}, \texttt{caveat\_issues} and
\texttt{scientific\_issues} --- emits one
\texttt{METRIC\ \textless{}name\textgreater{}=\allowbreak{}\textless{}count\textgreater{}}
line per audit, and returns exit status 0 only when every audit reports
zero issues. Individual checks that this test enforces include the
presence of scope and caveat language in the review documents
(\longtok{readme\_\allowbreak{}missing\_\allowbreak{}key\_\allowbreak{}results\_\allowbreak{}scope\_\allowbreak{}header},
\longtok{readme\_\allowbreak{}missing\_\allowbreak{}heldout\_\allowbreak{}language},
\longtok{checklist\_\allowbreak{}missing\_\allowbreak{}preliminary\_\allowbreak{}key\_\allowbreak{}results\_\allowbreak{}label}), the
absence of overstated checkpoint claims
(\longtok{misleading\_\allowbreak{}checkpoint\_\allowbreak{}availability\_\allowbreak{}claim},
\longtok{submission\_\allowbreak{}overstates\_\allowbreak{}checkpoint\_\allowbreak{}release}), the presence of
every reviewer caveat (\texttt{heldout\_scope},
\longtok{tool\_\allowbreak{}eval\_\allowbreak{}protocol}, \texttt{codegen\_subset},
\texttt{capacity\_confound}, \texttt{moe\_diagnostics},
\longtok{budget\_\allowbreak{}and\_\allowbreak{}splits}, \texttt{replication\_release}), the
absence of banned marketing formulations, the survival of the technical
caveats into the anonymous build
(\longtok{anonymous\_missing:\textless{}label\textgreater{}} for the
seven needles), and the presence of the audit call and its skip flag in
the export script. Verifies FR-16, NFR-4 and NFR-7.

\textbf{TC-7 --- Structural checks on the evaluation script and its
result files.} For each held-out evaluation script, assert by parsing
the source --- not by reading its documentation --- that it exposes a
seed argument, locks the dataset split to the test split, samples with a
non-zero temperature, consumes a checkpoint path, and writes the
declared metadata keys including the split, the sampling flag, the seed
and the model source. For each result file, assert that the schema
version is 2, that the evaluation status is one of completed or failed,
that the required configuration and summary fields are present, that the
recorded split is the test split, that the attempted count equals the
sum of correct and incorrect, and that a failed evaluation carries a
failure reason and a completed evaluation carries non-zero attempts.
Verifies FR-8 and NFR-4.

\textbf{TC-8 --- Build the report from source.} Compile the main report
and its supplementary through the full LaTeX chain and flag an empty
journal field in the bibliography. Verifies FR-16.

\textbf{TC-9 --- Dry-run the completion gate against a declared lane.}
Present a provider package to the gate and require it to validate the
declaration without executing anything: the expected outcomes are
acceptance with the declared evidence class, or a specific refusal
naming the missing element. A package must never be reported as a
benchmark result by this procedure. Verifies FR-13, FR-14 and NFR-7.

\textbf{TC-10 --- Registry query and null discipline.} Query the
registry for a labelled entry and its delta record, and confirm that a
field absent from the source is returned as unreported rather than as a
zero value, and that a same-label pair with divergent behaviour triggers
the flip-risk verdict. Verifies FR-9, FR-10, FR-11 and NFR-10.

\section{Requirements Traceability
Matrix}\label{requirements-traceability-matrix}

\textbf{Requirements traceability matrix.}

\begin{longtable}[]{@{}
  >{\raggedright\arraybackslash\hspace{0pt}}p{(\linewidth - 6\tabcolsep) * \real{0.3333}}
  >{\raggedright\arraybackslash\hspace{0pt}}p{(\linewidth - 6\tabcolsep) * \real{0.3333}}
  >{\raggedright\arraybackslash\hspace{0pt}}p{(\linewidth - 6\tabcolsep) * \real{0.3333}}@{}}
\caption{Requirements traceability matrix: each functional requirement
mapped to the artefacts that implement it and the audit check that
verifies it.}\tabularnewline
\toprule\noalign{}
\begin{minipage}[b]{\linewidth}\raggedright
Requirement
\end{minipage} & \begin{minipage}[b]{\linewidth}\raggedright
Implemented in
\end{minipage} & \begin{minipage}[b]{\linewidth}\raggedright
Verified by
\end{minipage} \\
\midrule\noalign{}
\endfirsthead
\toprule\noalign{}
\begin{minipage}[b]{\linewidth}\raggedright
Requirement
\end{minipage} & \begin{minipage}[b]{\linewidth}\raggedright
Implemented in
\end{minipage} & \begin{minipage}[b]{\linewidth}\raggedright
Verified by
\end{minipage} \\
\midrule\noalign{}
\endhead
\bottomrule\noalign{}
\endlastfoot
FR-1 Configuration object &
\longtok{platform\_local/unified/launcher.py}; run configurations under
\longtok{platform\_hybrid/experiments/} & TC-3 \\
FR-2 Back-end-agnostic execution &
\longtok{platform\_local/trl\_integrations/};
\longtok{platform\_hybrid/experiments/results/framework\_comparison.json}
& TC-3; parity shortfall recorded in §3.2 \\
FR-3 Per-step ZVF & Training-loop telemetry (Phase-1 design) & TC-1 \\
FR-4 Per-step telemetry & Telemetry writer;
\longtok{platform\_hybrid/experiments/results/p5p8/controller\_\allowbreak{}cf\_\allowbreak{}per\_\allowbreak{}step.tsv}
& TC-1, TC-4 \\
FR-5 Raw tensor persistence & Persisted per-group reward tensors & TC-1,
TC-2 \\
FR-6 Sweeps & Sweep runner & TC-3 \\
FR-7 White-box gradients & Single-GPU training path & Not automatable in
the audit suite; checks in Chapter 5 §5.5 (per-layer gradient
profile) \\
FR-8 Provenance record &
\longtok{zvf-program/flagship/e1\_\allowbreak{}e14\_\allowbreak{}completion\_\allowbreak{}gate.py}; run
provenance records & TC-5, TC-7, TC-9 \\
FR-9 Minimum-report manifest &
\longtok{platform\_hybrid/experiments/results/p5p8/minreport\_\allowbreak{}audit\_\allowbreak{}summary.json}
& TC-10 \\
FR-10 Registry query & \longtok{platform\_hybrid/registry/schema.json},
\texttt{query.py}, \texttt{entries/} & TC-10 \\
FR-11 Variant deltas, stack-diff &
\longtok{platform\_hybrid/registry/provenance/} & TC-10 \\
FR-12 Adaptive controller &
\longtok{platform\_hybrid/experiments/results/p5p8/controller\_\allowbreak{}cf\_\allowbreak{}summary.json}
& TC-4 \\
FR-13 Lane harness & \longtok{zvf-program/e1\_wave10/},
\longtok{zvf-program/e5\_successor27/}, \longtok{zvf-program/e2\_core/},
\longtok{zvf-program/e13\_balrog/},
\longtok{zvf-program/e6e9/check\_\allowbreak{}quota\_\allowbreak{}status.py}; runners actually used
on 2026-09-27:
\longtok{outputs/finish\_\allowbreak{}pending\_\allowbreak{}2026-09-27/\textless{}lane\textgreater{}/code/}
& TC-9 \\
FR-14 Completion gate &
\longtok{zvf-program/flagship/e1\_\allowbreak{}e14\_\allowbreak{}completion\_\allowbreak{}gate.py} & TC-9 \\
FR-15 Scope separation &
\longtok{outputs/E1\_\allowbreak{}E14\_\allowbreak{}FINAL\_\allowbreak{}RESULTS\_\allowbreak{}2026-09-19.md} & TC-6, TC-9 \\
FR-16 Blind-review packaging &
\longtok{platform\_local/run\_\allowbreak{}all\_\allowbreak{}audits.py},
\longtok{export\_\allowbreak{}guard\_\allowbreak{}audit.py} & TC-6, TC-8 \\
NFR-1 Attribution & Method sections and comparison artefacts & TC-3,
TC-10 \\
NFR-2 Reproducibility & Seed management in the runner & TC-3 \\
NFR-3 Recomputability & Documented analysis procedures & TC-1, TC-2,
TC-4 \\
NFR-4 Denominator discipline & Caveat and claim-strength checks & TC-6,
TC-7 \\
NFR-5 Isolation & Process-level execution & Operationally checked; not
in the audit suite \\
NFR-6 Budget & Spend receipts & TC-9 \\
NFR-7 Fail-closed behaviour & Gate constants and suite exit status &
TC-6, TC-9 \\
NFR-8 Least privilege & Credential-gate receipts & TC-9 \\
NFR-9 Auditability & Tracked execution paths & TC-9 \\
NFR-10 External traceability & Hash equality and trust-root validation &
TC-9, TC-10 \\
\end{longtable}

Two rows deserve comment. FR-7 has no automated check because the
white-box gradient path is a measurement procedure rather than an
invariant; its evidence is the result data itself. NFR-5 is likewise
checked operationally rather than by the suite, because process
isolation is a property of how runs are launched and is not visible in
any post-hoc artefact.

\section{Hardware Requirements}\label{hardware-requirements}

\textbf{Hardware requirements and their status.}

\begin{longtable}[]{@{}
  >{\raggedright\arraybackslash\hspace{0pt}}p{(\linewidth - 6\tabcolsep) * \real{0.3333}}
  >{\raggedright\arraybackslash\hspace{0pt}}p{(\linewidth - 6\tabcolsep) * \real{0.3333}}
  >{\raggedright\arraybackslash\hspace{0pt}}p{(\linewidth - 6\tabcolsep) * \real{0.3333}}@{}}
\caption{Hardware requirements and their status.}\tabularnewline
\toprule\noalign{}
\begin{minipage}[b]{\linewidth}\raggedright
Compute surface
\end{minipage} & \begin{minipage}[b]{\linewidth}\raggedright
Requirement
\end{minipage} & \begin{minipage}[b]{\linewidth}\raggedright
Status
\end{minipage} \\
\midrule\noalign{}
\endfirsthead
\toprule\noalign{}
\begin{minipage}[b]{\linewidth}\raggedright
Compute surface
\end{minipage} & \begin{minipage}[b]{\linewidth}\raggedright
Requirement
\end{minipage} & \begin{minipage}[b]{\linewidth}\raggedright
Status
\end{minipage} \\
\midrule\noalign{}
\endhead
\bottomrule\noalign{}
\endlastfoot
Managed training back-end (Tinker) & Sampling, training and
checkpointing for the \textasciitilde35B-total, \textasciitilde3B-active
MoE actor in bfloat16 with a pinned LoRA; per-layer gradients are not
exposed, which is why FR-7 is scoped to small models & Met \\
White-box GPU & One CUDA device with at least 24 GB (an L4-class cloud
notebook), enough for the 1.5B configuration and a reduced-length 3B
re-test & Met \\
Serverless GPU (Modal) & Large scoring sweeps and the hosted sampler
bridge; the bridge is a precondition for the lanes that use it & Met \\
Evaluation-lane cloud capacity & Some native environments exceeded
available execution capacity; hosted environments and a runtime adapter
were needed & Historical infrastructure limitation \\
Orchestration workstation & Orchestration, analysis, figures, the audit
suite and the report toolchain & Met \\
\end{longtable}

\emph{Sources: outputs/E1\_\allowbreak{}E14\_\allowbreak{}FINAL\_\allowbreak{}RESULTS\_\allowbreak{}2026-09-19.md;
platform\_tinker/reports/esa\_phase1/Phase1\_\allowbreak{}Project\_\allowbreak{}Report\_\allowbreak{}ZVF.tex;
withheld historical operational chronology; withheld historical campaign
chronology;
outputs/PES\_\allowbreak{}Phase2\_\allowbreak{}Review\_\allowbreak{}2026-09-12/finish/e9\_completion/status\_2026-09-19.json.}

\section{Software Requirements}\label{software-requirements}

The runtime is Python 3.11 with PyTorch and the Hugging Face stack
(Transformers, PEFT, Datasets) for model handling, tokenisation and
adapters; the RL back-ends are TRL, veRL through its HybridFlow path,
and OpenRLHF, driven through the shared launcher, with the Tinker client
SDK for the managed path. Experiment tracking is a weights-and-biases
project; telemetry is stored as JSONL and as persisted tensors; analysis
is carried out by the documented scripts named in the Phase-1
specification.

Semester 4 adds five software components to the requirement set. First,
the audit suite: \longtok{platform\_local/run\_\allowbreak{}all\_\allowbreak{}audits.py} together
with the nine audit modules it imports and the shared data model in
\texttt{utils/audit\_utils.py}, which defines the issue, result and
suite types and the \texttt{METRIC} line format; the suite requires a
Python interpreter with no third-party dependencies beyond the standard
library for its structural checks. Second, the completion gate and the
per-lane harnesses under \texttt{zvf-program/}, which require Ed25519
signature verification and JSON Schema validation against the published
provider-package, provider-grant and trust-root schema versions. Third,
the registry: a schema, an entry store and a query interface, together
with the provenance records under \texttt{platform\_hybrid/registry/}.
Fourth, the report toolchain: a LaTeX distribution capable of the
multi-pass build with BibTeX, plus a markdown assembler that
concatenates the chapter files and, when available, renders the document
through pandoc. Fifth, the judgement-layer harness under
\texttt{zvf-program/jev\_lab/}, which records its own receipts and is
used only for classification, faithfulness checking and value ordering.

\section{Constraints and Assumptions}\label{constraints-and-assumptions}

\textbf{Rewards are binary and verifiable.} The specification assumes a
programmatic checker --- exact match for the mathematics tasks and unit
tests for the code subset. Nothing in the requirement set addresses
learned reward models, and none of the diagnostics should be read as
applying to them.

\textbf{Held-out sets are small.} Several Phase-1 evaluation sets have
\(n\) between 8 and 20, and the controlled Semester-4 comparisons use
200 to 500 items (Table 1.1 gives the n behind each headline), with the
consequence that learning-gain magnitudes are noise-limited by
construction. This is a constraint on what the results can mean, not a
defect of the analysis: the small-\(n\) figures are reported with their
denominators and are not used to support comparative claims.

\textbf{The managed back-end is a black box at the layer level.}
Per-layer gradients are not exposed on the managed path, so claims about
the relationship between reward spread and gradient magnitude are
confined to the small-model white-box configuration and are directional
only.

\textbf{Group-size differences must not be attributed to sampling
noise.} The comparison protocols share a deterministic prompt schedule,
so that method-to-method differences in ZVF reflect training dynamics
rather than independent draws of prompts.

\textbf{Framework parity is incomplete.} The requirement in FR-2 covers
two seven-library rosters, but only two of the seven cross-launcher
entries produced completed runs in this release, and the remaining
comparison entries are dry-run placeholders. Comparative framework
claims are therefore constrained to the rosters and the completed runs
actually recorded, and the repository states the dry-run status
explicitly.

\textbf{The scale study reports descriptive limits rather than a scaling
law.} The cross-scale analysis spans more than seventy runs across seven
libraries and five model families, and its central result is a negative
one: no reliable positive cross-scale slope is recovered and no
saturation exponent is identifiable. The specification inherits that
boundary --- the requirement is to report identifiability limits, not to
assert a law.

\textbf{External dependencies are assumed unavailable until proven
otherwise.} Payload grants, hosted evaluations and quota increases are
treated as absent until a response arrives; a lane without the required
external benchmark scope is closed as externally blocked with a reopen
condition rather than left pending indefinitely.

\textbf{Arithmetic and string lookups stay in code.} Whatever judgement
layer is used for triage, classification or value ordering,
verification-type checks are performed by deterministic code. This rule
was banked after the judgement layer returned an incorrect answer to a
parity probe in a healthy window, and it constrains the architecture: no
requirement in this chapter may be verified by a language model where a
deterministic check is possible.

\textbf{Named variants are labels, not specifications.} The assumption
that a method name identifies a method is false on the available
evidence: entries sharing a label diverge substantially in measured
behaviour, which is why the registry's delta records and flip-risk
verdicts are specified as requirements rather than as documentation
niceties.



\part{Design and Implementation}


\chapter{Proposed Methodology}\label{proposed-methodology}

This chapter sets out the methodology of the work rather than its
results. It states four things in order: the design commitment to
attribution and the four de-confound pillars that operationalise it
(§4.1); the per-step telemetry that makes a run's internal state
observable at the granularity at which the mechanism acts (§4.2); the
formalisation of the Zero-Variance Fraction as a descriptive diagnostic
rather than a metric (§4.3); the eight-study portfolio P1--P8 and the
question each study asks (§4.4); and the design of the E1--E14
evaluation campaign, including its single trained actor, pinned native
evaluators, frozen sampling, receipt model, terminal-state taxonomy and
decontamination handling (§4.5). Throughout, the reporting discipline is
part of the method: figures are tied to named evidence, unavailable
primary artefacts are distinguished from surviving summaries (Appendix B
§B.11), and lanes without an admissible result are recorded as blocked (Figure~\ref{fig:pillars}).


\fig{fig_four_pillars}
{The four de-confound pillars}
{The four de-confound pillars. Each holds the whole stack fixed and varies exactly one factor. In Pillar~2 all nine rows of the method panel, vanilla GRPO included, are a declared simulation projection (Chapter~6).}
{fig:pillars}



\section{The Design Commitment to
Attribution}\label{the-design-commitment-to-attribution}

The premise of the project is that an RL-for-LLM result is not
identified by the name of its algorithm. Two runs both labelled GRPO may
differ in loss form, reference-policy handling, sampler, precision,
base-checkpoint identity, data order and group-size schedule, and the
label alone does not say which of those differences produced the
outcome. The project's response is a design commitment rather than a
post-hoc covariate adjustment: hold the whole experimental stack fixed
and vary exactly one factor at a time, so that an observed difference is
attributable to the factor that was varied rather than to the several
things that changed alongside it. This commitment is the methodological
spine of the whole thesis, and it has a practical corollary that shapes
every experiment reported later: a study is defined by the single axis
it varies and is only interpretable to the extent that every other axis
was genuinely pinned (Figure~\ref{fig:attribution}).


\fig{fig_attribution}
{The attribution design}
{The attribution design. Exactly one factor of the training stack is varied while the remainder is held fixed, so a measured difference cannot be attributed to an unstated framework difference.}
{fig:attribution}



The commitment is implemented in a unified launcher,
\longtok{platform\_local/unified/launcher.py}, which consumes a single
configuration object and dispatches it to thin per-framework adapters
--- \texttt{trl\_integrations/}, \texttt{verl/},
\texttt{platform\_modal/openrlhf/}, \texttt{platform\_hybrid/skyrl/} and
\texttt{platform\_tinker/atropos/} --- with the reference GRPO loop
implemented once in \longtok{platform\_tinker/tinkerrl/grpo.py}. The
reward grader and the decoding path are shared code. That sharing is the
load-bearing part of the design: because the grader and the sampler are
the same objects across back-ends, a difference in outcome between two
frameworks is attributable to the framework, and not to a divergent
reward parser or a differently configured decoder masquerading as an
algorithmic difference. The commitment also bounds what a result may
claim. A completed run supports a statement about that stack, at that
revision, with that grader; it does not license a statement about
``GRPO'' in the abstract, and the project's reporting standard (§4.4,
P5) is built around that distinction.

Four pillars carry the commitment, each holding the remainder of the
stack fixed while varying one factor:

\textbf{The four de-confound pillars.}

\begin{longtable}[]{@{}
  >{\raggedright\arraybackslash\hspace{0pt}}p{(\linewidth - 6\tabcolsep) * \real{0.3333}}
  >{\raggedright\arraybackslash\hspace{0pt}}p{(\linewidth - 6\tabcolsep) * \real{0.3333}}
  >{\raggedright\arraybackslash\hspace{0pt}}p{(\linewidth - 6\tabcolsep) * \real{0.3333}}@{}}
\caption{The four de-confound pillars, the single factor each isolates,
and the design that holds the remainder fixed.}\tabularnewline
\toprule\noalign{}
\begin{minipage}[b]{\linewidth}\raggedright
Pillar
\end{minipage} & \begin{minipage}[b]{\linewidth}\raggedright
Factor isolated
\end{minipage} & \begin{minipage}[b]{\linewidth}\raggedright
Design
\end{minipage} \\
\midrule\noalign{}
\endfirsthead
\toprule\noalign{}
\begin{minipage}[b]{\linewidth}\raggedright
Pillar
\end{minipage} & \begin{minipage}[b]{\linewidth}\raggedright
Factor isolated
\end{minipage} & \begin{minipage}[b]{\linewidth}\raggedright
Design
\end{minipage} \\
\midrule\noalign{}
\endhead
\bottomrule\noalign{}
\endlastfoot
1 · PPO versus GRPO, same stack & Advantage estimator: group baseline
versus learned value head & Same model, runner, clipped surrogate and
step budget on a saturated 0.5B addition task; the arms also differ in
prompts per step (16 × 8 versus 128) \\
2 · Zero-Variance Fraction & Signal starvation and its confounders &
Per-step ZVF logged on every run; the nine-row method panel, vanilla
GRPO row included, is a declared simulation projection (§6.3.2) \\
3 · Trainability and group size & Contrast density against the variance
knob & G ∈ \{2, 4, 8, 16\} × seeds, held out \\
4 · Held-out generalisation & Whether GRPO adds real capability &
Dr.~GRPO versus GRPO, pre-to-post, McNemar \\
\end{longtable}

Pillar 1 asks whether a reported GRPO advantage survives when the
advantage estimator is the main thing that changes. The comparison is
designed on \texttt{Qwen/Qwen2.5-0.5B} two-digit addition over 40 steps
and five seeds paired by seed, evaluated on 200 prompts per arm drawn
fresh from the training generator rather than from a disjoint split (the
two arms' item sets largely differ; §6.2.3), and the exact loss form of
both arms is in Table 4.2. Because that task is near ceiling, the design
can bound an estimator effect but cannot separate ``no effect'' from
``no headroom''; the result is reported in Chapter 6 §6.2.3. The
Qwen3-8B row in the project's own effect-size table is excluded from
Pillar 1 as a provenance conflict and is not treated as an estimate.
Pillar 2 turns a run's internal signal state into a first-class
measurement; it is formalised in §4.3. Pillar 3 varies the group size G
while holding model, data, reward and schedule fixed, and is the axis on
which the preference-density argument of P3 rests. Pillar 4 asks the
question the other three cannot answer --- whether any of this changes
held-out behaviour --- by comparing a pre-RL checkpoint with its post-RL
successor on data the training loop never saw.

The headline configuration of the internal report sources is LoRA rank
32, Adam with β₁ = 0.9 and β₂ = 0.95, learning rate 1 × 10⁻⁴, and the
seed set \{42, 123, 456, 789, 1024\}. It is not held fixed across the
programme. The canonical dispatch matrix uses rank 16 and learning rate
1 × 10⁻⁶, the stack probe uses 1 × 10⁻⁵, the P4 scale extension uses
rank 4, and the Tinker runners default to 3 × 10⁻⁵; Table 4.2 lists the
values each study actually used. The instrumented 79-run roster
described in that source spans the Qwen3, Qwen3.5, Llama-3.x,
DeepSeek-V3.1 and Nemotron-120B families, from 0.6B to approximately
671B total parameters, including mixture-of-experts anchors at the
frontier end of that range; the wider P1 cross-scale roster adds a 1T
anchor (Kimi-K2) and so spans 0.6B to 1T (Chapter 6 §6.1; source:
platform\_hybrid/paper/sections/p1\_abstract.tex). The statistical
apparatus is declared in the shared methods source but was not
pre-registered (Chapter 9 §9.2; the analysis log records 206
iterations). It consists of Welch t-tests and Mann--Whitney U tests for
two-sample contrasts; Cohen's d with Hedges--Olkin intervals; 95\%
bootstrap confidence intervals with 10,000 resamples; rliable aggregate
metrics (interquartile mean, median and optimality gap); and Bonferroni
and Benjamini--Hochberg corrections over a declared test family. Power
was computed rather than assumed: at 80\% power the minimum detectable
effect for an equal-n comparison of five seeds against five is d\_min =
2.024, improving to d\_min = 1.530 when a five-seed group is compared
against a pooled group of fifteen. Those figures are a standing reminder
of how coarse a five-seed design is, and the project reports several
null results that are null precisely because the design could not have
detected anything smaller.

Two limitations belong in this section rather than in a later
discussion. First, the framework roster is broader than the framework
evidence: although seven RL libraries are named in the study design,
only Tinker and TRL produced completed runs, and the veRL and OpenRLHF
entries in \texttt{framework\_comparison.json} are dry-run placeholders
rather than measurements. Cross-library claims are therefore narrower
than the roster implies. Second, the variance decomposition that
motivates the whole attribution commitment is itself four separate
marginal one-way ANOVAs --- library/framework, training algorithm, model
family and model scale --- each run against the same 42 experiments with
valid performance observations. On a corpus where the factors are
collinear a marginal η² absorbs the variance of each factor's
correlates, so the design reads the decomposition only as evidence that
framework and algorithm vary a great deal across this corpus, which is
what motivates holding them fixed --- not as an estimate of any factor's
independent contribution, which a marginal one-way analysis on a
non-orthogonal design cannot identify. The values are reported in
Chapter 6 §6.1. The historical twenty-test multiple-comparison family
that accompanies it pre-dates the Tier-A/B/C partition and is retained
as a reproducibility audit trail, explicitly not as a corrected survival
claim.

\section{Per-Step Telemetry}\label{per-step-telemetry}

Attribution is only as good as the run's observability. The intended
per-step telemetry record contains the Zero-Variance Fraction; gradient
utilisation; the collapse decomposition that separates all-correct from
all-wrong zero-variance groups; and, alongside them, reward, loss,
entropy, completion length and the KL divergence to the reference policy
when one exists and that measurement is exposed; actual availability is
runner-specific. Records are written to JSONL and mirrored to Weights \&
Biases, so a run can be analysed locally without a network dependency
and reviewed remotely without a local environment. The reason ZVF and
its decomposition are logged as separate fields rather than collapsed
into a single starvation number is that the same aggregate fraction can
describe opposite situations: a group of completions that are all
correct and a group that are all wrong both contribute zero within-group
variance, and the remedy for each is different. Logging the
decomposition preserves that distinction at the point of collection
instead of trying to recover it afterwards (Figure~\ref{fig:telemetry}).


\fig{fig_telemetry}
{Per-step telemetry recorded by the unified harness, and the destinations each record flows to}
{Per-step telemetry recorded by the unified harness, and the destinations each record flows to.}
{fig:telemetry}



The quantity being logged is the group-relative advantage actually used
by the runner:

\[A_i = \frac{r_i - \mu_g}{\sigma_g + \varepsilon},\qquad \mu_g = \frac{1}{G}\sum_j r_j,\qquad \sigma_g = \sqrt{\frac{1}{G}\sum_j (r_j - \mu_g)^2}\]

This form is taken from the Phase-1 report. Throughout this thesis
\(\sigma_g\) is the population standard deviation, with 1/G as written
above. Every GRPO-style runner in Table 4.2 computes it that way. Two
places use the unbiased 1/(G − 1) form instead: the reference
\texttt{zvf()} of §4.3, whose \texttt{ddof=1} sample variance leaves
exact ties unchanged but can change classification at a positive
tolerance, and the batch-wide whitening of the PPO arm in the same-stack
run. The ZVF-definition figure (Figure \ref{fig:zvfdef}) also shows the
unbiased form, which gives \(\sigma_g \approx 0.535\) for its example
group where the population form gives 0.5. When \(\sigma_g = 0\) every
\(A_i\) is zero, so the whole group contributes no gradient through the
advantage term. Whether anything else acts depends on the executed loss:
KL regularisation and auxiliary or critic losses can remain nonzero.
Clipping alone does not create a policy-surrogate gradient when the
advantage is zero. The zero-gradient statement is therefore confined to
the reward-relative surrogate in verified group-centred implementations;
it is not a claim that the whole optimiser step vanishes. Table 4.2
distinguishes verified absence, configuration evidence and unknown
managed settings.

Two further design choices distinguish this telemetry from conventional
training curves. The first is that the raw per-group reward tensors are
retained, not only the summaries computed from them. This makes a
published number auditable by recomputation: any downstream diagnostic,
such as \longtok{p2\_\allowbreak{}collapse\_\allowbreak{}analysis.py}, can reproduce a reported
statistic from the stored tensors rather than trusting a logged summary.
The convention has been exercised: study P2 recomputes three stored
GSM8K group tensors (Qwen3-8B sampling runs, G = 8) from their raw
values rather than reading the figures back from a log line, and the
recomputed values are reported in Chapter 6 §6.3.1. The second is that
the logging layer is treated as part of the apparatus and is expected to
fail. An early Weights \& Biases integration used the library's
thread-unsafe global run object from multiple worker threads, which
produced repeated ``Run is finished'' summary-write failures; the fix
was to give each run its own process in \texttt{campaign.py}. The defect
is recorded because telemetry that silently drops writes is
indistinguishable from telemetry showing a healthy run, and the remedy
--- process-level isolation per run --- is a structural property of the
campaign rather than a patch.

\section{Formalising the Zero-Variance
Fraction}\label{formalising-the-zero-variance-fraction}

We distinguish the exact-tie fraction, \(\mathrm{ZVF}_0\), from the
tolerance-based diagnostic, \(\mathrm{ZVF}_{\varepsilon}\). The first
counts groups with identical observed rewards; the second also counts
near-ties. Unless explicitly labelled exact-tie, ZVF here uses the
following sample-variance threshold. Formally, for a batch \(B_t\) of
prompt groups at step \(t\), each group \(g\) contributing \(K\) sampled
completions with rewards \(r_{g,1}, \ldots, r_{g,K}\),

\[\mathrm{ZVF}_t = \frac{1}{|B_t|} \sum_{g \in B_t} \mathbf{1}\!\left[\widehat{\mathrm{Var}}(r_{g,1}, \ldots, r_{g,K}) \le \varepsilon\right]\]

where \(\widehat{\mathrm{Var}}\) is the unbiased sample variance and ε =
10⁻⁶ for rewards normalised to {[}0, 1{]}. The exact-tie formulation
used in some internal sources instead writes
\(\mathrm{ZVF}_t = \frac{1}{|P_t|} \sum_{p \in P_t} \mathbf{1}[\mathrm{Var}_{c \sim C(p)}\, r(c) = 0]\),
with \(C(p)\) the set of sampled completions for prompt p.~These
formulations agree for the stored binary rewards at G = 8 and ε = 10⁻⁶,
but are not interchangeable for perturbed or continuous rewards. In
particular, the micro-jitter test changes exact ties while leaving this
tolerance-based diagnostic unchanged (§6.3.1). The reference
implementation is deliberately small, so that the diagnostic cannot
acquire hidden behaviour:

\begin{verbatim}
def zvf(rewards_2d: np.ndarray, eps: float = 1e-6) -> float:
    # rewards_2d shape: [num_groups, K]
    var_per_group = rewards_2d.var(axis=-1, ddof=1)
    return float((var_per_group <= eps).mean())
\end{verbatim}

Gradient utilisation is then defined as its complement,
\(\mathrm{GU}_t = 1 - \mathrm{ZVF}_t\). The complement is a descriptive
accounting identity, not a claim about optimiser behaviour. Its scope
must be read against the runner actually used, so the loss form of every
study is listed in Table 4.2, read from the runner code rather than from
the internal reports' prose. Three facts in the table govern how ZVF is
interpreted. First, an exact-tie group has zero reward-relative
surrogate gradient only where the runner actually centres advantages
within that group. Near-ties counted by \(\mathrm{ZVF}_{\varepsilon}\)
need not have exactly zero advantages. TRL's dump records β = 0.04 and
an on-device reference model, its launcher leaves β unspecified, the PPO
roster inherits adaptive-KL defaults, and the managed dump has β = null.
None of these is evidence of zero executed KL; unknown remains unknown.
Even verified zero-KL runs can update through optimiser state,
regularisation or auxiliary losses. Second, the inspected custom-loss
Tinker reference runners are REINFORCE with a group-mean baseline (the
actor's historical source binding is incomplete; Appendix D): the loss
is \(-A_i\) times the summed token log-probability, averaged over
completions, with no importance ratio and no clipping, and this thesis
calls them that rather than GRPO in the strict sense. Third, the runners
labelled GRPO are not one loss: the same-stack and Dr.~GRPO studies use
a clipped ratio with two inner epochs, the Tinker runners do not, and
the Phase-1 four-method run divides by no σ even in its ``grpo'' arm.
The scope must consequently be checked per row, including auxiliary
losses, masks and source-binding gaps, rather than inferred from the
algorithm label (Figures~\ref{fig:zvfdef} and~\ref{fig:gu}).


\fig{fig_zvf_definition}
{Definition of the Zero-Variance Fraction}
{Definition of the Zero-Variance Fraction. One prompt produces $G$ completions with scalar rewards; when the within-group reward standard deviation is zero the advantage is identically zero and the sample contributes no reward-relative surrogate gradient. Exact ties and the sample-variance threshold are distinguished in the text; they agree for the binary rewards illustrated here. The composition bar is measured on Qwen3-8B, held-out GSM8K slice, $G=8$, $T=1.0$, 200 problems $\times$ 3 seeds; all-correct (vertical hatch) and all-wrong (cross-hatch) together make up the ZVF.}
{fig:zvfdef}



\fig{fig_gradient_utilisation}
{Gradient utilisation, $1-\zvf$, as a function of group size $G$}
{Gradient utilisation, $1-\zvf$, as a function of group size $G$. Qwen2.5-0.5B, arithmetic-correctness task, 40 steps, 16 prompts, $G\in\{2,4,8,16\}$, seeds 42/123/456 ($n=3$ per $G$; \longtok{groupsize\_\allowbreak{}zvf\_\allowbreak{}sweep.json}). (a) Open circles are single seeds; squares are seed means with 95\% $t$-intervals ($df=2$). The line is a two-parameter power law, $\zvf = 0.922\,G^{-0.137}$, fitted to the four means (2 residual degrees of freedom), so it is descriptive rather than a validated law; $G=32^{\ast}$ was not trained and the dotted segment is extrapolation. (b) The absolute gain per doubling is nearly constant ($+0.074$, $+0.073$, $+0.059$), so the diminishing return is relative: $46\%\to31\%\to19\%$ of the previous level. Utilisation stays below 0.37 at $G=16$.}
{fig:gu}




\textbf{Loss form by study, read from the runner code.}

\begingroup\scriptsize\setlength{\tabcolsep}{3pt}
\begin{longtable}[]{@{}
  >{\raggedright\arraybackslash\hspace{0pt}}p{(\linewidth - 22\tabcolsep) * \real{0.0909}}
  >{\raggedright\arraybackslash\hspace{0pt}}p{(\linewidth - 22\tabcolsep) * \real{0.0909}}
  >{\raggedright\arraybackslash\hspace{0pt}}p{(\linewidth - 22\tabcolsep) * \real{0.0909}}
  >{\raggedright\arraybackslash\hspace{0pt}}p{(\linewidth - 22\tabcolsep) * \real{0.0909}}
  >{\raggedright\arraybackslash\hspace{0pt}}p{(\linewidth - 22\tabcolsep) * \real{0.0909}}
  >{\raggedright\arraybackslash\hspace{0pt}}p{(\linewidth - 22\tabcolsep) * \real{0.0909}}
  >{\raggedright\arraybackslash\hspace{0pt}}p{(\linewidth - 22\tabcolsep) * \real{0.0909}}
  >{\raggedright\arraybackslash\hspace{0pt}}p{(\linewidth - 22\tabcolsep) * \real{0.0909}}
  >{\raggedright\arraybackslash\hspace{0pt}}p{(\linewidth - 22\tabcolsep) * \real{0.0909}}
  >{\raggedright\arraybackslash\hspace{0pt}}p{(\linewidth - 22\tabcolsep) * \real{0.0909}}
  >{\raggedright\arraybackslash\hspace{0pt}}p{(\linewidth - 22\tabcolsep) * \real{0.0909}}@{}}
\caption{Loss form by study, read from the runner code. M =
\protect\longtok{platform\_hybrid/experiments/modal/}, T =
\protect\longtok{platform\_hybrid/experiments/tinker-runs/}, C =
\protect\longtok{zvf-program/colab-experiments/}. ``n/a'' means the study
does no training; ``not recoverable'' means the repository does not hold
the code or configuration that ran.}\tabularnewline
\toprule\noalign{}
\begin{minipage}[b]{\linewidth}\raggedright
Study / run family
\end{minipage} & \begin{minipage}[b]{\linewidth}\raggedright
Runner
\end{minipage} & \begin{minipage}[b]{\linewidth}\raggedright
Ratio / clip
\end{minipage} & \begin{minipage}[b]{\linewidth}\raggedright
KL / reference
\end{minipage} & \begin{minipage}[b]{\linewidth}\raggedright
Token mask
\end{minipage} & \begin{minipage}[b]{\linewidth}\raggedright
Advantage
\end{minipage} & \begin{minipage}[b]{\linewidth}\raggedright
Critic
\end{minipage} & \begin{minipage}[b]{\linewidth}\raggedright
LoRA r
\end{minipage} & \begin{minipage}[b]{\linewidth}\raggedright
lr
\end{minipage} & \begin{minipage}[b]{\linewidth}\raggedright
G
\end{minipage} & \begin{minipage}[b]{\linewidth}\raggedright
Steps
\end{minipage} \\
\midrule\noalign{}
\endfirsthead
\toprule\noalign{}
\begin{minipage}[b]{\linewidth}\raggedright
Study / run family
\end{minipage} & \begin{minipage}[b]{\linewidth}\raggedright
Runner
\end{minipage} & \begin{minipage}[b]{\linewidth}\raggedright
Ratio / clip
\end{minipage} & \begin{minipage}[b]{\linewidth}\raggedright
KL / reference
\end{minipage} & \begin{minipage}[b]{\linewidth}\raggedright
Token mask
\end{minipage} & \begin{minipage}[b]{\linewidth}\raggedright
Advantage
\end{minipage} & \begin{minipage}[b]{\linewidth}\raggedright
Critic
\end{minipage} & \begin{minipage}[b]{\linewidth}\raggedright
LoRA r
\end{minipage} & \begin{minipage}[b]{\linewidth}\raggedright
lr
\end{minipage} & \begin{minipage}[b]{\linewidth}\raggedright
G
\end{minipage} & \begin{minipage}[b]{\linewidth}\raggedright
Steps
\end{minipage} \\
\midrule\noalign{}
\endhead
\bottomrule\noalign{}
\endlastfoot
Tinker reference loop; actor preset reconstruction (Appendix D) &
\longtok{platform\_tinker/tinkerrl/grpo.py} & none (REINFORCE, group-mean
baseline) & none & full sequence & (r − μ)/(σ + 10⁻⁸), 1/G & none & 32 &
1e-5 (actor); 3e-5 default & 4 (actor); 8 default & 40 (actor) \\
P1 cross-scale Tinker anchors &
T/\longtok{scripts/tinker\_\allowbreak{}parallel\_\allowbreak{}runner.py},
T/\longtok{campaign\_v2.py},
T/\longtok{scripts/world\_\allowbreak{}class\_\allowbreak{}experiments.py}; per-anchor lineage for
the frontier anchors not recoverable & none in candidate source;
execution unknown & none in candidate source; execution unknown & full
sequence in candidate source & (r − μ)/σ, 1/G in candidate source & none
& 32 & 3e-5 & 8 & 20--30 logged \\
P2 recomputed tensors (0.130 / 0.190 / 0.155) &
\longtok{platform\_hybrid/experiments/tinker\_\allowbreak{}direct\_\allowbreak{}eval.py} (sampling
only, Qwen3-8B) & n/a & n/a & n/a & ZVF from population variance & n/a &
n/a & n/a & 8, T = 1.0 & 0 \\
P2 Phase-1 four-method batch (0.72--0.77) &
T/\longtok{scripts/n2\_\allowbreak{}reward\_\allowbreak{}tensor\_\allowbreak{}20260704.py} & Tinker
\longtok{importance\_sampling} ratio against sampler log-probs; no clip &
none & completion only & r − μ, no σ division (all arms) & none & 32 &
1e-5 & 8 & 40 \\
P3 Qwen2.5-0.5B arithmetic sweep &
M/\longtok{modal\_\allowbreak{}groupsize\_\allowbreak{}zvf\_\allowbreak{}sweep.py} & none & none & completion
only & (r − μ)/(σ + 10⁻⁶), 1/G & none & 16 (q, v) & 1e-4 & 2, 4, 8, 16 &
40 \\
P3 Qwen3-8B matched-budget panel & T/\longtok{live\_\allowbreak{}zvf\_\allowbreak{}probe.py}
(schema match; no launcher link) & none in candidate source; execution
unknown & none in candidate source; execution unknown & full sequence in
candidate source & (r − μ)/(σ + 10⁻⁸), 1/G in candidate source & none &
4 & 3e-5 & 2 / 16 & 160 / 20 \\
P3 Qwen3.5-4B openings &
\longtok{platform\_hybrid/experiments/openings/groupsize\_zvf.py} & none
& none & full sequence & (r − μ)/σ, 1/G; zero-variance groups skipped &
none & 8 & 3e-5 & 2--16 & 8 \\
P4 arithmetic, GRPO vs Dr.~GRPO & M/\longtok{modal\_\allowbreak{}drgrpo\_\allowbreak{}vs\_\allowbreak{}grpo.py}
& token-level ratio, clip 0.2, 2 epochs & none & completion only & GRPO
(r − μ)/(σ + 10⁻⁶), 1/G; Dr.~GRPO r − μ & none & 16 & 1e-4 & 8 & 40 \\
P4 GSM8K-CoT, Qwen2.5-1.5B-Instruct &
M/\longtok{modal\_\allowbreak{}drgrpo\_\allowbreak{}gsm8k\_\allowbreak{}cot.py} & as above & none & completion
only & as above & none & 16 & 1e-5 & 8 & 30 \\
P4 Qwen3-8B uncapped extension & T/\longtok{live\_\allowbreak{}zvf\_\allowbreak{}probe.py}
(\texttt{-\allowbreak{}\/\allowbreak{}-\allowbreak{}loss\ drgrpo}) & none & none & full sequence & GRPO 1/G;
Dr.~GRPO r − μ & none & 4 & 3e-5 & 8 & 30 \\
Same-stack PPO vs GRPO & M/\longtok{modal\_\allowbreak{}samestack\_\allowbreak{}ppo\_\allowbreak{}grpo.py} &
both arms: sequence-level ratio on length-normalised log-prob, clip 0.2,
2 epochs & none & completion only & GRPO (r − μ\_g)/(σ\_g + 10⁻⁶), 1/G;
PPO r − V, whitened over the batch with 1/(n − 1) & PPO: linear value
head on the prompt-end hidden state, 0.5 · MSE loss; no GAE (one-step) &
16 (q, v) & 1e-4 & GRPO 16 prompts × 8; PPO 128 prompts × 1 & 40 \\
P7 Colab four-arm audit & C/\longtok{e3\_\allowbreak{}open\_\allowbreak{}audit.py} & sequence-level
ratio, clip 0.2 (DAPO 0.2 / 0.28), 2 epochs & none & completion only &
population variance; Dr.~GRPO r − μ; zero-variance groups dropped & none
& full fine-tune & 2e-6 & 4, adaptive up to 10 & 10 \\
P7 Tinker adaptive-G & T/\longtok{scripts/qp7\_\allowbreak{}adaptive\_\allowbreak{}g\_\allowbreak{}20260704.py}
& Tinker \longtok{importance\_sampling} ratio; no clip & none &
completion only & r − μ & none & 32 & 1e-5 & 4 → 6 → 8 & 16 \\
Tinker side of the Tinker/TRL pair & launcher not recoverable; config
dump \longtok{framework\_\allowbreak{}config\_\allowbreak{}dumps/tinker\_\allowbreak{}qwen3\_\allowbreak{}8b\_\allowbreak{}gsm8k.yaml} &
not recoverable & unknown (managed; dump β null) & not recoverable & not
recoverable & none & 16 (dump) & 1e-5 & 8 & 30 \\
TRL side of the Tinker/TRL pair & M/\longtok{modal\_\allowbreak{}grpo\_\allowbreak{}trl.py} (TRL
\longtok{GRPOTrainer}) & TRL defaults, not set in code & code
unspecified; dump β = 0.04, on-device reference & TRL completion mask &
TRL default scaling & none & 32 in code, 16 in dump & 1e-5 & 8 & 30 \\
P1-roster PPO baselines & M/\longtok{modal\_\allowbreak{}ppo\_\allowbreak{}campaign.py} (TRL
\longtok{PPOTrainer}) & TRL PPO defaults (clip, GAE, adaptive KL) &
adaptive-KL defaults; executed settings unrecorded & TRL default & TRL
whitening & value head & 32 & 1e-5 & n/a & not recorded \\
\end{longtable}
\endgroup

Three properties make the definition more than a restatement of the
training loop. First, it is non-tautological under a simple null: if
completions were independent Bernoulli draws with per-prompt success
probability \(p_g\), then
\(\mathbb{E}[\mathrm{ZVF}_t] = \mathbb{E}_{p_g}[p_g^K + (1-p_g)^K]\), a
quantity that depends on the model's accuracy distribution and the group
size rather than on the observed batch. The identity
\(\mathrm{pass@}G - p^G = 1 - \mathrm{ZVF}\) follows algebraically under
the same i.i.d.-Bernoulli model and is therefore reported as an
algebraic remark rather than as evidence. Second, it is checked against
that null on instrumented runs: a Qwen2.5-0.5B sweep over G ∈ \{2, 4, 8,
16\}, three seeds and 40 steps --- twelve runs and 480 logged steps ---
compares per-step observed ZVF with the Bernoulli-null prediction
\(p^G + (1-p)^G\). Because the null for each step is computed from the
same step's mean reward, a high correlation is partly mechanical, so the
comparison is designed as a shape check rather than a validation, and
the baseline used is the per-step average of the closed form (results in
Chapter 6 §6.3.3). Third, at larger scale the same instrument is run on
Qwen3-8B against GSM8K test problems through the Tinker managed API at G
= 8, temperature 1.0, over 200 problems and three seeds (sampling only,
no training; these are the same three stored tensors whose ZVF P2
recomputes), with the zero-variance groups decomposed into all-correct
and all-wrong and compared with the same null. Deviations from the null
in either direction are not read as evidence about dependence between
completions.

That the quantity is portable across libraries is a separate, weaker
claim, and the project states its limits explicitly. A normaliser at
\longtok{platform\_modal/scripts/zvf\_\allowbreak{}compute\_\allowbreak{}cross\_\allowbreak{}framework.py} maps
each framework's reward and mask fields onto the common definition ---
TRL's \texttt{rewards} with \texttt{batch\_size\ /\ group\_size} and
\texttt{completion\_mask}; Tinker's \texttt{rollout.reward} with
\texttt{rollout.group\_id} and \texttt{rollout.eos\_mask}; OpenRLHF's
\texttt{reward\_score} with \texttt{prompt\_index} and
\texttt{attention\_mask}; veRL's \texttt{data\_source/reward} with
\texttt{uid} and \texttt{response\_mask} --- and emits JSONL records of
the form \texttt{\{step,\allowbreak{}\ zvf,\allowbreak{}\ batch\_\allowbreak{}\allowbreak{}reward\_\allowbreak{}\allowbreak{}mean,\allowbreak{}\ n\_\allowbreak{}groups,\allowbreak{}\ K\}}.
The scope correction attached to that parser is that it normalises
reward logs and does not make its runner canonical GRPO; because the
project does not bundle raw per-group rollout JSONL for every framework,
no fresh four-way pointwise ZVF overlay is claimed, only the narrower
implementation-level portability claim. The threshold ε is fixed in
advance, and its sensitivity is reported with the P2 results in Chapter
6 §6.3.1. Boundary conditions are declared where the diagnostic does not
apply: dense or continuous reward signals, K = 1, SFT-saturated
baselines and non-outcome-reward RL all fall outside the setting in
which the definition is meaningful.

The methodological point of this section is that ZVF is a diagnostic and
not a metric. The design commits to that position before the results
rather than inferring it from them, and it takes four precautions.
First, because the statistic is unsigned, a high value is read as
consistent with either a solved or an unsolved task and never as a score
in one direction; the 368-run audit that shows both extremes is reported
in Chapter 7 §7.1 and §7.3. Second, a pooled association between ZVF and
performance is always broken down by task type before it is read,
because the task mix of a pool can set the sign of the association
(results in Chapter 6 §6.3.4). Third, the association with final
accuracy and with final held-out outcome is tested, but ZVF is not
promoted to a causal or incrementally predictive statistic whatever that
test returns. Fourth, no step-level p-value is computed on the 480
autocorrelated steps of the instrumented sweep, because treating
correlated within-run steps as independent samples would manufacture
significance. One further boundary is recorded for the same reason: a
variance-mitigation table in the internal report sources is a declared
simulation projection and not evidence, and its rows are read as
synthetic rather than as independent library deployments. The result of
these choices is a diagnostic that is useful for describing where within
a batch the learning signal is and is not, and that is explicitly not
proposed as a quantity to be maximised, ranked or reported as a headline
score.

\section{The Eight-Study Portfolio}\label{the-eight-study-portfolio}

Semester 3 delivered the shared code and the initial experimental
foundation; the Semester 4 programme --- the research direction, the
expanded experiments, the analysis and the P1--P8 internal reports ---
is the present author's solo continuation, and the roster records 206
analysis iterations behind the numbers reported in this thesis. The
portfolio is organised as eight studies, each written up as an
unpublished internal report, each asking one question and each carrying
its own claim boundary; the titles quoted below are the internal
reports' working titles. This section states the questions and designs;
the results are in Chapters 6--7 and, for P8, Appendix F. The questions
and boundaries are stated here because they are methodological
commitments: a study's question determines which comparisons it may
make, and the boundary states what it declines to conclude.

P1, ``Scaling Laws for GRPO Post-Training of Language Models: A
Cross-Library, Cross-Scale Study'', asks whether GRPO post-training
outcomes scale with model scale in a way that can be fitted and
extrapolated. Its evidence base is more than seventy runs across seven
RL libraries and five model families from 0.6B to 1T parameters, fitted
with saturating forms and tested for a cross-scale slope. The frontier
anchors --- Qwen3-235B-MoE, DeepSeek-V3.1, Kimi-K2 and Nemotron-120B ---
were run through the Tinker managed API in a single-seed,
descriptive-only regime, so the report's claim boundary is a descriptive
limits-and-identifiability audit rather than a positive scaling law;
matched multi-seed cross-scale runs are an explicit non-claim. Results
are in Chapter 6 §6.2.

P2, ``The Zero-Variance Fraction: A Descriptive Diagnostic for Signal
Starvation in GRPO'', asks how much of a training batch carries usable
group-relative signal, and whether that fraction predicts the final
outcome. The study characterises ZVF descriptively, treats the identity
of §4.3 as algebraic rather than empirical, tests the association with
final held-out outcome without promoting ZVF to a predictive statistic,
and marks the variance-mitigation table as a simulation projection; the
retired roster item that previously carried a causal framing now lives
under \longtok{archive/absorbed/R02\_\allowbreak{}main\_\allowbreak{}zvf/}. Results are in Chapter
6 §6.3.

P3, ``Group Size in GRPO: Contrast Density and the Bridge to DPO'', asks
how trainability varies with the group size G and whether GRPO admits an
all-pairs preference interpretation. Its measured sweep covers G ∈ \{2,
4, 8, 16\} on the near-ceiling arithmetic task, with seed-level
equivalence tests against G = 2 and a one-way test of held-out accuracy
across G. No directly trained matched-budget G = 32 cell exists, so
every G = 32 training-performance value in P3 comes from an explicitly
labelled reconstructed token-budget grid and is held as a hypothesis
rather than a measurement. The claim boundary is that the data may not
be read as a universal optimum, and the all-pairs preference reading
resembling DPO is offered as an interpretation and not as a proven
equivalence. Results, including a conditional utility audit over 505
tasks, are in Chapter 6 §6.4.

P4, ``Length Bias and Held-Out Generalization in GRPO and Dr.~GRPO'',
asks whether the verbosity trap associated with GRPO appears in the
measured regime and whether either method changes held-out accuracy. Its
controlled comparison runs GRPO against Dr.~GRPO on a near-ceiling
arithmetic task and in a short-horizon GSM8K chain-of-thought regime
with a 200-token generation cap, alongside a multi-seed Qwen3-8B pre-RL
versus post-RL held-out control. The 200-token cap bounds the claim:
under a cap that already controls length, absence of the verbosity
signature would be the expected outcome rather than a positive finding.
Results are in Chapter 6 §6.5.

P5, ``Report the Stack, Not the Label: RL-for-LLM Results Are
Stack-Conditioned'', asks how much of a reported result is attributable
to the stack rather than to the algorithm label, and what a reportable
stack must contain. Three exhibits carry the argument: a nominally
matched-configuration backend audit, a same-label ``DAPO'' comparison
between true dynamic sampling on an open trainer and an asymmetric-clip
surrogate on a closed stack, and a twelve-cell head-to-head of four
labelled algorithms on one fixed stack; their results are in Chapter 7
§7.1. The study proposes an eight-item minimum reportable stack (seven
run-manifest fields and one evaluation item) --- loss form;
reference-policy and KL handling; sampler, backend and precision
including base-checkpoint identity by revision or hash; per-step ZVF and
gradient-utilisation trajectory; group-size schedule; a disjoint
held-out split; decontamination plus a parser-robustness probe; and
held-out pass@k curves alongside pass@1 --- and formalises an
RL-reproducibility threat model with flip-risk grades R0--R5. This is
the standard this thesis is written to.

P6, ``GRPO-Registry: A Machine-Readable Catalog of Group-Relative RL
Stacks and Their Variant Deltas'', asks whether stacks can be compared
mechanically rather than through prose. It defines stack records
carrying the seven manifest fields of the MIN-REPORT-RL standard, a
per-component variant-delta status, three variant-delta records verified
against the DAPO, Dr.~GRPO and GSPO papers, twenty seed stack entries,
and a reference CLI with a 0--100 MIN-REPORT auditor badge and a
\texttt{stackdiff} command that maps a pair of entries to an R0--R5
flip-risk verdict. The study is explicitly a resource study rather than
a second ranking, and it cites P5's seventeen-fold exhibit only to
motivate the schema. Results are in Chapter 7 §7.2.

P7, ``From Diagnostic to Controller: A Signal-Starvation Theory of GRPO
and an Adaptive Group-Size Intervention'', asks whether the diagnostic
of §4.3 can be turned into a control signal. It models usable learning
signal per prompt as \(S = p(1-p)(1-h_G(p))\) with
\(h_G(p) = p^G + (1-p)^G\), yielding three testable consequences: that
the expected ZVF indicator equals \(p^G + (1-p)^G\); that larger G
lowers ZVF; and that gradient magnitude tracks p(1 − p). Each
consequence is paired with a measurement that could falsify it: the
368-run audit, a model-by-group-size grid, and a toy-scale
open-backward-pass experiment whose result is admitted as directional
only. The study promotes pairwise-contrast density, PCD = (G − 1)/G ·
E{[}p(1 − p){]}, as the candidate control signal, subject to an
exact-tie micro-jitter sensitivity test, which does not falsify
tolerance-based ZVF, and then audits an adaptive group-size controller
against fixed recipes and DAPO-style dynamic sampling on a four-arm open
trainer and in a counterfactual controller trace. Results are in Chapter
7 §7.3. The methodological commitment that governs this study is stated
at its head: adaptive group size is not promoted, and no result
establishes prospective superiority over static G = 16, Dr.~GRPO or
dynamic sampling.

P8 occupies the eighth slot with a study outside the RL programme, ``LLM
vs.~XGBoost in Credit-Card Fraud: Sensor and Scribe, Not Scorer''. Its
question is whether a large language model should be used as the fraud
scorer or as a component around one, and it is designed as a comparison
of a boosted-tree model with a Qwen3.5-4B scorer on held-out splits,
with the scorer's cost as a triage component recorded alongside its
accuracy. Results are in Appendix F. The roster evolved after the
semester-4 title list was fixed, and the present position is recorded as
it stands: the fraud study is parked outside the ZVF thesis, contributes
no RL evidence, and is retained for its measurement discipline; the
active P8 slot in the current report roster is an exploratory workshop
artifact. No result from P8 is used anywhere in this thesis to support
an RL claim.

\section{The Design of the E1--E14 Evaluation
Campaign}\label{the-design-of-the-e1e14-evaluation-campaign}

The evaluation campaign is the applied half of the programme and is
deliberately narrow in a way that makes its results interpretable: it
evaluates one actor against fourteen benchmark suites that were not part
of its training data under a receipt, budget and decontamination regime,
instead of training many actors and comparing them. The single actor is
the base model \texttt{Qwen/Qwen3.6-35B-A3B} at commit
\texttt{995ad96e…}, with adapter \texttt{campaign-adapter} at commit
\texttt{withheld-adapter-revision…}, in bfloat16. How that adapter was
trained is reconstructed in Chapter 8 §8.1.1. ``Held-out'' is used here
in a narrow sense: the actor was trained only on API-Bank-RLVR and
SWE-Gym, and none of the fourteen suites is in that training data. It
does not mean that every suite was decontamination-checked. The content
check below covers LAB-Bench, AgentDojo and Omni-MATH, the E11 receipt
records its own verified decontamination, and the E10 receipt itself
declares \texttt{heldout\_claim:\ false}. One actor fixes the object of
study: every reported per-suite figure describes the same frozen policy,
so a difference between two suites is a property of the suite and not of
a model swap. The design consequence is that the campaign cannot support
model-comparison claims at all, and it makes none (Figure~\ref{fig:campaign}).


\fig{fig_campaign_design}
{Design of the E1--E14 native-benchmark evaluation campaign}
{Design of the E1--E14 native-benchmark evaluation campaign. One frozen actor serves all fourteen lanes, and each suite is graded by its own native evaluator at a pinned revision. Each lane lists its original-contract suite and, where one exists, its replacement scope side by side; Figure~\ref{fig:lanes} names the replacement scope where one exists. Figures and terminal states come from the 2026-09-19 ledger (\longtok{E1\_\allowbreak{}E14\_\allowbreak{}FINAL\_\allowbreak{}RESULTS\_\allowbreak{}2026-09-19.md} and \texttt{finish/Pending\_Experiments.md}; eleven named deterministic checks pass, scope in App.~B), updated with the six replacement scopes scored on 2026-09-27 ($\star$; \longtok{finish\_\allowbreak{}pending\_\allowbreak{}2026-09-27/<lane>/result.json}), which are separate receipts and are never pooled with a lane's original-contract figure. E14's $2271/4428 = 51.29\%$ counts the 2 unjudged rows as failures; the official scorer reports $51.31\% = 2271/4426$ on the judged rows.}
{fig:campaign}



Evaluators are the suites' own scorers, used natively rather than
reimplemented, and they are pinned by revision. The surviving native
runners live under \texttt{zvf-program/flagship/}, including
\longtok{e1\_\allowbreak{}swe\_\allowbreak{}bench\_\allowbreak{}pro\_\allowbreak{}full\_\allowbreak{}eval.py},
\longtok{e11\_\allowbreak{}verilog\_\allowbreak{}eval\_\allowbreak{}decontamination.py},
\longtok{eval\_\allowbreak{}pavlov\_\allowbreak{}sdab.py}, \longtok{eval\_\allowbreak{}pavlov\_\allowbreak{}xlam.py},
\longtok{mle\_\allowbreak{}bench\_\allowbreak{}eval.py} and \longtok{eval\_\allowbreak{}apex\_\allowbreak{}agents.py}. The
evaluation-side revisions pinned for the decontamination audit are
LAB-Bench \longtok{998a8e0a40cf116c80e1b0e7a805ebb5fb9fa838}, AgentDojo
\longtok{089ed468cf3ed0322acc66b0211f26d9d90dbf60} and Omni-MATH
\longtok{40ba231d8f16e29ecd40e6407e2c8640145a8f62}, with the
training-side envelope pinned at API-Bank-RLVR
\longtok{bf67c42626f02c305514b1df16dcabc5fc616333} (597 rows, MIT) and
SWE-Gym \longtok{bb94ed9e39bbeb96a7fcbfb533b80f25a7fd59cb} (2,438 train
rows, MIT). Using the native scorer is a substantive choice and not a
convenience: when the Omni-MATH lane produced two rows whose judged
outputs were truncated before any equivalence verdict, the native
scorer's omission of missing-verdict rows was adjudicated as correct
official behaviour rather than a parser bug, and the lane closed with no
re-judge because re-judging would deviate from the official scorer and
generate new evidence with no effect on the score.

For the training-selection protocol, sampling is frozen rather than
merely recorded: temperature 0.1, top-p 0.95, one sample per row, a
maximum of 128 response tokens, with the service seed fixed where the
provider supports it and otherwise recorded in the receipt. Prompt and
response limits are 1,200 and 128 tokens respectively, and the same
source requires the strict verifier implementation to be hashed before
paid work, with a perfect call defined as a verifier score of exactly
1.0 and no parser, target, split or verifier change permitted
mid-campaign. Freezing the sampler addresses the failure mode in which a
benchmark score quietly improves because decoding became more favourable
between runs; it also fixes the comparison between the base reference
and the adapter, which is evaluated side by side on untouched manifests.
The 128-token setting belongs to that protocol only. The E1--E14 lanes
were decoded under each suite's own settings, which range from 1,024 to
8,192 output tokens at temperatures 0 to 0.2, and are recorded per lane
in the serving table of Chapter 8 §8.1.2. The actor's own training run
also departed from the protocol in several fields (§8.1.1).

The receipt model is freeze-before-paid-work, and it is the part of the
campaign design that carries the most methodological weight. Before any
paid phase a run must record an immutable base-model revision, tokenizer
revision, container or environment digest and adapter revision; the
selection slice and the disjoint holdout are re-created from a pinned
dataset revision and sealed, are write-once, are not mounted in
training, and the word ``holdout'' is reserved for a suite or manifest
only once its split, task-ID, license, runtime and decontamination
receipts are complete. The selection slice is used for arm decisions and
is therefore unsuitable for confirmatory claims; the final evaluation
manifests are never used to choose an arm, alter a prompt, or re-run a
failed suite with a more favourable split. Tracking is verified before
spend rather than after: an online Weights \& Biases run must be
initialised and its Tinker run identifier recorded before the first paid
step, every step records loss, reward, configuration, seed, split
identity and cumulative cost, and every checkpoint is published to
Hugging Face with a unique repository or revision URL, with a failure of
any of these stopping the run fail-closed. If a receipt is missing, the
phase is blocked rather than silently shortened. At handoff the protocol
requires four separate evidence classes to be reported apart: completed
receipts; operational and provenance receipts; prospective gates and
budget projections; and admissible final results, the last existing only
if every gate passed.

That last distinction has teeth in practice, and the campaign's
disposition record shows it. The portfolio-level protocol defines a
predeclared equal-domain macro average over sixteen domains, a
conjunctive coverage gate, an integrity gate, a global improvement gate
requiring the lower bound of a one-sided 95\% interval for the portfolio
difference to exceed zero, a domain no-regression guard with a
predeclared five-percentage-point tolerance, and a safety guard bounding
any violation increase by two percentage points --- and it states that a
cost cap may not be met by dropping domains or suites, and that a suite
or domain which is missing, underpowered, contaminated or outside the
cost cap must be labelled incomplete or blocked, ``not negative and not
positive''. The portfolio training phase was never launched, so the
campaign does not report a portfolio result; the recorded outcome is a
set of per-lane evaluations of the one adapter, and the protocol's own
instruction that a partial subset is not a portfolio result and may not
be relabelled as one is honoured throughout.

Lane dispositions are recorded in a seven-state taxonomy rather than as
pass or fail: COMPLETE, for lanes with verified full-coverage results
and no open gate (E8 public, E10 benign, E11); COMPLETE\_\allowbreak{}TERMINAL\_\allowbreak{}NOTE,
for a lane complete under the native protocol with a recorded terminal
note (E14 replacement); CLOSED\_EXTERNAL, for lanes with no local or
paid path, closed by written closure records (E3, E7, E8-private,
E10-private, E12, E14-private); PENDING\_QUOTA, for lanes gated on cloud
quota resolution alone (E6, E9); AMENDMENT\_\allowbreak{}ACCEPTED\_\allowbreak{}LAUNCH\_\allowbreak{}PENDING
(E2, E13); REBUILD\_\allowbreak{}READY\_\allowbreak{}LAUNCH\_\allowbreak{}PENDING (E1 wave10, E5 successor27);
and CLOSED\_PARTIAL (E4). Those are the states of the 19 September
ledger. The three pending states were emptied on 27 September, when the
E1, E2, E5, E6, E9 and E13 replacement scopes were run to completion
(§8.4). This report assigns them COMPLETE on coverage, with two
qualifications recorded against the label: E6 still has 108 tasks
awaiting a judge, and E1's coverage spans two contracts reported
separately. The taxonomy exists so that a lane which has not been
measured is not silently absorbed into a percentage. Carried alongside
it is the project's evidence rule, stated in the ledger of record and
repeated in every results artefact: original-contract and
replacement-scope numbers are never pooled. The rule is not cosmetic.
The original-contract E1 lane and its replacement scope are different
suites with different task populations, and a table that added them
would produce a number that corresponds to no benchmark.

Decontamination is handled as a declared deterministic procedure with
its own stated limits. The check hashes complete UTF-8 fields and then
normalised fields with SHA-256; for fields of at least eight tokens it
hashes each consecutive eight-token sequence and computes exact set
Jaccard similarity for every training and evaluation pair sharing a
shingle; and it tests whether an entire normalised evaluation field is a
token-bounded contiguous substring of a training field. Against the
pinned training envelope, the per-suite inventories were 1,967 LAB-Bench
tasks over 13,113 text fields with 5,105 eligible for eight-token near
checks, 97 AgentDojo tasks over 71,966 fields with 9,382 eligible, and
4,428 Omni-MATH tasks over 13,284 fields with 9,394 eligible; the count
of near pairs at or above 0.80 was zero in all three suites, and the
maximum eight-gram Jaccard similarity was 0 for LAB-Bench and AgentDojo
and 0.13414634146341464 for Omni-MATH, below the declared 0.80
threshold, with all 1,338 candidate pairs sharing an eight-gram
examined. Short exact matches of zero to three normalised tokens were
retained and enumerated rather than dropped --- the longest examples
being seven characters in the LAB and Omni inventories and at most
seventeen raw characters in AgentDojo, such as ``Paris'' and
``Conference Room A'' --- and the audit states plainly that these
collisions do not demonstrate duplication of a task but do make an
unqualified ``zero overlap'' claim incorrect. The audit's scope limits
are recorded with equal care: it does not reconstruct which training
records were consumed: 512 is the size of the training pool recorded in
the E11 decontamination receipt
(\longtok{selected\_\allowbreak{}training\_\allowbreak{}row\_\allowbreak{}count}), from which the 40-step run
drew 80 prompts at two per step, and the transform and consumed-row
inventory that would identify them were not recorded; it establishes a
content comparison against the verified historical source envelope
rather than proving global pretraining cleanliness, semantic
disjointness, character-edit disjointness or image cleanliness; ASCII
normalisation loses mathematical punctuation and non-ASCII text; and
fields shorter than eight tokens fall outside the near-match and
containment rule. The audit itself is receipted: it re-verifies Weights
\& Biases run \texttt{withheld-id} against
\longtok{training\_\allowbreak{}identity\_\allowbreak{}verified.json}, matches the 23,733 original
source-field document hashes to the historical E11 receipt, adds 2,507
decoded nested API-Bank strings to reach 26,240 source documents, and
records 43,977,748 bytes downloaded against a 500 MiB cap, with the
checker not having enumerated the roughly 2.58 billion possible document
pairs.

The campaign records incomplete execution explicitly. Some September
lanes remained blocked by infrastructure or readiness conditions and
were completed on 27 September after changes to the execution
environment. The per-lane protocol deviations and results remain
separate (§8.4). Private account state, access routes and operational
approvals are omitted from the public derivative.

The campaign also carries its integrity accounting in the same register.
Eleven of eleven deterministic checks passed on 2026-09-19, covering
exact divisions, receipt presence, sealed-artifact hash equality and
lost-source absence probes. A judgment layer was added for
classification, faithfulness and value-ordering questions and then
partly invalidated by an external service outage: both
blocker-classification batteries returned degenerate distributions after
the incident, their receipts are retained as incident evidence, and the
per-lane blocker column therefore rests on document evidence and not on
model judgment. A design rule was banked from the healthy window and is
followed in this thesis: arithmetic and string lookups stay in code, and
model judgment is reserved for classification, faithfulness and value
ordering, because a parity probe was answered incorrectly even while the
service was healthy.

What the campaign does not claim follows directly from the design. It
does not claim improvement over any baseline anywhere, because no lane
carries a test establishing it. The one trained-versus-base comparison
is the paired small-scale rerun of Table 8.B, with 5 to 151 identical
items per lane, and its outcome is reported in Chapter 8; the campaign's
role in this thesis is to fix the evaluation protocol, the receipt
discipline and the disposition taxonomy that any future improvement
claim would have to satisfy. It does not claim that the fourteen suites
constitute a portfolio result. It does not pool original-contract with
replacement-scope figures, and it does not average across suites. Its
one trained actor means it says nothing about model comparisons. And
where the record is incomplete --- a lost execution source, an
unavailable external benchmark scope --- the incompleteness is written
into the ledger rather than resolved by inference.

\section{C1 Frozen Policy Diagnostic and Reference
Review}\label{c1-frozen-policy-diagnostic-and-reference-review}

C1 completed 4,096 fixed-order screening questions with eight
independently seeded responses per question. Eligibility required all
eight outputs to be parseable, uncapped and assigned zero exact-match
reward by frozen parser v2. After the full screen, the first 64 of 87
eligible questions received 24 additional responses each. Selection was
truncated deterministically; the 23 other eligible questions were not
followed up. The primary endpoint is whether at least one of the 24
fresh responses receives parser-defined exact-match reward one. No
question was removed or relabeled after review. {[}S1, S13{]}

The new workload contains 32,768 screening responses and 1,536 fresh
follow-up responses, totaling 34,304 new draws. A disjoint historical
smoke cohort contributes 1,024 retained draws imported once, giving
35,328 recorded unique draws. The historical smoke is an operational
quality record, not part of the primary denominator. The 512 initial
responses for the selected 64 questions are already included within
screening, and must not be counted again. {[}S1, S10{]}

The model and tokenizer were Qwen/Qwen3-8B at revision
\longtok{b968826d9c46dd6066d109eabc6255188de91218}, with BF16 and
thinking disabled. The dataset was openai/gsm8k, main configuration,
train split, revision \longtok{740312add88f781978c0658806c59bc2815b9866}.
Generation used temperature 0.7, top-p 0.8, top-k 20, min-p 0, a
1,024-token generation cap and 4,096-token context. The recorded
execution used one NVIDIA A100-SXM4-80GB, vLLM 0.11.0, PyTorch 2.8.0 and
Transformers 4.57.1, with eager execution enabled and prefix caching
disabled. Fixed request seeds do not imply bitwise reproducibility
across hardware or allocations. {[}S12, S13{]}

Screening had 504 parser failures and 19 cap hits among 32,768 outputs.
By design the selected initial groups had neither. Follow-up had ten
parser failures and four cap hits among 1,536 outputs. Those quality
counts do not certify semantic correctness. The primary denominator
remains 64 regardless of later quality flags. Descriptive prefixes of
eight, sixteen and twenty-four fresh draws yield 9/64, 10/64 and 11/64
parser events. They are nested observations, not independent comparisons
or evidence of an optimal rollout budget. {[}S1{]}

The Wilson interval uses a question-level binomial working model within
the selected cohort. It is not design-based or population coverage: both
outcome-dependent eligibility and fixed-order truncation restrict the
estimand. The observed 87/4,096 eligibility count similarly describes
this fixed screen, not benchmark-wide mathematical failure prevalence.

The first review solved the questions before comparing reference
answers, then inspected the required model outputs. A second reference
review was separately masked from the first judgments and model outputs.
Final reconciliation compared the exact question, reference and
annotations, retaining disagreements rather than forcing consensus. All
stages were assistant-assisted. Separate masking and repeat review do
not provide independent human validation or rule out shared model
biases. {[}S2, S6{]}

The first output audit covers all 64 cases, all 512 initial responses
and 143 required fresh responses: 655 responses in total, including all
38 primary-success outputs. The final mechanical coverage receipt
verifies those identities and stored logic; it is not an independent
mathematical adjudication. Six initially missing reasoning-caveat fields
were clarified in a separately attributed supplement without altering
the original fields, scores or verdicts. No coverage or annotation
defect remains unresolved in the final receipt. {[}S7{]}

For the eleven event questions, reconciliation read all 88 initial
outputs and all 38 primary-success fresh outputs in full. All 126
explicit final-number extractions are faithful and uncapped. Two
unit-sensitive non-events add sixteen inspected initial outputs. This is
not a claim to have validated reasoning in every one of the 1,536 fresh
responses. The first audit identified five frozen-v2 extraction false
negatives and no extraction false positives within its reviewed set;
these selected-case observations do not establish population parser
precision. {[}S2, S3, S5, S7{]}

The conservative witness rule requires a reference value supported in
both rounds without material ambiguity, eight faithful uncapped initial
claims independently judged wrong, and at least one faithful uncapped
fresh primary-success claim independently judged numerically correct.
Sound final arithmetic and sound visible reasoning are recorded
separately. No post-review filtering changes the 64-question
denominator.

Source labels S1-S13 are resolved in Appendix G. This design is distinct
from the historical V2 diagnostic.



\chapter{Implementation Details}\label{implementation-details}

This chapter describes what was built. Its organising claim is that the
engineering contribution of the project is a \textbf{dispatch and
evidence layer} rather than a training algorithm: the algorithmic
content of GRPO is inherited from the literature, but the work of making
six reinforcement-learning frameworks and six compute backends
addressable through one Tinker-style interface --- and of making every
number that comes out of them traceable to a sealed, hash-pinned
artefact --- is where the implementation effort sits. The chapter
proceeds from the abstraction (the adapter architecture), through the
concrete training loop, to the two engineering episodes that most shaped
the codebase (the parallel logging defect and the Colab white-box path),
and finally to the evaluation harness and the recovery
re-implementations undertaken after an execution-source loss in the
final phase of the project (Figure~\ref{fig:receipts}).


\fig{fig_receipt_pipeline}
{The evidence chain from sealed request through execution and native grading to a deterministic ledger check}
{The evidence chain from sealed request through execution and native grading to a deterministic ledger check. A break anywhere in the chain means no number is reported for that lane. Upper row: the chain; the 11/11 PASS badge is the 2026-09-19 deterministic ledger check of eleven named claims, not of every figure (\longtok{LEDGER\_\allowbreak{}CODE\_\allowbreak{}CHECK\_\allowbreak{}2026-09-19.json}). Lower row: the E1 wave-10 recovery, where the lost execution source stops the chain.}
{fig:receipts}



\section{Two Rosters, One Dispatch
Surface}\label{two-rosters-one-dispatch-surface}

A reader approaching the project's comparison material will encounter
two different framework rosters, and conflating them is the single most
likely misreading of the implementation. The internal reports' shared
methods note states the distinction explicitly: the cross-launcher
LLM-RL scaffold referenced in the framework-gap discussion is a
seven-library roster --- Tinker, TRL, SkyRL, veRL, OpenRLHF, Atropos,
and HF reference launchers --- that targets LLM post-training launchers
specifically, while the cross-RL-library evidence reported elsewhere
uses a different roster. Critically for this chapter, the same note
records that ``only Tinker and TRL produced completed runs in this
release; veRL and OpenRLHF entries in the
\texttt{framework\_comparison.json} artefact are dry-run placeholders''.
The veRL and OpenRLHF entries in that artefact carry peak and
last-ten-average reward values, but they are the output of the harness
resolving a launch plan, not the output of a training run, and they must
never be pooled with real results or read as framework comparison
evidence.

Atropos occupies a different position again. Rather than being a cell in
the dispatch matrix, it is integrated through a dedicated package under
\texttt{platform\_tinker/atropos/}, a \texttt{tinker\_atropos} module
containing \texttt{config.py}, \texttt{api.py}, \texttt{trainer.py},
\texttt{types.py}, \texttt{dataset.py}, \texttt{stats\_utils.py}, and an
\texttt{environments/} directory with per-task environment servers. Its
integration model is the Atropos managed-server pattern: a
three-terminal workflow in which \texttt{run-api} starts the Atropos
API, \texttt{launch\_training.py\ -\/-config\ configs/default.yaml}
starts Tinker training against a \longtok{TINKER\_\allowbreak{}API\_\allowbreak{}KEY}, and an
environment server --- a built-in one or any external Atropos
environment pointed at the same config file --- serves rollouts. The
README's quickstart example runs a fifty-step training run on
Llama-3.1-8B-Instruct against a GSM8K environment, and the environment
contract is that the \texttt{env} and \texttt{openai} sections are
consumed by the environment server while the \texttt{tinker} section is
consumed by the trainer and ignored by the environment. This separation
is what allows an arbitrary existing Atropos environment to be reused
against the Tinker backend without modification to the environment.

\section{The Adapter Architecture: What versus
Where}\label{the-adapter-architecture-what-versus-where}

The dispatch surface is built on one orthogonal decomposition, stated in
the backends module docstring: ``A \texttt{Backend} knows \textbf{where}
to run one experiment cell (local GPU, Modal, Colab, vast.ai, GCP, HF
Spaces); the framework dimension (trl/tinker/verl/openrlhf/skyrl) is the
\textbf{what}. \texttt{plan()} resolves a cell into a concrete
\texttt{LaunchPlan} without spending compute, so \texttt{-\/-dry-run}
and the matrix test work uniformly''. The \texttt{LaunchPlan} dataclass
carries \texttt{backend}, \texttt{framework}, \texttt{command},
\texttt{driver\_file}, \texttt{output}, \texttt{env}, and
\texttt{notes}, with a \texttt{format()} method for printing; the
default \texttt{\_execute} runs
\texttt{subprocess.run(plan.command,\allowbreak{}\ shell=\allowbreak{}True,\allowbreak{}\ check=\allowbreak{}True)}. Because
plan resolution is separated from execution, the entire matrix can be
validated without a GPU, an API key, or a dollar of spend --- a property
that the dry-run placeholders discussed above depend on, and which is
also the reason those placeholders must be labelled as such (Figure~\ref{fig:arch}).


\fig{fig_architecture}
{Six RL frameworks behind one Tinker-style API}
{Six RL frameworks behind one Tinker-style API. Five of the six are in the dispatch matrix (5 frameworks $\times$ 6 compute backends); only two executed (Tinker, TRL: last-10 mean training reward on Qwen/Qwen3-8B GSM8K, from \texttt{framework\_comparison.json}), two are dry-run launch plans and two were not executed. One configuration drives every framework, but the executed comparison bundled a different base checkpoint (Chapter~7), so outcome differences are reported with the stack rather than attributed to it.}
{fig:arch}



The matrix itself ships as data.
\longtok{platform\_local/unified/matrix.json} declares schema
\longtok{tinkerrl-framework-backend-matrix-v1} with thirty cells --- five
frameworks (trl, tinker, verl, openrlhf, skyrl) crossed with six
backends (local, modal, colab, vast, gcp, hfspaces), five cells per row
and per column --- and a \texttt{canonical\_experiment} block that pins
the single configuration every cell is meant to instantiate:
\texttt{Qwen/Qwen3-8B} on GSM8K with GRPO, thirty training steps,
\texttt{num\_generations} 8, per-device batch size 2, gradient
accumulation 8, learning rate \texttt{1e-06}, LoRA rank 16 and alpha 32
over all linear modules, \texttt{beta} 0.0, \texttt{epsilon} 0.2,
maximum completion length 1024, and seed 211. Each cell is a concrete
plan naming a framework, a backend, the launch command, the driver file
and an output location; for example, the \texttt{trl}/\texttt{local}
cell runs
\texttt{python\ -m\ platform\_local.unified\ -\/-backend\ local\ -\/-framework\ trl}
and dispatches in-process to
\longtok{platform\_local/trl\_integrations/trainer.py}.

The launcher implements the dispatch. \texttt{UnifiedLauncher} holds
three registries: \texttt{FRAMEWORKS} mapping \texttt{skyrl},
\texttt{tinker}, \texttt{verl}, \texttt{openrlhf}, and \texttt{trl} to
human-readable labels; \texttt{BACKENDS} mapping \texttt{local},
\texttt{modal}, \texttt{colab}, \texttt{vast}, \texttt{gcp}, and
\texttt{hfspaces} to labels; and \texttt{ALGORITHMS} covering
\texttt{grpo}, \texttt{ppo}, \texttt{reinforce}, \texttt{dapo}, and
\texttt{dpo}. Dispatch is outer-on-backend and inner-on-framework, which
is why the per-framework handlers can each be written against the
framework's own idioms. \texttt{\_run\_tinker} shells to
\texttt{python\ -\allowbreak{}m\ platform\_\allowbreak{}tinker.tinkerrl.grpo\_\allowbreak{}cli\ -\allowbreak{}\/\allowbreak{}-\allowbreak{}preset\ \textless{}spec.task\textgreater{}};
\texttt{\_run\_openrlhf} delegates to
\texttt{platform\_\allowbreak{}modal/\allowbreak{}openrlhf/\allowbreak{}trainer.py:\allowbreak{}run\_\allowbreak{}\allowbreak{}openrlhf\_\allowbreak{}\allowbreak{}training(\{\}\allowbreak{},\allowbreak{}\ output\_\allowbreak{}dir=\allowbreak{}"\allowbreak{}./\allowbreak{}checkpoints/\allowbreak{}openrlhf"\allowbreak{})},
which drives
\texttt{openrlhf.cli.train\_\allowbreak{}\allowbreak{}ppo\_\allowbreak{}\allowbreak{}ray\ -\allowbreak{}\/\allowbreak{}-\allowbreak{}advantage\_\allowbreak{}estimator\ group\_\allowbreak{}norm};
\texttt{\_run\_trl} generates a resumable GRPO training script through
\longtok{trl\_integrations.trainer.generate\_\allowbreak{}trl\_\allowbreak{}train\_\allowbreak{}script} into a
temporary directory and executes it; \texttt{\_run\_verl} imports
\texttt{verl.config} and \texttt{verl.trainer} directly; and
\texttt{\_run\_skyrl} requires a \texttt{SKYRL\_CHECKOUT} environment
variable and emits
\texttt{uv\ run\ -\/-extra\ gpu\ -\/-extra\ tinker\ -m\ skyrl.tinker.api}.
For the \texttt{local} and \texttt{colab} backends the launcher runs
\texttt{dispatch\_framework()} in-process, a path the code comments
identify as the one that breaks the Colab entry's self-recursion --- an
implementation detail worth recording because it indicates that the
backend and framework dimensions are not perfectly orthogonal in
practice, and the launcher carries framework-specific knowledge
precisely at the points where they intersect.

Fairness across frameworks is addressed by configuration auditing rather
than by assertion. The
\longtok{platform\_hybrid/experiments/framework\_\allowbreak{}config\_\allowbreak{}dumps/}
supplement ships one exact hyperparameter dump per framework for the
canonical Qwen3-8B + GSM8K matched-configuration run, and the
accompanying README recomputes the field partition from those dumps: of
44 tracked leaf fields, 25 are identical across all four frameworks
(LoRA rank, alpha, target modules and dropout; GRPO group size, rollouts
per group and max new tokens; sampling temperature, top-p, top-k and
do-sample; optimizer name, learning rate, weight decay, betas and
epsilon; reward type; and the task, seed and source-git-sha metadata),
while 11 are Tinker-managed --- serialised as
\texttt{null\ \ \#\ managed\_\allowbreak{}\allowbreak{}by\_\allowbreak{}\allowbreak{}tinker} in the Tinker dump --- covering
\texttt{kl.beta}, \texttt{kl.surrogate}, importance-sampling
granularity, reference-model offload and recompute, tokens per optimizer
step, reward broadcast precision, gradient checkpointing, micro-batch
size, gradient accumulation, and tensor-parallel size. The same README
records two disclaimers that matter more than the identical count:
\texttt{model} and \texttt{model\_sha256} are \emph{not} identical,
because the Tinker run uses \texttt{Qwen/Qwen3-8B-Base} while TRL, veRL
and OpenRLHF use the instruction-tuned \texttt{Qwen/Qwen3-8B}, and
\texttt{model\_sha256} is a placeholder rather than a verified checksum.
Remaining fields are framework-specific and documented rather than
harmonised --- KL beta is 0.04 for TRL and veRL against 0.02 for
OpenRLHF, importance-sampling granularity is token-level for TRL and
veRL against sequence-level for OpenRLHF, and reference-model placement
differs across all four (\texttt{on\_device}, \texttt{managed},
\texttt{separate\_process}, and \texttt{sharded\_ray}). The appendix's
original ``byte-identical'' framing is retracted on the strength of this
partition; the accurate statement is that a majority of tracked fields
are matched, a known Base-versus-Instruct confound survives, and eleven
fields are not under user control at the Tinker API level at all.

\section{The Tinker GRPO Loop}\label{the-tinker-grpo-loop}

The reference implementation of the training loop lives in
\longtok{platform\_tinker/tinkerrl/grpo.py} and is exposed through
\texttt{grpo\_cli.py}. The loop follows the standard four-stage shape
recorded in the Phase-1 report under ``The Tinker GRPO Loop'': sample a
group of \texttt{G} completions per prompt, compute group-relative
advantages, form the loss, and take an optimiser step. Advantages are
normalised within the group to mean zero and unit standard deviation,
using the population σ (\texttt{normalize\_rewards}, ε = 10⁻⁸). The loss
is not the clipped, KL-regularised surrogate of the original GRPO paper.
It is REINFORCE with a group-mean baseline,

\[L(\theta) = -\frac{1}{N} \sum_i A_i \sum_t \log \pi_\theta(y_{i,t} \mid x, y_{i,<t})\]

averaged over the N completions in the batch, with no importance ratio,
no clipping, no KL term and no reference policy. The sampler is
refreshed only at checkpoint boundaries (every \texttt{save\_every}
steps), so rollouts between refreshes come from slightly stale weights
with no importance correction. The implementation uses the Tinker SDK's
\texttt{ServiceClient} with \longtok{create\_\allowbreak{}lora\_\allowbreak{}training\_\allowbreak{}client},
\longtok{forward\_\allowbreak{}backward\_\allowbreak{}custom} with
\longtok{loss\_\allowbreak{}type\_\allowbreak{}input="logprobs"} so that per-token
log-probabilities flow into a custom loss, \texttt{optim\_step} with
\texttt{AdamParams}, and \texttt{sample()} for rollouts; resumption uses
\longtok{create\_\allowbreak{}training\_\allowbreak{}client\_\allowbreak{}from\_\allowbreak{}state\_\allowbreak{}with\_\allowbreak{}optimizer}.

The single most consequential property of this loop for the project's
research questions is the degenerate case. When the group's reward
standard deviation \texttt{σ\_g} is zero, every advantage \texttt{A\_i}
is zero, and because the loss has no KL or auxiliary term, the group
contributes exactly zero gradient. This is the mechanism behind the
zero-variance-fraction diagnostic that runs through the project; it is
recorded in the Phase-1 report as the observation that motivated the
collapse analysis, and it is the reason \texttt{zero\_loss} and
\texttt{zero\_reward} are first-class outputs of the training CLI rather
than derived quantities. Zero reward is counted as an average reward
equal to zero; zero loss is counted as a step loss below \texttt{1e-6}.

The CLI exposes eight presets --- \texttt{tooluse\_synth},
\texttt{tooluse\_baseline}, \texttt{tooluse\_heldout},
\texttt{tooluse\_xlam}, \texttt{gsm8k}, \texttt{math100},
\texttt{pavlov\_xlam}, and \texttt{pavlov\_portfolio} --- each resolving
a dataset factory and a reward function through a
\texttt{DATASET\_FACTORIES} registry and a \texttt{REWARD\_MAP}
(\texttt{ExactMathReward} for GSM8K, \texttt{StrictToolCallReward} for
\texttt{pavlov\_xlam}, \texttt{PavlovNonXLAMReward} for
\texttt{pavlov\_portfolio}). Tracking and campaign gates are enforced
before a run starts, through \texttt{config.validate\_tracking()} and
\longtok{config.validate\_\allowbreak{}campaign\_\allowbreak{}gate()}, and the CLI refuses to
launch without a \longtok{TINKER\_\allowbreak{}API\_\allowbreak{}KEY}. Each seed's result is
printed as a structured record containing \texttt{run\_id},
\texttt{sampler}, \texttt{avg\_first5}, \texttt{avg\_last10},
\texttt{peak\_reward}, \texttt{zero\_loss}, \texttt{zero\_reward},
\texttt{heldout\_reward}, and \texttt{reward\_trace}, so that a run's
evidence is emitted at the point of execution rather than reconstructed
afterwards. The eight presets replaced an earlier collection of
standalone \texttt{grpo\_*.py} scripts, with \texttt{legacy\_main}
retained to translate the old flag surface.

\section{The Parallel Runner and a Real Logging
Defect}\label{the-parallel-runner-and-a-real-logging-defect}

Sweep execution is handled by \texttt{parallel\_sweep.py} and
\texttt{campaign.py}, with the runner's design shaped directly by a
defect found during development. The defect is documented in the
group-size sweep findings and is worth quoting in full, because it is
the clearest illustration in the project of why logging infrastructure
deserves the same scrutiny as training code:

\begin{quote}
The threaded runner trains all G in parallel correctly, BUT
\texttt{wandb} uses a thread-unsafe global \texttt{run} --- with 4
threads calling \texttt{wandb.init(reinit=True)}, the G8/G16
\textbf{summary} writes failed (``Run is finished'') after another
thread finished its run. Per-step \texttt{wandb.log} mostly survived;
the summary dict didn't. \textbf{Fix for next parallel batch:} isolate
each run in its own PROCESS (multiprocessing / subprocess per config),
not a thread --- or use \texttt{wandb-core} service mode with explicit
per-run objects. Training itself was unaffected (all step data is in the
run log).
\end{quote}

Two aspects of this record are instructive. First, the failure was
silent with respect to training: gradients were correct, steps
completed, and per-step metrics were written. Only the run summary ---
the aggregated dictionary from which a results table would be built ---
was lost, which is precisely the class of failure that corrupts
downstream reporting while leaving the primary artefact intact. Second,
the remedy adopted was structural rather than retry-based: isolating
each run in its own process removes the shared mutable global entirely,
whereas retrying the summary write would have left the race in place.
The finding also notes the alternative of \texttt{wandb-core} service
mode with explicit per-run objects, which would preserve in-process
execution at the cost of depending on service-mode semantics that the
project chose not to rely on.

The same findings file documents the group-size sweep that exposed the
defect, and its numbers are reported here with the qualifications the
source attaches to them. The sweep ran \texttt{G\ ∈\ \{2,\ 4,\ 8,\ 16\}}
on Qwen3.5-4B for six steps on GSM8K at seed 0 under the W\&B group
\texttt{p3-groupsize-sweep}, with total token counts of 9,663, 19,964,
34,084, and 68,055 respectively. At \texttt{G\ =\ 2} the run showed a
50\% batch-level collapse with a \texttt{+0.000} gain; at
\texttt{G\ =\ 4} a 17\% collapse with a \texttt{+0.125} gain. The source
flags the \texttt{G\ =\ 4} gain explicitly as supported by one example
at n = 8 on a single seed and therefore \emph{suggestive rather than
proven}, and records that the independent verification of the sweep
returned the verdict SUSPECT. These are not results this thesis claims;
they are the observations that motivated the group-size investigation
and that the later controlled work supersedes, and they are reported in
the register in which they were recorded.

\section{The Colab White-Box Path (Per-Layer Gradient
Profile)}\label{the-colab-white-box-path-per-layer-gradient-profile}

The per-layer gradient profile --- numbered P1 in the Phase-1 series,
and distinct from the P1 cross-scale study of Chapter 6 --- asks a
question the managed Tinker backend cannot answer. Per-layer gradient
norms require access to the model's internals during the backward pass,
and the managed service exposes no such interface. The implementation
therefore leaves the managed path entirely and runs on a Colab L4 GPU in
\longtok{platform\_hybrid/experiments/openings/p1\_\allowbreak{}layer\_\allowbreak{}profile.py},
using Qwen2.5-1.5B-Instruct with LoRA adapters on the q/k/v/o
projections, five GRPO-style advantage-weighted steps, and per-layer
LoRA gradient norms recorded at each step. The white-box path is thus a
deliberate architectural exception: the project accepts the loss of
managed-scale training in exchange for gradient visibility, and confines
the exception to a question that only gradient visibility can address.

The first-pass result appeared strong. Across 28 layers, the top-k layer
set identified at step 1 overlapped the final top-k set perfectly, with
a measured overlap of 1.0; the top 25\% of layers accounted for a 0.476
share of gradient mass; layer 0 carried a gradient norm of 63 against
6--17 elsewhere, and a secondary mid-late band appeared at layers
17--23. A scaled re-test disciplined this result. The re-run used
\longtok{p1\_\allowbreak{}layer\_\allowbreak{}profile\_\allowbreak{}scaled.py} on Qwen2.5-3B-Instruct, with 24
real GSM8K problems, ten steps, \texttt{G\ =\ 4}, 36 layers, and seeds
\{0, 1\}. The step-1-to-final top-k overlap fell to 0.111 --- chance
level --- and the top-25\% concentration measured 0.39 with a standard
deviation of 0.013. The source's verdict is unambiguous: the original
predictability finding ``was an artifact of the toy
1.5B/hardcoded-arithmetic setup''. A later four-seed pass at the same
configuration (seeds \{0, 1, 2, 3\}) supersedes the two-seed figure and
gives a mean overlap of 0.083 ± 0.048, again approximately chance, with
concentration 0.39.

A companion freeze test asked whether the layer profile could be
exploited: if a small band of layers dominates the gradient, freezing
the rest should preserve the benefit at lower cost. It did not. The full
28-layer configuration trained 2.18M parameters and produced a mean
held-out gain of \texttt{+0.0625} (seed 0 \texttt{+0.125}, seed 1 0.0).
Freezing to a hot band comprising layer 0 plus layers 12--27 --- 17
layers, 1.32M parameters, or 0.61 times the full count --- produced
\texttt{+0.000}. The source concludes that ``aggressive freezing loses
the benefit\ldots{} The freeze lever does not work at this scale'', and
records that the comparison is underpowered at n = 8 held-out examples
over 2 seeds. Both P1 outcomes are negatives: the apparent structure did
not survive scaling, and the intervention it suggested did not survive
testing. They are reported as findings rather than suppressed, because a
layer-importance claim that collapses under a scaled re-test is exactly
the kind of result the project's evidence discipline exists to surface.

\section{Software Modules}\label{software-modules}

The module inventory recorded in the Phase-1 report remains the accurate
map of the codebase. The unified launcher,
\longtok{platform\_local/unified/launcher.py}, maps one configuration
onto a chosen backend; \longtok{platform\_tinker/tinkerrl/grpo.py} is the
reference GRPO loop against the managed backend;
\longtok{platform\_local/trl\_integrations/} and the \texttt{verl/}
integration hold the per-framework adapters, with
\texttt{platform\_modal/openrlhf/} and \texttt{platform\_hybrid/skyrl/}
completing the roster, and \texttt{platform\_tinker/atropos/} carrying
the Atropos integration described above. Analysis code is separated from
training code by intent: \longtok{p2\_\allowbreak{}collapse\_\allowbreak{}analysis.py} recomputes
zero-variance fraction and collapse statistics from stored reward
tensors, so the diagnostic is reproducible from evidence rather than
from a live run; \longtok{p1\_\allowbreak{}layer\_\allowbreak{}profile.py} implements the white-box
gradient probe; \texttt{parallel\_sweep.py} and \texttt{campaign.py}
drive sweeps; and \texttt{p8\_detector.py} is an exploratory
method-trace classifier evaluated by row-wise cross-validated AUROC.

The p8 detector deserves an explicit caveat in an implementation
chapter. It is an exploratory artefact evaluated by row-wise
cross-validation, and its AUROC is a measure of separability on the
traces it was given, not a benchmark result against an external suite;
it should not be read alongside the E-lane scores reported later in this
thesis. Beyond these modules, the \texttt{platform\_local/unified/}
package carries \texttt{canonical.py} for the canonical configuration
and \texttt{peft\_utils.py} for adapter handling, and the backends
package implements one module per backend (\texttt{local.py},
\texttt{modal.py}, \texttt{colab.py}, \texttt{vast.py}, \texttt{gcp.py},
\texttt{hfspaces.py}) behind the common \texttt{base.py} interface.

\section{The E1--E14 Evaluation
Harness}\label{the-e1e14-evaluation-harness}

Evaluation is governed by a completion gate,
\longtok{zvf-program/flagship/e1\_\allowbreak{}e14\_\allowbreak{}completion\_\allowbreak{}gate.py}, schema
\texttt{e1-e14-completion-gate-v1}. The gate maintains five evidence
classes --- \texttt{COMPLETE\_EXACT}, \texttt{PARTIAL\_EXACT},
\texttt{PARTIAL\_RECOVERY}, \texttt{BLOCKED\_EXTERNAL}, and
\longtok{LOCAL\_\allowbreak{}SETUP\_\allowbreak{}READY\_\allowbreak{}PROVIDER\_\allowbreak{}INPUT\_\allowbreak{}REQUIRED} --- and maps
lane source statuses onto them, with unknown states failing closed
rather than defaulting to a classification. Its \texttt{LANE\_SPECS}
table binds each lane to a replacement scope, a task count, and the
runner or adapter that produces its evidence: E2 FrontierSWE at 17 tasks
(\longtok{pavlov\_\allowbreak{}frontier\_\allowbreak{}swe\_\allowbreak{}eval\_\allowbreak{}adapter.py},
\longtok{frontier\_\allowbreak{}swe\_\allowbreak{}eval.py}), E3 SDAB at 80, E4 BankerToolBench at
100 (\longtok{run\_\allowbreak{}harbor\_\allowbreak{}via\_\allowbreak{}tinker\_\allowbreak{}bridge.py}), E5 APEX-Agents at
480, E7 BinaryAudit at 46 (\longtok{modal\_\allowbreak{}non\_\allowbreak{}e11\_\allowbreak{}readiness.py}), E9
MLE-bench at 75 (\longtok{e9\_\allowbreak{}mle\_\allowbreak{}bench\_\allowbreak{}streaming.py}), E12 AppBench at
6, and lanes E6, E8, E10, E13 and E14 with their respective runners. The
gate's own header states its scope precisely: it neither invokes an
adapter nor accepts a package as a benchmark result (Figure~\ref{fig:harness}).


\fig{fig_harness}
{Evaluation harness internals}
{Evaluation harness internals. The harness is fail-closed: when a precondition cannot be satisfied the lane records \textsc{blocked} rather than substituting an adjacent benchmark. Outcomes and the lane tally are taken from the 2026-09-19 ledger of record (\texttt{Pending\_Experiments.md}), updated with the six replacement scopes scored on 2026-09-27. Outcome A separates original-contract scores from replacement-scope results and gives coverage and score separately; Outcome B lists original-scope suites that are closed external, then the one lane whose scope is not settled. E14 is reported as $2271/4428 = 51.29\%$ with the two unjudged rows counted as failures (the official scorer reports 51.31\% over 4426 judged rows).}
{fig:harness}



The results the gate classifies are recorded in
\longtok{outputs/E1\_\allowbreak{}E14\_\allowbreak{}FINAL\_\allowbreak{}RESULTS\_\allowbreak{}2026-09-19.md}, verified by
eleven named deterministic checks on 2026-09-19 (scope in Appendix B)
and bound to surviving receipts. The single Tinker-trained actor is base
\texttt{Qwen/Qwen3.6-35B-A3B} at commit \texttt{995ad96e…} with adapter
\texttt{campaign-adapter} at commit \texttt{withheld-adapter-revision…},
in bfloat16. The document's opening constraint governs how everything
below may be read: original-contract and replacement-scope numbers are
never pooled, and per the project evidence rule no cross-suite average
is formed at any point.

The per-lane results the gate classified are reported once, in Chapter 8
(Tables 8.A--8.C), with each lane's scope and terminal state in Table
C.3; they are not repeated here.

One further result from this phase belongs in the implementation record
because of what it demonstrates about the harness. Under separate
authorization, the E4 BankerToolBench suite was re-run from scratch on
the base model through the Modal bridge's \texttt{/v1/responses} path at
\texttt{-\/-n-concurrent\ 2}. All 100 of 100 trials completed, with one
\texttt{NonZeroAgentExitCodeError}, and the mean reward was 0.0 ---
every trial in the zero bucket. The earlier report attributes the
mechanism to agents terminating after 2--14 steps (median about 6)
without producing output in \texttt{deliverables/}, causing the verifier
to short-circuit to zero on a ``No deliverables found'' condition, and
that 27 or more trajectories ended in role-token degeneration of the
form \texttt{user\ user\ assistant\ …} repeated. This bridge path set no
stop sequences, which produces exactly this degeneration, but the
available comparison does not identify the bridge contribution to the
full zero (Section 8.5) of the cap over 3h29m, against 16.35M prompt and
2.18M completion tokens. Reporting a mean reward of 0.0 with its
mechanism attached, rather than reporting a low score and moving on, is
the standard this chapter's evidence rules are meant to enforce.

\section{The Bridge and Sampler Path}\label{the-bridge-and-sampler-path}

Runners reach a Tinker-trained model through
\longtok{zvf-program/flagship/modal\_\allowbreak{}tinker\_\allowbreak{}openai\_\allowbreak{}bridge.py}, a
Modal-deployed OpenAI-compatible bridge whose own header states that
``the deployed bridge is infrastructure, not a benchmark result. Its
default paid-inference budget is zero''. The bridge declares app name
\texttt{campaign-serving-bridge}, model alias \texttt{campaign-actor}
over model id \texttt{Qwen/Qwen3.6-35B-A3B}, and pins its sampler at
{[}operational identifier withheld{]}, with the Hugging Face side pinned
by commit \texttt{withheld-adapter-revision}. Runtime limits are 16,384
maximum completion tokens against a 65,536-token model context, with an
8,192-token context prefix reserved; \longtok{DEFAULT\_\allowbreak{}MAX\_\allowbreak{}USD} is
\texttt{Decimal("0.00")}, so every paid path must be opted into
explicitly at launch. The pinned dependency set is \texttt{fastapi}
0.116.1, \texttt{huggingface-hub} 1.27.0, \texttt{tinker} 0.30.0,
\texttt{transformers} 5.5.4, and \texttt{wandb} 0.21.0.

The architectural point of the bridge is that it converts a Tinker
sampler into an OpenAI-compatible endpoint. A harness written against
the OpenAI chat-completions API --- which most agentic benchmark
harnesses are --- can therefore be pointed at a Tinker-trained actor
without modification, and the E4 rerun described above used exactly this
path in its \texttt{/v1/responses} form. The same pattern appears in the
actor-side configuration of the E1 recovery wave, where the actor is
addressed as an OpenAI-compatible chat-completions service at
\texttt{http://127.0.0.1:18015/v1} (container port 8000) under model
name \texttt{campaign-actor-bf16}, with \texttt{max\_tokens} 8192,
temperature 0, seed 809, top-p 0.95, and
\texttt{chat\_\allowbreak{}\allowbreak{}template\_\allowbreak{}\allowbreak{}kwargs\ \{"\allowbreak{}enable\_\allowbreak{}thinking"\allowbreak{}:\allowbreak{}\ false\}}. The
bridge and the local actor endpoint are two deployments of the same
interface contract, and keeping that contract stable across deployment
modes is what allows one harness to serve both.

\section{\texorpdfstring{Recovery Re-implementations under
\texttt{zvf-program/}}{Recovery Re-implementations under zvf-program/}}\label{recovery-re-implementations-under-zvf-program}

In the final phase of the project an execution-source loss --- deletion
of the historical temporary execution workspace directory --- removed
the finish-era recovery and build code, confined in scope to the E1
wave10 controller, the E9 restored dependencies, and the E6 validator,
while the replacement-lane native runners survived under
\texttt{zvf-program/flagship/}. The project's lost-source policy is that
sealed requests, validation receipts and test logs pin any faithful
re-implementation. Four such re-implementations are worth describing,
because together they show what the evidence discipline looks like when
it is expressed as code rather than as convention. None of them produced
a benchmark number; their value is that their safety properties are
testable offline.

\textbf{The four recovery re-implementations.}

\begingroup\scriptsize\setlength{\tabcolsep}{3pt}
\begin{longtable}[]{@{}
  >{\raggedright\arraybackslash\hspace{0pt}}p{(\linewidth - 10\tabcolsep) * \real{0.2000}}
  >{\raggedright\arraybackslash\hspace{0pt}}p{(\linewidth - 10\tabcolsep) * \real{0.2000}}
  >{\raggedright\arraybackslash\hspace{0pt}}p{(\linewidth - 10\tabcolsep) * \real{0.2000}}
  >{\raggedright\arraybackslash\hspace{0pt}}p{(\linewidth - 10\tabcolsep) * \real{0.2000}}
  >{\raggedright\arraybackslash\hspace{0pt}}p{(\linewidth - 10\tabcolsep) * \real{0.2000}}@{}}
\caption{The four recovery re-implementations. Schedule offsets, hashes,
reservation identifiers and HMAC domain strings are recorded in the
cited source files rather than repeated here.}\tabularnewline
\toprule\noalign{}
\begin{minipage}[b]{\linewidth}\raggedright
Artefact
\end{minipage} & \begin{minipage}[b]{\linewidth}\raggedright
Scope
\end{minipage} & \begin{minipage}[b]{\linewidth}\raggedright
Safety properties enforced in code
\end{minipage} & \begin{minipage}[b]{\linewidth}\raggedright
Offline tests
\end{minipage} & \begin{minipage}[b]{\linewidth}\raggedright
Status
\end{minipage} \\
\midrule\noalign{}
\endfirsthead
\toprule\noalign{}
\begin{minipage}[b]{\linewidth}\raggedright
Artefact
\end{minipage} & \begin{minipage}[b]{\linewidth}\raggedright
Scope
\end{minipage} & \begin{minipage}[b]{\linewidth}\raggedright
Safety properties enforced in code
\end{minipage} & \begin{minipage}[b]{\linewidth}\raggedright
Offline tests
\end{minipage} & \begin{minipage}[b]{\linewidth}\raggedright
Status
\end{minipage} \\
\midrule\noalign{}
\endhead
\bottomrule\noalign{}
\endlastfoot
E1 wave10 (\longtok{zvf-program/e1\_wave10/}) & Rebuilds the sealed v6
controller for 16 SWE-bench Multilingual tasks (9 gson, 5 terraform, 2
immutable-js) & Acts only on the sealed request hash and its embedded
task list; requests built only through the lane's native request
builder; v6 process-group contract with unconditional finalisation;
\longtok{GENERATION\_FAILED} reserved for completed responses without a
valid diff; write-once receipts; no retry, resume or replay & 10/10
(\longtok{test\_\allowbreak{}process\_\allowbreak{}group.py} 4/4,
\longtok{test\_\allowbreak{}cleanup\_\allowbreak{}regressions.py} 6/6), matching the lost suite &
Reservation recorded \longtok{RESERVED\_\allowbreak{}NOT\_\allowbreak{}DISPATCHED} \\
E2 driver (\longtok{zvf-program/e2\_core/driver.py}) & Schedules the
CORE-Bench target-helper GPU chain on GCP (one \longtok{n1-standard-4}
with a T4) & Validation-only by default; refuses \longtok{-\/-execute}
without the required execution-readiness evidence; exact early deletes;
no retry & 12/12 (schedule arithmetic, deletes, cutoff,
\longtok{UNVERIFIED}/\texttt{LATE} reporting, IAM gate) & Unexecuted
design artifact \\
E5 successor27 (\longtok{zvf-program/e5\_successor27/controller.py}) &
Rebuilds the expired sealed packet for 27 never-started Tau3 banking
tasks (97 in scope) & Pinned model revision and fresh run identifiers,
with the earlier namespace barred; concurrency 3; fixed host, actor and
dispatch deadlines; offline \texttt{verify}, \texttt{plan} and
\texttt{gate} subcommands & 52/52; \texttt{verify} and \texttt{plan}
return \longtok{PLAN\_\allowbreak{}READY\_\allowbreak{}NOT\_\allowbreak{}EXECUTED}, \texttt{gate} and
\texttt{launch} refuse & Needs a bound authorisation receipt and a
native runtime adapter \\
E13 supervisor (\longtok{zvf-program/e13\_balrog/hosted\_supervisor.py})
& Hosted supervisor for the BALROG lane; design artefact, not deployed &
Absolute deadline D = t0 + 3600 s that is never rebased; short
HMAC-authenticated leases that reject replays; renewal and receipt
signatures in separate domains; live execution unimplemented and guarded
by \longtok{NeverExecutedError} & 17/17 & Not deployed \\
\end{longtable}
\endgroup

\emph{Sources: zvf-program/e1\_wave10/README.md;
zvf-program/e2\_core/driver.py;
zvf-program/e5\_successor27/controller.py;
zvf-program/e13\_balrog/hosted\_supervisor.py;
outputs/E1\_\allowbreak{}E14\_\allowbreak{}FINAL\_\allowbreak{}RESULTS\_\allowbreak{}2026-09-19.md.}

The companion E2 \texttt{launch\_gate.py} is deliberately trivial: it
checks only that the amendment acceptance record and the surviving
CORE-Bench harness entry points exist, which keeps the gate's claim
small enough to be true.

Read together, these four artefacts are the implementation answer to a
question the results chapters raise: how a project maintains evidential
integrity when its execution environment is partly lost and partly
unreachable. The answer taken here is to re-derive the controllers from
the sealed requests and the surviving test records, to keep every
re-implementation offline-testable, and to make the no-retry, no-replay,
hash-pinned, fail-closed properties explicit in code so that they can be
demonstrated rather than asserted. The cost of that discipline is
visible: four controllers, three of them complete and one a deliberate
stub, produced no new benchmark numbers in this phase. That is the
intended trade.

None of the four controllers ran the lanes in the end. On 27 September
2026 a single lead authorisation superseded the sealed-reservation and
owner-binding gates they enforce, and each lane was run with a thin
runner that calls the benchmark's native loop and grader directly. The
runners live under
\longtok{outputs/finish\_\allowbreak{}pending\_\allowbreak{}2026-09-27/\textless{}lane\textgreater{}/code/},
and each lane's \texttt{result.json} lists the bypass as a deviation
(§8.4). These fail-closed properties describe holding the lanes, not
measuring them.



\part{Results}


\chapter{Results I: Core GRPO Studies
(P1--P4)}\label{results-i-core-grpo-studies-p1p4}

This chapter reports the four core GRPO studies that make up the first
half of the P1--P8 internal reports: the cross-library and cross-scale
behaviour of GRPO training reward (P1), the Zero-Variance Fraction as a
diagnostic of signal starvation (P2), the effect of group size on
trainability (P3), and the length-bias and held-out-generalisation
behaviour of GRPO against its Dr.~GRPO variant (P4). Each study is
reported here in the form the evidence supports, which in every case is
more conservative than the hypothesis that motivated it. Three of the
four pillars return a negative or an explicitly bounded result, and the
fourth returns a diagnostic whose measured association with the outcome
it is meant to predict is weaker than the literature would suggest (Figure~\ref{fig:scaling}).


\fig{fig_scaling_null}
{Cross-scale behaviour of reward across the studied model range: no reliable monotone scaling trend is recovered}
{Cross-scale behaviour of reward across the studied model range: no reliable monotone scaling trend is recovered. (a) Twelve GRPO anchors from 4B to 1T parameters, one seed each, listed in the key with their number of logged steps. $\bar R$ is the mean over a run's logged steps of the per-step training reward (the fraction of rollouts rewarded, on a 0--1 scale; no further normalisation). Whiskers are $\pm1$~SD of the per-step reward within the run, clipped to $[0,1]$; they describe step-to-step noise, not a confidence interval. Open markers have fewer than ten logged steps. The dashed line and grey band are the constant model selected by AICc (pooled mean 0.641) and the $\pm\sigma_{\mathrm{step}}$ noise floor (0.096). OLS slope $+0.081\pm0.129$ per decade ($R^2=0.038$, permutation $p=0.54$); Spearman $\rho(\log_{10}N,\bar R)=+0.149$ ($p=0.64$). (b) OLS slope per decade of $N$ for seven summary metrics, with 95\% CIs ($\pm1.96$~SE) and permutation $p$ over the twelve anchors. Each slope is in its own metric's units (reward, squared reward or probability), so slopes are not comparable across rows; every interval includes zero. Sources: experiments/results/scaling\_\allowbreak{}law\_\allowbreak{}extended\_\allowbreak{}frontier.tsv and scaling\_\allowbreak{}law\_\allowbreak{}power\_\allowbreak{}law.tsv.}
{fig:scaling}



\section{Scope, Sources and Evidence
Boundary}\label{scope-sources-and-evidence-boundary}

The four pillars share one methodological commitment, which is the
organising principle of the whole project: every comparative claim is
made with the entire training stack held fixed except the one factor
under test. Where that was not possible, the result is reported as an
association rather than a contrast, and the confound is named. The
reader should treat this chapter as the primary result surface for
P1--P4; the corresponding LaTeX sources are the canonical artefacts, and
every figure quoted below is taken from one of them.

Each study draws on a different evidence base, and the differences
matter. P1 pools 70+ runs across seven reinforcement-learning libraries
and five model families spanning 0.6B to 1T parameters, but the
frontier-scale anchors within that pool are single-seed runs executed
through a managed training and inference API, so the cross-scale
analysis is descriptive and its uncertainty is wide. P2's measured
magnitudes come from a controlled five-seed panel plus a matched-budget
rollout budget check; a portion of its comparative table is a declared
simulation projection rather than a set of independent library
deployments, and that distinction is preserved below. P3's measured
sweep covers G ∈ \{2, 4, 8, 16\}; every G = 32 and above number in the
pillar is a joint-fit reconstruction, not a trained arm. P4's controlled
comparison is sixteen runs across two algorithms and two tasks; its
roster-level audit is drawn from a separate, larger pool that mixes
models, tasks and libraries.

The corpus-level motivation for the fixed-stack commitment (Chapter 4
§4.1) is a variance decomposition computed as four separate marginal
one-way ANOVAs, each run against the same 42 experiments with valid
performance observations: library/framework η² = 0.546, training
algorithm η² = 0.558, model family η² = 0.471, model scale η² = 0.322.
These four values sum to 1.897. A set of shares that sums above one is
not a partition of a common variance budget, and on a corpus where the
factors are collinear a marginal η² absorbs the variance of each
factor's correlates, so the values show that framework and algorithm
vary a great deal across this corpus and do not estimate either factor's
independent contribution. The historical twenty-test multiple-comparison
family that accompanies the decomposition silently includes single-seed
rows whose Welch statistics are mathematically undefined, and reduces to
a single surviving comparison under the restricted family; it is
retained as a reproducibility audit trail and not as a corrected
survival claim.

Throughout this chapter, ``statistically indistinguishable'' means that
a formal equivalence or difference test failed to separate the arms at
the stated level, not that the arms were demonstrated to be equivalent.
Where an equivalence test was run and returned a null, that is stated
explicitly.

\section{P1 --- Cross-Library and Cross-Scale GRPO
Behaviour}\label{p1-cross-library-and-cross-scale-grpo-behaviour}

The question P1 sets out to answer is whether the benefit of GRPO
post-training scales with model size. The framing is deliberate: the
hypothesis of interest is a Chinchilla-style power law in parameter
count N, and the P1 report's declared claim boundary is that it reports
a descriptive limits and identifiability audit rather than a positive
scaling law. The short answer is that no reliable scaling law is
recovered from the reward curves, and that the evidence base is too
small and too heterogeneous to falsify a plausible scaling hypothesis
even if one were present.

\subsection{The anchor pool and the fitted saturation
forms}\label{the-anchor-pool-and-the-fitted-saturation-forms}

The fitted anchor pool spans roughly 2.4 orders of magnitude in
parameter count. This is a subset of the full 0.6B--1T roster: for the
cross-scale fits, only anchors with a usable per-step reward trace and a
known parameter count were retained. Two properties of that pool limit
what the fits can support, and both are stated here rather than left to
the reader. First, the anchors are not measured on a common footing:
five of the twelve are interrupted runs whose recorded values are means
over three to five logged steps, against twenty to thirty for the
remaining seven, so the dependent variable is a short-horizon snapshot
for part of the pool. Second, scale is collinear with training recipe
across these anchors rather than independent of it: the learning rate,
LoRA rank and related settings change together and only at the scale
boundary, so a fit through the pool cannot separate an absent scale
effect from recipe changes offsetting scale changes. The null reported
below is therefore a statement about this pool as configured, and not an
estimate of how reward scales with parameter count in general. The
frontier-scale anchors --- Qwen3-235B-MoE, DeepSeek-V3.1 at the 685B
class, Kimi-K2 at 1T, and Nemotron-120B --- were obtained through the
Tinker managed training and inference API in a single-seed,
descriptive-only regime, since full base weights are not released for
those anchors and the project does not own frontier hardware.

Each trace was fitted with a saturating-exponential form
\(R(t) = R_{\max}(1 - e^{-\lambda t})\), with the saturation rate λ
constrained by the optimiser to the interval {[}0, 10{]}. The fitted
results are the first place the negative result becomes visible. Across
the five characterised anchors, the ordinary-least-squares trend slope
on the raw reward trace, β\_OLS, satisfies \textbar β\_OLS\textbar{} ≤
0.0037 in every case --- that is, the traces are flat to within a few
thousandths of a reward unit per step. Qwen3.5-4B fits a mean reward of
0.817 with a peak of 1.000 and λ pinned at the bound of 10.0; Qwen3-8B
fits 0.285 with a peak of 0.625; Llama-3.1-8B-Instruct fits 0.869 with a
peak of 1.000; DeepSeek-V3.1 fits 0.844 with a peak of 1.000. Four of
these five traces have λ pinned at the fit boundary --- that is, the
optimiser wanted a faster saturation than the constraint permitted ---
and the P1 report's verdict is that for those traces no single
saturation exponent is identifiable at all.

The exception is Nemotron-120B, which fits \(R_{\max} = 0.182\) with
\(\lambda = 0.99\), well inside the bound, and which the pipeline
classifies as a collapse rather than a saturation. The autopsy table
records a peak at step 3, a fitted \ensuremath{\hat{R}}\_max of 0.875, a late-phase mean
\ensuremath{\bar{R}}\_late of 0.208, a drop of −0.667 from peak to late phase, a
zero-reward step fraction of 0.55, and a post-peak OLS slope of −0.0036
per step. Across the rest of the pool, the zero-reward step fraction is
at most 0.067. The Nemotron run is therefore better described as a
distinct collapse phase than as a point on a smooth trend, and this is
the reason the pillar's headline object is a local, stack-conditioned
taxonomy of endpoint ceilings rather than a functional form in N.

\subsection{The cross-scale test and its
nulls}\label{the-cross-scale-test-and-its-nulls}

Three cross-scale statistics were computed against log₁₀ N: the slope of
the mean reward, the slope of the peak reward, and the slope of the
reward variance. None is distinguishable from zero. The mean-reward
slope is −0.024 per decade with a bootstrap mean of −0.117 and a 95\%
interval of {[}−0.796, +0.313{]}, Spearman ρ = −0.07 (a later
statistical audit,
\longtok{results/berkeley/adding\_\allowbreak{}error\_\allowbreak{}bars\_\allowbreak{}summary.json},
re-bootstraps the same anchors to a slope of −0.023 and {[}−0.796,
+0.323{]}; neither interval excludes zero); the peak-reward slope is
+0.037 with a bootstrap mean of −0.010 and an interval of {[}−0.623,
+0.197{]}, ρ = +0.22; the variance slope is −0.002 with an interval of
{[}−0.063, +0.027{]}, ρ = −0.14. A residual-resampling bootstrap with
1,000 draws applied to the Nemotron anchor returns \ensuremath{\bar{R}}\_max = 0.310, \ensuremath{\bar{λ}} =
3.26, P(λ ≥ 9.5) = 0.277 and a t₈₀ of 17.87 steps with an interval of
{[}0.16, 161.4{]} --- an interval wide enough to contain almost any
behavioural claim. A 70/30 anchor holdout returned slope estimates of
+0.0000, +0.0000, −0.0000, −0.0000 and +0.0016 across the five holdout
configurations. Fitting a power law over the full twelve-anchor set from
4B to 1T returns p ≥ 0.46 for every exponent tested.

Two model-selection checks were run to see whether the saturation story
was at least the best available description of the traces. Fitting
constant, linear and saturating forms to each trace, AIC and BIC select
the constant model on every single trace, with a ΔAIC of at most 2 over
the runner-up. That is, a flat line describes these reward curves at
least as well as any saturating curve does. A changepoint test with a
permutation null estimated change locations \ensuremath{\hat{τ}} of 27, 28, 5, 15 and 18
steps across five traces, with permutation p-values of 0.493, 0.789,
0.742, 0.909 and 0.137; none is significant. An
effective-plateau-horizon analysis at tolerance ε = 0.20 found that only
two of the five traces --- Qwen3.5-4B and Llama-3.1-8B-Instruct ---
reach a horizon at all, both at 27 steps, because the others never
satisfy the tolerance.

A compute-adjusted proxy, which folds an estimated token budget into the
scale axis, returns a slope of +482 ± 490 per decade. The point estimate
is positive and the interval is an order of magnitude wider than the
estimate, so the compute axis does not rescue the scaling story either.

\subsection{Where the geometric template
fails}\label{where-the-geometric-template-fails}

The three-phase learning template used throughout the project --- a slow
start, a period of rapid improvement, and a plateau, with collapse as a
degenerate tail --- was tested directly against the anchor pool as an
internal consistency check. Only 2 of the 12 anchors match the
template's geometric prediction. The template itself survives as a
descriptive vocabulary: it names the failure modes the pipeline detects,
and the collapse criteria (peak ≥ 0.4, late mean below 40\% of peak,
zero-reward fraction ≥ 0.30) are used in the pillar's flagging
throughout. What does not survive is the claim that these traces realise
the template as a progression.

One further contrast deserves note because it is our own measurement
rather than a synthesis. Under a same-stack comparison of PPO and GRPO
on a saturated Qwen2.5-0.5B two-digit addition task, evaluation accuracy
is 0.990 for GRPO and 0.992 for PPO, and the paired algorithm difference
is −0.002 with p = 0.374 (paired t on n = 5 seeds, each seed's held-out
accuracy on 200 items; source:
platform\_hybrid/paper/sections/frontier\_\allowbreak{}synthesis\_\allowbreak{}scaling.tex;
recomputed from
platform\_hybrid/experiments/results/samestack\_\allowbreak{}ppo\_\allowbreak{}grpo.json). The 200
evaluation items are fresh random draws from the training generator, not
a disjoint split. The pairing is by seed only, not by item: each arm
drew its evaluation set from the shared random stream after training,
and because GRPO consumes 16 problem draws per step against PPO's 128,
the two arms of a seed were evaluated on almost entirely different
items. Replaying the stream for the five seeds, the GRPO and PPO
evaluation sets share only 6, 6, 4, 7 and 3 of 200 items, and the
evaluation items that had also been drawn for training number 12, 16,
13, 15 and 13 of 200 for GRPO against 90, 86, 80, 83 and 74 for PPO,
roughly 7\% against 40\% (corrected; see Table E.1; replay of
\longtok{platform\_hybrid/experiments/modal/modal\_\allowbreak{}samestack\_\allowbreak{}ppo\_\allowbreak{}grpo.py}).
The paired difference is therefore a difference of per-seed accuracies
on different item sets with unequal train--evaluation overlap, not an
item-paired contrast. At a 0.98--1.00 ceiling this does not change the
null reading, but it is a further reason the null cannot be read as
equivalence. This is the project's closest de-confound: the runner,
clipped surrogate and step budget are shared, but the arms also differ
in prompts per step (16 prompts × 8 completions for GRPO against 128
distinct prompts for PPO; Table 4.2). The two estimators are
statistically indistinguishable, and because both sit at the ceiling the
null cannot separate no effect from no headroom. They are not shown to
be equivalent. The five per-seed differences are −0.005, +0.005, −0.005,
0.000 and −0.005, each one item of 200, at a 0.98--1.00 accuracy
ceiling. A paired TOST with a ±0.005 margin gives p = 0.104, so
equivalence at that margin is not established, although it is at ±0.010
(p = 0.008). The exact sign-flip permutation p on the same five pairs is
0.625. It is reported here rather than in P2 or P3 because it is the
algorithmic counterpart of the P1 null --- the cross-scale story is not
the only axis on which GRPO's advantage fails to separate from its
alternatives. The project's historical effect-size table also carries a
Qwen3-8B PPO-versus-GRPO row. It is marked a provenance conflict and not
estimable, and its Welch test, recorded in the historical twenty-test
family at p = 0.7605, is carried as a provenance record rather than as
an estimate, so no statistic is quoted for it here.

\subsection{A withdrawn architectural
contrast}\label{a-withdrawn-architectural-contrast}

An earlier version of this analysis reported a MoE-versus-dense gap
across six pairs of anchors, with a mean gap of +0.338 and a permutation
p = 0.023. That claim has been withdrawn. The defect is architectural:
the Nemotron-120B anchor is a LatentMoE, not a dense model, so the
original six-versus-six allocation misclassified it. With the
classification corrected, the same test returns a one-sided permutation
p ≈ 0.178. A leave-Nemotron-out recomputation returns a gap of +0.278
with p ≈ 0.049, and this too is withdrawn, because removing the
misclassified anchor from one side of the comparison does not repair the
design. The withdrawal is recorded here because the project's convention
is that retired claims remain visible with the reason for retirement
attached.

A related cross-anchor statistic is retained, with its own caveat.
Across the twelve anchors, seven have \textbar Δ\_1T\textbar{} ≤ 0.2,
where Δ\_1T measures how far the anchor's endpoint deviates from the
three-phase template's terminal prediction; the largest deviation is
Qwen3.5-27B at +0.562 and the largest negative is Nemotron at −0.50. A
Cohen's κ analysis of the categorical phase assignments is also reported
in the source: across three labelling methods, the heuristic and AIC
labels agree at κ = +1.000 on the five-anchor set (both label every
trace ``plateau'') and κ = +0.625 on the twelve-anchor set, while the
changepoint labels agree with neither (κ between −0.029 and +0.077). The
supportable summary is that categorical capability structure carries
more signal than log₁₀ N, which is what the P1 report's abstract says
and what the numbers support.

\section{P2 --- The Zero-Variance
Fraction}\label{p2-the-zero-variance-fraction}

The second pillar asks a mechanistic question: when GRPO stops learning,
is the cause visible in the reward statistics of the groups it samples?
The diagnostic proposed is the Zero-Variance Fraction (Figure~\ref{fig:zvflib}).


\fig{fig_zvf_by_library}
{Mean Zero-Variance Fraction by method in a simulation projection (a) and by measured experiment family (b); neither panel compares backend libraries}
{Mean Zero-Variance Fraction by method in a simulation projection (a) and by measured experiment family (b); neither panel compares backend libraries. (a) Nine variance-mitigation methods in a simulation projection ($G=8$, math-verifiable task, 5 seeds each; hatched bars mark projected, not model-training, values): bar $=$ mean of per-seed mean ZVF, whisker $=\pm1$~SD across seeds ($\le 0.008$). The projection ranks no methods; GRPO (cross-hatched) is the reference at 0.481, and in 3 of its 5 seeds the per-step collapse flag fires. (b) Measured runs: gsm8k $=$ Qwen3-8B, three seeds of 200 GSM8K prompts at $G=8$ (mean $\pm1$~SD); drgrpo $=$ Qwen2.5-0.5B arithmetic, GRPO and Dr.\ GRPO, 5 seeds each (10 runs, mean $\pm1$~SD); gsize $=$ Qwen2.5-0.5B group-size sweep, mean over $G\in\{2,4,8,16\}$ with a dashed whisker spanning the per-$G$ values (a range over $G$, not seed dispersion); tool-32B (Qwen3-32B) and tool-8B (Llama-3.1-8B-Instruct) are single runs (open diamonds, no interval) in which every group has zero reward variance. Source: platform\_hybrid/experiments/results/zvf\_\allowbreak{}by\_\allowbreak{}library.tsv and the per-seed files it cites.}
{fig:zvflib}



\subsection{Definition and
measurement}\label{definition-and-measurement}

For a prompt x sampled from the training distribution, GRPO emits G
completions and scores each with a verifier. Centring those rewards
within the group produces the advantage signal. If all G completions
receive the same reward --- all correct, or all incorrect --- every
centred reward is zero and the group contributes no contrastive signal
to that step's update. The exact-tie fraction records how often this
occurs; the default sample-variance tolerance in §4.3 also counts
near-ties. They agree for the binary rewards below. Gradient utilisation
is defined descriptively as GU = 1 − ZVF.

What a zero-variance group does to the update depends on the loss form.
The source states the general case: in an objective with a KL penalty or
auxiliary loss, a zero-variance group does not make the total gradient
zero, and ZVF measures only the fraction of groups contributing zero
centred reward contrast. An exact-tie group contributes zero to a
verified group-centred policy surrogate, but the full update can include
KL, critic or other auxiliary losses and optimiser state. Table 4.2
contains both nonzero-KL configuration evidence and unknown managed
settings; neither may be relabelled zero. Near-ties counted by the ε =
10⁻⁶ diagnostic need not have exactly zero policy gradients. ZVF
therefore describes reward contrast, not measured useful compute or the
whole update.

The statistic is validated by recomputation. Three stored GSM8K group
tensors (Qwen3-8B, G = 8, sampling only, 200 problems per seed) were
re-read from disk and their ZVF recomputed from the raw per-group
rewards; the recomputed values are 0.130, 0.190 and 0.155, pooled
0.1583, matching the reported values exactly. This is the mechanical
check that the quantity in the tables is the quantity the code computes.
The earlier ≤0.4\% ε-sensitivity statement did not distinguish exact
ties from thresholded variance and is replaced here by a reproducible
sweep over these same 600 groups. With independent seeded δ
\textasciitilde{} U(0,10⁻⁴) jitter, exact ties fall from 95/600 to
0/600. At sample-variance ε = 10⁻⁹, 71/600 groups remain below
threshold; at ε = 10⁻⁸ through 10⁻², including the declared 10⁻⁶, all
95/600 remain. Unperturbed binary rewards retain 95/600 across this
range. At G = 8, jitter alone has sample variance at most 2.86 × 10⁻⁹,
so it cannot cross 10⁻⁶ in an originally tied group. PCD changes by
+2.984 × 10⁻⁷. The example diagnoses exact-tie sensitivity, not failure
of tolerance-based ZVF; it does not establish robustness to larger
perturbations or other reward scales.

\subsection{Measured magnitudes across
libraries}\label{measured-magnitudes-across-libraries}

The headline measurement is a nine-row panel run at K = 8 with five
seeds and a 100-step rollout budget, spanning vanilla GRPO and eight
named variance-mitigation methods. Vanilla GRPO records a mean ZVF of
0.481 against a mean last-10 accuracy of 0.379, with zero of five seeds
collapsing, zero of five drifting, and five of five plateauing. The
remaining rows are: AERO 0.220 ZVF with 0.399 accuracy; CPPO 0.295 with
0.392; NGRPO 0.300 with 0.387; SCAFGRPO 0.317 with 0.409; MCGRPO 0.146
with 0.322; GIFT 0.114 with 0.326; AREAL 0.121 with 0.328; and ES 0.114
with 0.316. The four lowest-ZVF rows each record one of five seeds
drifting and four of five plateauing.

The variance-mitigation rows require a caveat that the source states
explicitly and that this chapter preserves: they are a declared
simulation projection used for sensitivity analysis, and they ``should
be read as synthetic, not as independent real-library deployments''. The
caveat covers the whole panel. The vanilla GRPO row and its
collapse/drift/plateau counts come from the same dry-run file,
\longtok{experiments/results/variance\_mitigation.tsv}, as the other
eight rows, and Chapter 2 §2.3 already records that GRPO baseline row as
synthetic. None of the nine rows is a measured deployment. Anyone citing
the nine-row ordering as evidence that a named method mitigates failure
would be citing a projection.

Two further results sharpen the picture, and the first mixes measured
and projected rows. First, across 60 runs the late-phase ZVF of GRPO
runs sits in the range 0.87--0.97, while early-phase ZVF for
variance-mitigation GRPO and for the largest group size sits in the
range 0.04--0.13; smaller group sizes start higher, between 0.25 and
0.61. The lag-1 autocorrelation of the ZVF trace is approximately 0.94,
meaning the quantity moves smoothly rather than jittering, which is what
makes it usable as a warning indicator at all. Of those 60 runs, 45 are
the five-seed variance-mitigation projection rows above (corrected; see
Table E.1). The measured part is the remainder, which includes the
12-run group-size sweep. The source does not report the ranges or the
autocorrelation separately for the measured subset, so these figures
describe a pool that is mostly synthetic.

Second, and more directly relevant to practice, the diagnostic rises
\emph{after} the reward plateau rather than before it. Measuring
first-passage times, ZVF crosses its warning threshold only 1 to 3 steps
ahead of the plateau. The recommended operating rule from the pillar is
therefore not predictive but custodial: report ZVF alongside mean
reward, and treat a rise of more than 0.3 sustained over ten consecutive
rollout steps as a warning. The project reports this limitation against
its own proposed diagnostic rather than in a footnote.

\subsection{The all-correct wall and the mastery--incapacity
aliasing}\label{the-all-correct-wall-and-the-masteryincapacity-aliasing}

The mechanism behind large ZVF values on easy tasks is a saturation
effect that the project calls the all-correct wall. A matched-budget
panel was run at 2,560 rollouts per arm --- G = 2 over 160 steps against
G = 16 over 20 steps --- on Qwen3-8B with GSM8K, LoRA rank 4 and batch
size 8. The late G = 2 arms reach a mean training reward of
approximately 1.0 with ZVF between 0.75 and 1.0; the G = 16 arms hold
ZVF at or below 0.25 at the same rollout budget. The reading is that
when the policy solves the task, nearly every group becomes all-correct,
nearly every centred reward is zero, and the estimator loses its
contrast even though the reward signal is at its maximum. Signal
starvation in this regime is driven by ease, not by difficulty --- which
is the opposite of the intuition that motivates variance-mitigation
work.

This produces the pillar's most important structural caveat. ZVF is
unsigned and symmetric: an all-wrong group (\(K_x = 0\)) and an
all-correct group (\(K_x = G\)) score identically. The statistic
therefore cannot distinguish mastery from incapacity, and this aliasing
is the explanation the project gives for the weak association between
ZVF and final held-out accuracy. Under an i.i.d.-Bernoulli model the
closed form is \(\mathrm{ZVF}(p_x, G) = h_G(p_x) = p_x^G + (1-p_x)^G\),
and it follows in one line. If the G rewards of a group are independent
Bernoulli(p) draws, the group has zero variance exactly when all G
rewards are equal. That happens when all are 1, with probability
\(p^G\), or all are 0, with probability \((1-p)^G\). The two events are
disjoint, so \(P(\text{zero variance}) = p^G + (1-p)^G\). At
\(p = \tfrac12\) this is \(h_G = 2^{1-G}\), so \(\ln h_G\) falls by
\(\ln 2\) for each unit increase in G. The project's source document for
this closed form and the aliasing analysis is headed as reasoning
contributions and proposed theorems from ChatGPT Pro Extended and Gemini
Deep Think, ``NOT experiments this paper ran''. That provenance is kept
on record, and AI assistance is declared in the front matter, but the
result does not rest on it: the derivation above is elementary and is
the argument this thesis relies on. What the project itself measured is
the empirical ZVF against the closed form, evaluated at the per-step
training reward rate because ZVF is a training-batch quantity
(corrected; see Table E.1). Recomputed from the per-step training
rewards in
\longtok{platform\_hybrid/experiments/results/groupsize\_\allowbreak{}zvf\_\allowbreak{}sweep.json}
(3 seeds × 40 steps per arm), the per-step-average null is 0.828, 0.709,
0.597 and 0.522. The measured mean ZVF is 0.838, 0.764, 0.691 and 0.631,
so the residuals are +0.010, +0.055, +0.094 and +0.109: measured ZVF
sits \emph{above} the null at every G. Gradient utilisation rises across
the same range from 0.162 to 0.369. The direction of that gap depends
entirely on which baseline is used, and the project's own formalisation
names the correct one. Evaluating the closed form at a single pooled or
held-out rate overstates the Bernoulli prediction, because \(h_G\) is
convex and the per-step reward rate varies; the right baseline is the
per-step average of the closed form, which is \(0.60\) at \(G = 8\)
rather than the \(0.33\) obtained at the pooled mean. Against the
per-step-average baseline the measured ZVF at \(G = 8\) (\(0.69\))
\emph{exceeds} the null by approximately \(0.09\), and the source is
explicit that most of that apparent gap is the Jensen effect rather than
correlated completions and that it therefore ``do{[}es{]} not read it as
strong evidence that the independence null is a lower bound''. This
chapter uses only the per-step-average baseline. It is a weak
consistency check on shape rather than evidence about dependence.

The formalisation's own summary of the same sweep, twelve runs and 480
logged steps, reports per-step observed ZVF correlating at r = 0.92 with
the Bernoulli-null prediction \(p^G + (1-p)^G\), with mean ZVF of 0.84,
0.76, 0.69 and 0.63 at G = 2, 4, 8 and 16 respectively. That correlation
is a shape check, not a validation: the null for each step is computed
from the same step's mean reward, so a high r is partly mechanical. At
larger scale the same instrument was run on Qwen3-8B against GSM8K test
problems through the Tinker managed API at G = 8, temperature 1.0, over
200 problems and three seeds (sampling only, no training; these are the
three stored tensors of §6.3.1), yielding a mean ZVF of 0.16 composed of
12.7\% all-correct groups, 3.2\% all-wrong groups and 84\% mixed groups;
the raw correlation with the mechanical null is approximately r = 0.93
with rank ρ ≈ 0.65, and the i.i.d. null over-predicts (approximately
0.28 against the observed 0.16). The direction of the deviation from the
null is therefore opposite in the two runs: above the null on the 0.5B
sweep, below it on this Qwen3-8B sampling run. The source does not state
which baseline produced the 0.28, and neither deviation is read as
evidence about dependence between completions.

\subsection{Prognostic value}\label{prognostic-value}

The pillar's central claim is not that ZVF is large but that it is
informative, and here the results split cleanly. Pooling the measured
runs, the naive correlation between mean ZVF and last-10 accuracy is r =
+0.22, which is not significant. The residual table on disk pools n = 17
per-experiment rows, and at n = 17 the Fisher-z 95\% CI is {[}−0.29,
+0.63{]} (corrected; see Table E.1). After mean reward is regressed out
of last-10 accuracy, the pooled residual correlation with ZVF is r =
+0.205 (Fisher 95\% CI {[}−0.31, +0.62{]}, p = 0.43, n = 17
experiments), and its interval includes zero. The source's own
conclusion from that analysis is that ``we therefore do not treat ZVF as
an incremental predictor above reward mean'', and this chapter keeps
that conclusion. Neither the raw nor the residualised diagnostic is
shown to predict graded outcome on measured data.

Pooling across task types shows why a pooled association cannot be read
as a score. Over N = 15 experiments spanning five model families, the
association between ZVF and performance is Pearson r = −0.769 (p =
0.0008) with Spearman ρ = −0.784 (p = 0.0005), but it is confounded by
task type: tool-use experiments reach ZVF = 100\% at approximately 0\%
reward while GSM8K experiments average ZVF = 8.5\%. Within GSM8K the
association weakens to r = +0.40 (p = 0.09, N = 19), and then to r =
+0.13 (p = 0.79) when restricted to the seven model-varying runs. The N
= 19 is a separate and larger pool of Tinker GSM8K runs, not a subset of
the N = 15 pooled experiments, so the two analyses are not nested. A
statistic whose sign is set by which tasks are included in the pool is
not a score. On the instrumented runs ZVF also does not track final
accuracy (r ≈ 0.09, n = 12).

Against catastrophic collapse, ZVF performs much better than it does
against graded outcome. Mean ZVF correlates with catastrophic collapse
at Spearman ρ = 0.56 (Fisher 95\% CI {[}0.19, 0.79{]}, n = 23 pooled
cells) and point-biserial 0.62. Its correlation with final held-out
accuracy is ρ = 0.27, and that is not significant: the registry's
bootstrap 95\% CI is {[}−0.37, 0.88{]} and the Fisher CI is {[}−0.16,
0.61{]}, p = 0.21, n = 23. The 23 cells pool 80 heterogeneous local
rows, and the source does not state how many of them are
variance-mitigation projection rows. A three-channel risk index built on
the diagnostic --- magnitude, critical slowing down via a rolling lag-1
autocorrelation with window 15, and drift slope, blended 0.3/0.5/0.2
with a max-fusion variant --- was evaluated over a 52-row
cross-experiment panel comprising 45 variance-mitigation seeds, 2
tool-use collapse anchors and 5 scaling-law collapse anchors, with 11
failures and 41 safe runs. Cross-experiment AUROCs are 0.663 for
magnitude, 0.758 for critical slowing down, 0.579 for drift, 0.763 for
the blended risk index and 0.929 for the max-fusion index;
within-methods, the same values fall to 0.073, 0.335, 0.293, 0.396 and
0.805 respectively. None of these AUROCs is a measurement on real
training runs, and this chapter reports them as a synthetic-panel
exercise. All 45 variance-mitigation rows come from the declared
simulation projection (§6.3.2). The seven collapse anchors enter with
channel values fixed by the analysis script, not measured: mean ZVF 1.0
or an imputed 0.95, lag-1 0.95 or 0.85, and slope 0 or 0.015. The 41
safe rows are therefore all synthetic. The max-fusion index was selected
after the per-axis results were seen. For the 0.929 figure, a
Hanley--McNeil 95\% CI with 11 positives and 41 negatives is {[}0.82,
1.00{]}; that interval describes this constructed panel only. No AUROC
is computed on the measured rows alone, because they contain too few
failures to estimate one. The gap between the cross-experiment and
within-method columns is a warning about the index's transportability:
most of the cross-experiment discrimination is coming from the fact that
collapse clusters by experimental family, not from a within-family
predictive signal. The declared thresholds are a risk above 0.55 to
treat as failure-positive, below 0.30 as safe, and the interval between
as requiring investigation.

The pillar's own conclusion names three load-bearing limitations: the
diagnostic is mechanically coupled to reward sparsity, group size and
baseline accuracy; its strong association holds only for catastrophic
collapse and not for the graded held-out outcome; and the unsigned
statistic aliases mastery with incapacity. A proposed replacement, the
pairwise contrast density, is introduced in the synthesis document with
a target association of ρ ≳ 0.45; the project states that this bar is
not yet met by the replacements and that the head-to-head protocol
remains future work.

\section{P3 --- Group Size and Contrast
Density}\label{p3-group-size-and-contrast-density}

The third pillar studies G, which plays two roles simultaneously in
GRPO: it is a variance-reduction knob, like a Monte-Carlo baseline, and
it is a preference-density knob, since centring rewards within a group
induces implicit pairwise contrasts --- nonzero only for mixed-reward
groups. Practitioners often treat larger G as strictly better by analogy
to variance reduction, which conflates the two roles and ignores cost.
The pillar's conclusion is that the measured evidence does not justify a
group-size optimum and that the dependence of trainability on G is
non-monotone (Figure~\ref{fig:groupsize}).


\fig{fig_group_size}
{Group size $G$: one single-seed sweep, one three-seed sweep and one reconstruction}
{Group size $G$: one single-seed sweep, one three-seed sweep and one reconstruction. (a) Qwen3-8B on GSM8K through Tinker, one seed, 30 steps: peak and last-10 in-training reward (not held-out); $G=8$ is highest, but with no seed replication no optimum can be established. (b) Qwen2.5-0.5B on arithmetic, 3 seeds, 40 steps: held-out accuracy and mean train reward, bars $\pm1.96$~SE (approximate 95\% CI). The held-out apex at $G=8$ (0.990) overlaps $G=4$ (0.988) and $G=16$ (0.978); overlap is descriptive, not a test, and the arms are statistically indistinguishable. (c) Mean zero-variance fraction on the same sweep, falling by 0.207 from $G=2$ to $G=16$. (d) Reconstructed retention of $G=4$ relative to $G=32$ against token budget (dashed line and open markers: reconstructed values, not matched-budget training runs; bars are the file's conditional reconstruction intervals). The point estimate reaches the Wu et al.\ 0.976 retention level at 1 of 4 budgets ($T=1$M), but the equivalence test at that margin is not passed at any budget. Sources: paper/expected\_results.json (a); experiments/results/group\_\allowbreak{}size\_\allowbreak{}effect.tsv (b--d).}
{fig:groupsize}



\subsection{The measured sweep}\label{the-measured-sweep}

The primary sweep measures G ∈ \{2, 4, 8, 16\} on a near-ceiling
two-digit addition task with Qwen2.5-0.5B, three seeds and 40 steps,
evaluated on 200 fresh draws from the training generator. Held-out
accuracy (200 fresh problems per run) traces an inverted U with its apex
at G = 8. Values are mean ± SE over n = 3 seeds: 0.9900 ± 0.0029 (SD
0.0050) at G = 8, 0.9817 ± 0.0044 (SD 0.0076) at G = 2, 0.9883 ± 0.0017
(SD 0.0029) at G = 4 and 0.9783 ± 0.0060 (SD 0.0104) at G = 16. At n = 3
seeds the arms are statistically indistinguishable: a one-way ANOVA
across the four arms gives F(3, 8) = 1.82, p = 0.22, recomputed from
\longtok{platform\_hybrid/experiments/results/groupsize\_\allowbreak{}zvf\_\allowbreak{}sweep.json}.
The minimum detectable effect at 3 versus 3 seeds is d = 3.07, so only
very large differences could have been detected. A separate check run at
B = 4000 with seed 20260702 records G = 2 and G = 16 at 0.9817 and
0.9783, a difference of +0.0033 (G = 2 minus G = 16) and an interval of
{[}−0.008, +0.015{]}; G = 4 and G = 8 add +0.010 and +0.012
respectively, and G = 4 versus G = 32 was not measured on that sweep.

The retention figures are the cleanest way to state the same result. G =
2 retains 100.3\% of G = 16. A later unified synthesis at the same
configuration reports ``retention'' ratios and signal-to-noise ratios
with intervals. The ``retention'' is a ratio of last-10 \emph{training}
reward, not of held-out accuracy, so this chapter calls it a
training-reward ratio. The point values are G = 2 1.0000, G = 4 0.9941,
G = 8 0.9960 and G = 16 1.0068, with SNR 1.46, 1.76, 1.94 and 2.16. All
four ratios sit above the Wu et al. \citep{wu2025ittakestwo} 97.6\%
reference bar. At the seed level (paired TOST on per-seed last-10
training reward against G = 2, margin ±0.02, n = 3 seeds), the TOST
p-values are 0.012 at G = 4, 0.043 at G = 8 and 0.056 at G = 16
(corrected; see Table E.1). Equivalence to G = 2 at that margin is
therefore supported, narrowly, for G = 4 and G = 8, and is \textbf{not}
established for G = 16.

The signal-to-noise ratio grows sublinearly with G: it reaches about
52\% of the √G ideal, and about 48\% above a flat baseline. Reward
itself is statistically flat across the sweep, in the range
0.971--0.984. The marginal ZVF gain from G = 8 to G = 16 is 0.6906 →
0.6312, a fall of 0.059, which is smaller than the G = 2 → G = 4 fall of
0.8380 → 0.7635, a fall of 0.075. Diminishing returns on the contrast
axis are visible in the measured data, which is what the reconstruction
later formalises.

The pillar's own framing of this regime is important and is stated in
its abstract: near the ceiling, variance is compressed, so the finding
is ``no detectable G effect under saturation'', not evidence of
equivalence on harder, unsaturated tasks. That is, the flat retention
curve is a property of an easy task, and the pillar explicitly declines
to generalise it. The measured equivalence tests were run up to G = 16,
and at the seed level they establish equivalence only up to G = 8; there
is no directly trained, matched-budget G = 32 cell in the primary sweep.
The retention-versus-budget axis is where the picture changes: R(1M) =
0.976, R(4M) = 0.833, R(16M) = 0.750 and R(64M) = 0.727, fitted by
\(R(T) = 0.729 + 0.247\,e^{-T/5.76\mathrm{M}}\) with \(R^2 = 0.948\) and
an asymptote of 0.729. The four points are themselves reconstructed
(§6.4.2), and a three-parameter curve through four points leaves one
residual degree of freedom. The fit is therefore illustrative only, and
no confidence interval is reported for its asymptote (corrected; see
Table E.1). The Wu reference line and this reconstructed curve cross at
T = 1M and diverge by 25 percentage points at T = 64M. The operative
conclusion is scoped to the 0.5B arithmetic task: no G is
distinguishable there, so G = 2 is the cheapest setting at statistically
indistinguishable accuracy. It is not a recommendation for harder tasks,
and it does not conflict with the single-seed Qwen3.5-4B observation of
§6.4.3 that G = 2 under-trains with a 50\% batch collapse, which is a
different model and task.

\subsection{The same-stack four-axis audit and the reconstructed
budgets}\label{the-same-stack-four-axis-audit-and-the-reconstructed-budgets}

The strongest de-confound for P3 is a four-axis factorization audit of
the same training stack. On the algorithm axis, the residual effect is
small: η² = 0.023, from the same five-seed same-stack contrast as
§6.2.3. The GRPO−PPO difference is −0.002, paired by seed only: the two
arms of each seed were evaluated on different items, with PPO's
evaluation items overlapping its training draws far more (§6.2.3;
corrected; see Table E.1). The factorization file reports a Monte-Carlo
permutation p of 0.62; the exact sign-flip p over the five pairs is
0.625, and the paired t p of 0.374 in §6.2.3 is the primary test for
this contrast. Equivalence on the algorithm axis is not established,
because the TOST at ±0.005 in §6.2.3 does not establish it (corrected;
see Table E.1). On the group-size axis the file reports η² = 0.54 and
0.63, with Cohen's d = −1.47 for the G = 2 → 16 comparison. The 0.54 is
\longtok{eta2\_\allowbreak{}G\_\allowbreak{}last10}, a share of last-10 \emph{training reward}, not
of terminal held-out accuracy; 0.63 is for gradient norm. Recomputed on
the 12 sweep runs (3 seeds per G), the training-reward share gives F(3,
8) = 3.08, p = 0.090. The corresponding held-out share is η² = 0.41,
F(3, 8) = 1.82, p = 0.22. The source itself grades the G-axis result
only SUGGESTIVE. The supportable reading is that the group-size axis
shows a large but not statistically significant share of training-reward
variance at n = 3 seeds per arm, and that the algorithm label is not
shown to matter on this stack.

A Miller-style \citep{miller2024errorbars} error-bars audit covered the
pillar's four headline contrasts. Two of them are held here at
hypothesis-generator status rather than reported as precise (corrected;
see Table E.1): the gradient-utilisation ratio at T = 1M, with a point
estimate of 5.03 and no interval that n = 4 budgets can support, and the
negative retention-versus-training-token slope, which is fitted to the
four \emph{reconstructed} budgets of the grid below, not to trained
arms. The SNR-versus-√G slope has an interval of {[}+0.148, +0.583{]},
which contains the ideal value of +0.500 and is therefore graded NULL.
The native Wu-style equivalence claim that G = 2 ≈ G = 16 is also graded
NULL. This chapter therefore reports none of the four P3 headlines as
precise, and the pillar's own summary is that the data do not justify a
universal optimum.

Everything at or above G = 32 in the pillar comes from an explicitly
reconstructed token-budget grid, retained only as a hypothesis
generator. The reconstruction table gives, for G = 4 versus G = 32, a T
= 1M comparison of 0.41 against 0.42 with a difference of −0.01, a
97.6\% retention ratio and a ``generalizes'' verdict of yes, a
point-estimate criterion of retention at least 90\% --- the equivalence
test at the 97.6\% margin, with a reconstruction interval of {[}0.873,
1.092{]} at T = 1M, is not passed at any of the four budgets; at T = 4M,
0.55 against 0.66 with −0.11, 83.3\% and no; at T = 16M, 0.63 against
0.84 with −0.21, 75.0\% and no; at T = 64M, 0.64 against 0.88 with
−0.24, 72.7\% and no. The source labels these as conditional
reconstruction intervals, not fresh-run confidence intervals. The T =
64M accuracy row across G = 4, 8, 16, 32 and 64 reads 0.64, 0.80, 0.85,
0.88 and 0.87 --- the apex moves with budget, which is precisely the
pattern that makes a universal G optimum incoherent, and precisely the
pattern the project declines to claim as measured. A fit of the
critical-budget form \(|\Delta(T)| = a(1 - e^{-T/\tau})\) gives â = 0.30
and \ensuremath{\hat{τ}} ≈ 12.1M tokens, with a 5\% critical budget of 2.21M tokens and an
interval of {[}1.29M, 4.77M{]}. A difficulty-gated adaptive rule derived
from this analysis reports a mean selected group size of 7.0 and a
projected 56.2\% rollout saving against always using G = 16, with 5 of
20 cells under-predicting by more than 5 percentage points. These are
projections, and the chapter marks them as such.

A separate conditional utility audit over 505 tasks, scoring each group
size by (1 − ZVF)/√G, selects G = 4 over G = 5 by 0.00177 with a 95\%
bootstrap interval of {[}0.00082, 0.00270{]}. The margin is positive but
small, and it is conditional on cohort and objective.

\subsection{The curriculum and matched-token null
results}\label{the-curriculum-and-matched-token-null-results}

Two interventions that P3's design space naturally suggests were run and
both returned nulls. They are reported here rather than elsewhere
because they are group-size interventions, and because their negative
results are load-bearing for the chapter's overall argument.

The first is a difficulty curriculum that filters collapsed groups
during training. The intervention works as designed on its own terms: it
drives the zero-gradient waste fraction from 0.50 to 0.00, eliminating
wasted rollouts completely. It does so at approximately 4.81 times the
sampling cost, and its held-out gain is identical to the baseline's ---
both arms move from 0.80 to 0.85 on a held-out set of n = 20 (one seed,
eight steps), a gain of +0.05 in each case, which is one item of 20. At
a matched 30k-token budget the curriculum ties the baseline at +0.028
against +0.028, and at five times the cost it loses. The verdict
recorded is Null. An important provenance note attaches to this result:
an independent verification pass conducted on 6 July 2026 flagged the
curriculum claim, along with a single-seed G = 4 ``win'' of +0.125, as
SUSPECT and overclaimed, before multi-seed re-testing confirmed both as
noise. The group-size sweet-spot result is itself a null for the same
reason --- the +0.125 advantage at G = 4 rested on a single held-out
example and did not survive multi-seed re-testing, and G = 2
under-trains with a 50\% batch collapse and a +0.0 gain.

The second null family concerns loss formulations intended to remove the
length bias. On the held-out metric, the standard sum formulation is
nominally the best of the three tested: +0.125 against +0.000 for
mean-normalisation and +0.042 for the surprise formulation. These are
measured on held-out sets of n = 12--20 per arm, so each difference
corresponds to about one or two items and none is distinguishable from
noise. The length change under the standard arm is recorded as Δlen =
−16. The verdict is Null. A fourth null, on a step-1 predictive
layer-freeze, is reported in the same table because it was proposed as a
cheap way to stabilise group-relative training: it achieves a per-step
overlap of 1.0 on a 1.5B hardcoded-arithmetic toy, but only 0.11 ---
approximately chance, on both seeds --- at 3B on real GSM8K, while the
concentration statistic still holds at approximately 0.39. A later
four-seed pass at 3B supersedes that two-seed figure with a mean overlap
of 0.083 ± 0.048 (mean ± SD, n = 4 seeds), also approximately chance.
Verdict: Null, and a reversal of its own toy-scale positive under
scaling.

The pattern across all four nulls is consistent and worth stating
plainly: the interventions that target the mechanism the diagnostic
identifies do not improve held-out outcomes at matched cost. That is a
negative result about the diagnostic's prescriptive power as much as
about the interventions themselves, and it is one of the reasons P2 is
reported as a monitoring tool rather than an optimisation objective.

\section{P4 --- Length Bias and Held-Out
Generalisation}\label{p4-length-bias-and-held-out-generalisation}

The fourth pillar tests the most widely repeated specific criticism of
GRPO: that it rewards verbosity, inflating completion length in a way
that Dr.~GRPO corrects. The pillar's headline is a scoped null, and the
scope is the point (Figure~\ref{fig:lengthbias}).


\fig{fig_length_bias}
{The post-hoc peak-then-decay rule (peak before 65\% of logged steps and terminal reward below $0.90\times$ peak) on the roster and on the controlled comparison}
{The post-hoc peak-then-decay rule (peak before 65\% of logged steps and terminal reward below $0.90\times$ peak) on the roster and on the controlled comparison. (a) The four of eleven roster GRPO runs that carry the flag (single seed each; legend gives peak$\rightarrow$terminal reward and t/p $=$ terminal$\div$peak; stars mark peaks). (b)--(d) The controlled sixteen-run comparison, seed mean (line) and min--max across seeds (band); dashed vertical line at 65\%, dotted horizontal line at the arm's last-10 mean. (b) Arithmetic, Qwen2.5-0.5B, 40 steps, 5 seeds per arm: every seed peaks at 1.0 and stays there, so the rule fires on 0 of 10 runs. (c, d) GSM8K-CoT, Qwen2.5-1.5B-Instruct, 30 steps, 3 seeds per arm: seed peaks (stars $=$ flagged, open circle $=$ not flagged) fall at 33--63\% of training and the last-10 mean (0.26 GRPO, 0.25 Dr.\ GRPO) is far below every peak, so the rule fires on 3 of 3 GRPO and 2 of 3 Dr.\ GRPO runs. Completion length falls in all sixteen runs, so the pillar's length-bias flag (length rising while reward stays flat or falls) is zero everywhere: the peak-then-decay rule marks noisy reward in both algorithms, not length inflation. Flags recomputed from the per-step logs in drgrpo\_\allowbreak{}vs\_\allowbreak{}grpo.json and drgrpo\_\allowbreak{}gsm8k\_\allowbreak{}cot\_\allowbreak{}full.json.}
{fig:lengthbias}



\subsection{The controlled null}\label{the-controlled-null}

Sixteen controlled runs were analysed: ten on an arithmetic task with
Qwen2.5-0.5B over 40 steps, and six on GSM8K chain-of-thought with
Qwen2.5-1.5B-Instruct over 30 steps, split evenly between GRPO and
Dr.~GRPO. The arithmetic half is saturated: evaluation accuracy is 0.987
for GRPO and 0.992 for Dr.~GRPO (paired p = 0.35, five seeds), and
completions average 4.7 tokens, so its accuracy and length nulls are
nulls at a ceiling. Only the GSM8K-CoT half has headroom. At these
30--40-step horizons, neither algorithm exhibits the verbosity-trap
signature on either task: the pillar's length-bias flag, which fires
when completion length grows while reward stays flat or falls, is zero
on every one of the sixteen runs. That flag is not the roster audit's
peak-then-decay rule of §6.5.2. Recomputed from the per-step logs, the
peak-then-decay rule (peak before 65\% of training, last-10 mean reward
below 90\% of peak) fires on none of the ten arithmetic runs but on five
of the six GSM8K-CoT runs, three GRPO and two Dr.~GRPO, because the
reward trace on that task is noisy and peaks early; length still falls
in all five.

The summary correlations are as follows, given as mean ± SD across seeds
(n = 5 arithmetic, n = 3 GSM8K-CoT). On arithmetic, the rank correlation
between training step and completion length is −0.405 ± 0.176 for GRPO
and −0.307 ± 0.181 for Dr.~GRPO; on GSM8K-CoT it is −0.622 ± 0.301 and
−0.355 ± 0.255. The step-versus-reward correlations are +0.709 ± 0.082
and +0.753 ± 0.074 on arithmetic, and +0.046 ± 0.162 and +0.025 ± 0.091
on GSM8K-CoT. The length-versus-reward correlations are −0.419 ± 0.155
and −0.293 ± 0.182 on arithmetic, and −0.501 ± 0.258 and −0.618 ± 0.118
on GSM8K-CoT. Mean reward over the last ten steps is 0.988 for both
algorithms on arithmetic, and 0.265 for GRPO against 0.252 for Dr.~GRPO
on GSM8K-CoT. The direction of the length trend is negative everywhere:
all sixteen length slopes are negative, with the arithmetic GRPO slope
at −0.029 tokens per step and the GSM8K GRPO slope at −0.481, and the
count of runs with a positive horizon is zero. On the harder task,
GRPO's final-ten-step mean completion length is 184.4 tokens against
Dr.~GRPO's 189.3 --- a difference of +2.7\% in the direction opposite to
the critique.

The pillar's abstract states the scope of this null explicitly. The
claim is bounded by the 200-token cap used in that regime, and the
summary sentence is that ``the verbosity-trap signature is absent in
this regime''. The pillar is not claiming that GRPO never inflates
length; it is claiming that within a capped, short-horizon regime it
does not, and it names the length-confounded regime in which a Dr.~GRPO
advantage should appear as future work.

\subsection{The roster-level verbosity-trap
audit}\label{the-roster-level-verbosity-trap-audit}

Set against the controlled null is a descriptive audit over the whole
experiment roster, and here the picture differs. The audit defines four
statistics per trace with at least five steps: an instability index
σ(r\_tail)/\textbar μ(r\_tail)\textbar{} over the final 50\% of
training; a monotonicity score, the fraction of consecutive step-pairs
with strictly improving reward; a regression severity, the largest
normalised drop from peak; and a binary verbosity-trap flag, set when
the peak occurs before 65\% of training and the terminal reward falls
below 90\% of peak.

On that definition, the eleven GRPO runs in the roster have a mean
instability index of 0.267 ± 0.335, against 0.785 ± 0.297 (mean ± spread
across runs; the source does not label the ±, and its size relative to
the mean implies an SD) for the seventeen PPO-labelled runs. Four of the
eleven GRPO runs, 36\% (Wilson 95\% CI {[}15\%, 65\%{]}, one run per
configuration), carry the flag on this definition, including both
Qwen3-8B runs and the Nemotron-120B run. That count is reported here as
a description of the roster, not as evidence of length bias, because the
flag does not isolate the behaviour it is named for. Three properties of
the rule prevent the stronger reading. Its thresholds are post hoc:
neither the 65\% peak position nor the 0.90 terminal-to-peak ratio is
pre-registered anywhere in the repository, and the rule cannot be traced
to a commit that predates the traces it judges. It is not specific to
length: the metric it tests is in-training reward, and the project's own
source reads the flag as ``a generic instability detector'' rather than
a length measure, so a run can carry it without any change in response
length. And its base rate is high: applied to the wider trace pool
rather than this eleven-run subset, the rule fires on the majority of
runs --- including PPO-labelled runs and the Dr.~GRPO control in the one
matched comparison --- so 36\% sits at or below the rate the rule
produces on traces with no decay at all. The project's own length data
point the other way: under the 200-token cap, length drifts from 194.4
to 183.5 tokens, i.e.~shorter. The run-level audit therefore does not
establish length inflation, and the pillar's negative finding rests on
the controlled comparison below rather than on this roster count. The
Nemotron run records the highest instability index in the pool at 1.194,
with a regression severity of 1.0 and a terminal-to-peak reward ratio of
0.57. The two DeepSeek-V3.1 GRPO runs, by contrast, show markedly lower
instability, with a mean of 0.143. The PPO rate of 12/17 has a Wilson
95\% CI of {[}47\%, 87\%{]}, which overlaps the GRPO interval, so the
two rates are not distinguished either.

Two caveats on that comparison must travel with it. First, the PPO pool
is dominated by classic-RL library runs (SB3, CleanRL, Tianshou) that
failed to train at all, with last-10 rewards below 5\%; a near-zero
reward trajectory mechanically produces a high instability index.
Restricted to LLM-native PPO, the PPO Llama instability is 0.105, lower
than GRPO Llama's 0.228, and the source therefore explicitly declines to
claim that GRPO is more or less stable than PPO at the algorithmic
level. Second, the audit's reward metric is the in-training reward, not
a held-out one, and the audit pool is single-seed per configuration:
each Tinker configuration was run once and is reported descriptively.
The only multi-seed matched base-versus-RL control on Qwen3-8B is the
five-seed comparison discussed below.

An earlier draft of this report's summary attributed to P4 the claim
that Dr.~GRPO is the only algorithm without the verbosity-trap pattern.
That claim could not be located in any P4 source. What the sources state
is that 36\% of the audit's GRPO runs carry the flag; the audit does not
report a per-algorithm flag rate for Dr.~GRPO, and in the controlled
sixteen-run comparison the length-bias flag is zero for both algorithms
while the peak-then-decay rule fires on both (three of three GRPO and
two of three Dr.~GRPO runs on GSM8K-CoT). The finding recorded here is
therefore the sourced one: the peak-then-decay pattern is not specific
to GRPO, it appears in both arms of the controlled comparison without
any length inflation, and the project makes no claim that Dr.~GRPO is
uniquely free of it. The narrowest related statement the sources do
support is a length-drift one --- that Dr.~GRPO is the only algorithm
whose length keeps drifting upward after the initial compression phase
--- which is a different statistic and is not the verbosity-trap flag.

\subsection{Elevated analyses}\label{elevated-analyses}

Because the controlled 30--40-step null is a null, the pillar runs a
sequence of progressively more sensitive analyses to establish whether
the trap is present but invisible at that horizon. Four of these deserve
report, and all four are exploratory. They were run across iterations
12, 20 and 24 of the analysis log after the controlled null was known.
The seed counts are small: n = 5 arithmetic and n = 3 GSM8K seeds, and
three seeds admit only 10 distinct bootstrap resamples. Several of the
intervals resample autocorrelated steps or windows within runs rather
than seeds. The p-values below are uncorrected unless stated, and
roughly fifteen tests are reported across the three iterations.

A sliding-window trap-onset analysis with window W = 10 relaxes the flag
from ``peak before 65\% of training'' to ``local peak followed by decay
within a window'', and finds traps in 2 of 10 arithmetic runs (20\%), 4
of 6 GSM8K runs (67\%) and 6 of 16 runs overall (37.5\%) --- every one
of which is a local re-expansion rather than a terminal divergence. This
is the analysis that most nearly matches the roster-level trap rate. It
is a point estimate of trap prevalence in the controlled data, not an
upper bound (corrected; see Table E.1). The Wilson 95\% CIs are
{[}5.7\%, 51\%{]} for 2/10, {[}30\%, 90\%{]} for 4/6 and {[}18.5\%,
61\%{]} for 6/16, with runs as the unit.

A paired bootstrap on the Dr.~GRPO minus GRPO difference in rank
correlations gives, on arithmetic, Δρ(step, len) = +0.098 with an
interval of {[}+0.073, +0.119{]} --- excluding zero, which suggests that
Dr.~GRPO compresses length relative to GRPO on the easy task
(seed-paired bootstrap, n = 5); Δρ(step, reward) = +0.043 with an
interval of {[}−0.024, +0.112{]}, which is not significant; and Δρ(len,
reward) = +0.126 with an interval of {[}+0.045, +0.201{]}, again
excluding zero. On GSM8K the same three quantities are +0.267 {[}+0.164,
+0.419{]}, −0.020 {[}−0.069, +0.065{]} and −0.117 {[}−0.265, +0.015{]}.
The exploratory reading is that Dr.~GRPO's length attenuation is
task-dependent, and that on the harder task it deepens rather than
weakens the length--reward coupling.

A decile-binned coupling analysis on GSM8K shows the mechanism directly:
completions in the first length decile, spanning 153--180 tokens, earn a
mean reward of 0.351, while completions in the eighth decile, spanning
194--195 tokens, earn 0.233. Longer completions are rewarded less, which
is the opposite of the verbosity-trap prediction. A 100-step
cross-validation on the arithmetic metrics file returns ρ(step, len) =
+0.41 and ρ(step, reward) = +0.27, but the absolute length range across
that run is 0.08 tokens against a cap of 5 tokens --- a positive
correlation over a range too small to be a length effect.

The length--ZVF coupling analysis at iteration 12 is the most
mechanistically interesting of the elevated results, because its sign
flips between tasks. On arithmetic, ρ(len, zvf) is −0.449 {[}−0.584,
−0.324{]} for GRPO and −0.291 {[}−0.428, −0.139{]} for Dr.~GRPO; on
GSM8K it is +0.286 {[}+0.231, +0.348{]} and +0.187 {[}+0.148, +0.225{]}.
The Dr.~GRPO minus GRPO difference is +0.158 {[}+0.075, +0.239{]} on
arithmetic and −0.099 {[}−0.132, −0.041{]} on GSM8K. Conditioning ZVF on
length bins makes the arithmetic structure concrete: the shortest bin
has ZVF 0.92 and the longest 0.62, a drop of −0.30 {[}−0.34, −0.26{]},
whereas on GSM8K the corresponding values are 0.30 and 0.41, a rise of
+0.11 {[}+0.10, +0.13{]}. On the easy task, shorter completions are more
often all-same-reward; on the hard task, longer ones are.

The mechanism analysis at iteration 20 fits, per run, a regression of
step-level mean length on step-level mean reward. On arithmetic,
Dr.~GRPO reduces the slope magnitude: GRPO −3.65 ± 0.08 against Dr.~GRPO
−3.22 ± 0.15 (mean ± SD over 5 seeds). The paired differences below are
recomputed from
\longtok{platform\_hybrid/experiments/results/length\_\allowbreak{}bias\_\allowbreak{}mechanism\_\allowbreak{}per\_\allowbreak{}run.tsv}
as Dr.~GRPO minus GRPO, paired by seed, n = 5 (corrected; see Table
E.1). Each is given with a paired-t 95\% CI, the paired-t p and the
exact sign-flip permutation p; the permutation p cannot fall below
0.0625 at n = 5. The OLS slope difference is +0.43 {[}+0.18, +0.68{]},
paired t p = 0.009, permutation p = 0.0625. The Theil--Sen slope
difference is +1.11 {[}+0.33, +1.88{]}, p = 0.017, permutation p =
0.0625. The ρ(L, R) difference is +0.10 {[}−0.02, +0.22{]}, p = 0.080.
The lag-1 autocorrelation is essentially unchanged, −0.01 {[}−0.27,
+0.26{]}, p = 0.94. On GSM8K (n = 3) the OLS difference reverses sign,
−4.41 {[}−37.7, +28.9{]}, p = 0.63, far too wide to interpret. Across
the fourteen exploratory tests of iterations 20 and 24,
Benjamini--Hochberg correction leaves the arithmetic OLS and Theil--Sen
differences at q = 0.116 each, the smallest q in the family, so no test
in it survives correction (corrected; see Table E.1). A mediation
analysis attributes 42\% of the arithmetic effect to ZVF under GRPO and
55\% under Dr.~GRPO, against 28\% and 12\% on GSM8K.

The drift analysis at iteration 24 closes the sequence. Windowed ρ drift
on arithmetic is +0.326 for GRPO with p = 0.035, decomposing into an
early −0.330 and a late −0.003, against +0.187 for Dr.~GRPO with p =
0.399; on GSM8K the corresponding values are +0.019 with p = 0.905 and
−0.084 with p = 0.769. The sign-flip analysis of late windows (w = 10)
gives mean late non-negative fractions of 0.527 and 0.382 for arithmetic
GRPO and Dr.~GRPO and 0.048 and 0.000 for GSM8K GRPO and Dr.~GRPO.
Tested against one half with an exact two-sided sign-flip permutation
over seeds, the p-values are 0.688, 0.625, 0.250 and 0.250 (corrected;
see Table E.1). The unit is the seed, not the within-run window: each
seed contributes one fraction, so n = 5 on arithmetic and n = 3 on
GSM8K, where 0.250 is the smallest p the test can return. The GSM8K
Dr.~GRPO cell, in which not a single late window flips sign on any of
the three seeds, therefore reaches only that floor. The source's p
\textless{} 0.001 for it was a one-sample t-test on three identical
values, whose zero variance gives t = −∞. Recomputed over the same
fourteen exploratory tests, Benjamini--Hochberg leaves no test below
0.05: the smallest q is 0.116 (the arithmetic OLS and Theil--Sen
differences), and it stays 0.116 for every fourteen-test selection from
the sixteen iteration-20 and iteration-24 tests
(\longtok{platform\_modal/scripts/length\_\allowbreak{}bias\_\allowbreak{}iter24.py};
\longtok{platform\_hybrid/experiments/results/length\_\allowbreak{}bias\_\allowbreak{}iter24\_\allowbreak{}signflip.tsv}).
The iteration-24 drift p-values above are unchanged. The iteration-20
and iteration-24 analyses are therefore exploratory throughout, with no
multiplicity-surviving result.

\subsection{Held-out generalisation and the scale
extension}\label{held-out-generalisation-and-the-scale-extension}

The pillar's held-out results are the part of the chapter closest to a
conventional capability claim, and they are reported with corresponding
care. The only multi-seed matched base-versus-RL control on Qwen3-8B is
a five-seed comparison on Qwen3-8B-Instruct, which moves from 82.0\% to
83.3\% with p = 0.26. Stated as a mean improvement over the pre-RL
checkpoint, the gain is +0.013 with a paired p = 0.256 at n = 5 ---
small and not statistically significant, and a matched pre-RL control
for the Qwen3-8B-Base checkpoint is not available. In summary, on this
near-ceiling configuration, GRPO does not produce a held-out improvement
that can be separated from noise at the available seed budget, and that
GRPO and Dr.~GRPO produce held-out gains that are indistinguishable from
each other at that budget.

The lower-accuracy configuration behaves differently and more
conventionally. On Qwen2.5-1.5B-Instruct with GSM8K-CoT, both algorithms
do produce significant pre-to-post held-out gains: +6.2 percentage
points for GRPO and +5.0 for Dr.~GRPO (n = 3 seeds × 200 held-out items
per arm). The unit of analysis is the seed. A one-sample t on the three
per-seed deltas gives p = 0.0051 for GRPO (deltas +0.055, +0.060,
+0.070) and p = 0.0099 for Dr.~GRPO (+0.045, +0.060, +0.045). As
supporting detail, per-seed McNemar exact p-values are 0.071, 0.036 and
0.016 for GRPO and 0.122, 0.017 and 0.122 for Dr.~GRPO; no pooled
McNemar p-value is used, because the same 200 items recur across seeds
(corrected; see Table E.1). The McNemar detail also overstates what is
attributable to training, because the greedy pre-RL evaluation is not
deterministic: the six runs evaluate the same base model on the same 200
items before any training, yet their stored pre-RL correctness vectors
take two distinct forms, 40/200 and 41/200, which disagree on 7 of the
200 items, including the GRPO and Dr.~GRPO runs of the same seed
(recomputed from
\longtok{platform\_hybrid/experiments/results/drgrpo\_\allowbreak{}gsm8k\_\allowbreak{}cot\_\allowbreak{}full.json}).
About 3.5\% of items can therefore flip with no training at all, so some
discordant pairs are evaluation noise and each per-seed delta carries
roughly ±1 item of measurement noise. The seed-level t-tests above
absorb this noise in their between-seed variance; the per-seed McNemar
p-values do not, and are supporting detail only. The two algorithms'
gains do not differ (paired by seed, p = 0.25) (recomputed from
platform\_hybrid/experiments/results/drgrpo\_\allowbreak{}gsm8k\_\allowbreak{}cot.json; source:
platform\_hybrid/paper/sections/p4\_\allowbreak{}results\_\allowbreak{}intro.tex). This is the
project's clearest evidence that the pipeline trains something real; it
is reported alongside the near-ceiling null precisely because the
contrast between the two configurations is what the pillar's
bounded-null interpretation rests on. Across the same comparison, the
pooled mean completion length drifts from 194.4 to 183.5 tokens, that is
downward.

A scale extension tests whether the null survives when the cap is
lifted. Qwen3-8B with LoRA rank 4 was run on GSM8K for 30 steps at G = 8
and batch size 4, three seeds per algorithm, with the completion cap
raised fivefold to 1,024 tokens. All six arms compress: GRPO moves 1004
→ 905, 981 → 944 and 996 → 900 tokens, and Dr.~GRPO moves 999 → 931, 972
→ 902 and 1000 → 878 --- declines of between 3.8\% and 12.2\% in every
arm. Late ZVF is 0.45, 0.70 and 0.72 for Dr.~GRPO against 0.47, 0.47 and
0.55 for GRPO, and step-0 completions are approximately 98\% of the cap.
The step-0 saturation is the caveat that limits this extension: with
completions already at 98\% of the ceiling at the first step, the run
starts truncated, and length compression under a binding cap is not the
same phenomenon as length control in an uncapped regime. The pillar
reports the extension as consistent with the null, not as an independent
confirmation of it.

Finally, the P4-facing headline in the Miller-style error-bars audit is
the iteration-136 arithmetic late-training-efficiency contrast. Late
efficiency is training-reward gain per unit of length change. The
contrast is Dr.~GRPO minus GRPO, paired by seed, n = 5, with mean
difference +0.020. Recomputed from
\longtok{length\_\allowbreak{}bias\_\allowbreak{}iter136\_\allowbreak{}step\_\allowbreak{}coupling.tsv}, the effect size is
\(d_z\) = +2.68, the mean paired difference divided by the SD of the
differences. The exact one-sided Wilcoxon signed-rank p is 0.031, which
is 1/32, its floor at n = 5 (corrected; see Table E.1). The paired t
that matches \(d_z\) gives t(4) = 5.96, two-sided p = 0.004. The
two-sided exact sign-flip p is 0.0625, again the n = 5 floor: all five
pairs point the same way, yet the non-parametric test cannot reach 0.05
two-sided. The contrast was one of eight directional hypotheses in the
iteration-136 family (four per task). Under Holm correction across the
eight, the family's own Wilcoxon p becomes 0.25. This chapter therefore
grades it \emph{suggestive}, not decisive (corrected; see Table E.1): a
consistent direction across five seeds on a training metric, with
two-sided non-parametric p = 0.0625 at n = 5. The audit's three NULL
verdicts fall on Pillar-1 and Pillar-3 headlines rather than on this
one. The pillar also records its own limitations plainly: a
single-benchmark held-out evaluation, a single-seed top-10-checkpoint
frontier comparison, weak length identifiability, and marginal paired
tests.

\subsection{Same-stack rerun on the unsaturated
configuration}\label{same-stack-rerun-on-the-unsaturated-configuration}

The same-stack PPO/GRPO and group-size nulls of §6.2.3 and §6.4.1 sit at
evaluation accuracy near 0.99, where a null cannot separate no effect
from no headroom. To test the contrasts where headroom exists, they were
rerun on 2 October 2026 on the configuration of §6.5.4:
Qwen2.5-1.5B-Instruct, GSM8K with a boxed-answer correctness reward and
a 200-token chain-of-thought cap, LoRA rank 16 on the attention
projections, learning rate 1 × 10⁻⁵, 30 steps, two inner epochs of the
clipped surrogate with per-response length normalisation, on one A10G
per run. Three arms share that stack and a fixed count of 64 generations
per step, but differ in advantage estimation, critic training and
distinct-prompt exposure: GRPO with G = 8 (8 prompts × 8 completions),
GRPO with G = 2 (32 × 2), and PPO with a learned value head on the
prompt-end hidden state (64 × 1, value loss 0.5 × MSE, batch-normalised
advantage). Each run evaluates the same 200 GSM8K \emph{test} items
greedily before and after training, so the evaluation set is disjoint
from the training split. Five paired seeds were run per arm; the unit is
the seed, and the analysis, including a ±2-point equivalence margin, was
fixed in the script before launch.

\textbf{Same-stack rerun on Qwen2.5-1.5B-Instruct GSM8K-CoT, five paired
seeds, 200 test items.}

\begingroup\scriptsize\setlength{\tabcolsep}{3pt}
\begin{longtable}[]{@{}
  >{\raggedright\arraybackslash\hspace{0pt}}p{(\linewidth - 8\tabcolsep) * \real{0.2500}}
  >{\raggedright\arraybackslash\hspace{0pt}}p{(\linewidth - 8\tabcolsep) * \real{0.2500}}
  >{\raggedright\arraybackslash\hspace{0pt}}p{(\linewidth - 8\tabcolsep) * \real{0.2500}}
  >{\raggedright\arraybackslash\hspace{0pt}}p{(\linewidth - 8\tabcolsep) * \real{0.2500}}@{}}
\caption{Same-stack rerun on Qwen2.5-1.5B-Instruct GSM8K-CoT;
paired-\(t\) intervals over five seeds. Pre-training accuracy is
identical across seeds because evaluation is greedy on a fixed item
set.}\tabularnewline
\toprule\noalign{}
\begin{minipage}[b]{\linewidth}\raggedright
Arm
\end{minipage} & \begin{minipage}[b]{\linewidth}\raggedright
Accuracy before → after
\end{minipage} & \begin{minipage}[b]{\linewidth}\raggedright
Change (95\% CI)
\end{minipage} & \begin{minipage}[b]{\linewidth}\raggedright
Last-10 training reward
\end{minipage} \\
\midrule\noalign{}
\endfirsthead
\toprule\noalign{}
\begin{minipage}[b]{\linewidth}\raggedright
Arm
\end{minipage} & \begin{minipage}[b]{\linewidth}\raggedright
Accuracy before → after
\end{minipage} & \begin{minipage}[b]{\linewidth}\raggedright
Change (95\% CI)
\end{minipage} & \begin{minipage}[b]{\linewidth}\raggedright
Last-10 training reward
\end{minipage} \\
\midrule\noalign{}
\endhead
\bottomrule\noalign{}
\endlastfoot
GRPO, G = 8 & 0.200 → 0.245 & +0.045 {[}+0.033, +0.057{]} & 0.270 \\
GRPO, G = 2 & 0.200 → 0.250 & +0.050 {[}+0.026, +0.074{]} & 0.275 \\
PPO, learned value head & 0.200 → 0.195 & −0.005 {[}−0.033, +0.023{]} &
0.196 \\
\end{longtable}
\endgroup

Both GRPO arms improve on every seed, which reproduces the +6.2-point
GRPO gain of §6.5.4 at a similar size. The PPO arm shows no detectable
held-out gain: its training reward drifts down and its post-training
accuracy is indistinguishable from the starting point. After training,
GRPO at G = 8 exceeds PPO by +0.050 {[}+0.015, +0.085{]} (paired \(t\),
p = 0.016), and the difference is positive for all five seeds, which
puts the exact two-sided sign-flip p at its n = 5 floor of 0.0625. GRPO
at G = 8 and at G = 2 differ by −0.005 {[}−0.038, +0.028{]} (p = 0.69);
that interval is wider than the ±0.02 margin, so equivalence is not
established either (TOST p = 0.13).

The PPO result has a narrow reading. The PPO arm is the minimal one that
matches the GRPO arms' stack: a linear value head trained from zero for
30 steps at the policy's learning rate, with no critic warm-up, no
generalised advantage estimation and no tuning. Its advantage is
therefore close to batch-centred reward with a noisy baseline, so what
the contrast shows is that on this fixed stack and budget the
group-relative baseline shows a held-out gain and this minimal
learned-critic baseline shows no detectable gain. It is a
single-configuration result with five seeds, not a verdict on PPO with a
tuned critic, and the 64-prompt PPO batch also sees eight times as many
distinct prompts per step as GRPO at G = 8. The paired-t result depends
on a small-sample parametric approximation; the exact two-sided
sign-flip result does not reach 0.05. Both describe seed variation on
the same fixed 200-item evaluation slice, not uncertainty over GSM8K
items or model/task families. Fixed generation counts are not measured
equality of generated tokens or total compute. Mean training completions
are 190.16 (G8), 191.62 (G2) and 194.16 (PPO) tokens under cap 200, with
PPO last-ten length 196.19; non-ceiling accuracy does not establish a
non-binding length regime. For group size, the rerun removes the
accuracy-saturation objection and still finds no detectable effect
between G = 2 and G = 8 at this budget, without establishing that none
exists.

\section{Cross-Pillar Observations and
Limitations}\label{cross-pillar-observations-and-limitations}

Three observations span the four pillars and are worth stating together,
because the chapter's overall contribution is better captured by their
pattern than by any single number.

First, the project's strongest results are its de-confounds, and its
de-confounds mostly return nulls. The matched same-stack PPO-versus-GRPO
contrast returns a seed-paired difference of −0.002 with p = 0.374 (n =
5 seeds, arms evaluated on largely different items, at a near-ceiling
accuracy, not shown equivalent by TOST at ±0.005); the four-axis
factorization audit attributes η² = 0.023 to the algorithm axis and η² =
0.54 of last-10 training-reward variance to the group-size axis (F(3, 8)
= 3.08, p = 0.09, n = 12 runs); the controlled length-bias comparison
returns zero length-bias flags across sixteen runs. In each case the
stack was held fixed except the factor under test (and, in the PPO arm,
the number of distinct prompts per step), and in each case the expected
algorithmic advantage failed to appear. Each of these comparisons was
run on the saturated 0.5B arithmetic task, at evaluation accuracy
0.98--0.99. The group-size axis carries the largest variance share, but
at n = 3 seeds per arm it does not reach significance on either training
reward or held-out accuracy. It is also a hyperparameter rather than an
algorithm choice. The chapter's empirical conclusion is therefore
conditional on saturation: on a 0.5B arithmetic task at evaluation
accuracy near 0.99, no algorithm-label effect was detectable, which
cannot distinguish ``no effect'' from ``no headroom''; the sampling
configuration is the more plausible lever but was not shown to matter at
this seed budget. Rerun on the unsaturated Qwen2.5-1.5B GSM8K-CoT
configuration (§6.5.5), the picture changes for one contrast: GRPO at G
= 8 exceeds a minimal learned-critic PPO by +0.050 {[}+0.015, +0.085{]}
(paired-t p = 0.016; exact sign-flip p = 0.0625), while the minimal PPO
arm shows no detectable gain on this budget; G = 2 and G = 8 remain
indistinguishable without being shown equivalent.

Second, the diagnostics and interventions are less prescriptive than
their motivating literature suggests, and the project says so. ZVF
measures signal starvation reliably. It associates with catastrophic
collapse (ρ = 0.56, 95\% CI {[}0.19, 0.79{]}, n = 23 cells), and its
association with graded held-out accuracy is not distinguishable from
zero (ρ = 0.27, 95\% CI {[}−0.37, 0.88{]}, n = 23), because it is
unsigned and therefore cannot distinguish mastery from incapacity. The
four interventions built on it --- the difficulty curriculum, the
group-size sweet spot, the length-bias loss formulations and the step-1
layer freeze --- all return nulls at matched cost, and two of them (the
curriculum and the G = 4 win) were flagged as SUSPECT by an independent
verification pass before multi-seed re-testing reclassified them as
noise. Reporting an intervention as a null after an independent reviewer
flagged it as overclaimed is, in this project's convention, part of the
result.

Third, the single-seed regime is the dominant shared limitation. The
cross-scale anchors are single-seed and descriptive. The roster-level
verbosity-trap audit is single-seed per configuration, and its
in-training reward metric is not a held-out one. Only the P3 sweep, the
P4 controlled comparison, the five-seed Qwen3-8B held-out control and
the P4 scale extension operate at n ≥ 3, and the Qwen3-8B control is the
only multi-seed matched base-versus-RL comparison on Qwen3-8B; the
Qwen2.5-1.5B GSM8K-CoT pre/post runs and the small-sample pairs of Table
8.B are the other base-versus-trained comparisons. Every headline in
this chapter is therefore stated at the strength that its own evidence
base supports, which is in several cases weaker than the strength at
which it would need to be stated to generalise beyond the named stack.
The reconstructions (G = 32 and above in P3, the T = 64M grid, the
difficulty-gated saving projection) are held at hypothesis-generator
status for the same reason: they are joint-fit extrapolations, not
trained arms, and the project declines to let them carry a
recommendation.

Taken together, the four pillars establish no general algorithm ranking:
the unsaturated PPO contrast is baseline-specific and test-sensitive,
while the group-size contrasts remain unresolved at these seed budgets.
The method's most widely repeated specific failure mode --- length
inflation --- is absent in the regimes where it was tested directly and
is not established by the roster flag, which is not specific to length.
That is a narrower and more cautious set of conclusions than the P1--P4
internal reports set out to establish, and the chapter reports it as
such.



\chapter{Results II: Reporting, Registry and Controller
(P5--P7)}\label{results-ii-reporting-registry-and-controller-p5p7}

The three internal reports covered in this chapter form a chain rather
than a set of independent results. P5 makes a claim about the epistemic
status of reinforcement-learning results for language models and
proposes a checkable reporting standard to match it. P6 turns the
standard into a machine-readable artifact with a validator, a badge and
a diff tool. P7 takes the diagnostic that the standard requires authors
to publish --- the zero-variance fraction (ZVF) --- and asks whether it
is mechanistically predictable and ready to control rollout allocation.
Each report is explicit about where its claim stops, and this chapter
reports their evidence in the same register.

\section{A Minimum Reporting Standard for Stack-Conditioned Results
(P5)}\label{a-minimum-reporting-standard-for-stack-conditioned-results-p5}

P5's position is that results in this field are
\emph{stack-conditioned}, and that reporting the framework or algorithm
label without the stack is therefore insufficient to support the claims
that are routinely made from it. The evidence for that position is not a
single demonstration but a family of levers, each showing that varying
one component of the stack, with the label held fixed, moves or flips
the reported outcome or its telemetry.

The cleanest lever is a managed 12-cell head-to-head: four
group-relative labels (GRPO \citep{shao2024deepseekmath}, DAPO
\citep{yu2025dapo}, Dr.~GRPO \citep{liu2025understanding}, GSPO
\citep{zheng2025gspo}) run on one stack --- Tinker-managed, Qwen3.5-4B,
GSM8K, three seeds per arm. Mean last-10 training reward spans 0.710 to
0.744, a band of 0.034, with per-arm values of 0.744 (DAPO), 0.742
(Dr.~GRPO), 0.723 (GRPO) and 0.710 (GSPO); mean ZVF over the same arms
is 0.578, 0.500, 0.567 and 0.511. The labels move telemetry by more than
they move outcome, and neither movement is large.

The largest lever in the corpus is an outcome flip of 5.0\% to 85.6\%
final training reward under a fixed label when the backend is changed.
That figure must be read with its provenance caveat attached, because
the same inspection that surfaced it also found that the nominal backend
comparison bundled a different base checkpoint: Qwen3-8B-Base on the
managed Tinker run against the instruction-tuned Qwen3-8B on TRL. The
number is therefore evidence that item 3 of the standard must record
backend \emph{and} checkpoint identity; it is not a causal estimate of
the backend alone. A second same-label lever points the other way: the
label ``DAPO'' produced mean ZVF 0.00 on an open trainer running true
dynamic sampling and 0.58 on a closed stack running an asymmetric-clip
surrogate.

These exhibits are summarised by a variance decomposition, and its
design has to be read with it. The algorithm term comes from the
four-method same-stack N2 run (GRPO, AERO, GIFT, AREAL; one seed; 40
steps each, so 4 groups of 40 per-step observations, n = 160, df = 3
between and 156 within): η² = 0.0454 on ZVF, with a within-method
(stratified) step-bootstrap CI of {[}0.0094, 0.1476{]}, and η² = 0.0075
{[}0.0019, 0.0790{]} on mean reward (corrected; see Table E.1). The
source interval was drawn from the pooled steps, so it described a null
band rather than the estimate's uncertainty. The corrected ZVF interval
therefore does not support the reading that the algorithm explains at
most about 6\% of ZVF variance: an algorithm share of up to roughly 15\%
is compatible with these data, and the source's DECISIVE grade on this
term falls to SUGGESTIVE. Even the corrected intervals are too narrow,
because the steps are autocorrelated within a run and those 156
within-group degrees of freedom overstate the independent information
(recomputed by
\longtok{platform\_modal/scripts/p5p8/p5\_\allowbreak{}headline\_\allowbreak{}cis\_\allowbreak{}full.py}; output
\longtok{platform\_hybrid/experiments/results/p5p8/p5\_\allowbreak{}headline\_\allowbreak{}cis\_\allowbreak{}full.tsv}).
No group-size term is reported from this decomposition: the group-size
file, \longtok{groupsize\_\allowbreak{}zvf\_\allowbreak{}sweep.tsv}, holds one row per G level (4
cells, one observation per cell, df = 0 within), on which any one-way
decomposition returns η² = 1 by construction (corrected; see Table E.1).
The 98-cell mega-campaign decomposition is an unbalanced partial grid
(98 of 120 cells over model family, task slice, G, temperature and seed,
with two seeds). Its four stack axes account for 72.9\% to 92.7\% of the
variance in the five logged metrics; the seed axis rounds to 0.000, so
no stack-to-seed ratio is quoted as an effect size (corrected; see Table
E.1).

The standard that P5 proposes has eight items --- seven run-manifest
fields and one evaluation item --- and each item earned its place by
exhibiting a documented lever of this kind. Item 1 is the loss form,
reported as a tuple of ratio level, clip low and high, token mask,
advantage normalisation and dynamic-sampling status --- the same label
with and without dynamic sampling produced opposite ZVF telemetry (0.00
against 0.58). Item 2 is the reference policy and KL handling, which
distinguishes family members that are otherwise conflated and whose
omission is unrecoverable after the fact. Item 3 is the sampler, backend
and precision, including base-checkpoint identity, motivated by the 17×
exhibit above. Item 4 is the per-step ZVF and gradient-utilisation
trajectory, because a single mean aliases opposite regimes: across a
368-run audit drawn from a 790-run corpus, mean ZVF is 0.97 and 0.95 for
two low-reward models (Llama-3.2-1B and Llama-3.2-3B), falls to 0.25 and
0.29 at mid-range rewards, and rises again to 0.87 and 0.96 for two
near-saturated models. Item 5 is the group-size schedule, since group
size is a signal-geometry knob rather than a throughput knob; measured
mean ZVF for Qwen3-32B falls from 0.33 at G=4 to 0.18 at G=8 and 0.08 at
G=16, and for Llama-3.1-8B from 0.76 to 0.65 to 0.47, while the two
boundary cases compress the effect as the theory predicts. Item 6
requires a held-out split disjoint from the reward environment. Item 7
requires a contamination check and a parser-robustness probe: adding
micro-jitter δ \textasciitilde{} U(0,10⁻⁴) collapsed the
\emph{exact-tie} fraction from 95/600 (0.1583) to 0/600 while pairwise
contrast density changed negligibly. Recalculation with the declared
sample-variance tolerance ε = 10⁻⁶ instead leaves ZVF at 95/600 before
and after jitter (§6.3.1). The example demonstrates sensitivity of exact
ties to rounding, not failure of the tolerance-based diagnostic.

Item 8 is deliberately different in kind: held-out pass@k curves at k ∈
\{1, 8, 32\} belong to the evaluation report rather than to the run
manifest. The lever here is a matched base-versus-post-RL panel --- four
adapters at G ∈ \{2, 4, 16, 32\}, one stack, the same problems --- on
which pass@1 deltas scatter from −0.5 to +1.9 points (under which
reading the G=2 arm would be called a regression) while the pass@32
frontier improves by +1.5 to +2.5 points in all four arms. Neither
reading reaches significance at a single training seed, which is itself
the argument: without pass@k curves and stated noise bands the same
panel supports any narrative the reader prefers.

P5 also ships the machinery that makes the standard enforceable rather
than advisory. An auditor scores a manifest against the eight items and
emits a 0--100 weighted badge; on an n=103 audit the weighted total was
55.0 with a badge span of {[}35.8, 56.7{]} and standard deviation 3.1,
awarding zero gold, zero silver, 99 bronze and 4 wood. Coverage,
however, is where the negative result lies. On 98 manifests, every one
of the seven manifest fields is present as a key 100\% of the time, yet
item 2 is \emph{validated} in 0 of 98 cases (0.0\%), and all twelve of
its sub-fields are present at 0 in each case; nine of the sub-fields
across items 1 and 2 --- advantage normalisation, clip high, clip low,
dynamic sampling, KL coefficient, KL estimator, ratio level, reference
snapshot and token mask --- sit at 0.0\% with zero information content.
Twenty-three sub-fields examined across the population show no field at
a full null rate, but twenty of the twenty-three carry at most two
distinct values, and the twenty entries collapse to ten distinct 23-bit
fingerprints.

Two further analyses bound what the standard can be expected to buy. The
first is a discriminative-entropy audit: four of the seven manifest
fields are vacuous on this corpus --- loss form, reference policy and
KL, sampler/backend/precision, and decontamination notes each carry too
little task-discriminative information to distinguish runs --- while the
per-step ZVF path and the group-size schedule carry the information. A
field-sufficiency regression agrees: a random forest on the disclosed
stack fields, scored by out-of-fold R² over the 98 cells (8-fold
cross-validation, 8 fold seeds), predicts ZVF at R² = 0.832 {[}0.772,
0.879{]} and mean reward at R² = 0.993 {[}0.990, 0.996{]}. The
label-only comparator is R² = 0.000 by construction, because all 98
cells share one algorithm label, so it is a floor rather than a
competing model. The second is a manifest self-sufficiency test, and it
splits sharply by metric. The manifest carries essentially all of the
ZVF signal (R² = 0.776 {[}0.730, 0.846{]} against 0.773 {[}0.735,
0.834{]} for the cells, ΔR² = −0.003 {[}−0.045, +0.034{]}) but only a
minority of the reward signal (R² = 0.280 {[}0.219, 0.443{]} against
0.743 {[}0.688, 0.820{]}, ΔR² = +0.463 {[}+0.309, +0.564{]}); all 98
cells are therefore labelled \longtok{honest\_\allowbreak{}but\_\allowbreak{}vacuous} on the
three-axis triangulation.

The most useful negative in P5 is the placebo-replacement bound. The
four-of-seven vacuity pattern admits two readings: the schema chose the
wrong items, or any stack-descriptor schema would be vacuous on a
single-stack corpus. Projecting six candidate replacement items onto the
live corpus refutes the first reading. Four of the six --- KL
coefficient, clip range low, advantage normalisation and mini-batch size
--- project to exactly one unique value on the 98-cell corpus with a
bootstrap CI excluding any positive information gain, because the corpus
never varies them; temperature schedule would add 0.995 bits and lift
the total from 11.41 to 12.41 bits, and a model-family label is
redundant with an existing column. The variance partition confirms it:
stack axes explain η² = 0.947 of ZVF variance across 25 unique buckets,
with 5.3\% seed residual and an algorithm term of zero by construction.
The schema is not the binding constraint; the corpus is. On the live
corpus the fingerprint over the seven manifest fields carries 11.4 bits
in total, four fields carry exactly zero, and Hamming distance over the
fingerprint is associated with the absolute difference in ZVF (Spearman
ρ = +0.44), PCD (ρ = +0.39) and mean reward (ρ = +0.28) across all 4,753
cell pairs, while the badge itself is constant across cells and cannot
discriminate seed-pair dispersion (ρ = +0.000, n = 48 pairs). Because
pairs that share a cell are not independent, significance comes from a
Mantel test that permutes cell labels jointly on the distance matrix
(9,999 permutations): p \textless{} 10⁻³ for all three outcomes, and the
association holds when seed replicates are averaged into 50
configurations (ρ = +0.47, +0.43, +0.33). The Hamming distance takes
only four values and is in effect a function of group size and task, so
the test shows that cells differing on those fields differ in outcome,
not that the fingerprint adds information beyond them. The source
paper's ZVF and PCD values (ρ = +0.529 and +0.541) could not be
reproduced: its own analysis script gives +0.407 and +0.336 on a
2,000-pair sample, and the full-matrix values are those above
(corrected; see Table E.1).

\section{The GRPO-Registry (P6)}\label{the-grpo-registry-p6}

P6's premise is that the deltas defining GRPO variants --- DAPO's
asymmetric clip, dynamic sampling and token-level loss; Dr.~GRPO's
dropped length and standard-deviation normalisation; GSPO's
sequence-level importance ratio \citep{zheng2025gspo} --- live in paper
appendices and framework defaults rather than in any machine-readable
form. The registry is the response, and it ships as three artifacts at
\texttt{registry/}: \texttt{schema.json}, \texttt{entries/*.json} and
\texttt{query.py}.

The schema defines two record types under a single \texttt{oneOf}, both
carrying \texttt{schema\_version}. The stack record's core is
\texttt{min\_report}, one sub-object per minimum-report item, with
base-checkpoint identity required inside the \texttt{sampler\_backend}
block precisely because the audit that motivated item 3 found a managed
backend silently pinning a different base checkpoint. Two design
decisions carry weight. The first is the distinction between an
unreported field and a field reported as absent: a null means the author
said nothing, whereas an explicit \texttt{false}, \texttt{0.0} or
\texttt{"none"} means the author made a checkable claim --- ``a stack
that says `we use no KL term' is making a checkable claim; a stack that
says nothing is not.'' The second is that item 8, the held-out pass@k
requirement, is attached under outcomes as an evaluation-report field
rather than folded into the run manifest, preserving the eight-item
standard's split between what describes a run and what describes a
result. Per-component delta status is recorded on the stack side in
\longtok{variant\_\allowbreak{}deltas\_\allowbreak{}applied}, with each entry's status drawn from
\{implemented, surrogate, absent, unknown\}. The shipped entry for the
managed DAPO run makes the point concretely: \texttt{clip\_higher} is
recorded as \texttt{surrogate} with the note that the asymmetric clip
was set through the user-facing config while the underlying managed loss
form is unverifiable, \texttt{dynamic\_sampling} is recorded as
\texttt{absent}, and \longtok{token\_\allowbreak{}level\_\allowbreak{}loss} is \texttt{unknown}
because it is managed by the runtime. An optional \texttt{ci\_method}
object (method, bootstrap count, seed, confidence level, source) is
populated on the seven derived entries with \texttt{n\_boot\ =\ 2000}
and a percentile 95\% interval, and left null on the remaining
twenty-four.

Three variant-delta records were verified against their source papers:
DAPO with five components (clip higher at ε\_low = 0.2 and ε\_high =
0.28, dynamic sampling, token-level loss, overlong reward shaping, KL
removed); Dr.~GRPO with two removals (the 1/\textbar o\_i\textbar{}
response-length normalisation in the loss and the division by the
within-group reward standard deviation in the advantage); and GSPO with
a length-normalised sequence-likelihood importance ratio and
sequence-level clipping.

The seed population is twenty stack records in five families, all of
which pass validation. Family A holds four framework config dumps
(Qwen3-8B, GSM8K, seed 42) scoring badges 74, 60, 60 and 62; family B
holds the four open Colab trainer arms at toy scale, scoring 96, 96, 96
and 96; family C is the managed head-to-head at 57, 62, 62 and 60;
family D is a same-stack four-method run at 67 each; and family E is a
five-seed risk-index batch at 43 each. The finding worth stating plainly
is why the ordering looks the way it does. The config dumps were
assembled for a reviewer-requested fairness audit and score only 57 to
74: they nail items 2, 3 and 5, expose item 1 only partially, and cover
items 6 and 7 not at all. The open Colab arms score 96 not because the
runs are better but because an open backward pass, coupled with explicit
declarations that no KL term and no decontamination check were used,
makes every field reportable. Reproducibility supplements of the kind
commonly shipped with RL papers do not automatically satisfy MIN-REPORT.

The health audit is the registry's own evidence that it works. Every
record shipped with P6 parses: 31 of 31 pass, comprising 20 stack
records and 11 variant-delta records, with zero broken delta pointers.
The delta-mix table is the sharpest single result in P6: the label
``DAPO'' is claimed by ten entries, of which four implement a delta, one
records a surrogate, three record it as absent and two record it as
unknown --- the same label covering four distinct implementations (Table
\texttt{tab:p6-delta-mix}). Per-item coverage on the twenty stack
records is 100\% for the held-out split (40/40), 95\% for
sampler/backend (76/80), 90\% for telemetry (54/60), 68\% for the
group-size schedule (41/60), 48\% for reference policy and KL (29/60),
32\% for the loss form (38/120) and 20\% for decontamination (8/40)
(Table \texttt{tab:p6-coverage}). The per-openness breakdown inverts the
conventional expectation: the open subset records sampler/backend at
91.7\% and telemetry at 83.3\% with reference KL at 58.3\%, while the
managed subset records sampler/backend and telemetry at 100\% but
reference KL at 33.3\% and the loss form at only 22.9\%.

The registry was stress-tested rather than asserted. A batch of
adversarial perturbations with B = 4000 bootstrap resamples --- 13
categories × 31 entries × 10 mutations = 4,030 generated, of which 800
were no-ops for their entry and were skipped, leaving 3,230 applied ---
moved the overall robustness statistic from 0.864 {[}0.756, 0.878{]}
before the corresponding schema bumps to 1.000 {[}1.000, 1.000{]} after
them, closing three silent-failure categories: a missing identifier
pattern check (0.645 to 1.000), MIN-REPORT leaves declared in
\texttt{properties} but not in \texttt{required} (0.000 to 1.000, fixed
by promoting key presence to a schema invariant), and free-form delta
references (0.000 to 1.000, fixed by replacing \texttt{delta\_id} with
an enum generated from the entry filenames). A separate consistency
check cross-references the variant-delta records against the MIN-REPORT
blocks: of 31 triples from 13 entries, 7 are MATCH, 0 are MISMATCH, 0
are MISSING\_REPORT, 5 are SURROGATE\_OBS and 19 are NOT\_APPLICABLE,
giving an implementation-honesty match rate of 7/7 = 1.000 with a Wilson
95\% CI of {[}0.65, 1.00{]} (a percentile bootstrap of 7/7 is degenerate
at {[}1, 1{]}).

The empirical payoff of the schema is the same-stack measurement.
Re-executing four GRPO-family methods (GRPO, AERO, GIFT, AREAL) on a
single identical stack --- managed sampler, G = 8, seed 0, 40 steps, 16
prompts per step --- with the full per-prompt reward tensor logged at
every step, produces reward deltas of magnitude ≤ 2 percentage points.
Two of the three show within-run intervals that exclude zero on the side
of a \emph{decrease} (AERO at −0.014 {[}−0.023, −0.005{]} and AREAL at
−0.020 {[}−0.032, −0.008{]}). Those intervals are paired bootstrap
intervals over the last 10 of 40 steps of a single seed-0 run per
method, so they carry no seed variance and are not replication. GIFT's
delta is +0.016 {[}−0.007, +0.040{]}. Exact sign-flip tests over the
same 10 paired steps, Holm-adjusted across the three contrasts, give p =
0.047 for AERO, 0.047 for AREAL and 0.23 for GIFT; when prompt-level
variation is included (160 paired prompt-steps), the AERO decrease is no
longer distinguishable from zero (Holm p = 0.20) (corrected; see Table
E.1). The magnitude thesis survives. Whether any difference exceeds
seed-to-seed noise was not tested, so neither the claim that all
differences are within noise nor its negation is established. The
largest same-stack effects are not on reward at all: AERO and AREAL
lengthen completions by approximately 31 tokens and tighten the length
coefficient of variation, and GIFT's loss shift of −1.67×10⁴ is a
baseline constant. A later registry extension measures the anti-herding
residual
\(\delta_{\mathrm{div}} = \mathrm{ZVF}^{\mathrm{iid}} - \mathrm{ZVF}^{\mathrm{obs}}\)
on the same tensors, reporting a stack-mean of 0.049 with a range of
{[}0.039, 0.053{]}; only GIFT shows a borderline-significant reduction
against GRPO (p = 0.066, CI {[}−0.021, +0.001{]}), which resolves a
synthesis range of {[}0.13, 0.23{]} downwards --- the anti-herding
diversity bonus is real but smaller than the qualitative synthesis had
guessed.

The reference CLI, \texttt{registry/query.py}, implements four verbs
using only the standard library: \texttt{list};
\texttt{query\ -\/-item\ reference\_kl} returning FULL, PARTIAL or
UNREPORTED; \texttt{badge}; and \texttt{stackdiff}, which maps a pair of
entries onto the R0--R5 ladder and reports the highest rung triggered
--- R0 identical manifests, R1 nuisance-only, R2 runtime differences, R3
loss-form or KL-handling differences, R4 differing variant-delta sets,
R5 unauditable because a differing item is null on at least one side.
The badge formula weights the seven items equally and is outcome-blind:
\(\mathrm{badge}(s) = \lfloor 100 \cdot \tfrac{1}{7} \sum_i c_i(s) \rfloor\).
\texttt{stackdiff} prints a warning at R4 and above that outcome
differences must not be attributed to the shared algorithm label, and
the P6 report is careful to describe the ladder as ``a risk scale, not
an effect-size scale.'' Applied to the shipped entries,
\texttt{stackdiff} flags the open-versus-closed ``DAPO'' pair as R5 from
manifests alone, and across the twenty seed entries the MIN-REPORT badge
spans 43--96. Two boundaries deserve restating. The registry reports
stated provenance and conflicts; it does not reconcile unresolved
exact-ID or model conflicts, nor does it validate nominal algorithm
comparisons. And the badge is coarse relative to the corpus --- twenty
entries collapse to ten distinct fingerprints, and twenty of
twenty-three sub-fields carry at most two values.

\section{From Diagnostic to Controller
(P7)}\label{from-diagnostic-to-controller-p7}

P7 asks a narrower question than P5 and P6. Given that the zero-variance
fraction is worth reporting, is it mechanistically predictable, and is
it ready to control rollout allocation? The P7 report's own framing is
that this does not overturn the descriptive restraint of the earlier
work --- it \emph{tests} ZVF's predictive value through falsifiable
predictions rather than assuming it --- and it flags at the outset that
the weakest link in the chain, the causal connection from within-group
contrast to gradient magnitude, rests on toy-scale directional evidence
only.

The theory is compact. Fix a prompt with latent per-rollout success
probability p and a group of G conditionally independent Bernoulli
rewards. Write \(h_G(p) = p^G + (1-p)^G\) for the probability that the
group is degenerate, and model the usable signal per prompt as
\(S(p,G) = p(1-p)(1-h_G(p))\). The first factor is the contrast a
non-degenerate group can carry; the second is the probability that any
signal survives centring. Both factors vanish at the endpoints and peak
at p = ½, so this proposed signal proxy is largest at intermediate
latent success probability and vanishes at the endpoints; it is not a
measured law of learning speed. Three consequences follow, each paired
with the measurement that could have falsified it. T1 states that the
expected ZVF indicator is exactly \(h_G(p)\). T2 states that \(h_G\) is
strictly decreasing in G for 0 \textless{} p \textless{} 1, with
exponential decay rate −ln max(p, 1−p). T3 states that with unnormalised
group-centred advantages the mean absolute advantage in a group is
exactly 2\ensuremath{\hat{p}}(1−\ensuremath{\hat{p}}).

The checks have different evidential status. Binning by the same group's
empirical success rate is an algebraic identity, not an independent
validation of latent-p predictions. For T1, binning the per-group reward
tensors of the GSM8K runs by empirical success rate gives a ZVF
indicator averaging 1.000 at \ensuremath{\hat{p}} ∈ \{0, 1\} and 0.000 at every interior
bin; at run level, measured ZVF sits \emph{above} \(h_G(\bar p)\) at
every G, which the P7 report attributes to the Jensen gap --- the
binomial model evaluates \(h_G\) at the mean \(p\) while the prompt
population can mix different latent success probabilities. The Jensen
gap explains only part of it. On the repository sweep the measured ZVF
also sits above the per-step-average null, which already absorbs the
across-step variation in p: 0.838 against 0.828 at G = 2, 0.764 against
0.709, 0.691 against 0.597 and 0.631 against 0.522 at G = 16, recomputed
from \longtok{groupsize\_\allowbreak{}zvf\_\allowbreak{}sweep.json}. This is the same baseline used
in Chapter 6 §6.3.3 and Chapter 4 §4.3. For T2, two independent grids
agree. The 368-run audit shows monotone decline for every model with an
interior success rate (Qwen3-32B 0.33 → 0.18 → 0.08; Llama-3.1-8B 0.76 →
0.65 → 0.47), with the two boundary rows behaving as the theory flags
them: Llama-3.2-1B sits near p ≈ 0 where \(h_G \approx 1\) regardless of
G (0.98 → 0.97), and Qwen3-4B-Instruct sits near saturation where the
decay rate is tiny and seed noise dominates, producing the single
non-monotone cell (0.89 → 0.84 → 0.91). The repository-internal sweep at
G ∈ \{2, 4, 8, 16\} shows the same decline, with mean ZVF falling from
0.838 to 0.631 while held-out accuracy stays essentially flat near 0.98
throughout --- raising G buys signal, not, on this near-saturated task,
outcome. T3 cannot be checked on the closed stack at all, because it
does not expose per-group gradients; the open-backward-pass measurement,
corr(‖∇‖, \ensuremath{\hat{p}}(1−\ensuremath{\hat{p}})) = +0.71 on a 0.5B model with synthetic arithmetic, is
reported strictly as directional evidence.

Because ZVF aliases mastery with incapacity, the P7 report promotes the
pairwise contrast density PCD = ((G−1)/G)·E{[}p(1−p){]} as the signal to
control on. The micro-jitter experiment in §7.1 motivates recording the
metric definition: the exact-tie fraction falls from 0.158 to 0.000
while PCD is nearly unchanged, but \(\mathrm{ZVF}_{10^{-6}}\) remains
0.158. This experiment alone does not establish that PCD is a better
control signal than tolerance-based ZVF. The controller's design follows
from the theory: a \texttt{zvf-triage} callback watches per-step ZVF
with PCD as an aliasing guard, and an adaptive-G controller escalates
group size on ZVF spikes, changing only the sampler budget so that it
composes with any group-relative objective.

The four-arm open-trainer audit is the first place where the
intervention is measured against alternatives on the same data, budget
and seeds. Against a 120-rollout baseline, GRPO gains +0.500 held-out at
mean ZVF 0.25; Dr.~GRPO gains +0.575 at 0.27 on the same 120 rollouts;
DAPO-style dynamic sampling gains +0.550 while driving ZVF to 0.00 at
174 rollouts, a +45\% cost; and adaptive-G gains +0.575 while lowering
mean ZVF only modestly to 0.23 at 186 rollouts. Three readings are
stated. Telemetry behaves as the mechanisms dictate. Dynamic sampling
buys perfect telemetry at a 45\% rollout premium and does not exceed the
cheaper fixed recipe. Adaptive-G spends comparable overhead and ties the
best held-out result, concentrating its spend on flagged steps --- and
the tie is small-n and single-task, so the description is ``competitive
with, not superior to, the best fixed recipe.''

The controller simulation then asks a counterfactual question on the
real reward tensors: for each of the 2,560 prompt-step decisions in the
four-method run, when would each controller have fired, what group size
would it have chosen, and how much contrast would it have restored? Its
zero-rescue result is conditional on the plug-in estimator, rather than
evidence that fresh rollouts cannot recover contrast. Across all four
methods and all five thresholds of the point-estimate rule, the
controllers save \emph{zero} prompts. Of GRPO's 640 prompt-step pairs,
461 have identical observed rewards --- 427 all-correct and 34 all-wrong
--- with \ensuremath{\hat{p}} at exactly 0 or 1, and the counts are near-identical for the
other methods (AERO 461, GIFT 493, AREAL 452). Substituting \ensuremath{\hat{p}} ∈ \{0, 1\}
into \(h_G\) forces the estimated degeneration probability to one for
every G. Resampling the same eight identical observed rewards to size G
= 16 also cannot create an unobserved outcome: the resulting 0/1,867
rescue count is mechanical. Neither calculation identifies latent p as 0
or 1. These traces contain predominantly all-correct groups, but
independent fresh responses at larger G from the historical panel are
needed to estimate how often new outcomes restore contrast there. The
separate frozen-checkpoint diagnostic in §7.4 observes new responses
under a different, explicitly fixed contract; it does not reproduce this
historical panel. The operating range where the trigger fires
selectively is τ ∈ {[}0.70, 0.80{]}, 12 to 19 fires per 40 steps at a
cost ratio of 1.30 to 1.48, and the cost differences are retrospective
budget accounting under that estimator.

In the same trace the escalation asymmetry is stark: 1,723 of 1,867
escalation events (92.3\%) occur on all-correct groups and only 144 on
all-wrong groups. A symmetric ZVF trigger mostly spends on prompts whose
eight observed completions were all correct; that is not proof of
certain success on future draws. The P7 report is explicit that deleting
those fires is a counterfactual proposal only, because performance
preservation has not been measured.

Two refinements sharpen the picture and then complicate it. The Bayesian
controller replaces the point estimate with a Beta(1,1) posterior and a
mid-range probability m(k, G) = Pr{[}0.05 ≤ p ≤ 0.95 \textbar{} k, G{]},
escalating when m exceeds a posterior threshold. At τ\_post = 0.60 it
fires on the same 466.75 degenerate prompt-steps as the point-estimate
rule at τ = 0.50 --- with a 95\% bootstrap CI of {[}454.25, 485.00{]}
--- but at strictly lower cost, 8,854 rollouts against 10,240, a 14\%
saving; and all 466.75 of those mean fires are on observed all-equal
groups, so the plug-in rescue calculation again reports zero; this is
not a posterior-predictive zero-rescue guarantee. Above τ\_post = 0.65
it never fires, because the maximum m over degenerate observations is
0.630.

The posterior-predictive eight-controller comparison produces a
model-based metric reversal with bootstrap CIs attached. Its
contrast-restoration estimates depend on the assumed posterior and are
not observed outcomes from fresh rollouts; these intervals do not
validate that assumption. On the cost axis the Pareto winner is
Dualformer-Auto, an automatic difficulty-gated de-escalation rule, at a
mean cost ratio of 0.657 {[}0.599, 0.738{]} against 1.000 for fixed G =
8 --- a 34\% rollout saving, achieved because the prompts it shrinks to
G = 2 would have been degenerate at G = 8 anyway. On a different axis,
contrast restored per 1,000 extra rollouts, the ordering reverses:
\texttt{zvf-triage} at τ = 0.50 restores 90.10 units per 1,000 extra
rollouts {[}89.40, 90.77{]}, against 86.68 for τ = 0.70, 84.00 for
Dualformer-Auto and 80.00 for the Bayesian rule, and all four pairwise
differences have CIs excluding zero. The Bayesian controller ---
Pareto-dominant on the earlier axis --- is dominated here, because its
principled prior fires on the hardest-to-restore degenerate prompts. The
regime decomposition explains the mechanism rather than leaving it as a
ranking: \texttt{zvf-triage} at τ = 0.50 fires on 100\% of steps in
every regime on this panel, so it is operationally equivalent to
``escalate every step,'' and its apparent advantage is mechanical
exposure rather than a smarter rule. Its restoration decomposes 64.76\%
degenerate, 13.69\% boundary and 21.56\% mid-range; the non-degenerate
35.24\% share is what the Bayesian rule sacrifices.

Two further analyses bound what this family of shrink-only rules can
achieve. A per-prompt hindsight-optimal calculation scans G' down from a
base of 8 and finds a bimodal optimum: G* = 8 for 72.9\% of the 2,560
prompt-steps, which are observed all-equal groups assigned zero rescue
by the plug-in model, and G* = 2 for the remaining 27.1\%, which are
already non-degenerate at G = 8 and are cheapest to shrink. Mean G* is
6.376, giving a cost ratio of 0.797 {[}0.776, 0.818{]} --- a 20.3\%
saving that is pure economy, since mean contrast restored per prompt is
0.000 and mean contrast lost to shrinking is 0.130 ZVF units. This
hindsight optimum is not a ceiling on cost (corrected; see Table E.1):
Dualformer-Auto's observed 0.66 to 0.68 on the four-bin rule is
\emph{cheaper}. The two optimise different objectives. The hindsight
calculation keeps every already-degenerate prompt at G* = 8 and only
shrinks non-degenerate ones, whereas Dualformer-Auto also shrinks
prompts that would have been degenerate at G = 8. The 0.797 is therefore
the best cost ratio under a no-shrinking-of-saturated-prompts
constraint, not a lower bound on cost for any adaptive rule. A
finite-pool correction then shows that even this bound is loose: because
a finite pool never collides more than independent sampling, the exact
hypergeometric contrast-preservation probability exceeds the binomial
one everywhere, by +0.0494 {[}+0.0481, +0.0506{]} on the 693 relevant
prompt-steps, which means the binomial model over-predicts starvation
and the ``zero contrast restored'' accounting understates true preserved
contrast by roughly 34.2 ZVF units.

Cross-base replication is where the retreat from ``controller'' to
``audit'' is earned. On a second panel of five independent seeds, the
trigger fires 4.20 times per seed {[}3.00, 5.40{]} at τ = 0.70, with the
per-step rate falling from 0.48 to 0.28 as the panel's overall ZVF
drops; what replicates is the sign and the operating-point logic, not a
fixed firing magnitude. The Bayesian branch is entirely silent on this
panel, firing 0.0 times per seed, because every step's implied mid-range
probability lies in {[}0.95, 0.999{]} --- a real negative finding that
scopes the Bayesian controller to distributions containing boundary
steps. And the aggregate firing rate is stable while the \emph{location}
is not: at τ = 0.70 with a de-escalation threshold of 0.95, the
cross-seed Jaccard overlap of the fire sets is 0.133, only 13.3\% of
steps see all five seeds choose the same group size, and the mean
cross-seed Kendall-τ of per-step ZVF trajectories is −0.060. The firing
rate is seed-stable; the firing location is not, and the training-step
index is a meaningless axis for the ZVF signal across seeds. The
replayed unification confirms the direction: on this panel
Dualformer-Auto is the only controller that saves compute (+14\%
{[}0.10, 0.18{]}) while \texttt{zvf-triage} spends more than baseline
(−28\% {[}−0.36, −0.20{]}), with disjoint CIs.

Five design rules close the P7 report, and they read as constraints
discovered rather than principles asserted. Gate on PCD and alarm on
ZVF: escalate only when ZVF is high \emph{and} PCD indicates interior
difficulty, because high ZVF with near-zero PCD and high reward is
consistent with mastery --- where \(h_G(p) \approx 1 - G\varepsilon\)
near the boundary and added rollouts buy order ε --- while high ZVF with
near-zero reward is consistent with incapacity; finite observations
alone do not identify either endpoint, so the proposed choice between
data intervention and resampling still needs prospective validation.
Escalate multiplicatively in small steps, since at p = ½ each additional
rollout halves the degenerate fraction; the audited 4 → 6 → 8 schedule
at multipliers of 1.5 and 1.33 captured most of the achievable reduction
at a total cost of +55\% rollouts, whereas G = 16 costs four times the
rollouts of G = 4 for a ZVF drop from 0.76 to 0.63 and no held-out gain
on a near-saturated task. Trigger on spikes rather than levels, because
ZVF is temporally sticky with a lag-1 autocorrelation of approximately
0.94 and drifts upward over training, so a level threshold fires either
never or permanently. Never compare controllers across stacks; ship the
seven manifest fields with item 5 carrying the escalation log. And treat
ZVF = 0 as a purchase rather than a prize --- dynamic sampling shows
perfect telemetry is always available at a 45\% rollout premium, and on
this task it was not the best use of the marginal rollout.

The chapter's summary of P7 is the report's own. No result establishes
prospective superiority over static G = 16, Dr.~GRPO, or dynamic
sampling; adaptive group size is explicitly \emph{not} promoted to a
recommended policy; and all intervention evidence is single-task,
small-n, or retrospective. The report's contribution is a retrospective
controller audit together with a pre-registered prospective test plan,
and the bakeoff that would close the question --- a matched-budget,
multi-seed held-out comparison against the strong static baseline --- is
future work rather than a claim. In the same spirit, the report flags
its own August 2026 evidence boundary: unsigned ZVF alone cannot select
between retiring, distilling, retrying, resampling or warm-starting a
prompt, so any operator policy remains a prospective evaluation contract
rather than a demonstrated benefit.

\section{Historical V2 Frozen-Policy Fresh-Rollout
Diagnostic}\label{historical-v2-frozen-policy-fresh-rollout-diagnostic}

\textbf{Format-induced contrast is distinct from numeric contrast.} A
prospectively frozen diagnostic completed on 2 October 2026 found that
33 of 461 initially homogeneous eight-response groups became mixed after
24 additional responses: 7.16\%, with a prompt-level Wilson 95\%
interval of 5.14-9.88\%. Twenty of those first reward changes were
parser failures whose visible final number still agreed with the gold
answer. The remaining 13/461 (2.82\%) were assistant-audited
numeric-correctness contrasts: two wrong-to-correct recoveries and
eleven correct-to-wrong changes. Thus the raw 7.16\% measures variation
in this diagnostic's binary reward, not a useful learning-signal rate or
an accuracy gain.

The main cohort contained 512 GSM8K \textbf{train} prompts, each with 32
ordered responses from a frozen public Qwen/Qwen3-8B checkpoint in BF16
non-thinking mode. All 16,384 main responses completed. A disjoint
32-prompt smoke cohort was excluded from every main estimate. Each draw
was a fresh response to the same original prompt, with a unique request
seed, rather than a continuation of a previous answer. The primary
endpoint was predeclared: among prompts whose first eight binary rewards
all agreed, count those with both rewards in the combined 32. The strict
parser used exact rational matching; parser failures received reward
zero, while cap-hit outputs retained their exact-match reward and were
separately flagged. Neither kind of output was removed from the primary
denominator. Appendix B §B.12 gives the pinned revisions, decoder,
receipt chain and audit method.

\begingroup\scriptsize\setlength{\tabcolsep}{3pt}
\begin{longtable}[]{@{}
  >{\raggedright\arraybackslash\hspace{0pt}}p{(\linewidth - 8\tabcolsep) * \real{0.2000}}
  >{\raggedleft\arraybackslash}p{(\linewidth - 8\tabcolsep) * \real{0.2667}}
  >{\raggedleft\arraybackslash}p{(\linewidth - 8\tabcolsep) * \real{0.2667}}
  >{\raggedleft\arraybackslash}p{(\linewidth - 8\tabcolsep) * \real{0.2667}}@{}}
\caption{Fresh reward mixedness under the frozen diagnostic contract.
The all-correct numerator includes all 20 format-induced first changes;
its other 11 cases are numeric failures. Intervals use prompts as the
unit and a binomial working model, not a guarantee of population
coverage under arbitrary dependence.}\tabularnewline
\toprule\noalign{}
\begin{minipage}[b]{\linewidth}\raggedright
Initial eight-response group
\end{minipage} & \begin{minipage}[b]{\linewidth}\raggedleft
Became mixed / eligible
\end{minipage} & \begin{minipage}[b]{\linewidth}\raggedleft
Rate
\end{minipage} & \begin{minipage}[b]{\linewidth}\raggedleft
Wilson 95\% interval
\end{minipage} \\
\midrule\noalign{}
\endfirsthead
\toprule\noalign{}
\begin{minipage}[b]{\linewidth}\raggedright
Initial eight-response group
\end{minipage} & \begin{minipage}[b]{\linewidth}\raggedleft
Became mixed / eligible
\end{minipage} & \begin{minipage}[b]{\linewidth}\raggedleft
Rate
\end{minipage} & \begin{minipage}[b]{\linewidth}\raggedleft
Wilson 95\% interval
\end{minipage} \\
\midrule\noalign{}
\endhead
\bottomrule\noalign{}
\endlastfoot
All homogeneous (primary) & 33 / 461 & 7.16\% & 5.14-9.88\% \\
All-wrong (descriptive) & 2 / 9 & 22.22\% & 6.32-54.74\% \\
All-correct (descriptive) & 31 / 452 & 6.86\% & 4.87-9.57\% \\
\end{longtable}
\endgroup

The main corpus had 290/16,384 parser failures (1.77\%) and 18/16,384
cap hits (0.11\%), each below the predeclared 5\% quality threshold.
These global quality checks do not make a parser-induced reward change a
numeric error. Separate assistant reviewers reviewed all 33 initial
groups and first-changing responses in disjoint partitions, solving the
arithmetic and applying the frozen rubric; this was model-assisted
qualitative review, \textbf{not human annotation}. A coordinator
separately checked both recoveries and one bread-rise wording ambiguity.
That question's gold answer, 20 hours, assumes serial rises; it is
ambiguous rather than definitely wrong and was already excluded from the
numeric subset for a formatting failure. The two recoveries changed
eight initial answers of 1152 to the supported value
\(6(150+6\times30+6\times12)=2412\), and eight initial answers of 12 to
the supported coin total \(2+5+3.25+2+0.4+0.35=13\).

The predeclared clean-all-32 sensitivity was 12/439 (2.73\%). Because
eligibility requires all 32 outputs to be parseable and uncapped, it
selects on post-initial outcomes and cannot replace the primary
estimate. A \textbf{post-hoc exploratory} check found no later
parseable, uncapped numeric change in any of the 20 parser-first groups;
it leaves the frozen first-event audit unchanged and cannot predict
further draws.

This parser effect is \textbf{diagnostic-contract-specific}. The
repository's existing GSM8K test evaluator uses a more permissive
extractor, including numeric conclusions embedded in prose. These
strict-parser events are therefore not evidence of reward failures in
the historical evaluator or training runs. The 13 conservative numeric
witnesses show that eight equal observed rewards need not imply a fixed
numeric outcome on fresh draws for this checkpoint, task and decoder.
They do not reconstruct a historical checkpoint, demonstrate training or
controller benefit, identify a group-size optimum, or estimate held-out
capability. The earlier failed V1 engineering smoke remains separate and
unpooled.

\section{C1 Results with Reference and Unit
Boundaries}\label{c1-results-with-reference-and-unit-boundaries}

Benchmark reference errors, ambiguous questions and equivalent units
limit what this diagnostic can establish. Its frozen parser-defined
result is \textbf{11/64 questions (17.19\%; Wilson 95\% interval
9.88-28.21\%)}. Assistant-assisted reconciliation corroborates
\textbf{three numeric recovery witnesses on the same 64-question
denominator}. This 3/64 count (4.69\%) describes the support found under
a conservative review rule; it does not replace the frozen endpoint or
estimate a population recovery rate. The other eight parser events
include three convention-dependent interpretations, four materially
ambiguous cases and one match to an incorrect reference answer.
{[}S1-S4{]}

The study contains no training update. It does not measure useful
reinforcement-learning gradients, reasoning improvement, controller
superiority, held-out accuracy or general model capability. The prompts
are public GSM8K training questions and are fresh only to this campaign.
Historical V2 evidence remains separate and unpooled.

\begin{longtable}[]{@{}
  >{\raggedright\arraybackslash\hspace{0pt}}p{(\linewidth - 6\tabcolsep) * \real{0.3000}}
  >{\raggedleft\arraybackslash}p{(\linewidth - 6\tabcolsep) * \real{0.4000}}
  >{\raggedright\arraybackslash\hspace{0pt}}p{(\linewidth - 6\tabcolsep) * \real{0.3000}}@{}}
\caption{C1 results and interpretation. The final three rows partition
the eight events outside the conservative witness count.
{[}S1-S4{]}}\tabularnewline
\toprule\noalign{}
\begin{minipage}[b]{\linewidth}\raggedright
Quantity
\end{minipage} & \begin{minipage}[b]{\linewidth}\raggedleft
Result
\end{minipage} & \begin{minipage}[b]{\linewidth}\raggedright
Interpretation
\end{minipage} \\
\midrule\noalign{}
\endfirsthead
\toprule\noalign{}
\begin{minipage}[b]{\linewidth}\raggedright
Quantity
\end{minipage} & \begin{minipage}[b]{\linewidth}\raggedleft
Result
\end{minipage} & \begin{minipage}[b]{\linewidth}\raggedright
Interpretation
\end{minipage} \\
\midrule\noalign{}
\endhead
\bottomrule\noalign{}
\endlastfoot
Eligible in the fixed screen & 87 / 4,096 & Eight parser-zero, valid and
uncapped outputs \\
Selected for fresh draws & 64 / 87 & First 64 in frozen order \\
Frozen primary event & 11 / 64 & 17.19\%; Wilson 95\% 9.88-28.21\% \\
Conservative numeric witnesses & 3 / 64 & 4.69\%; support count, not a
replacement endpoint \\
Convention-dependent event readings & 3 cases & Retained review
disagreement \\
Materially ambiguous events & 4 cases & Requested quantity or premises
unresolved \\
Wrong-reference event & 1 case & Matching the reference reverses a
correct initial answer \\
\end{longtable}

The three corroborated witnesses are Santana presents (C00, source rank
676), Suzanne bread (C07, rank 912), and Hank reading (C38, rank 2699).
Santana's supported answer is \((4+7)-3=8\) presents; Suzanne's is
\(3\times2\times52/12=26\) loaves; Hank's is
\(5(30+60)+2\times2(30+60)=810\) minutes. Each has eight reviewed,
genuinely wrong initial final answers and at least one independently
correct fresh final answer. {[}S3{]}

Santana also shows why final-answer correctness cannot validate
reasoning. Its earliest successful fresh response, sample index 12,
returns eight through an incorrect argument that no first-half birthday
gifts were needed. Later successful indices 15 and 23 use the valid
\(11-3\) calculation. Hank's reference final value is supported even
though its reference explanation mislabels subtotals. These distinctions
remain visible in the event ledger. All sample and selected indices in
this addendum are zero-based. {[}S3{]}

The checkpoint case (C35, rank 2593) gives the opposite result.
Checkpoints at miles one and twenty-five enclose twenty-four miles
across three gaps, so the separation is eight miles. All eight initial
answers of eight are numerically correct. Four fresh outputs match the
erroneous reference value six by dividing by four. The frozen analysis
counts the event exactly as specified; semantic review does not call it
mathematical recovery. {[}S3{]}

The citations, raft and fundraiser cases have favorable ordinary
interpretations, but retain first-versus-second review disagreements.
They are described individually rather than combined with the three
corroborated witnesses into another adjusted recovery estimate. The
remaining four cases concern aggregation of walking distance, the
cake-and-ham probability event, the pet-owner denominator and Reggie's
missing initial marble count. A gold match does not resolve those
questions. {[}S3{]}

The original second-review categories are retained: sixteen clear
wrong-reference judgments, nineteen materially ambiguous questions, two
inconsistent questions, fifteen supported references and twelve
references supported with conventional assumptions. Refinement separates
four correct final answers with flawed reference explanations from the
supported categories, leaving the following mutually exclusive view.
{[}S2, S4{]}

\begin{longtable}[]{@{}lr@{}}
\caption{Refined labels preserve the original classifications and
distinguish final-value errors from explanation defects. They are
second-review judgments, not a consensus benchmark error
count.}\tabularnewline
\toprule\noalign{}
Refined second-review category & Questions \\
\midrule\noalign{}
\endfirsthead
\toprule\noalign{}
Refined second-review category & Questions \\
\midrule\noalign{}
\endhead
\bottomrule\noalign{}
\endlastfoot
Clear wrong reference final value & 16 \\
Materially ambiguous & 19 \\
Inconsistent question & 2 \\
Correct final value with flawed reference explanation & 4 \\
Supported & 12 \\
Supported with conventional assumptions & 11 \\
Total & 64 \\
\end{longtable}

A narrower cross-review analysis corroborates fifteen wrong-final-value
judgments under the stated or ordinary mathematical model. The towel
case explicitly assumes equal towel weights. Separately, two
discrete-count questions have inconsistent premises and erroneous
reference relations, three correct-final references have corroborated
text defects, and two materially ambiguous questions contain explicit
reference faults. These are \textbf{22 disjoint cases with specified
reference defects, not 22 wrong labels}. The four correct-final cases in
the refined second review and the three corroborated correct-final
defects are different sets: a notation defect newly identified for
Sonia/Joss is not corroborated in both rounds. {[}S2, S4{]}

Thirteen first/second category disagreements are retained. The coin
question illustrates the reason: gross deposits total 225 cents and the
balance after withdrawal is 200 cents. The second review judges the
recorded 200 wrong for the requested deposits, while the first retains a
gross-versus-net target ambiguity. It is excluded from the narrower
fifteen. Ordinary mathematical conventions are not defects merely
because every physical assumption is unstated; genuinely different
quantities or plausible readings are described case by case. None of
these audit counts estimates whole-GSM8K error prevalence because the
cohort was selected for eight parser-zero outputs and then truncated.
{[}S2, S4{]}

Equivalent representations also matter. For Jackson ingredients (C36),
all eight initial values 0.24 are equivalent to the reference 24 percent
under the same conventional independent-selection probability
interpretation; formal odds would be 6:19. The scalar parser extracted
0.24 faithfully. For Gunner plugs (C62), 200 pairs equal 400 individual
plugs. Four initial responses preserve pair units, three remain
unit-ambiguous and one incorrectly calls 200 the number of individual
plugs. Neither case supports a claim of eight genuine mathematical
errors. The frozen scalar references and rewards remain unchanged.
{[}S5{]}

\subsection{Separate Prospective Parser
Validation}\label{separate-prospective-parser-validation}

The undeployed parser v3 candidate was frozen at 02:43:30 UTC on 3
October 2026, before the earliest target batch at 02:59:39 UTC. The
fixed target was the last 256 screening questions, eight responses each,
yielding 2,048 future outputs. This was a predefined future subset of an
ongoing study, not preregistration before the entire C1 campaign.
Historical candidate tuning is separate. C1's primary score remains
based on frozen parser v2. {[}S8, S9{]}

All 2,025 previously valid full parse objects were preserved. Four
previously rejected outputs at four distinct prompts became accepted:
one bold answer heading, two display-box blocks and one hash-prefixed
XML answer block. Two separately blinded, output-only assistant
annotations agreed on the same unambiguous scalar in all four changed
cases. No changed acceptance was ambiguous or incorrectly extracted
under those reviews. Nineteen outputs remained rejected. {[}S8, S9{]}

This supports only a narrow pass on the four observed extraction
changes. It does not establish semantic precision for the 2,025
unchanged outputs, validate the nineteen rejects, measure task
correctness, or estimate population false-positive rates. The responses
are clustered within prompts. Four examples do not exhaust their wrapper
classes or malformed-input behavior. The candidate remains undeployed
and is not used to recompute C1 recovery. {[}S9{]}

All source labels refer to Appendix G.



\chapter{The E1--E14 Evaluation
Campaign}\label{the-e1e14-evaluation-campaign}

This chapter reports the evaluation of a single Tinker-trained actor
against fourteen benchmark suites outside its recorded training-suite
list, recorded lane by lane under sealed reservations, immutable
receipts and explicit terminal states. Item-level decontamination is not
established. Lanes that did not finish and lanes that could not start
are reported too, because the methodological aim is to make the boundary
between what was measured, what was only reported, and what can still be
verified explicit. The current repository does not preserve every cited
artifact, so it cannot support an unqualified claim that every figure
has a surviving, hash-pinned evidence chain. Lane-by-lane detail that
does not change any figure is collected in Appendix D (Figure~\ref{fig:lanes}).


\fig{fig_lane_status}
{Terminal status of all fourteen E1--E14 lanes: the 2026-09-19 ledger of record, updated with the six replacement scopes scored on 2026-09-27 (double border, $\star$)}
{Terminal status of all fourteen E1--E14 lanes: the 2026-09-19 ledger of record, updated with the six replacement scopes scored on 2026-09-27 (double border, $\star$). Each cell names the suite whose scope the state refers to, and status is carried by glyph, word, border style and hatching as well as colour. Cells show the replacement scope where one exists (E1 SWE-bench Multilingual, E2 CORE-Bench, E5 Tau3, E6 WebArena, E8 LAB-Bench public, E9 MLDevBench, E10 AgentDojo benign, E13 BALROG, E14 Omni-MATH) and the original contract otherwise; the original-contract suites and their results (e.g.\ E1 SWE-bench Pro $2/731$, E2 FrontierSWE, E5 APEX-Agents, E9 MLE-bench) are shown in the campaign-design figure (Chapter~4) and are never pooled with replacement-scope numbers. *E14 closes terminal-complete: $2271/4428 = 51.29\%$ with the two unjudged rows counted as failures (official scorer: 51.31\% over 4426 judged rows).}
{fig:lanes}



Two reporting rules govern everything below, and they are worth stating
before any number appears. First, original-contract results and
replacement-scope results are never pooled and never averaged across
suites; a replacement scope is a different suite, a different interface
and often a different grader, and combining such figures would produce a
number with no referent. Second, no figure in this chapter is quoted to
a precision beyond what its source receipt records, and no comparison
against a baseline, a sibling suite, or the untrained base model is
asserted for the campaign lanes of Sections 8.2--8.6. Those lanes have
no matched base-model arm. The one paired baseline comparison in this
chapter is the small-scale trained-versus-base rerun of Section 8.7
(Table 8.B): 14 lanes, 5 to 151 identical items each, one vLLM engine
for both arms. It finds no significant difference in any lane, and this
chapter therefore claims no improvement of any kind over the base model.
Where a lane's record itself contains an unresolved discrepancy, the
discrepancy is stated.

The consolidated ledger of record for the campaign is
\longtok{outputs/E1\_\allowbreak{}E14\_\allowbreak{}FINAL\_\allowbreak{}RESULTS\_\allowbreak{}2026-09-19.md}, which
supersedes the working tables of 2026-09-05 and 2026-09-12; the per-lane
terminal-state ledger,
\texttt{withheld\ historical\ campaign\ chronology}, was recovered from
the original Mac checkout on 3 October 2026 with updates through 21
September, while remaining absent from the frozen repository; and
selected arithmetic of the ledger was historically re-verified by
deterministic code checks recorded in
\longtok{outputs/verification/LEDGER\_\allowbreak{}CODE\_\allowbreak{}CHECK\_\allowbreak{}2026-09-19.json},
which reports eleven of eleven named checks passing with zero failures
(Appendix B lists their scope; they do not cover every figure in this
chapter).

\textbf{Evidence availability (private recovery audit, 3 October 2026).}
Six original files covering all five previously missing basenames were
recovered from the Mac checkout: the E10 native receipt, E8 diagnostics,
training-identity verification, E14 terminal note and two distinct
\texttt{Pending\_Experiments.md} snapshots. All remain absent from the
frozen repository, but their exact UTF-8 bytes and SHA-256 hashes are
preserved in the withheld source archive described in Appendix B §B.11.
A separately retained source audit records bounded read-only checks of
underlying E8/E10/E14 artifacts on the Mac: 29 exact reference-path hash
matches, 1,979 E8 archive/local member checks, 525 E10 archive-member
checks and E14 stored-disposition inspection. Exported E8 grade records
and E10 selected trace records permitted independent tally checks in the
retained private evidence audit; those underlying rows are not included
in this public document package; E14 exports aggregate traversal results
and two exclusion records, not all underlying rows. This verifies stored
evidence and arithmetic within that scope, without rerunning models or
scorers, rejudging correctness, or reconstructing the missing training
logs. It supersedes the earlier missing-primary disclosure for these
files only; the frozen-checkout absence record remains valid. E14's
recovered note contains a denominator inconsistency, and the recovered
finish ledger does not substantiate the later E4 rerun; both are
disclosed below.

\section{The evaluated actor and the one-actor
design}\label{the-evaluated-actor-and-the-one-actor-design}

A single actor was evaluated across the whole campaign. The base model
is \texttt{Qwen/Qwen3.6-35B-A3B} at revision
\longtok{995ad96eacd98c81ed38be0c5b274b04031597b0}, and the trained
parameter set is the LoRA adapter \texttt{campaign-adapter} at commit
\texttt{withheld-adapter-revision}, served in bfloat16 under the
identifier \texttt{campaign-actor-bf16}. The public package does not
contain the original execution route or operational repository access
identities. It therefore cannot reproduce the original sampler execution
from these materials alone; no current provider-retention or
account-status claim is made here.

The one-actor design is deliberate. No lane received a separate
checkpoint, per-suite prompt engineering beyond its native harness, or
tuning against its own evaluation set: the recovered training-identity
record lists only API-Bank-RLVR and SWE-Gym (§8.1.1), while the
recovered E10 receipt explicitly reports
\longtok{TRAINING\_\allowbreak{}INVENTORY\_\allowbreak{}ABSENT} and makes no held-out claim. The
E10 contract's declaration of AgentDojo training is explicitly marked as
not evidence of training consumption; neither record proves
content-level decontamination. Differences between suites therefore
reflect the interaction of one fixed policy with fourteen different
interfaces, observation spaces, tool conventions and graders, and they
are not evidence about what a specialised actor could achieve. The
campaign also cannot support any statement about seed variance.
Inference conditions were frozen per lane and recorded in that lane's
own receipt (§8.1.2).

\subsection{How the actor was trained}\label{how-the-actor-was-trained}

The adapter is attributed to a single Tinker training run (W\&B run
\texttt{withheld-id}). The recovered
\longtok{training\_\allowbreak{}identity\_\allowbreak{}verified.json} records that run ID,
aggregate dataset revision, start time, base Git commit and hashes of
\texttt{config.yaml}, \texttt{wandb-metadata.json} and
\texttt{output.log}. Its status is
\longtok{TRAINING\_\allowbreak{}CONFIG\_\allowbreak{}VERIFIED\_\allowbreak{}NOT\_\allowbreak{}CONTENT\_\allowbreak{}DECONTAMINATION}, and
it states that the dirty source diff was not uploaded. The underlying
W\&B files and local run receipt have not been recovered in this
package, so their bytes cannot yet be matched afresh to those recorded
hashes. Appendix D separates these surviving identity fields from the
still-reported configuration details in Table 8.1.

\textbf{Reported training record of the E1--E14 actor.} The
identity-verification artifact is recovered; underlying local run files
remain unavailable in this package. Current-code evidence is labelled
separately.

\begin{longtable}[]{@{}
  >{\raggedright\arraybackslash\hspace{0pt}}p{(\linewidth - 6\tabcolsep) * \real{0.3333}}
  >{\raggedright\arraybackslash\hspace{0pt}}p{(\linewidth - 6\tabcolsep) * \real{0.3333}}
  >{\raggedright\arraybackslash\hspace{0pt}}p{(\linewidth - 6\tabcolsep) * \real{0.3333}}@{}}
\caption{Reported training record of the E1--E14 actor, with the
source-availability limitations above.}\tabularnewline
\toprule\noalign{}
\begin{minipage}[b]{\linewidth}\raggedright
Field
\end{minipage} & \begin{minipage}[b]{\linewidth}\raggedright
Value
\end{minipage} & \begin{minipage}[b]{\linewidth}\raggedright
Evidence
\end{minipage} \\
\midrule\noalign{}
\endfirsthead
\toprule\noalign{}
\begin{minipage}[b]{\linewidth}\raggedright
Field
\end{minipage} & \begin{minipage}[b]{\linewidth}\raggedright
Value
\end{minipage} & \begin{minipage}[b]{\linewidth}\raggedright
Evidence
\end{minipage} \\
\midrule\noalign{}
\endhead
\bottomrule\noalign{}
\endlastfoot
Launch &
\texttt{python\ -\allowbreak{}m\ platform\_\allowbreak{}tinker.tinkerrl.grpo\_\allowbreak{}cli\ -\allowbreak{}\/\allowbreak{}-\allowbreak{}preset\ pavlov\_\allowbreak{}portfolio},
no overrides & \longtok{wandb-metadata.json} \\
Base model & \longtok{Qwen/Qwen3.6-35B-A3B} at revision
\longtok{995ad96e…} & \longtok{config.yaml}; run receipt \\
Loss & REINFORCE with a group-mean baseline:
\(A_i = (r_i - \mu_g)/(\sigma_g + 10^{-8})\) with population \(\sigma\);
loss \(-A_i \sum_t \log \pi(y_{i,t})\), averaged over completions. No
importance ratio, no clipping, no KL term, no reference model &
\longtok{platform\_tinker/tinkerrl/grpo.py} (\longtok{normalize\_rewards},
\longtok{make\_\allowbreak{}grpo\_\allowbreak{}loss\_\allowbreak{}fn}) \\
Sampler refresh & Only at checkpoints: steps 1--20 sampled from the
step-0 weights and steps 21--40 from the step-20 weights, with no
importance correction & \texttt{grpo.py} \\
Reward & \longtok{PavlovNonXLAMReward}: API-Bank rows scored by
\longtok{StrictToolCallReward}, SWE-Gym rows by \longtok{PatchReward} (1.0
for an exact diff, otherwise partial credit capped at 0.99) &
\longtok{platform\_tinker/tinkerrl/grpo\_cli.py}, \texttt{grpo.py}
(current code; see Appendix D) \\
Training data & Suites \longtok{api\_\allowbreak{}bank\_\allowbreak{}rlvr\_\allowbreak{}train} and
\longtok{swe\_\allowbreak{}gym\_\allowbreak{}train} (API-Bank-RLVR at \longtok{bf67c426…}, SWE-Gym
at \longtok{bb94ed9e…}); dataset revision \longtok{bd1f0db6…}; training
pool of 512 rows & Recovered training-identity record binds the suite
IDs and aggregate revision; run receipt, \longtok{config.yaml} and E11
decontamination receipt are the reported sources for remaining fields \\
Declared evaluation suites & 14 primary evaluation suites (the original
E1--E14 contracts, including \longtok{verilog\_eval});
\longtok{heldout\_\allowbreak{}suite\_\allowbreak{}ids} empty; \longtok{evaluate\_heldout} false &
run receipt \\
Group size, batch, steps & G = 4; 2 prompts per step; 40 steps; 80
prompt draws and 320 completions in total & \longtok{config.yaml}; run
receipt \\
Optimiser & lr 1 × 10⁻⁵; Adam β₁ 0.9, β₂ 0.95, ε 10⁻⁸ &
\longtok{config.yaml} \\
LoRA & rank 32, alpha 32, all linear layers & \longtok{config.yaml};
adapter config \\
Decoding during training & prompt 1,536 and response 384 tokens;
temperature 0.7, top-p 0.95; seed 809 & \longtok{config.yaml} \\
Reward trace & mean 0.025 over the first 5 steps and 0.2425 over the
last 10; reward zero on 22 of 40 steps; loss zero on 35 of 40 steps &
run receipt; \longtok{output.log} \\
Training time & started 2026-08-09; runtime 1,934 s & historical
training metadata; underlying files remain unavailable \\
Checkpoints & Scientific schedule reported at step 0, step 20, step 40
and final; the final adapter follows step 40 without an additional
optimizer step & Historical scientific training record;
execution-specific repository paths omitted \\
Execution provenance & Exact private execution identities and access
paths are withheld; the public package cannot replay the original
execution route & Public-source limitation, Appendix I \\
\end{longtable}

Four items remain unavailable or were never created --- the exact
training code, which 80 training prompts were drawn, the Tinker sampler
weights, and any held-out training-time evaluation --- and the run
departs from the planned Pavlov protocol; both are recorded in Appendix
D so that no reader fills them in.

What this means for the chapter is direct. The actor received 40 steps
of sparse-reward, critic-free policy-gradient training, with zero loss
on 35 of them. The fourteen suite results are properties of this
checkpoint, not evidence about what GRPO-style training contributes. The
campaign is reported as an exercise in evaluation governance. The only
trained-versus-base evidence is the small-sample paired comparison of
Table 8.B.

\subsection{Serving conditions per
lane}\label{serving-conditions-per-lane}

The E4 zero (§8.5) prompted identification of a serving bridge that sent
no stop sequences. The defect can cause failed agent trajectories, but
its contribution to the full 100-trial zero is not identified by the
available comparison. The same bridge,
\longtok{zvf-program/flagship/modal\_\allowbreak{}tinker\_\allowbreak{}openai\_\allowbreak{}bridge.py}, served
several other lanes, so Table 8.2 records for every lane how the actor
was served and whether that defect could apply. Its
\texttt{/v1/responses} route hard-codes \texttt{stop=None} and thinking
on. Its \texttt{/v1/chat/completions} route sends stop sequences only if
the client supplies them, and the agent clients used here (opencode,
mini-swe-agent and the APEX client) do not. The deployed bridge version
(sha256 \texttt{f59c08b9…}) matches neither committed version of the
file.

\textbf{Serving conditions of the original campaign runs, by lane.}

\begingroup\scriptsize\setlength{\tabcolsep}{3pt}
\begin{longtable}[]{@{}
  >{\raggedright\arraybackslash\hspace{0pt}}p{(\linewidth - 16\tabcolsep) * \real{0.1250}}
  >{\raggedright\arraybackslash\hspace{0pt}}p{(\linewidth - 16\tabcolsep) * \real{0.1250}}
  >{\raggedright\arraybackslash\hspace{0pt}}p{(\linewidth - 16\tabcolsep) * \real{0.1250}}
  >{\raggedright\arraybackslash\hspace{0pt}}p{(\linewidth - 16\tabcolsep) * \real{0.1250}}
  >{\raggedright\arraybackslash\hspace{0pt}}p{(\linewidth - 16\tabcolsep) * \real{0.1250}}
  >{\raggedright\arraybackslash\hspace{0pt}}p{(\linewidth - 16\tabcolsep) * \real{0.1250}}
  >{\raggedright\arraybackslash\hspace{0pt}}p{(\linewidth - 16\tabcolsep) * \real{0.1250}}
  >{\raggedright\arraybackslash\hspace{0pt}}p{(\linewidth - 16\tabcolsep) * \real{0.1250}}@{}}
\caption{Serving conditions of the original campaign runs, by
lane.}\tabularnewline
\toprule\noalign{}
\begin{minipage}[b]{\linewidth}\raggedright
Lane
\end{minipage} & \begin{minipage}[b]{\linewidth}\raggedright
Serving path
\end{minipage} & \begin{minipage}[b]{\linewidth}\raggedright
Stop sequences
\end{minipage} & \begin{minipage}[b]{\linewidth}\raggedright
Template / thinking
\end{minipage} & \begin{minipage}[b]{\linewidth}\raggedright
Max tokens
\end{minipage} & \begin{minipage}[b]{\linewidth}\raggedright
Temp.
\end{minipage} & \begin{minipage}[b]{\linewidth}\raggedright
Parse failure / truncation
\end{minipage} & \begin{minipage}[b]{\linewidth}\raggedright
Could the stop defect apply?
\end{minipage} \\
\midrule\noalign{}
\endfirsthead
\toprule\noalign{}
\begin{minipage}[b]{\linewidth}\raggedright
Lane
\end{minipage} & \begin{minipage}[b]{\linewidth}\raggedright
Serving path
\end{minipage} & \begin{minipage}[b]{\linewidth}\raggedright
Stop sequences
\end{minipage} & \begin{minipage}[b]{\linewidth}\raggedright
Template / thinking
\end{minipage} & \begin{minipage}[b]{\linewidth}\raggedright
Max tokens
\end{minipage} & \begin{minipage}[b]{\linewidth}\raggedright
Temp.
\end{minipage} & \begin{minipage}[b]{\linewidth}\raggedright
Parse failure / truncation
\end{minipage} & \begin{minipage}[b]{\linewidth}\raggedright
Could the stop defect apply?
\end{minipage} \\
\midrule\noalign{}
\endhead
\bottomrule\noalign{}
\endlastfoot
E1 SWE-bench Pro & Mixed: 476 Tinker sampler, 255 Modal vLLM (merged
PEFT) & none; relies on EOS & chat template, thinking off & 8,192 & 0.2
& 14 generation failures, 4 artefacts lost; 40/476 Tinker and 22/255
vLLM responses at the cap; 593 of 713 submitted patches fail
\texttt{git\ apply} as corrupt patches (hunk-header line counts), so
tests ran on the unpatched base & Not through the bridge. The Tinker
path also sent no stops, but no response contains role-token lines and
cap rates match across backends \\
E2 FrontierSWE & bridge (Harbor/opencode) & not recorded; client sets
none & not recorded & not recorded & not recorded & not recorded &
Yes \\
E3 SDAB & no model output & --- & --- & --- & --- & --- & n/a \\
E4 BankerToolBench & bridge (opencode); base rerun via
\longtok{/v1/responses} & none & template default (thinking on) & not
recorded & bridge default 0.2 & pass3 archive missing & Yes; defect
documented, aggregate causal effect unmeasured \\
E5 APEX-Agents & bridge (actor and judge) & not recorded & not recorded
& not recorded & not recorded & 3 unscored failures, 1 interrupted &
Yes \\
E6 WebArena & no model output & --- & --- & --- & --- & --- & n/a \\
E7 BinaryAudit & bridge (mini-swe-agent) & not recorded & not recorded &
16,384 (output limit) & not recorded & attempt ended on an HTTP 402
bridge-budget error & Yes, but the attempt ended on a budget error \\
E8 LAB-Bench & Modal vLLM, merged BF16 & none; vLLM EOS & thinking off &
1,024 & 0 & 1,259/1,967 unparsed (64.0\%); 1,086/1,967 length-truncated
(55.2\%) & No \\
E9 MLE-bench & streaming arm via bridge \longtok{/chat/completions};
merged arm via vLLM & none in either arm & thinking off & 4,096 & 0.2 &
not recorded & Yes for the streaming arm (single-turn); no for the
merged arm \\
E10 AgentDojo & Modal vLLM, merged BF16 & none; vLLM EOS & thinking on &
4,096 & 0 & 1 length finish in 418 recorded response bodies (407 main
run, 11 continuation) & No \\
E11 VerilogEval & Tinker sampler, called directly & none & template
default (thinking on) & 1,024 & 0.2 & 150/312 extraction failures
(48.1\%), attributed to the 1,024-token cap with thinking on rather than
to answer format; finish reasons not recorded (mean response ≈ 1,002 of
1,024 tokens) & Not through the bridge; unknown, because per-item
outputs are not kept locally \\
E12 AppBench & no model output & --- & --- & --- & --- & --- & n/a \\
E13 BALROG & Modal vLLM, merged BF16 & none; vLLM EOS & thinking on &
8,192 & 1.0 & 78 length finishes against 189 stop finishes & No \\
E14 Omni-MATH & Modal vLLM, merged BF16 & none; vLLM EOS & thinking off
& 2,048 & 0 & 2,982/4,428 responses at the 2,048-token cap (67.3\%),
34.5\% correct against 85.9\% uncapped; 2 unjudged & No \\
\end{longtable}
\endgroup

\emph{Sources: per-lane receipts under
\longtok{outputs/modal\_\allowbreak{}e1\_\allowbreak{}e14/2026-08-16/},
\longtok{outputs/public\_\allowbreak{}portfolio\_\allowbreak{}2026-09-05/},
\longtok{outputs/e2\_\allowbreak{}frontier\_\allowbreak{}swe/}, \longtok{outputs/e5\_\allowbreak{}apex\_\allowbreak{}agents/},
\texttt{outputs/e7\_binaryaudit/}, \longtok{outputs/e9\_\allowbreak{}mle\_\allowbreak{}bench/}, and
\longtok{outputs/PES\_\allowbreak{}Phase2\_\allowbreak{}Review\_\allowbreak{}2026-09-12/finish/e13\_control/};
bridge bindings in
\longtok{outputs/modal\_\allowbreak{}e1\_\allowbreak{}e14/2026-08-16/non\_\allowbreak{}e11\_\allowbreak{}readiness/}.}

The recovered E8 diagnostic now supports its aggregate failure counts,
and the recovered E10 receipt supports its declared inference limits.
Other entries still depend on unrecovered raw responses or finish-era
sources, so recovery of these two receipts does not verify every serving
and failure count in the table. Two consequences follow. First, the
lanes served through the bridge (E2, E4, E5, E7 and the E9 streaming
arm) may be exposed to the E4 mechanism, but exposure does not
demonstrate that it caused their failures or quantify its contribution
to E4. Second, the reported format and length failures mean that E8, E11
and E14 headline figures are end-to-end harness scores; §8.2
distinguishes their sensitivity figures and missing evidence. For E11
and E14 the dominant failure is the generation budget, and for E1 it is
patch formatting: those three scores measure the harness configuration
as much as the actor (§8.2, §8.3). A note on the scope of the
stop-sequence mechanism is in Appendix D.

\section{Suites with strictly complete
scopes}\label{suites-with-strictly-complete-scopes}

The historical ledger and its check log report three strictly complete
scopes: E8, E10 and E11. E8 and E10 are replacement scopes. E11 is an
original-contract full-suite result that also meets the
strict-completeness check, because the suite evaluated is VerilogEval
itself; the ledger lists it both among its replacement scopes and among
its original-contract results, but it is one run and one result. A
fourth, E14, is recorded terminal-complete under its native protocol in
the recovered terminal note, subject to its denominator inconsistency
and the raw-output limits below. The six replacement scopes completed on
27 September postdate that check and are reported separately in §8.4 (Figures~\ref{fig:e8} and~\ref{fig:e11}).


\fig{fig_e8_categories}
{LAB-Bench public split, per-category accuracy, all 1{,}967 questions graded across eight categories}
{LAB-Bench public split, per-category accuracy, all 1{,}967 questions graded across eight categories. Each bar is correct/total for one category (labels give the counts). Accuracy is strongly category-dependent, so no pooled cross-category average is drawn or claimed. CloningScenarios is a historically reported zero (0/33). The original diagnostic file was recovered and byte-verified on 3 October (Appendix B); its aggregate counts agree with these bars.}
{fig:e8}



\fig{fig_e11_decomposition}
{VerilogEval decomposition: two native framings of 156 tasks, 67/156 code-completion and 62/156 spec-to-RTL, summing to $129/312 = 41.35\%$ pass@1}
{VerilogEval decomposition: two native framings of 156 tasks, 67/156 code-completion and 62/156 spec-to-RTL, summing to $129/312 = 41.35\%$ pass@1. The dashed top bar is the sum of the two rows, not a third measurement. Graded by the suite's native harness (receipt \longtok{e11\_\allowbreak{}full\_\allowbreak{}receipt.json}); the $129/311 = 41.48\%$ value that excludes \longtok{Prob099\_\allowbreak{}m2014\_\allowbreak{}q6c} is a non-canonical sensitivity only.}
{fig:e11}




\textbf{E8 --- LAB-Bench, public split.} The recovered
\longtok{outputs/public\_\allowbreak{}portfolio\_\allowbreak{}2026-09-05/labbench\_\allowbreak{}final\_\allowbreak{}diagnostics.json}
records all 1,967 expected items evaluated across eight categories, with
450 correct and 1,517 incorrect responses. Those counts give 450/1967 =
0.2288 (Wilson 95\% CI {[}0.211, 0.248{]}, unit = question, n = 1,967).
It also records 1,259 unparsed responses
(\longtok{native\_\allowbreak{}answer\_\allowbreak{}parse\_\allowbreak{}none}) and 1,086 length-truncated
finishes. The supplementary validation found the native receipt at its
recorded hash and checked 1,979 archive members against both declared
and local hashes. Its 1,967 unique exported grade records independently
corroborate the correct/incorrect totals, category counts and 1,086
length finishes. The 1,259 parse-none count remains a field of the
byte-preserved diagnostic: the exported grade records contain no
parsed-answer field for a separate recount. The reported score is an
end-to-end harness score with substantial formatting and length
limitations. The conditional descriptive sensitivity remains 450/708 =
0.636 (Wilson 95\% CI {[}0.600, 0.670{]}), not the headline or a direct
measure of domain knowledge; its response-level classification has not
been independently regraded. Per-category figures and evidence limits
are in Appendix D.

\textbf{E10 --- AgentDojo, benign utility scope.} The recovered
\longtok{outputs/public\_\allowbreak{}portfolio\_\allowbreak{}2026-09-05/agentdojo\_\allowbreak{}native\_\allowbreak{}receipt.json}
records status \longtok{COMPLETE\_\allowbreak{}NATIVE\_\allowbreak{}BENIGN\_\allowbreak{}UTILITY}, 97 completed
episodes, 88 utility passes, AgentDojo v1.2.2 and decontamination status
\longtok{TRAINING\_\allowbreak{}INVENTORY\_\allowbreak{}ABSENT}. Its score is 88/97 = 0.9072
(Wilson 95\% CI {[}0.833, 0.950{]}, unit = episode, n = 97). The receipt
binds the benchmark version but contains no upstream commit field; the
previously attributed commit
\longtok{089ed468cf3ed0322acc66b0211f26d9d90dbf60} is therefore not
verified by this receipt. Its scope is native default benign utility
only: \texttt{security\_evaluation} and \texttt{heldout\_claim} are
false, while \texttt{attack} and \texttt{defense} are null. Completion
and utility are distinct: 97 completed episodes do not mean 97 correct
episodes. The supplementary validation checks 525 archived member pins
and maps all 97 selected episodes to unique exact-hash traces (94
original and 3 explicit continuation), whose stored utility fields
independently tally to 88 passes and zero handled errors. It excludes
extra historical partial traces from that mapping. This is stored-trace
corroboration, not a rerun of AgentDojo or an independent judgement of
task correctness. The earlier 94-episode attempt and reported collector
repair retain their separate missing-source limitations in Appendix D.

\textbf{E11 --- VerilogEval, two native framings.} All 312 pinned prompt
instances were evaluated under the upstream NVlabs/verilog-eval harness,
but they represent 156 underlying designs, each appearing in
code-completion (\texttt{code-complete-iccad2023}) and
specification-to-RTL framing. The descriptive result remains 129 passes
of 312 = 0.4135, decomposed into 67/156 = 0.4295 for code-completion and
62/156 = 0.3974 for specification-to-RTL. Because paired framings share
a design, a Wilson interval treating all 312 outcomes as independent is
inappropriate. Resampling the 156 design pairs together gives a 95\%
percentile cluster-bootstrap interval of {[}0.343, 0.487{]} (200,000
replicates, seed 20261002). This describes design-composition
uncertainty under an exchangeable-design assumption for this checkpoint
and these fixed sampled generations; it is not training-seed,
decoding-seed or protocol uncertainty. The observed full
pinned-benchmark fraction itself remains descriptive. The surviving
receipt records 150 extraction failures included in the denominator, so
129/312 is an end-to-end harness score: 150 of its 183 failures are
extraction failures. Those failures are better attributed to the
generation budget than to answer format (corrected; see Table E.1). The
run used a 1,024-token cap with the chat template's default thinking
mode on, and the receipt's 312,622 response tokens give a mean of about
1,002 per sample, so nearly every response hit the cap, most of them
before closing the reasoning block. On the 50 items of the 26 September
small-scale rerun (§8.7), the receipt's verdicts give 23/50, while the
same trained adapter with thinking off and a 4,096-token cap scores
33/50 with no extraction failures (11 items pass only in the second run
and 1 only in the first; exact McNemar p ≈ 0.006). That comparison also
changes engine and temperature (managed sampler at 0.2 against vLLM at
0) and covers 50 of 312 prompts, so it shows the direction and rough
size of the budget effect rather than a corrected score. The extractor
also accepted a module found inside an unterminated reasoning block, and
per-sample response text was not kept, so some of the 162 extracted
candidates, and possibly some passes, may be drafts from the reasoning
rather than final answers. The parseable-only descriptive sensitivity is
129/162 = 0.796, hard to interpret for that reason; no interval is given
because the receipt's aggregate extraction count does not provide the
per-design extraction indicators needed for a clustered conditional
interval. One prompt,
\longtok{verilog\_eval/spec-to-rtl/Prob099\_\allowbreak{}m2014\_\allowbreak{}q6c}, fails against
its own test bench under both simulators; excluding it yields 129/311 =
0.4148, but no immutable justification for the exclusion was captured,
so the raw 129/312 is canonical and 129/311 is a noncanonical
sensitivity only. Pair counts, the bootstrap method, harness revisions,
verdict rule and amendment are in Appendix D.

\textbf{E14 --- Omni-MATH replacement scope, receipt-recorded
terminal-complete.} The surviving ledger and historical arithmetic check
record 4,426 accepted rows of 4,428 and official accuracy
0.5131043831902395, consistent with 2,271/4,426. Retaining 2,271 passes
but counting both unjudged rows as failures gives 2,271/4,428 = 51.29
per cent (Wilson 95\% CI {[}0.498, 0.528{]}), 0.02 percentage points
below the judged-row rate. The recovered
\longtok{outputs/PES\_\allowbreak{}Phase2\_\allowbreak{}Review\_\allowbreak{}2026-09-12/finish/e14\_\allowbreak{}terminal\_\allowbreak{}note\_\allowbreak{}2026-09-19.json}
records truncated judge outputs before an Equivalence Judgement section,
two recorded skip reasons, no re-judge and terminal completion at 4,428
dispositions. The note nevertheless labels \texttt{2271/4428} as
\texttt{0.5131043831902395}, a denominator inconsistency: that decimal
equals 2,271/4,426, not 2,271/4,428 (corrected; see Table E.1). The
original note is preserved unchanged, and this thesis retains both
clearly labelled arithmetic conventions. The supplementary Mac-side
traversal found the stored disposition, judged-row and
new-judge-response files with 4,428 valid JSON rows apiece. It reports
4,426 accepted dispositions, 2,271 correct and two skipped, and exports
both exclusion records with missing Equivalence Judgement headings. It
also records the native score's explicit 4,426 denominator and the
separate all-disposition diagnostic. The full E14 rows are not exported
here, so the cloud check verifies the exported arithmetic and exclusions
rather than independently recounting all 4,428 rows or rejudging any
mathematics. The headline is strongly shaped by the 2,048-token actor
cap (corrected reading; see Table E.1). Recounting the stored
dispositions, the 2,982 responses that hit the cap are correct in 1,029
cases (34.5\%), against 1,242 of the 1,446 uncapped responses (85.9\%).
The split is confounded by difficulty, since harder problems both run
longer and fail more, so it does not identify what a larger budget would
score; it does mean 51.29\% should not be read as the actor's Omni-MATH
capability independent of its generation budget.

\section{Original-contract results}\label{original-contract-results}

Seven lanes have an original-contract figure, reported separately from
any replacement scope with its coverage caveats; two are full-suite
results. Lane detail is in Appendix D.

\begin{itemize}
\tightlist
\item
  \textbf{E1 --- SWE-bench Pro.} 2 of 731 tasks resolved, a pass@1 of
  2/731 = 0.00274, or 0.274 per cent (Wilson 95\% CI {[}0.08\%,
  0.99\%{]}, unit = task), over the complete 731-row test split of
  \texttt{ScaleAI/SWE-bench\_Pro}. The fourteen generation failures and
  four lost artifacts stay in the denominator. The 713 ``native
  evaluations'' are not 713 tests of a candidate fix (corrected reading;
  see Table E.1). Parsed with \texttt{git\ apply}, 593 of the 713
  submitted patches are rejected as corrupt, because their \texttt{@@}
  hunk headers state line counts that do not match the hunk bodies; the
  official entry script runs a plain \texttt{git\ apply} and then the
  tests, so for those 593 the tests ran against the unpatched
  repository. With \texttt{git\ apply\ -\/-recount}, 711 of 713 parse.
  Both resolved instances are among the 120 patches that parse. The
  2/731 therefore mostly measures whether the actor could write
  arithmetically consistent diff headers, not whether it could solve the
  issues, and it is not a measure of problem-solving ability.
\item
  \textbf{E2 --- FrontierSWE.} 1 of 17 tasks evaluated,
  replay-normalised score 0.8628; partial, because revision-bound
  authorisation for the remaining sixteen was never obtained.
\item
  \textbf{E4 --- BankerToolBench.} 1 of 100 tasks evaluated, recovery
  metric 0.3115 --- a verifier recovery metric on a single task, not the
  simple fraction 37/128 and not a suite score; the pass3/pass16 archive
  is missing.
\item
  \textbf{E5 --- APEX-Agents.} 7 of 480 tasks natively scored. The
  prefix mean of 0.050505 treats unscored attempts as zero and is not a
  native 480-task score; the seven scored tasks average 0.079365.
\item
  \textbf{E7 --- BinaryAudit.} One of 46 attempts recorded, verifier
  reward 0.0 after an agent error, no grade.
\item
  \textbf{E9 --- MLE-bench.} 40 of 75 competitions carry native grades,
  a coverage of 53.33 per cent, and the suite score is null; a separate
  merged-vLLM arm holds one valid grade (H\&M, 0.02132) and is not
  combined.
\item
  \textbf{E11 --- VerilogEval} is reported in §8.2; it is the second
  original-contract full-suite result, alongside E1, and not an
  additional one.
\end{itemize}

No lane here has a matched base-model run, so none licenses a claim that
the actor outperforms its untrained counterpart, a predecessor model, or
a published number.

\section{Replacement-scope lanes completed on 27 September
2026}\label{replacement-scope-lanes-completed-on-27-september-2026}

At the 19 September ledger of record, six replacement scopes were
started but incomplete: E1 (SWE-bench Multilingual, 35 of 300 graded),
E2 (CORE-Bench, 45 capsule setups, no grades), E5 (Tau3, 20 of 97
cleaned) and E13 (BALROG, 13 of 255 episodes) were launch-pending, and
E6 (WebArena, 0 of 812) and E9 (MLDevBench, 0 of 34) were
\texttt{PENDING\_QUOTA} because required execution capacity was
unavailable. All six were run to completion on 27 September (Figure~\ref{fig:replacement}).


\fig{fig_replacement_results}
{The six replacement-scope results scored on 2026-09-27, each with its 95\% confidence interval}
{The six replacement-scope results scored on 2026-09-27, each with its 95\% confidence interval. (a) Rates over all attempted items, with errors, timeouts and ungraded items counted as failures (Wilson intervals). For E6 WebArena the filled point is the lower bound $90/812$; the square is the rate over the 704 judged tasks and the triangle the value if all 108 unjudged tasks had passed. (b) E13 BALROG reports native progression, not a success rate: per-environment means with percentile-bootstrap intervals ($B = 10{,}000$), and the BALROG overall score, the unweighted mean of the six environment means, with its normal-approximation interval. Each row is a different suite with its own metric and denominator, so no cross-lane average is drawn, and none of these values is pooled with an original-contract figure. Source: \longtok{outputs/finish\_\allowbreak{}pending\_\allowbreak{}2026-09-27/<lane>/result.json}.}
{fig:replacement}



The later launch used a lead-issued protocol override rather than
independent approval. That governance limitation matters: the candidate
was both the experiment lead and the author of the override. The public
derivative omits operational authorization records and account limits;
the scientific departures from the native protocols remain explicit in
Table 8.D.

\Needspace{10\baselineskip}\addcontentsline{lot}{table}{\protect\numberline{8.C}Replacement-scope results, 27 September 2026}\textbf{Table
8.C.} Replacement-scope results, 27 September 2026. One run per item,
native grader. Errors, timeouts and ungraded items count as failures in
the denominator. Graded counts items with a native verdict. Intervals
are Wilson 95\% for rates; E13's uses BALROG's standard error, and its
score is the mean of six per-environment means.

{\def\LTcaptype{none} % do not increment counter
\begingroup\scriptsize\setlength{\tabcolsep}{3pt}
\begin{longtable}[]{@{}
  >{\raggedright\arraybackslash\hspace{0pt}}p{(\linewidth - 14\tabcolsep) * \real{0.0563}}
  >{\raggedright\arraybackslash\hspace{0pt}}p{(\linewidth - 14\tabcolsep) * \real{0.2958}}
  >{\raggedright\arraybackslash\hspace{0pt}}p{(\linewidth - 14\tabcolsep) * \real{0.0986}}
  >{\raggedright\arraybackslash\hspace{0pt}}p{(\linewidth - 14\tabcolsep) * \real{0.0845}}
  >{\raggedright\arraybackslash\hspace{0pt}}p{(\linewidth - 14\tabcolsep) * \real{0.1549}}
  >{\raggedright\arraybackslash\hspace{0pt}}p{(\linewidth - 14\tabcolsep) * \real{0.1549}}
  >{\raggedright\arraybackslash\hspace{0pt}}p{(\linewidth - 14\tabcolsep) * \real{0.1549}}@{}}
\toprule\noalign{}
\begin{minipage}[b]{\linewidth}\raggedright
Lane
\end{minipage} & \begin{minipage}[b]{\linewidth}\raggedright
Benchmark (scope)
\end{minipage} & \begin{minipage}[b]{\linewidth}\raggedright
Attempted
\end{minipage} & \begin{minipage}[b]{\linewidth}\raggedright
Graded
\end{minipage} & \begin{minipage}[b]{\linewidth}\raggedright
Metric
\end{minipage} & \begin{minipage}[b]{\linewidth}\raggedright
Result
\end{minipage} & \begin{minipage}[b]{\linewidth}\raggedright
95\% CI
\end{minipage} \\
\midrule\noalign{}
\endhead
\bottomrule\noalign{}
\endlastfoot
E1 & SWE-bench Multilingual (190 tasks not attempted before) & 190 & 57
& resolved rate & 1/190 = 0.005 & {[}0.001, 0.029{]} \\
E2 & CORE-Bench hard (all 45 capsules) & 45 & 45 & task accuracy & 27/45
= 0.600 & {[}0.455, 0.730{]} \\
E5 & Tau3 banking (all 97 tasks) & 97 & 72 & pass\^{}1 & 8/97 = 0.082 &
{[}0.042, 0.154{]} \\
E6 & WebArena (all 812 tasks) & 812 & 704 & success rate, lower bound &
90/812 = 0.111 & {[}0.091, 0.134{]} \\
E9 & ML-Dev-Bench (all 34 tasks) & 34 & 34 & success rate & 10/34 =
0.294 & {[}0.168, 0.462{]} \\
E13 & BALROG (all 255 episodes, 6 environments) & 255 & 255 &
progression & 26.1\% & {[}22.6, 29.6{]} \\
\end{longtable}
\endgroup
}

\emph{Sources:
\longtok{outputs/finish\_\allowbreak{}pending\_\allowbreak{}2026-09-27/\textless{}lane\textgreater{}/result.json}
for each lane.}

\Needspace{10\baselineskip}\addcontentsline{lot}{table}{\protect\numberline{8.D}Deviations from the fail-closed protocol, 27 September 2026}\textbf{Table
8.D.} Deviations from the fail-closed protocol, 27 September 2026.
Condensed from the \texttt{deviations} field of each lane's result file;
the full text is in that file.

{\def\LTcaptype{none} % do not increment counter
\begin{longtable}[]{@{}
  >{\raggedright\arraybackslash\hspace{0pt}}p{(\linewidth - 6\tabcolsep) * \real{0.0385}}
  >{\raggedright\arraybackslash\hspace{0pt}}p{(\linewidth - 6\tabcolsep) * \real{0.8077}}
  >{\raggedright\arraybackslash\hspace{0pt}}p{(\linewidth - 6\tabcolsep) * \real{0.1538}}@{}}
\toprule\noalign{}
\begin{minipage}[b]{\linewidth}\raggedright
Lane
\end{minipage} & \begin{minipage}[b]{\linewidth}\raggedright
Deviation from reference protocol/scaffold
\end{minipage} & \begin{minipage}[b]{\linewidth}\raggedright
Recorded in
\end{minipage} \\
\midrule\noalign{}
\endhead
\bottomrule\noalign{}
\endlastfoot
E1 & Actor at 32,768-token context (65,536 on 12 Sept), per-task
\longtok{max\_tokens} cap; source-context collector re-implemented; thin
runner \longtok{run\_remaining.py} bypassing paperwork gates, native eval
in Modal sandboxes; batches r06/r07 harvested from orphaned sandboxes
(exec stdout/returncode lost); two driver relaunches reusing, never
resampling, completed generations; context-overflow short-circuit &
\longtok{E1/result.json} \\
E2 & Per-task GCP VMs instead of native Azure shapes; thin
\texttt{bash}/\longtok{query\_image}/\texttt{finish} loop instead of the
CORE-Agent scaffold (native 8,100 s budget, 150-turn cap, 26k-token
window); \longtok{query\_image} answered by the same actor;
infrastructure aborts rerun from scratch, none because of score &
\longtok{E2/result.json} \\
E5 & Thin runner \longtok{code/run\_tau3.sh} replaces the fail-closed
\longtok{e5\_successor27} controller; dedicated 65k-context actor copy;
OpenAI roles through a spend-capped relay & \longtok{E5/result.json} \\
E6 & GCP VM instead of the AWS AMI; thin runner
\longtok{code/e6\_driver.py} mirroring native \texttt{run.py};
\longtok{max\_tokens} 384, few-shot system turns mapped to
user/assistant; retired native judge swapped to \longtok{openai/gpt-4.1},
which was unavailable (108 tasks ungraded, counted as failures); 1,500 s
per-task timeout; last 103 shopping-admin tasks sharded over three
containers of the official image & \longtok{E6/result.json} \\
E9 & GCP instead of AWS; backported one-line OpenHands context-window
fix; loopback relay forcing non-thinking and restoring dropped tool-call
text; \longtok{max\_\allowbreak{}input\_\allowbreak{}tokens} 28,000; one run per task &
\longtok{E9/result.json} \\
E13 & Non-thinking template (original run had thinking on); NetHack
capped at 2,000 steps (never bound); macOS arm64 instead of the Linux
x86 image, dynamics not bit-verified; killed or hung episodes rerun from
scratch with new seeds & \longtok{E13/result.json} \\
\end{longtable}
}

\emph{Sources:
\longtok{outputs/finish\_\allowbreak{}pending\_\allowbreak{}2026-09-27/\textless{}lane\textgreater{}/result.json},
field \texttt{deviations}.}

The six figures are on six different benchmarks with six different
metrics, and they are not averaged, pooled with the same lane's
original-contract result, or compared with a base model. Five lanes (E2,
E5, E6, E9, E13) now have a native score over their full replacement
scope. E6's is a lower bound: 108 tasks needing an LLM judge could not
be graded, and the true figure lies between 0.111 and 0.244. E1 has two
separately reported scores, 1/190 and 4/110 (the earlier waves, served
at a different context), that together cover all 300 tasks but are not
combined. E5's pass\^{}1 over the 72 tasks without an infrastructure
error, 8/72 = 0.111, is a sensitivity figure only, and E13's 26.1 per
cent is not compared with the 26-episode small-scale E13 rows of Tables
8.A and 8.B. Per-lane narratives are in Appendix D.

\section{The E4 rerun as a
diagnostic}\label{the-e4-rerun-as-a-diagnostic}

A fresh full BankerToolBench rerun was reported to have run the base
model with a zero-initialised LoRA snapshot through the Modal bridge
\texttt{/v1/responses} path after the lane closed
\texttt{CLOSED\_PARTIAL} on 19 September. The surviving
\longtok{outputs/e1\_\allowbreak{}e14\_\allowbreak{}small\_\allowbreak{}scale\_\allowbreak{}2026-09-26/E4/result.json}
(\longtok{prior\_\allowbreak{}base\_\allowbreak{}run}) repeats 100 trials, one error and a mean
reward of 0.0. The recovered terminal ledger,
\texttt{withheld\ historical\ campaign\ chronology}, includes updates
through 21 September and explicitly says the rerun was NOT launched at
that snapshot. It therefore cannot support the later completed-run
account. The raw job receipt,
\longtok{outputs/e4\_\allowbreak{}banker\_\allowbreak{}toolbench/official\_\allowbreak{}repo\_\allowbreak{}ff6db552/jobs/btb-banking-tasks-tinker-bridge/result.json},
remains unrecovered in this package. The 100-trial result is supported
here by the surviving \longtok{prior\_\allowbreak{}base\_\allowbreak{}run} secondary record; the
earlier detailed claims of 2--14 agent steps, no deliverables and at
least 27 role-token-degenerate trajectories are not freshly verified.
The surviving bridge finding documents a route that passes no stop
sequences and can run past end-of-turn tokens into invented turns. Under
a corrected shim, the same base model scored 0.4087 on task
\texttt{btb-19b3361c} in the six-task small-scale rerun (§8.7). That
contrast demonstrates a plausible failure mechanism, but it compares one
successful corrected task with 100 bridge trials while changing tasks,
serving path and caps. It does not identify how much of the aggregate
zero the bridge caused or exclude model limitations. The 0.0 is retained
as a confounded diagnostic, not a controlled model-capability estimate.
It is not set against the original single-task outcome, and no valid
full-suite BankerToolBench capability estimate exists for this actor
family. The fuller reported audit and evidence limitations are in
Appendix D (Figure~\ref{fig:e4}).


\fig{fig_e4_diagnostic}
{The E4 rerun diagnostic}
{The E4 rerun diagnostic. Across 100 trials agents terminate at a median of roughly six steps without producing deliverables, and more than a quarter end in role-token degeneration, consistent with generation running past the end-of-turn token because the serving bridge passed no stop sequences. The causal share of the zero reward is unidentified: the corrected-path check changes task, path and caps together. Left: per-trajectory step counts (\texttt{final\_metrics.total\_steps}), stacked by whether the final agent message carries at least two concatenated role tokens (29/100; the ledger records ``27+''). Right: where the verifier chain stopped (\texttt{verifier/test-stdout.txt}, \texttt{reward.json}), all bars on one scale. Base-model rerun of 2026-09-21, all 100 trial directories under \longtok{outputs/e4\_\allowbreak{}banker\_\allowbreak{}toolbench/}. These raw trial artefacts are absent from the reviewed checkout; the graphic preserves the historical summary rather than a fresh receipt-level verification.}
{fig:e4}



\section{Externally blocked lanes and terminal
states}\label{externally-blocked-lanes-and-terminal-states}

Seven original benchmark contracts remain outside the completed
evaluation scope: SDAB, BinaryAudit, LifeSciBench, AgentHarm, AppBench,
OpenReward and FrontierMath. Their required benchmark payload or hosted
evaluation scope is not available in this public package. Six are
recorded as externally closed; OpenReward proceeds only through the
separate BALROG replacement scope. These statuses are scope limits and
not graded model failures. Private access correspondence, operational
requests and account state are omitted. The following small-scale reruns
keep original and substitute benchmarks separate (Figure~\ref{fig:taxonomy}).


\fig{fig_terminal_taxonomy}
{Terminal-state taxonomy used to close every lane}
{Terminal-state taxonomy used to close every lane. Each state has a defined meaning and a named reopen condition where one exists. The seven states are the ledger's full vocabulary; 17 lane-states cover 14 lanes because E8, E10 and E14 carry separate original and replacement rows. States classify lanes, not results: an unrecognised state is recorded as blocked rather than absorbed into a percentage, and original-contract and replacement-scope numbers are never pooled. The completion gate's \texttt{PARTIAL\_EXACT}/\texttt{PARTIAL\_RECOVERY} labels are a separate evidence-class vocabulary (Chapter~5). Lane assignments are those of 2026-09-19; the six lanes in the three pending leaves were all run and scored on 2026-09-27 (Chapter~8).}
{fig:taxonomy}



\section{Small-scale reruns of all fourteen
lanes}\label{small-scale-reruns-of-all-fourteen-lanes}

This section adds a separate, smaller set of runs, made on 26 September
2026, to give every one of the fourteen lanes a measured number under
one fixed protocol, including the lanes that never produced a score.
None of these figures is pooled with, averaged against, or substituted
for a result reported earlier in this chapter.

The protocol fixes one actor family, a fixed item-selection seed
(20260926), temperature 0 and a non-thinking chat template, and scores
each lane with its native grader where one could be run. Lanes whose
original benchmark is externally blocked were run on a public
substitute, and each substitute's gap from the original suite is
recorded in its lane file. Items that errored or timed out are counted
as failures in the denominator.

Two arms were run. The first is the base model,
\texttt{Qwen/Qwen3.6-35B-A3B} with no adapter, sampled on Tinker (Table
8.A). The second is a paired comparison of the trained actor against the
same base model, both served by one vLLM configuration on Modal, on
identical items (Table 8.B). The pairing exists because the adapter's
effect on individual token probabilities was measured to be about the
same size as the numerical difference between the two serving engines on
the same base weights, so a trained-on-vLLM against base-on-Tinker
comparison would confound the adapter with the engine, and none is made.

\Needspace{10\baselineskip}\addcontentsline{lot}{table}{\protect\numberline{8.A}Base model on Tinker, small-scale subsets}\textbf{Table
8.A.} Base model on Tinker, small-scale subsets.

{\def\LTcaptype{none} % do not increment counter
\begingroup\scriptsize\setlength{\tabcolsep}{3pt}
\begin{longtable}[]{@{}
  >{\raggedright\arraybackslash\hspace{0pt}}p{(\linewidth - 12\tabcolsep) * \real{0.1667}}
  >{\raggedright\arraybackslash\hspace{0pt}}p{(\linewidth - 12\tabcolsep) * \real{0.1667}}
  >{\raggedright\arraybackslash\hspace{0pt}}p{(\linewidth - 12\tabcolsep) * \real{0.1667}}
  >{\raggedright\arraybackslash\hspace{0pt}}p{(\linewidth - 12\tabcolsep) * \real{0.1667}}
  >{\raggedright\arraybackslash\hspace{0pt}}p{(\linewidth - 12\tabcolsep) * \real{0.1667}}
  >{\raggedright\arraybackslash\hspace{0pt}}p{(\linewidth - 12\tabcolsep) * \real{0.1667}}@{}}
\toprule\noalign{}
\begin{minipage}[b]{\linewidth}\raggedright
Lane
\end{minipage} & \begin{minipage}[b]{\linewidth}\raggedright
Benchmark run
\end{minipage} & \begin{minipage}[b]{\linewidth}\raggedright
Scope
\end{minipage} & \begin{minipage}[b]{\linewidth}\raggedright
Metric
\end{minipage} & \begin{minipage}[b]{\linewidth}\raggedright
Result
\end{minipage} & \begin{minipage}[b]{\linewidth}\raggedright
95\% CI
\end{minipage} \\
\midrule\noalign{}
\endhead
\bottomrule\noalign{}
\endlastfoot
E1 & SWE-bench Pro public test split & original-public-subset & resolved
rate & 0/10 = 0.000 & {[}0.00, 0.28{]} \\
E2 & EffiBench & substitute & pass@1 & 23/30 = 0.767 & {[}0.59,
0.88{]} \\
E3 & Terminal-Bench v2.0 & substitute & task success rate & 4/8 = 0.500
& {[}0.22, 0.78{]} \\
E4 & BankerToolBench (official Harbor tasks) & original-public-subset &
mean verifier reward & 0.068 (n = 6; one nonzero task, 0.409) & --- (not
estimable) \\
E5 & tau2-bench & substitute & pass\^{}1 & 3/10 = 0.300 & {[}0.11,
0.60{]} \\
E6 & MiniWoB++ & substitute & episode success rate & 15/36 = 0.417 &
{[}0.27, 0.58{]} \\
E7 & BinaryAudit upstream cbd86c7 & original-public-subset & task
success rate & 0/6 = 0.000 & {[}0.00, 0.39{]} \\
E8 & LAB-Bench public MCQ & substitute & accuracy & 26/80 = 0.325 &
{[}0.23, 0.43{]} \\
E9 & MLE-bench lite/low-complexity split & original-public-subset &
any-medal rate & 1/5 = 0.200 & {[}0.04, 0.62{]} \\
E10 & AgentHarm test\_public & original-public-subset & harmful refusal
rate & 20/30 = 0.667 & {[}0.49, 0.81{]} \\
E11 & VerilogEval & original-public-subset & pass@1 & 35/50 = 0.700 &
{[}0.56, 0.81{]} \\
E12 & AppBench public tasks & original-public-subset & rubric-item pass
rate & 122/151 = 0.808 & {[}0.74, 0.86{]} \\
E13 & BALROG & substitute & mean progression (\%) & 23.1\% & --- \\
E14 & Omni-MATH test & substitute & accuracy & 38/100 = 0.380 & {[}0.29,
0.48{]} \\
\end{longtable}
\endgroup
}

\Needspace{10\baselineskip}\addcontentsline{lot}{table}{\protect\numberline{8.B}Trained actor versus base model on the same vLLM engine, identical items}\textbf{Table
8.B.} Trained actor versus base model on the same vLLM engine, identical
items.

{\def\LTcaptype{none} % do not increment counter
\begingroup\scriptsize\setlength{\tabcolsep}{3pt}
\begin{longtable}[]{@{}
  >{\raggedright\arraybackslash\hspace{0pt}}p{(\linewidth - 12\tabcolsep) * \real{0.1667}}
  >{\raggedright\arraybackslash\hspace{0pt}}p{(\linewidth - 12\tabcolsep) * \real{0.1667}}
  >{\raggedright\arraybackslash\hspace{0pt}}p{(\linewidth - 12\tabcolsep) * \real{0.1667}}
  >{\raggedright\arraybackslash\hspace{0pt}}p{(\linewidth - 12\tabcolsep) * \real{0.1667}}
  >{\raggedright\arraybackslash\hspace{0pt}}p{(\linewidth - 12\tabcolsep) * \real{0.1667}}
  >{\raggedright\arraybackslash\hspace{0pt}}p{(\linewidth - 12\tabcolsep) * \real{0.1667}}@{}}
\toprule\noalign{}
\begin{minipage}[b]{\linewidth}\raggedright
Lane
\end{minipage} & \begin{minipage}[b]{\linewidth}\raggedright
Items
\end{minipage} & \begin{minipage}[b]{\linewidth}\raggedright
Trained
\end{minipage} & \begin{minipage}[b]{\linewidth}\raggedright
Base
\end{minipage} & \begin{minipage}[b]{\linewidth}\raggedright
Trained − base
\end{minipage} & \begin{minipage}[b]{\linewidth}\raggedright
Test
\end{minipage} \\
\midrule\noalign{}
\endhead
\bottomrule\noalign{}
\endlastfoot
E1 & 10 & 0.000 & 0.000 & 0.000 & McNemar p = 1 (b/c 0/0) \\
E2 & 30 & 0.700 & 0.800 & -0.100 & McNemar p = 0.25 (b/c 0/3) \\
E3 & 8 & 0.625 & 0.375 & 0.250 & McNemar p = 0.5 (b/c 2/0) \\
E4 & 6 & 0.066 & 0.093 & -0.027 & one nonzero pair (−0.160); no
interval \\
E5 & 10 & 0.500 & 0.200 & 0.300 & McNemar p = 0.38 (b/c 4/1) \\
E6 & 36 & 0.444 & 0.472 & -0.028 & McNemar p = 1 (b/c 1/2) \\
E7 & 6 & 0.167 & 0.000 & 0.167 & McNemar p = 1 (b/c 1/0) \\
E8 & 80 & 0.312 & 0.312 & 0.000 & McNemar p = 1 (b/c 6/6) \\
E9 & 5 & 0.000 & 0.000 & 0.000 & McNemar p = 1 (b/c 0/0) \\
E9 (above median) & 5 & 0.000 & 0.200 & -0.200 & McNemar p = 1 (b/c
0/1) \\
E10 (refusal rate) & 30 harmful, 30 benign & 0.700 & 0.667 & 0.033 &
McNemar p = 1 (b/c 3/2) \\
E10 (harm score) & 30 harmful, 30 benign & 0.247 & 0.236 & 0.011 & 95\%
CI {[}-0.053, 0.078{]} \\
E10 (benign score) & 30 harmful, 30 benign & 0.780 & 0.794 & -0.015 &
95\% CI {[}-0.050, 0.012{]} \\
E11 & 50 & 0.660 & 0.660 & 0.000 & McNemar p = 1 (b/c 2/2) \\
E12 & 151 & 0.748 & 0.781 & -0.033 & McNemar p = 0.3 (b/c 5/10);
app-level sign-flip p = 0.25 (n = 6) \\
E13 & 26 & 25.510 & 23.630 & 1.875 & 95\% CI {[}-2.73, 7.09{]} \\
E14 & 100 & 0.430 & 0.440 & -0.010 & McNemar p = 1 (b/c 5/6) \\
\end{longtable}
\endgroup
}

These are small samples. Most lanes score between five and eighty items,
so a single item moves a lane's rate by several percentage points and
the confidence intervals are wide. In Table 8.B, b and c count items
solved only by the trained actor and only by the base model
respectively, and the CI is a paired bootstrap interval on the
difference. Four lane-specific caveats matter. For E12 the 151 rubric
items come from only six generated applications, so its McNemar p-value
overstates the evidence; the exact sign-flip test over the six per-app
differences gives p = 0.25, and the grader is an LLM judge that is the
base model itself, so the E12 rates are judge-dependent. For E4, only
\texttt{btb-19b3361c} scores above zero in either arm (trained 0.397,
base 0.557; 0.409 in Table 8.A), so no interval is valid. For E13, the
1.875 difference is the mean of per-environment differences computed
from unrounded per-item values (recomputed 1.8752), not 25.510 − 23.630.
For E9, the one competition that earned a medal in the Tinker run
(nomad2018) has a 46.8k-token prompt that exceeds the 32,768-token vLLM
limit and fails identically in both vLLM arms, so the vLLM medal rates
are not comparable with Table 8.A. The paired differences are not
combined across lanes, and none of the seventeen paired tests excludes
zero. At these sizes that is an inconclusive result, not evidence of no
effect: with 5 to 10 items, an exact McNemar test needs at least six
discordant items, all in one direction, to reach p \textless{} 0.05. No
equivalence test was run. Fuller E12 detail and the principal per-lane
caveats from each result file are in Appendix D.

\section{Summary of the campaign at the point of
writing}\label{summary-of-the-campaign-at-the-point-of-writing}

By the 19 September ledger the campaign had established two full-suite
results and one complete replacement scope with a score: E1 at 2/731 =
0.274 per cent on SWE-bench Pro, E11 at 129/312 = 41.35 per cent pass@1
(design-cluster bootstrap 95\% interval {[}34.3, 48.7{]}, n = 156
designs) across its two framings, and E8 at 450/1967 = 0.2288 ({[}0.211,
0.248{]}) on the LAB-Bench public split, alongside a complete
benign-utility evaluation of E10 at 88 of 97 episodes (0.907, {[}0.833,
0.950{]}), a terminal-complete E14 replacement scope at 2,271/4,428 =
51.29 per cent accuracy, and five original-contract partials (E2, E4,
E5, E7, E9), none of which is a suite score. Three of those scores are
dominated by the harness rather than the actor: 593 of E1's 713
submitted patches were corrupt for \texttt{git\ apply}, E11's 150
extraction failures follow from a 1,024-token cap with thinking on, and
E14's capped responses are correct at 34.5\% against 85.9\% uncapped
(§8.2, §8.3). On 27 September six further replacement scopes were
completed under a self-issued override whose deviations are recorded per
lane (Tables 8.C and 8.D). Seven original contracts remain externally
blocked, and no lane now waits on a launch or a cloud quota. The
campaign did not establish a multi-seed estimate, a held-out security
result, a BankerToolBench suite score, a graded result for the 108
WebArena tasks that need an LLM judge, or a baseline comparison for the
campaign lanes themselves; the only trained-versus-base comparison,
Table 8.B, is inconclusive in every lane. The E8 diagnostic, E10 native
receipt, training-identity verification, E14 terminal note and both
pending-state snapshots are recovered in the private package while
remaining absent from the frozen repository. Their receipt contents are
now inspectable, and bounded stored-payload checks corroborate E8/E10
counts and the E14 denominator distinction; training-source, E4
raw-evidence and semantic-reproduction limits remain as in §8.2 and
Appendix D. The available evidence chains, reproduction paths and gaps
are given in Appendix B; the next chapter draws the cross-study
synthesis and conclusions.



\part{Synthesis}


\chapter{Cross-Study Synthesis, Threats to Validity, Conclusions and
Future
Work}\label{cross-study-synthesis-threats-to-validity-conclusions-and-future-work}

\section{Cross-Study Synthesis}\label{cross-study-synthesis}

This report has presented two bodies of work built with different
instruments. The first is the P1--P8 corpus, which probes GRPO-style
post-training across libraries, model scales, group sizes and datasets
and ends in a reporting standard, a machine-readable registry and a
controller audit. The second is the E1--E14 campaign, a governance-bound
evaluation of one Tinker-trained actor against fourteen agentic and
reasoning suites outside its recorded training-suite list; the recovered
identity record does not establish content-level decontamination. They
are the same methodological commitment applied twice --- once to
\emph{what does GRPO do}, once to \emph{what does a GRPO-trained actor
score}. That commitment is attribution: every comparative claim holds
the entire training stack fixed except the one factor under test, and
where that was unachievable the confound is reported alongside the
number rather than absorbed.

The corpus's own synthesis, recorded at the close of Phase 1, is that on
this compute scale the \emph{measurement} layer is useful while most
\emph{algorithmic} levers are noise-limited. The dominant measured
phenomenon is advantage-signal starvation, with the Zero-Variance
Fraction at approximately 0.72--0.77 on the Phase-1 main configuration,
and it is easy-driven: the collapsed groups are mostly prompts the model
has already solved. That rules out the most tempting remedy, because
re-baselining the advantages merely converts the signal into a
difficulty proxy with no learning content, and the remedy that remains
--- a curriculum that removes the wasted rollouts --- returned a null.
The narrow systems observation survived better than any algorithmic one:
per-step telemetry separated the available method traces more cleanly
than reward-only features, though the classifier's integrity
generalisation was never tested.

The Semester 4 internal reports sharpen those threads rather than
reversing them. P2's diagnostics are recomputable from stored raw
tensors rather than read back from logged summaries: three stored GSM8K
group tensors reproduce ZVF at 0.130, 0.190 and 0.155, pooling to
0.1583, and the correlation between ZVF and final held-out outcome, ρ =
0.27, is not significant (95\% CI {[}−0.37, 0.88{]}, n = 23 pooled
cells). ZVF is therefore offered as a diagnostic and not a predictor.
Once mean reward is regressed out, the pooled measured residual is r =
+0.21 (95\% CI {[}−0.31, +0.62{]}, n = 17), so ZVF is not an incremental
predictor either (Chapter 6 §6.3.4). P4's matched comparison of GRPO
against Dr.~GRPO finds the two held-out gains indistinguishable, with
pooled mean completion length drifting from 194.4 to 183.5 tokens under
a 200-token cap that bounds what the comparison can say. P4's companion
error-bar audit grades the single P4-facing headline as DECISIVE: audit
H7, the iter-136 arithmetic late-training-efficiency contrast, paired n
= 5 seeds, \(d_z = +2.68\). Re-examined, it is \emph{suggestive} only
(Chapter 6 §6.5.4). Its ``one-sided parametric p = 0.031'' is the floor
of an exact one-sided Wilcoxon test at n = 5. The two-sided exact
sign-flip p is 0.0625, even though all five pairs point the same way.
The direction was one of eight hypotheses in its family, and Holm
correction across those eight raises the family's p to 0.25. The audit's
other three DECISIVE grades do not survive either. Two (H1, H2) rest on
an interval that excludes its own point estimate and on four
reconstructed budgets (§6.4.2; Table E.1), and one (H6) is an AUROC
computed on a synthetic and imputed panel (§6.3.4). Of the seven audited
headlines, then, none is decisive, one is suggestive and the rest are
null or hypothesis-generating. What survives is a consistent-direction
late-training efficiency contrast on a training metric, not a capability
claim.

P3 and P7 both treat group size as a lever and are best read as a pair,
because they disagree in emphasis. P3 finds that under saturation the
retention ratio between G = 2 and G = 16 is approximately 1.00. That
ratio is of last-10 training reward, and a seed-level TOST (n = 3,
margin ±0.02) does not establish G = 16 as equivalent to G = 2 (p =
0.056; Chapter 6 §6.4.1). P3 also finds that the effective
signal-to-noise ratio reaches only about 52\% of the √G ideal, and that
a 505-task conditional utility audit selects G = 4 over G = 5 by 0.00177
with a 95\% bootstrap interval of {[}0.00082, 0.00270{]} --- positive,
small, and explicitly cohort- and objective-conditional. P7 models
usable signal per prompt as \(S = p(1-p)(1-h_G(p))\) with
\(h_G(p) = p^G + (1-p)^G\), derives three falsifiable consequences, and
checks them: a 368-run Weights \& Biases audit across six model labels
(Appendix A, row P7-C1) reproduces the predicted U-shape in accuracy,
with ZVF 0.95--0.97 at the low-reward extreme and 0.87--0.96 at the
high-reward extreme and 0.25--0.29 at mid-range rewards of 0.35--0.50; a
model × G grid confirms the monotone effect at interior accuracy,
Qwen3-32B moving 0.33 → 0.18 → 0.08 across G = 4/8/16; and a toy-scale
open-backward-pass experiment finds corr(‖∇‖, p(1−p)) = +0.71 on a 0.5B
model, which the P7 report labels directional only. Where P3's sweep
could not find an optimum, P7 supplies a mechanism under which a shallow
one should exist. That is a reconciliation, not a resolution: the
mechanism is checked in telemetry, and its causal link to the gradient
remains toy-scale.

The same report is clear about the control it does not have. A
micro-jitter test collapses the exact-tie fraction from 0.158 to 0.000
while barely changing PCD, but the declared variance-threshold ZVF
remains 0.158 (§6.3.1). It therefore does not establish superiority of
PCD over tolerance-based ZVF. In a four-arm open-trainer audit, adaptive
group size matches the best fixed-recipe held-out gain, +0.575, at 186
rollouts, while DAPO-style dynamic sampling drives ZVF to 0.00 at +45\%
rollout cost --- an accuracy-versus-compute trade, not a controller win.
In the larger 2,560-observation controller trace, 1,723 of 1,867
escalation events, 92.3\%, occur on all-correct groups and only 144 on
all-wrong, so a symmetric ZVF trigger spends most of its budget on
observed all-correct groups. Eight successes do not prove latent p = 1,
and the plug-in simulation's zero rescue is imposed by its estimator.
Deleting those fires remains counterfactual because performance
preservation was not measured. Adaptive group size is therefore not
promoted to a recommended policy.

A separate frozen-checkpoint diagnostic now observes fresh responses
directly (§7.4; Appendix B §B.12): 33/461 initially homogeneous groups
became reward-mixed, 7.16\% {[}5.14\%, 9.88\%{]}. Twenty first changes
are strict-parser format failures with the same visible gold number;
only 13/461 (2.82\%) are assistant-audited numeric contrasts, comprising
two recoveries and eleven new failures. The format effect belongs to
this diagnostic contract, not the repository's more permissive
historical evaluator. The numeric witnesses reject certainty from eight
equal observations on this new checkpoint/task/decoder, but supply no
historical-checkpoint reproduction, training or controller benefit,
group-size optimum, or held-out capability estimate from the GSM8K train
split.

P5 and P6 supply the corpus's methodological claim, and they contain its
sharpest internal tension. P5's headline exhibits are a nominally
matched-configuration backend audit whose final training rewards span
5.0\% to 85.6\% --- a 17× same-label swing the P5 report declines to
assign to the backend, because the managed stack also silently pinned a
different base checkpoint (Qwen3-8B-Base on Tinker against the
instruction-tuned Qwen3-8B on TRL), and that undisclosed bundling is
itself part of the exhibit --- and a pair of runs carrying the same
``DAPO'' label in which an open trainer with true dynamic sampling
produced mean ZVF 0.00 while a closed stack running an asymmetric-clip
surrogate produced mean ZVF 0.58. Alongside them sits a third exhibit: a
12-cell head-to-head of four labelled algorithms on one fixed stack,
landing within a 0.034 band with mean last-10 reward 0.710--0.744, a
spread the P5 report calls comparable to seed-level variation. P5's
framing is that a label-only comparison can hide stack-level variation,
and that the converse also holds: when the stack is genuinely fixed, the
labels stop separating. P6 reads the same exhibit in the reverse
direction, using it to motivate a schema while explicitly declining to
re-claim it as novel empirical science; on the twenty seed entries the
MIN-REPORT auditor badge spans 43--96, and \emph{stackdiff} flags the
open-versus-closed ``DAPO'' pair as an R5 label-flip risk from manifests
alone.

Two gaps should be named. The first is that P7's controller audit and
P5's stack-attribution programme have not been joined: the controller is
evaluated on open trainers where the update rule is auditable, while the
standard is motivated by a closed-stack divergence that cannot be
audited. The second is P8, the credit-card fraud side-probe, absorbed
here as a measurement-discipline case study and nothing more. Its two
AUCs are not a like-for-like ranking: XGBoost's 0.7955 and the
Qwen3.5-4B SFT arm's 0.48268 on a positive-enriched 500-row held-out
evaluation sit on different held-out splits by construction.

The E1--E14 campaign is a different instrument and should be read as
one. Its figures have unequal source accessibility. The six original
files covering five missing basenames listed in Appendix B §B.11 are now
preserved in a withheld source archive, while remaining absent from the
frozen checkout. Supplementary read-only payload checks corroborate
stored E8/E10 counts and E14 disposition/denominator reporting; they do
not rerun a model or scorer, independently judge correctness, establish
content decontamination or recover the missing training logs. A single
Tinker-trained actor --- base \texttt{Qwen/Qwen3.6-35B-A3B} at commit
\texttt{995ad96e…}, with adapter \texttt{campaign-adapter} at commit
\texttt{withheld-adapter-revision…}, in bfloat16 --- was evaluated
against fourteen suites outside its recorded training-suite list, under
per-lane native evaluation contracts. The recovered AgentDojo receipt
binds version v1.2.2 but no upstream commit; the broader revision claims
retain their source-specific limits (§8.2). Three scopes closed strictly
complete: two replacement scopes, LAB-Bench public at 1,967/1,967 across
eight categories and AgentDojo benign utility with 97/97 episodes
completed and utility 88/97 = 0.907 (Wilson 95\% CI {[}0.833, 0.950{]}),
and one original-contract suite, VerilogEval, at 312/312 with 129/312 =
41.35\% pass@1 canonical (design-cluster bootstrap 95\% CI {[}34.3\%,
48.7\%{]}, resampling 156 paired designs (§8.2)). Omni-MATH closed
terminal-complete under its native protocol, with all 4,428 dispositions
recorded, 4,426 accepted, and an accuracy of 2,271/4,428 = 51.29\% with
the two unjudged rows counted as failures (the stored native-score
fields give 51.31\% over the 4,426 judged rows). The recovered terminal
note incorrectly associates that decimal with denominator 4,428; the
thesis retains the corrected arithmetic and preserves the source
unchanged (Table E.1; sources: recovered E14 terminal note and
\texttt{stored-payload\ source\ audit}, Appendix B §B.11). Two suites
carry full-coverage scores under their original contracts: SWE-bench Pro
at 2/731 = 0.274\% pass@1 (Wilson 95\% CI {[}0.08\%, 0.99\%{]}) and
VerilogEval, the same E11 result counted once. Both, and Omni-MATH, are
dominated by the harness configuration: 593 of 713 SWE-bench Pro patches
were corrupt for \texttt{git\ apply}, and the VerilogEval and Omni-MATH
scores are bound by their generation budgets (§8.2, §8.3). Six further
replacement scopes were run in full on 27 September 2026 (Table 8.C):
CORE-Bench hard at 27/45 = 0.600, Tau3 at 8/97 = 0.082, WebArena at
90/812 = 0.111 (a lower bound, with 108 judge-dependent tasks ungraded),
ML-Dev-Bench at 10/34 = 0.294, BALROG at 26.1\% progression over 255
episodes, and SWE-bench Multilingual at 1/190 on the tasks not attempted
before, kept separate from the earlier 4/110. Lane scope and status are
tabulated once, in Table C.3. The rest are partial or externally
blocked, each with a terminal state and a named gate.

What the campaign establishes is descriptive, and its scope matters. For
\emph{this} actor, at \emph{this} revision, under \emph{each suite's
own} evaluator and denominator, these are the scores. It is not a
post-training delta: no per-suite base-model comparator appears in the
campaign ledger at the campaign's own scopes, so no campaign number,
high or low, is evidence about what the GRPO step contributed. The one
paired comparison of trained actor and base model is the 26 September
small-scale rerun (Table 8.B): 14 lanes, 5 to 151 identical items each,
one serving engine. It finds no significant difference in any lane,
which at those sizes is inconclusive rather than evidence of no effect.
The surviving E4 \longtok{prior\_\allowbreak{}base\_\allowbreak{}run} secondary record reports a
BankerToolBench base-model rerun with mean reward 0.0 across 100 trials.
The recovered finish ledger explicitly says the rerun was not yet
launched at its 21 September snapshot, so it is not evidence for the
later completion. The bridge finding documents omitted stop sequences
and reports invented turns; the detailed historical claims of
termination after 2--14 steps without \texttt{deliverables/} output
remain unverified against the still-unrecovered raw job receipt and
trajectories. One corrected-path task achieved positive reward,
supporting a concrete harness defect, but task, serving path and caps
changed together; the fraction of the 100-trial zero attributable to
that defect is unidentified. The result describes that
model-plus-harness configuration (Section 8.5; sources:
\longtok{outputs/e1\_\allowbreak{}e14\_\allowbreak{}small\_\allowbreak{}scale\_\allowbreak{}2026-09-26/E4/result.json},
\longtok{prior\_\allowbreak{}base\_\allowbreak{}run};
\longtok{outputs/e1\_\allowbreak{}e14\_\allowbreak{}small\_\allowbreak{}scale\_\allowbreak{}2026-09-26/E4/BRIDGE\_\allowbreak{}STOP\_\allowbreak{}TOKEN\_\allowbreak{}FINDING.md}).

What governance bought is visible in the ledger itself: eleven of eleven
named deterministic checks passing on 2026-09-19 (scope in Appendix B),
and the recovered finish ledger preserving a historical snapshot
recording E7 as ``1/28 errored'' visibly beside the superseding table's
``1/46 attempted'', without inventing a reconciliation.

\section{Rigour and Threats to
Validity}\label{rigour-and-threats-to-validity}

\textbf{What rigour was exercised.} Diagnostics are recomputable from
stored raw tensors, and headline structural claims were re-verified by
eleven named deterministic code checks on 2026-09-19, all passing,
covering exact divisions, receipt presence, sealed-artifact hash
equality and lost-source absence probes. Held-out \emph{n} is reported
alongside every gain, and failed runs are counted rather than hidden:
failed attempts remain part of the stated limitations.

Table 9.1 lists the threats discussed below, grouped as internal,
construct and external, with the mitigation applied to each. Most
mitigations narrow what is claimed rather than repair the evidence.

\textbf{Threats to validity and the mitigation applied to each.}

\begin{longtable}[]{@{}
  >{\raggedright\arraybackslash\hspace{0pt}}p{(\linewidth - 6\tabcolsep) * \real{0.3333}}
  >{\raggedright\arraybackslash\hspace{0pt}}p{(\linewidth - 6\tabcolsep) * \real{0.3333}}
  >{\raggedright\arraybackslash\hspace{0pt}}p{(\linewidth - 6\tabcolsep) * \real{0.3333}}@{}}
\caption{Threats to validity and the mitigation applied to each. † marks
figures drawn from Tinker data only; replicating them requires Tinker
API access (\protect\longtok{LIMITATIONS\_\allowbreak{}AND\_\allowbreak{}IMPACT.md}
§6.3).}\tabularnewline
\toprule\noalign{}
\begin{minipage}[b]{\linewidth}\raggedright
Type
\end{minipage} & \begin{minipage}[b]{\linewidth}\raggedright
Threat
\end{minipage} & \begin{minipage}[b]{\linewidth}\raggedright
Mitigation
\end{minipage} \\
\midrule\noalign{}
\endfirsthead
\toprule\noalign{}
\begin{minipage}[b]{\linewidth}\raggedright
Type
\end{minipage} & \begin{minipage}[b]{\linewidth}\raggedright
Threat
\end{minipage} & \begin{minipage}[b]{\linewidth}\raggedright
Mitigation
\end{minipage} \\
\midrule\noalign{}
\endhead
\bottomrule\noalign{}
\endlastfoot
Internal & T1. Most Tinker configurations ran once, so those runs carry
no variance estimate or significance test; no result family reaches the
ten seeds Henderson et al.~(2018) recommend & Quantitative conclusions
rest on multi-seed runs (three to eight seeds, listed below);
single-seed Tinker figures are descriptive and marked † \\
Internal & T2. Horizons of 30--50 steps are early-training snapshots &
Every claim is scoped to the observed horizon; no asymptotic or
``converges to'' claim is made \\
Internal & T3. The E4 pass3/pass16 per-task archive is lost; only
\longtok{TINKER\_BRIDGE.json} survives & 0.3115 is reported as a
receipt-only verifier recovery metric on 1 of 100 tasks
(\longtok{CLOSED\_PARTIAL}), not as a suite score and not as the 37 of
128 verifier criteria met \\
Construct & T4. Each lane uses its own native denominator (2/731,
90/812, 4,428 \ldots) & No pooling within or across lanes; no comparison
or ranking across denominators \\
Construct & T5. No cross-suite aggregate exists & Deliberate: an
aggregate over native denominators has no referent \\
Construct & T6. E10's completed scope measures benign utility only & The
benign scope is named at every reference; no adversarial-robustness
claim \\
Construct & T7. The retained E11 result used an original managed sampler
route that is not reproducible from this public package & The lane is
receipt-bound to its sampler path; no backend-parity claim \\
External & T8. Tinker's server-side loss, normalisation, batching and
hardware are closed, and the SDK version moved (0.16.1 to 0.30.0) & The
stack and SDK version are reported with each result; no numerical-parity
claim \\
External & T9. Six original contracts lack third-party access (E3, E7,
E8, E10, E12, E14) & Each lane has a terminal close-out record with a
reopen condition; replacement scopes are reported under their own
names \\
\end{longtable}

\textbf{Statistical power: single-seed Tinker experiments.} Most Tinker
configurations ran once, because cost precluded repetition, and those
single-seed runs carry no variance estimates or significance tests. The
exceptions are stated where they occur. They are the three-seed-per-arm
managed head-to-head (Chapter 7 §7.1), the five-seed Qwen3-8B held-out
control (p = 0.26, Chapter 6 §6.5.4) and the Qwen3-8B scale extension.
Henderson et al.~(2018) recommend at least ten seeds, and no result
family reported here reaches that bar. The seed counts are: one seed for
most Tinker runs and all P1 anchors; three for the P3 sweep, the P4
GSM8K-CoT and scale-extension arms and the managed head-to-head; five
for the same-stack PPO/GRPO contrast, the P4 arithmetic arms and the
Qwen3-8B pre/post control; and eight for the variant audit (P11). The
limitations document says the ten-seed bar was met for the TRL baseline
and the GSM8K evaluation and then that the evaluation was cut to five
seeds; the run records in Appendix A show five seeds for both the TRL
Modal campaign (Table A.3) and the GSM8K pre/post evaluation. The 30B
MoE scaling analysis uses three seeds; the bootstrap confidence
intervals assume independent observations, so any positive correlation
from shared initialisation or data ordering would make the reported
intervals too narrow, and the true intervals wider; and secondary
analyses carry uncorrected p-values. Quantitative conclusions therefore
come from the open-source side, where every hyperparameter is auditable,
while Tinker figures are descriptive.

\textbf{Horizons of 30--50 steps.} The Tinker runs are early-training
snapshots; long-horizon effects such as reward hacking, catastrophic
forgetting and late-stage policy collapse are unlikely to manifest, and
the source conclusion is direct: ``Asymptotic claims are unsupported''.
This bounds results rather than vitiating them: P1's failure to recover
a scaling law across roughly 2.4 orders of magnitude is a bounded
negative, and P4's drift from 194.4 to 183.5 mean completion tokens is
bounded by the 200-token cap.

\textbf{Closed Tinker internals.} The Tinker SDK exposes a custom loss
path, a limited optimiser surface and standard LoRA hyperparameters, but
not the exact server-side GRPO loss formulation, the reward
normalisation or baseline-subtraction scheme, minibatch or
gradient-accumulation strategy, hardware configuration, or system-level
telemetry. This applies unevenly: the inspected custom-loss Tinker
reference loops define their loss in client code (REINFORCE with a
group-mean baseline, Table 4.2), while the actor's missing original
execution source limits verification of its historical loss path
(Appendix D), whereas the runners that use Tinker's built-in
\texttt{importance\_sampling} loss leave its server-side details
unrecorded. The consequence is best quoted: Tinker results measure ``the
\textbf{platform's} implementation of GRPO, not an abstract algorithmic
specification''. The illustrative case in the limitations document
compares a pooled last-10 training reward of 99.9\% for Tinker GRPO with
73.4\% for TRL GRPO on GSM8K. Those figures are training rewards, not
held-out accuracy, and they pool Tinker runs on 4B--235B models against
TRL on Qwen2.5-0.5B, so model size and implementation are confounded;
the source itself attributes part of the gap to the model difference.
The only same-model pair, Qwen3-8B, used the Base checkpoint on Tinker
and the Instruct checkpoint on TRL (Table A.2), so it does not remove
the confound either. The threat propagates into the campaign, whose
actor was trained on Tinker. The SDK is also a moving target: Phase 1
pinned \texttt{tinker==0.16.1}, while the September 2026 bridge required
\texttt{tinker==0.30.0} after the older SDK produced a
no-longer-supported error.

\textbf{Per-suite native denominators are not comparable.} Every E1--E14
result is expressed on its own suite's denominator, and those
denominators differ in size, in what they count and in what they
exclude. SWE-bench Pro is 2/731, whose denominator retains 14 generation
failures and 4 artifact losses alongside 713 native evaluations, 593 of
which ran on the unpatched base because the submitted patch was corrupt
for \texttt{git\ apply} (§8.3); APEX-Agents is 7/480; E9 MLE-bench is
40/75 competitions natively graded, stated as 53.33\% coverage with a
null suite score, because coverage and score are different quantities;
E1's replacement SWE-bench Multilingual is 1/190 on one contract and
4/110 on another, never summed to a 300-task figure; E6 WebArena is
90/812 with 108 tasks ungraded, so its denominator counts items that
were never graded; E13 BALROG is a mean of six per-environment means
rather than a rate over episodes; and E14's Omni-MATH records 4,428
dispositions of which 4,426 were accepted, with the headline diagnostic
computed over all 4,428 (the stored native score's 51.31\% uses the
4,426 judged rows; the recovered terminal note's contrary denominator
label is corrected in Table E.1). A 41.35\% on one denominator and a
0.274\% on another do not license a comparison, an average or a ranking,
and none is offered; the rule holds within lanes too, and the ledger
states it flatly: ``N/A means unverified; no pooled results''.

\textbf{No cross-suite aggregate.} No mean, median or composite index
over suites appears anywhere in this report or in the campaign materials
it draws on: an aggregate over native denominators would be a number
with no referent, mixing suites that measure different capabilities at
different difficulties under different scoring rules. The cost is that
the campaign has no single headline; the three strictly complete scopes,
the terminal-complete Omni-MATH scope, the six replacement scopes
completed on 27 September and the two original-contract scores are the
most the record supports.

\textbf{The E10 benign-utility scope says nothing about adversarial
robustness.} The E10 lane's completed evidence is a benign-utility
evaluation of 97/97 completed tasks with utility 88/97 = 0.907 under
AgentDojo's benign scope; its original contract, private AgentHarm
tasks, is externally blocked and carries a closure record with a named
reopen condition. The completeness result therefore carries no
adversarial-robustness measurement, and this report makes no such claim
for the trained actor. The lane has a dual-use dimension, though the
limitations analysis finds the marginal uplift of the released artefacts
small, because the base model's capability is overwhelmingly
pre-existing: the pre-RL Qwen3-8B-Instruct checkpoint already scores
82.0\% on held-out GSM8K, and full GRPO adds only +1.3 points to 83.3\%,
not significant at p = 0.26. The 82.0\% is the same Instruct checkpoint
before RL, not a separate base model.

\textbf{The retained E11 result came through the original sampler.} The
129/312 full-run receipt attributes its samples to the managed Tinker
sampler rather than a locally served copy. Private execution and
repository access identities are omitted, so this public package cannot
replay the original route. A different serving stack would be a protocol
change with unmeasured same-weights, different-sampler divergence. The
separate four-task step-40 pilot has different source and
decontamination limitations and is not substituted for the full run.

\textbf{Known evidence loss in the finish-era recovery source.} The
record is explicit: the historical source-loss incident recorded in
\texttt{withheld\ historical\ source-loss\ chronology} removed the
execution sources of the finish-era recovery and build code --- the E1
wave-10 controller, the E9 restored dependencies and the E6 validator
--- while sealed requests, validation receipts and test logs survive
with matching hashes, and the replacement lanes' native runners survive
under \texttt{zvf-program/flagship/}. The lost-source policy was applied
rather than waived: a faithful re-implementation was written into a
tracked path under \texttt{zvf-program/e1\_wave10/}, where it verifies
10/10 targeted regressions offline (\longtok{test\_\allowbreak{}process\_\allowbreak{}group.py}
4/4, \longtok{test\_\allowbreak{}cleanup\_\allowbreak{}regressions.py} 6/6). Two consequences
follow: a sealed request without its execution source is not a rerun, so
the E9 MLDevBench runtime recipe was never reconstructed (the 27
September run built a fresh environment on Google Cloud instead); and
the policy response --- execution sources henceforth live under tracked
paths --- cannot restore what was already lost. A related item: the Jev
(TypeSafe) judgment harness lost reliability from roughly 13:10Z on
2026-09-19, both blocker-classification batteries returning degenerate
distributions and a controlled yes-probe falling from about 1.0 to 0.3;
the receipts are retained as incident evidence and the batteries must be
re-run once the health check passes, so the ledger's
first-uncleared-blocker column stands on document evidence rather than
judgment output.

\textbf{Scope and confound inventory.} Several Phase-1 threats remain
live. Many early runs logged train-set rewards only, so time-series
comparisons spanning old and new runs must account for the
instrumentation change. Initial phases used LoRA at ranks 8--64
exclusively; historical baseline comparisons use default hyperparameters
from the Tinker cookbook and the TRL documentation; and no full
learning-rate × rank × batch × group-size grid was swept. Tinker, Modal,
Colab and the PES A100 node use different accelerators, memory
hierarchies and training stacks, so observed differences in reward
dynamics are not purely algorithmic, and the work was not
pre-registered. The model-scale range is 0.6B--1T, with the anchors
above 235B single-seed and descriptive (Chapter 6 §6.1); the limitations
record's ``reliable'' range of 0.6B--235B is the multi-run part of that
roster, so dense-70B and trillion-parameter MoE extrapolations are
unreliable; benchmark coverage is narrow and biased toward verifiable
rewards such as exact match, unit tests and schema validity, with no
human preference data, so findings should not be extended to
reward-model-based RLHF; and the work was carried out at one institution
in Bengaluru, so benchmark choices and evaluation design reflect that
context.

The resource accounting is itself a threat. The report estimates roughly
60 A100 GPU-hours of in-kind PES University compute and projects
environmental use of about 776 kWh and about 296 kg CO₂-eq; the dominant
uncertainty is the Tinker GPU-hour count, inferred because the platform
exposes no hardware telemetry, giving a pessimistic Tinker contribution
of about 540 kg against an optimistic one of about 130 kg. The
mitigation offered is reproducibility by artefact.

\section{Conclusions}\label{conclusions}

This report set out to build a multi-framework training harness for
GRPO-style post-training of large language models and to study that
post-training under a single methodological rule: hold the stack fixed,
change one factor, report what happens. The harness dispatches one
configuration to several back-ends, but only Tinker and TRL produced
completed runs, on different base checkpoints, so it is not a completed
multi-framework benchmark. The four research questions of §1.4 have the
following answers.

\begin{itemize}
\tightlist
\item
  \textbf{RQ1 (ZVF).} Yes, it is recomputable: the three stored Qwen3-8B
  tensors reproduce 0.130, 0.190 and 0.155 exactly. Its information
  about outcome is limited: its reported association with catastrophic
  collapse is stronger than with final held-out accuracy, and the latter
  is not distinguishable from zero (§6.3; Table 1.1). Its magnitude is
  regime-dependent, from 0.16 to 0.84 across the measured runs.
\item
  \textbf{RQ2 (algorithm labels on a fixed stack).} Unanswered at this
  budget. PPO versus GRPO, GRPO versus Dr.~GRPO and G ∈ \{2, 4, 8, 16\}
  all sit at evaluation accuracy 0.98--0.99 on 0.5B two-digit
  arithmetic, and nulls at ceiling are not evidence of equivalence: they
  cannot separate no effect from no headroom. The unsaturated
  GRPO/Dr.~GRPO comparisons (Qwen2.5-1.5B GSM8K-CoT; the eight-seed
  variant audit of §1.2) are inconclusive and establish no equivalence
  either (§6.2.3, §6.4, §6.5; Table 1.1). The same-stack rerun on the
  unsaturated Qwen2.5-1.5B GSM8K-CoT configuration answers part of it:
  GRPO at G = 8 exceeds a minimal learned-critic PPO by +0.050
  {[}+0.015, +0.085{]} (five paired seeds, paired-t p = 0.016; exact
  sign-flip p = 0.0625), a baseline-specific, suggestive gap against an
  untuned critic whose gain is not detectable at this budget, and G = 8
  minus G = 2 is −0.005 {[}−0.038, +0.028{]}, neither different nor
  equivalent at a ±0.02 margin (§6.5.5).
\item
  \textbf{RQ3 (manifest-level flip risk).} Yes, for the one pair tested:
  \texttt{stackdiff} flags the open-versus-closed ``DAPO'' pair as R5
  from manifests alone (§7.2). Whether it flags risk that matters on
  outcome was not tested beyond that pair.
\item
  \textbf{RQ4 (governance-first evaluation).} It yields receipt-bound
  per-suite scores and an explicit terminal state for every lane (Table
  C.3). It does not yield a post-training delta: the actor had 40
  training steps, and the paired Table 8.B comparison is inconclusive in
  every lane.
\end{itemize}

On the analytical side the contribution is a measurement layer and a set
of bounded negatives. The Zero-Variance Fraction is a reproducible
diagnostic of advantage-signal starvation, recomputable from stored
tensors, regime-dependent in magnitude (0.72--0.77 on the Phase-1 main
configuration, 0.16 to 0.84 across the measured runs; Table 1.1), with
an association with final held-out outcome that is not distinguishable
from zero (ρ = 0.27, 95\% CI {[}−0.37, 0.88{]}, n = 23). It is
descriptive, not predictive. Its theory has been made falsifiable, and
two of its three consequences reproduced: the U-shape in accuracy across
reward levels and the monotone group-size effect; the gradient link
remains toy-scale and directional. Group size therefore has a candidate
mechanism, but the measured three-seed sweep found its arms
statistically indistinguishable and establishes no optimum (Chapter 6
§6.4). Length bias is not established: the roster-level flag that marks
four of eleven GRPO runs --- peak before 65\% of training, terminal
reward below 90\% of peak --- has post hoc thresholds, is not specific
to length, and fires on most runs in the wider trace pool, so it is not
evidence of length bias; in the controlled comparison the same rule
fires on both GRPO and Dr.~GRPO (five of six GSM8K-CoT runs) while no
run shows length growing, completion length drifts shorter, and the two
are indistinguishable on held-out gain (Chapter 6 §6.5). No claim is
made that Dr.~GRPO is uniquely free of the pattern. The levers we could
test --- a naive curriculum, a matched-token-budget intervention, an
integer group-size sweep --- did not produce reliable held-out gains.
The earlier same-stack de-confounds (PPO versus GRPO, GRPO versus
Dr.~GRPO, and group size) were all measured on a saturated 0.5B
arithmetic task with evaluation accuracy near 0.99, so their nulls say
that no difference was detectable there, not that the methods are
equivalent; the 2 October unsaturated Qwen2.5-1.5B GSM8K-CoT rerun now
supplies PPO/GRPO and G = 2 versus G = 8 contrasts (§6.5.5). Its +0.050
GRPO G8-minus-PPO estimate is suggestive and baseline-specific (paired-t
p = 0.016; exact sign-flip p = 0.0625), while the group-size contrast
remains inconclusive. PPO was minimal and untuned, prompt exposure
differed, and the 200-token cap remained potentially binding. The
strongest original contribution remains P5, the reporting standard,
whose status the record states precisely: a targeted search of the GRPO,
model-reporting and provenance literature through July 2026 did not
identify an equivalent GRPO-specific machine-readable run datasheet with
rollout provenance, and this is ``an evidence-bounded novelty claim, not
proof of non-existence''. The earlier second bet --- predictive layer
freezing from step-1 behaviour --- is weakened rather than withdrawn: it
held at 1.5B and collapsed to chance at 3B across four seeds, at a mean
of 0.083 ± 0.048.

The E1--E14 campaign is an evaluation-governance contribution, not
evidence about GRPO: the actor received 40 steps of critic-free
policy-gradient training (§8.1.1), so its scores describe a checkpoint
rather than a training method. It delivered what a governance-first
campaign can deliver: complete, receipt-bound results on three suites
(two replacement scopes and one original contract), with Omni-MATH
terminal-complete; full-coverage original-contract scores on two (E1 and
E11); and an explicit, auditable terminal state for every remaining
lane, with a named gate and a named reopen condition instead of an
indefinite pending. Its discipline is negative, which is why it is
credible: original-contract and replacement-scope numbers are never
pooled, no cross-suite aggregate is computed, and no claim of
improvement over the base model is made for this campaign (the paired
trained-versus-base comparison, Table 8.B, is inconclusive in every
lane). The controlled Chapter 6 gains are separate, stack-specific
findings and are not inferred from these campaign scores. The strongest
complete positive result, 129/312 = 41.35\% pass@1 on VerilogEval at
full 312/312 coverage, is an end-to-end harness score (150 of its 183
failures are extraction failures, driven by a 1,024-token cap with
thinking on; on a 50-item subset the same adapter scores 33/50 with
thinking off at 4,096 tokens against 23/50 in this run) and is reported
as a measurement of one actor under one generation budget on one suite,
with operational source bindings retained outside this public
derivative.

The report's overall conclusion is about method rather than algorithms.
On this compute scale, and with these instruments, the measurements are
informative while the algorithmic levers are noise-limited; the stack is
a material part of the result, so a label without its stack is not a
reproducible experimental specification; and the most durable artefacts
produced here are the ones that make other people's results checkable
--- the ZVF/PCD diagnostic layer, the eight-item minimum reportable
stack (seven manifest fields and one evaluation item), the registry and
its CLI, and a campaign ledger with named receipts and gates, with
recovered primary receipts preserved separately from the frozen checkout
and remaining raw-evidence and provenance limits stated explicitly
(Appendix B §B.11).

\subsection{C1 adds numerical witnesses and exposes reference
limits}\label{c1-adds-numerical-witnesses-and-exposes-reference-limits}

The C1 extension makes the evidence boundary more precise. Its frozen
score is 11/64 on a deterministic first-64 subset, with a question-level
Wilson 95\% working-model interval of 9.88-28.21\%. Assistant-assisted
reconciliation corroborates three numerical recovery witnesses on that
unchanged denominator. Three further events retain convention-dependent
interpretations, four are materially ambiguous and one matches an
incorrect checkpoint reference after eight correct initial answers. The
frozen score is preserved, not retrospectively repaired. Selection on
parser zeros can enrich the reference faults and representation
mismatches that prevent a direct mathematical-recovery interpretation
(Appendix G, sources S1-S5).

These three witnesses show that eight observed wrong numerical answers
do not force future numerical answers to remain wrong for this frozen
checkpoint, task and decoder. They do not validate the reasoning in a
successful response: Santana's earliest successful number was reached by
an incorrect argument. They do not measure a gradient, training update,
held-out capability or controller benefit. Reference auditing and
unit-aware evaluation therefore belong in the measurement design before
a future intervention is interpreted. The full 64-case audit retains
thirteen category disagreements and provides no whole-benchmark
error-prevalence estimate. Historical V2 remains separate and unpooled.

The separately frozen parser-v3 candidate passes only its narrow
four-change extraction check in 2,048 prospectively selected future
outputs. Preservation of 2,025 old-valid objects and agreement on four
new scalar extractions do not establish task correctness or broad parser
precision. The candidate remains undeployed; the C1 endpoint stays on
parser v2 (Appendix G, sources S8-S9).

\section{Future Work}\label{future-work}

The work leaves three kinds of open question: evaluation lanes whose
measurement is incomplete, analytical claims that the present evidence
could not settle, and methodological tools that have been built but not
yet tested at scale.

\textbf{Closing the remaining evaluation gaps.} The six replacement
lanes that were launch-pending or quota-blocked on 21 September were run
to completion on 27 September 2026 with the reported lane outcomes
(Chapter 8 §8.4). Two measurement gaps remain in them. The 108 WebArena
tasks that need an LLM judge are ungraded because the judge was
unavailable, and the native judge model has since been retired, so a
regrade would record a judge substitution. SWE-bench Multilingual has
been attempted on all 300 tasks, but under two serving contracts, so a
single 300-task figure would need the first 110 tasks rerun under the 27
September contract (32,768-token context and the re-implemented source
collector). Six original-contract scopes are recorded
\texttt{CLOSED\_EXTERNAL} with written closure records: E3 (SDAB private
bundle), E7 (BinaryAudit private payload), E8-original (LifeSciBench
package), E10-original (AgentHarm private tasks), E12 (AppBench
deployment) and E14-original (FrontierMath hosted evaluation). A seventh
original contract, E13's OpenReward held-out games, is also blocked on
external access but has no closure record, because the lane continues
through its BALROG replacement scope (Table C.3). A new evaluation would
require the unavailable benchmark scope. A full BankerToolBench score
for E4 would be a new-runtime experiment; the recovered pending-state
ledger recorded a grader-only lower bound of approximately an
undisclosed operational estimate before launch. The later 100-trial
zero-mean rerun is reported separately by
\longtok{outputs/e1\_\allowbreak{}e14\_\allowbreak{}small\_\allowbreak{}scale\_\allowbreak{}2026-09-26/E4/result.json}
(\longtok{prior\_\allowbreak{}base\_\allowbreak{}run}) and remains a confounded diagnostic with its
raw job receipt unrecovered, not a suite capability estimate or a result
substantiated by that pending-state ledger (§8.5).

\textbf{Turning the campaign into a post-training measurement.} The
campaign reports scores for one actor, not the contribution of the
training step. The paired same-engine comparison of §8.7 (Table 8.B) is
the first base-model arm; scaling it from 5--151 items per lane to the
full native scopes, with more than one training seed and a longer
training run than the actor's 40 steps, would turn each lane from a
descriptive score into a post-training delta. The original execution
route is not reproducible from this public package. A future controlled
rerun on an explicitly documented serving stack would measure the
currently unmeasured sampler divergence.

\textbf{Settling the analytical questions.} Five Phase-1 directions
remain open, each answering a threat identified above:

\begin{itemize}
\tightlist
\item
  adequately powered evaluation, moving from the Phase-1 probes'
  held-out \emph{n} of 8--20 (the multi-seed controlled comparisons
  already use 200--500 items, Chapter 1 §1.5) to the full GSM8K test
  split or a large stratified subset at five or more seeds;
\item
  a token-budget-optimal curriculum with staleness bounds, tested
  multi-seed;
\item
  testing whether the mid-to-late band in the per-layer gradient profile
  supports a FLOP-saving freeze schedule, even though step-1
  predictability did not survive scaling;
\item
  hardening the Phase-1 method-trace classifier on grouped run-level and
  temporal splits rather than row-wise cross-validation;
\item
  consolidating the diagnostic findings and the reporting standard for
  publication.
\end{itemize}

Three more follow from the Chapter 6 results. The first is a
length-confounded regime, with the completion cap lifted and longer
horizons, in which a Dr.~GRPO length advantage would be expected to
appear if it exists. The second is the head-to-head protocol between ZVF
and the pairwise-contrast density, to test whether PCD meets its target
association. The third extends the unsaturated same-stack rerun of
§6.5.5: a PPO arm with a warmed-up, tuned critic and generalised
advantage estimation, so that the GRPO--PPO gap can be separated from
critic quality, and enough seeds or evaluation items to resolve the G =
2 versus G = 8 contrast at a ±0.02 margin.

\textbf{Testing the methodological tools.} Three extensions come from
the Semester-4 work:

\begin{itemize}
\tightlist
\item
  a pre-registered controller bakeoff against static G = 16 and naive
  boundary heuristics, the test P7 declines to claim in advance of;
\item
  a cross-framework divergence study that runs an identical task,
  reward, decoding and seed configuration across TRL, veRL, OpenRLHF and
  Tinker, turning P5's attribution argument from an audit of existing
  runs into a designed experiment;
\item
  dogfooding the registry schema on the campaign's own runs, so that
  every future benchmark result ships with a verifiable provenance
  record.
\end{itemize}

Two commitments carry forward, because the threats above are structural
rather than incidental. The first is the never-pool rule:
original-contract and replacement-scope numbers stay separate, and
per-suite native denominators are never averaged. The second is that
failures are reported with their mechanism, as the E11 sampler path, the
E4 zero, the lost execution sources and the seed-count discrepancy are
here.



\cleardoublepage
\phantomsection
\addcontentsline{toc}{chapter}{References}
\bibliographystyle{IEEEtran}
\bibliography{references}


\appendix


\part{Appendices}


\chapter{Historical Run Registry and Coverage
Reconciliation}\label{historical-run-registry-and-coverage-reconciliation}

This appendix records the historical run inventories used by the report
and the later reconciliation of those inventories. It is not an
exhaustive current account history. Provider logging identities,
independent scientific executions, checkpoint revisions, and
claim-to-run links are different units and must not be added as though
they were interchangeable. Run identifier, library and platform, base
model, algorithm, group size, dataset, seed, steps, and terminal status
are retained where a primary source records them; unknown fields remain
unknown. A claim without supporting primary evidence must be qualified
or withdrawn rather than inferred from an inventory count.

Historical inventories differ in scope and counting units.
Provider-ledger entries, opening-study records, claim links and
independent scientific trials are not interchangeable. This public
derivative omits all account-wide totals, project breakdowns and exact
provider identities. Only scientific study configurations, denominators,
outcomes and limitations are retained; it makes no complete-coverage
claim.

\section{Historical boundary and source
scope}\label{historical-boundary-and-source-scope}

The Semester-3 boundary is the public tag
\texttt{capstone-final-2026-04-25}, commit \texttt{21a99ef7} dated 23
April 2026. Attribution of inherited and individual work rests on commit
history, not provider-account ownership or activity counts. The study
tables below retain scientific configurations, seeds, steps, outcomes
and evidence tiers. Operational identifiers are withheld; locally named
scientific study units remain for interpretation.

\textbf{Historical source scope in the public derivative.}

\begin{longtable}[]{@{}
  >{\raggedright\arraybackslash\hspace{0pt}}p{(\linewidth - 4\tabcolsep) * \real{0.5000}}
  >{\raggedright\arraybackslash\hspace{0pt}}p{(\linewidth - 4\tabcolsep) * \real{0.5000}}@{}}
\caption{Historical source scope; no account-level project census is
included.}\tabularnewline
\toprule\noalign{}
\begin{minipage}[b]{\linewidth}\raggedright
Source category
\end{minipage} & \begin{minipage}[b]{\linewidth}\raggedright
Public interpretation
\end{minipage} \\
\midrule\noalign{}
\endfirsthead
\toprule\noalign{}
\begin{minipage}[b]{\linewidth}\raggedright
Source category
\end{minipage} & \begin{minipage}[b]{\linewidth}\raggedright
Public interpretation
\end{minipage} \\
\midrule\noalign{}
\endhead
\bottomrule\noalign{}
\endlastfoot
Historical provider-ledger snapshots & Source-accounting records, not
independent experiment counts; account breakdowns are omitted \\
Scientific study records & Reported configurations, seeds, steps,
outcomes and limitations appear in the following study tables \\
\end{longtable}

\section{Semester-3 shared benchmark
runs}\label{semester-3-shared-benchmark-runs}

The Semester-3 contribution was the multi-framework benchmark
infrastructure: a common harness, a fixed evaluation configuration, and
integrations for several independent RL post-training libraries. These
runs are the shared baseline measurements, and they are the ones for
which the registry is \emph{most} complete, because each was produced
under a declared configuration.

\textbf{Cross-framework comparison at the fixed benchmark
configuration.} Configuration for all four rows: model
\texttt{Qwen/Qwen3-8B} (Tinker row uses \texttt{Qwen/Qwen3-8B-Base}),
algorithm GRPO, group size G = 8, learning rate 1e-5, dataset
GSM8K{[}:500{]}, 30 steps, seed 42. Metric is \texttt{last10\_avg}, the
mean of the final ten reported training rewards. Source:
\longtok{platform\_hybrid/experiments/results/framework\_comparison.json},
generated by \longtok{aggregate\_\allowbreak{}framework\_\allowbreak{}comparison.py}.

\begingroup\scriptsize\setlength{\tabcolsep}{3pt}
\begin{longtable}[]{@{}
  >{\raggedright\arraybackslash\hspace{0pt}}p{(\linewidth - 18\tabcolsep) * \real{0.1111}}
  >{\raggedright\arraybackslash\hspace{0pt}}p{(\linewidth - 18\tabcolsep) * \real{0.1111}}
  >{\raggedright\arraybackslash\hspace{0pt}}p{(\linewidth - 18\tabcolsep) * \real{0.1111}}
  >{\raggedright\arraybackslash\hspace{0pt}}p{(\linewidth - 18\tabcolsep) * \real{0.1111}}
  >{\raggedright\arraybackslash\hspace{0pt}}p{(\linewidth - 18\tabcolsep) * \real{0.1111}}
  >{\raggedright\arraybackslash\hspace{0pt}}p{(\linewidth - 18\tabcolsep) * \real{0.1111}}
  >{\raggedright\arraybackslash\hspace{0pt}}p{(\linewidth - 18\tabcolsep) * \real{0.1111}}
  >{\raggedright\arraybackslash\hspace{0pt}}p{(\linewidth - 18\tabcolsep) * \real{0.1111}}
  >{\raggedright\arraybackslash\hspace{0pt}}p{(\linewidth - 18\tabcolsep) * \real{0.1111}}@{}}
\caption{Semester-3 cross-framework shared benchmark
runs.}\tabularnewline
\toprule\noalign{}
\begin{minipage}[b]{\linewidth}\raggedright
Run identifier
\end{minipage} & \begin{minipage}[b]{\linewidth}\raggedright
Stack / platform
\end{minipage} & \begin{minipage}[b]{\linewidth}\raggedright
Model
\end{minipage} & \begin{minipage}[b]{\linewidth}\raggedright
Algo
\end{minipage} & \begin{minipage}[b]{\linewidth}\raggedright
G
\end{minipage} & \begin{minipage}[b]{\linewidth}\raggedright
Dataset
\end{minipage} & \begin{minipage}[b]{\linewidth}\raggedright
Seed
\end{minipage} & \begin{minipage}[b]{\linewidth}\raggedright
Steps
\end{minipage} & \begin{minipage}[b]{\linewidth}\raggedright
Outcome / status
\end{minipage} \\
\midrule\noalign{}
\endfirsthead
\toprule\noalign{}
\begin{minipage}[b]{\linewidth}\raggedright
Run identifier
\end{minipage} & \begin{minipage}[b]{\linewidth}\raggedright
Stack / platform
\end{minipage} & \begin{minipage}[b]{\linewidth}\raggedright
Model
\end{minipage} & \begin{minipage}[b]{\linewidth}\raggedright
Algo
\end{minipage} & \begin{minipage}[b]{\linewidth}\raggedright
G
\end{minipage} & \begin{minipage}[b]{\linewidth}\raggedright
Dataset
\end{minipage} & \begin{minipage}[b]{\linewidth}\raggedright
Seed
\end{minipage} & \begin{minipage}[b]{\linewidth}\raggedright
Steps
\end{minipage} & \begin{minipage}[b]{\linewidth}\raggedright
Outcome / status
\end{minipage} \\
\midrule\noalign{}
\endhead
\bottomrule\noalign{}
\endlastfoot
\longtok{campaign\_\allowbreak{}v2\_\allowbreak{}w1\_\allowbreak{}qwen3-8b-base} & Tinker & Qwen/\allowbreak{}Qwen3-8B-Base &
GRPO & 8 & gsm8k-500 & 42 & 30 & last10 = 0.85625 --- \textbf{real
run} \\
\longtok{modal\_\allowbreak{}trl\_\allowbreak{}trl\_\allowbreak{}qwen3\_\allowbreak{}8b} & TRL on Modal H100 & Qwen/\allowbreak{}Qwen3-8B
& GRPO & 8 & gsm8k-500 & 42 & 30 & last10 = 0.05, 735.7 s ---
\textbf{real run} \\
\longtok{verl-dryrun} & verl & Qwen/\allowbreak{}Qwen3-8B & GRPO & 8 & gsm8k-500 & 42
& 30 & last10 = 0.5528 --- \textbf{seeded dryrun, not a training run} \\
\longtok{openrlhf-dryrun} & OpenRLHF & Qwen/\allowbreak{}Qwen3-8B & GRPO & 8 &
gsm8k-500 & 42 & 30 & last10 = 0.4785 --- \textbf{seeded dryrun, not a
training run} \\
\end{longtable}
\endgroup

The two \texttt{dryrun} rows are a documented threat to validity, not a
result: the generator falls back to a deterministic seeded trace from
\texttt{verl.VERLTrainer} / \texttt{OpenRLHFTrainer} when the framework
is absent from the sandbox, and the JSON carries the note ``Seeded
deterministic fallback \ldots{} (verl not installed in sandbox)''. Any
comparison reading them as measurements of verl or OpenRLHF is invalid.
The Tinker row also initialises from the \emph{base} checkpoint while
the open-framework rows use the instruction-tuned one, so even the two
real rows are not a clean framework contrast.

\textbf{Modal multi-framework campaign, 4 frameworks × 5 seeds = 20
runs.} Source: \longtok{modal\_\allowbreak{}results\_\allowbreak{}all.json} (repository root). All
rows: model \texttt{Qwen/Qwen2.5-0.5B}, dataset arithmetic/math, 125
training steps, NVIDIA L4, one run per seed. Seeds are 42, 123, 456,
789, 1024. Group size is not recorded in the artefact and is therefore
shown as ``---''.

\begingroup\scriptsize\setlength{\tabcolsep}{3pt}
\begin{longtable}[]{@{}
  >{\raggedright\arraybackslash\hspace{0pt}}p{(\linewidth - 18\tabcolsep) * \real{0.1111}}
  >{\raggedright\arraybackslash\hspace{0pt}}p{(\linewidth - 18\tabcolsep) * \real{0.1111}}
  >{\raggedright\arraybackslash\hspace{0pt}}p{(\linewidth - 18\tabcolsep) * \real{0.1111}}
  >{\raggedright\arraybackslash\hspace{0pt}}p{(\linewidth - 18\tabcolsep) * \real{0.1111}}
  >{\raggedright\arraybackslash\hspace{0pt}}p{(\linewidth - 18\tabcolsep) * \real{0.1111}}
  >{\raggedright\arraybackslash\hspace{0pt}}p{(\linewidth - 18\tabcolsep) * \real{0.1111}}
  >{\raggedright\arraybackslash\hspace{0pt}}p{(\linewidth - 18\tabcolsep) * \real{0.1111}}
  >{\raggedright\arraybackslash\hspace{0pt}}p{(\linewidth - 18\tabcolsep) * \real{0.1111}}
  >{\raggedright\arraybackslash\hspace{0pt}}p{(\linewidth - 18\tabcolsep) * \real{0.1111}}@{}}
\caption{Semester-3 Modal campaign runs.}\tabularnewline
\toprule\noalign{}
\begin{minipage}[b]{\linewidth}\raggedright
Run identifier
\end{minipage} & \begin{minipage}[b]{\linewidth}\raggedright
Stack / platform
\end{minipage} & \begin{minipage}[b]{\linewidth}\raggedright
Model
\end{minipage} & \begin{minipage}[b]{\linewidth}\raggedright
Algo
\end{minipage} & \begin{minipage}[b]{\linewidth}\raggedright
G
\end{minipage} & \begin{minipage}[b]{\linewidth}\raggedright
Dataset
\end{minipage} & \begin{minipage}[b]{\linewidth}\raggedright
Seed
\end{minipage} & \begin{minipage}[b]{\linewidth}\raggedright
Steps
\end{minipage} & \begin{minipage}[b]{\linewidth}\raggedright
Outcome / status
\end{minipage} \\
\midrule\noalign{}
\endfirsthead
\toprule\noalign{}
\begin{minipage}[b]{\linewidth}\raggedright
Run identifier
\end{minipage} & \begin{minipage}[b]{\linewidth}\raggedright
Stack / platform
\end{minipage} & \begin{minipage}[b]{\linewidth}\raggedright
Model
\end{minipage} & \begin{minipage}[b]{\linewidth}\raggedright
Algo
\end{minipage} & \begin{minipage}[b]{\linewidth}\raggedright
G
\end{minipage} & \begin{minipage}[b]{\linewidth}\raggedright
Dataset
\end{minipage} & \begin{minipage}[b]{\linewidth}\raggedright
Seed
\end{minipage} & \begin{minipage}[b]{\linewidth}\raggedright
Steps
\end{minipage} & \begin{minipage}[b]{\linewidth}\raggedright
Outcome / status
\end{minipage} \\
\midrule\noalign{}
\endhead
\bottomrule\noalign{}
\endlastfoot
\longtok{trl\_\allowbreak{}grpo\_\allowbreak{}math} & TRL / Modal L4 & Qwen/\allowbreak{}Qwen2.5-0.5B & GRPO &
--- & math & 42, 123, 456, 789, 1024 & 125 & final acc 0.735 / 0.810 /
0.620 / 0.740 / 0.765; 139.8--183.7 s \\
\longtok{sb3\_\allowbreak{}ppo\_\allowbreak{}math} & Stable-Baselines3 / Modal L4 &
Qwen/\allowbreak{}Qwen2.5-0.5B & PPO & --- & math & 42, 123, 456, 789, 1024 & 125 &
final acc 0.003 / 0.014 / 0.011 / 0.012 / 0.010; \longtok{steps\_\allowbreak{}to\_\allowbreak{}95}
null for all five \\
\longtok{cleanrl\_\allowbreak{}ppo\_\allowbreak{}math} & CleanRL / Modal L4 & Qwen/\allowbreak{}Qwen2.5-0.5B &
PPO & --- & math & 42, 123, 456, 789, 1024 & 125 & final acc 0.010 /
0.008 / 0.014 / 0.009 / 0.004; 25.4--32.4 s \\
\longtok{tianshou\_\allowbreak{}ppo\_\allowbreak{}math} & Tianshou / Modal L4 & Qwen/\allowbreak{}Qwen2.5-0.5B &
PPO & --- & math & 42, 123, 456, 789, 1024 & 125 & final acc 0.011 /
0.002 / 0.005 / 0.002 / 0.009; 1.67--3.30 s \\
\end{longtable}
\endgroup

\textbf{Group-size sweep, shared-harness runs.} Source:
\longtok{platform\_hybrid/experiments/results/groupsize\_\allowbreak{}zvf\_\allowbreak{}sweep.tsv}
(per-run data in \longtok{groupsize\_\allowbreak{}zvf\_\allowbreak{}sweep.json}). Model
\texttt{Qwen/Qwen2.5-0.5B} on a synthetic arithmetic-correctness task, 3
seeds per group size (42, 123, 456; 12 runs), 40 steps, 16 prompts per
step, one NVIDIA A10 on Modal, with held-out accuracy on 200 fresh
problems per run. An earlier version of this table labelled these runs
Qwen3-8B / Tinker / GSM8K. The raw per-run records contradict that
label, and the held-out figures below match the raw records exactly.
Held-out accuracy is the mean ± SE over seeds; \texttt{mean\ ZVF} is the
mean zero-variance fraction.

\begingroup\scriptsize\setlength{\tabcolsep}{3pt}
\begin{longtable}[]{@{}
  >{\raggedright\arraybackslash\hspace{0pt}}p{(\linewidth - 18\tabcolsep) * \real{0.1111}}
  >{\raggedright\arraybackslash\hspace{0pt}}p{(\linewidth - 18\tabcolsep) * \real{0.1111}}
  >{\raggedright\arraybackslash\hspace{0pt}}p{(\linewidth - 18\tabcolsep) * \real{0.1111}}
  >{\raggedright\arraybackslash\hspace{0pt}}p{(\linewidth - 18\tabcolsep) * \real{0.1111}}
  >{\raggedright\arraybackslash\hspace{0pt}}p{(\linewidth - 18\tabcolsep) * \real{0.1111}}
  >{\raggedright\arraybackslash\hspace{0pt}}p{(\linewidth - 18\tabcolsep) * \real{0.1111}}
  >{\raggedright\arraybackslash\hspace{0pt}}p{(\linewidth - 18\tabcolsep) * \real{0.1111}}
  >{\raggedright\arraybackslash\hspace{0pt}}p{(\linewidth - 18\tabcolsep) * \real{0.1111}}
  >{\raggedright\arraybackslash\hspace{0pt}}p{(\linewidth - 18\tabcolsep) * \real{0.1111}}@{}}
\caption{Semester-3 group-size sweep runs.}\tabularnewline
\toprule\noalign{}
\begin{minipage}[b]{\linewidth}\raggedright
Run identifier
\end{minipage} & \begin{minipage}[b]{\linewidth}\raggedright
Stack / platform
\end{minipage} & \begin{minipage}[b]{\linewidth}\raggedright
Model
\end{minipage} & \begin{minipage}[b]{\linewidth}\raggedright
Algo
\end{minipage} & \begin{minipage}[b]{\linewidth}\raggedright
G
\end{minipage} & \begin{minipage}[b]{\linewidth}\raggedright
Dataset
\end{minipage} & \begin{minipage}[b]{\linewidth}\raggedright
Seed
\end{minipage} & \begin{minipage}[b]{\linewidth}\raggedright
Steps
\end{minipage} & \begin{minipage}[b]{\linewidth}\raggedright
Outcome / status
\end{minipage} \\
\midrule\noalign{}
\endfirsthead
\toprule\noalign{}
\begin{minipage}[b]{\linewidth}\raggedright
Run identifier
\end{minipage} & \begin{minipage}[b]{\linewidth}\raggedright
Stack / platform
\end{minipage} & \begin{minipage}[b]{\linewidth}\raggedright
Model
\end{minipage} & \begin{minipage}[b]{\linewidth}\raggedright
Algo
\end{minipage} & \begin{minipage}[b]{\linewidth}\raggedright
G
\end{minipage} & \begin{minipage}[b]{\linewidth}\raggedright
Dataset
\end{minipage} & \begin{minipage}[b]{\linewidth}\raggedright
Seed
\end{minipage} & \begin{minipage}[b]{\linewidth}\raggedright
Steps
\end{minipage} & \begin{minipage}[b]{\linewidth}\raggedright
Outcome / status
\end{minipage} \\
\midrule\noalign{}
\endhead
\bottomrule\noalign{}
\endlastfoot
\longtok{groupsize\_g2} & Modal A10 (shared harness) & Qwen/\allowbreak{}Qwen2.5-0.5B
& GRPO & 2 & arithmetic, 200 held-out & 3 seeds & 40 & held-out acc
0.9817 ± 0.0044; mean ZVF 0.8380 \\
\longtok{groupsize\_g4} & Modal A10 (shared harness) & Qwen/\allowbreak{}Qwen2.5-0.5B
& GRPO & 4 & arithmetic, 200 held-out & 3 seeds & 40 & held-out acc
0.9883 ± 0.0017; mean ZVF 0.7635 \\
\longtok{groupsize\_g8} & Modal A10 (shared harness) & Qwen/\allowbreak{}Qwen2.5-0.5B
& GRPO & 8 & arithmetic, 200 held-out & 3 seeds & 40 & held-out acc
0.9900 ± 0.0029; mean ZVF 0.6906 \\
\longtok{groupsize\_g16} & Modal A10 (shared harness) & Qwen/\allowbreak{}Qwen2.5-0.5B
& GRPO & 16 & arithmetic, 200 held-out & 3 seeds & 40 & held-out acc
0.9783 ± 0.0060; mean ZVF 0.6312 \\
\end{longtable}
\endgroup

Additional exploratory records are retained in the historical source
corpus. Incomplete configuration fields prevent treating all inventory
entries as comparable scientific trials. This public derivative omits
account-level breakdowns.

\section{Semester-4 P1--P8 study runs}\label{semester-4-p1p8-study-runs}

The P1--P8 series is the individual Semester-4 contribution, and its
runs are indexed by \emph{claim} rather than by project. The audit
ledger at \longtok{platform\_hybrid/experiments/results/claim\_\allowbreak{}to\_\allowbreak{}run/}
maps each published claim to its supporting run identifiers, provider
references retained in the non-public source records, whether the model
named in the paper matches the model actually trained, the seed and step
counts, the held-out metric, and an evidence tier. The retained claim
map has eighteen scientific claim rows. Private link and
provider-inventory counts are omitted; the table retains the scientific
evidence tiers and study outcomes.

\textbf{P1--P8 claim-to-run ledger.} Tier vocabulary: A = ≥5 seeds and
≥100 steps; B = 3--4 seeds and 50--99 steps; C = completed/descriptive;
R = deterministic resource or schema claim; X = unlinked, contradicted,
or provenance-conflicted.

\begingroup\scriptsize\setlength{\tabcolsep}{3pt}
\begin{longtable}[]{@{}
  >{\raggedright\arraybackslash\hspace{0pt}}p{(\linewidth - 20\tabcolsep) * \real{0.0600}}
  >{\raggedright\arraybackslash\hspace{0pt}}p{(\linewidth - 20\tabcolsep) * \real{0.1100}}
  >{\raggedright\arraybackslash\hspace{0pt}}p{(\linewidth - 20\tabcolsep) * \real{0.1000}}
  >{\raggedright\arraybackslash\hspace{0pt}}p{(\linewidth - 20\tabcolsep) * \real{0.1300}}
  >{\raggedright\arraybackslash\hspace{0pt}}p{(\linewidth - 20\tabcolsep) * \real{0.1100}}
  >{\raggedright\arraybackslash\hspace{0pt}}p{(\linewidth - 20\tabcolsep) * \real{0.0500}}
  >{\raggedright\arraybackslash\hspace{0pt}}p{(\linewidth - 20\tabcolsep) * \real{0.0900}}
  >{\raggedright\arraybackslash\hspace{0pt}}p{(\linewidth - 20\tabcolsep) * \real{0.1000}}
  >{\raggedright\arraybackslash\hspace{0pt}}p{(\linewidth - 20\tabcolsep) * \real{0.2000}}
  >{\raggedright\arraybackslash\hspace{0pt}}p{(\linewidth - 20\tabcolsep) * \real{0.0500}}@{}}
\caption{Semester-4 P1-P8 study runs: each claim mapped to its linked
runs, stack, model and algorithm.}\tabularnewline
\toprule\noalign{}
\begin{minipage}[b]{\linewidth}\raggedright
Claim
\end{minipage} & \begin{minipage}[b]{\linewidth}\raggedright
Linked runs (representative)
\end{minipage} & \begin{minipage}[b]{\linewidth}\raggedright
Stack
\end{minipage} & \begin{minipage}[b]{\linewidth}\raggedright
Model
\end{minipage} & \begin{minipage}[b]{\linewidth}\raggedright
Algo
\end{minipage} & \begin{minipage}[b]{\linewidth}\raggedright
G
\end{minipage} & \begin{minipage}[b]{\linewidth}\raggedright
Dataset
\end{minipage} & \begin{minipage}[b]{\linewidth}\raggedright
Seeds / steps
\end{minipage} & \begin{minipage}[b]{\linewidth}\raggedright
Outcome
\end{minipage} & \begin{minipage}[b]{\linewidth}\raggedright
Tier
\end{minipage} \\
\midrule\noalign{}
\endfirsthead
\toprule\noalign{}
\begin{minipage}[b]{\linewidth}\raggedright
Claim
\end{minipage} & \begin{minipage}[b]{\linewidth}\raggedright
Linked runs (representative)
\end{minipage} & \begin{minipage}[b]{\linewidth}\raggedright
Stack
\end{minipage} & \begin{minipage}[b]{\linewidth}\raggedright
Model
\end{minipage} & \begin{minipage}[b]{\linewidth}\raggedright
Algo
\end{minipage} & \begin{minipage}[b]{\linewidth}\raggedright
G
\end{minipage} & \begin{minipage}[b]{\linewidth}\raggedright
Dataset
\end{minipage} & \begin{minipage}[b]{\linewidth}\raggedright
Seeds / steps
\end{minipage} & \begin{minipage}[b]{\linewidth}\raggedright
Outcome
\end{minipage} & \begin{minipage}[b]{\linewidth}\raggedright
Tier
\end{minipage} \\
\midrule\noalign{}
\endhead
\bottomrule\noalign{}
\endlastfoot
P1-C1 & 20 --- {[}operational identifier withheld{]} & Tinker (managed)
+ W\&B & 5 GSM8K anchors: Qwen3.5-4B, Qwen3-8B, Llama-3.1-8B-Instruct,
Nemotron-120B, DeepSeek-V3.1 & --- & --- & GSM8K & 1 per anchor (seed
42) / 20--30 & No held-out metric on the exact-linked runs; the claim
rests on training-reward traces & C \\
P1-C2 & 1 --- {[}operational identifier withheld{]} (HF-arbitrated) &
Tinker (managed) + HF adapter &
\longtok{nvidia/NVIDIA-Nemotron-3-Super-120B-A12B-BF16} (MoE, A12B
active) & --- & --- & \longtok{frontier\_gsm8k} & 1 (seed 42) / 20 & Peak
0.875, last-10 0.1625, zero-reward step fraction 0.55 (11/20); no
held-out metric & C \\
P2-C1 & 80 local rows pooled to 23 cells & Local analysis pool (no W\&B)
& Heterogeneous; Qwen3-8B at G = 8 among the listed rows & --- & --- &
GSM8K & up to 3 per source row / 30--600 & Mean ZVF versus collapse:
Spearman 0.56, point-biserial 0.62 (n = 23) & C \\
P2-C2 & 80 (as P2-C1) & as P2-C1 & as P2-C1 & --- & --- & GSM8K & as
P2-C1 & Spearman ρ = 0.27, bootstrap 95\% CI {[}−0.37, 0.88{]} (n = 23)
& C \\
P3-C1 & 12 --- {[}operational identifier withheld{]} & W\&B-managed open
trainer; no Tinker identity & Qwen/\allowbreak{}Qwen3.5-4B & --- & 2, 4, 8, 16 & ---
(saturated arithmetic regime; dataset not recorded) & 3 per G (seeds 0,
1, 2) / 8 & Mean held-out Δ: G=2 +0.0000, G=4 +0.0000, G=8 −0.0625, G=16
+0.0000; evaluation size not logged & C \\
P3-C2 & 5 --- {[}operational identifier withheld{]} & W\&B, sampling
only & --- (G = 32 cells) & --- & 32 & --- & 5 cells / 0 optimizer steps
& No matched G = 4/G = 32 training replication; the comparison cannot be
executed & X \\
P4-C1 & 6 --- \longtok{local:drgrpo\_\allowbreak{}gsm8k\_\allowbreak{}cot:grpo:seed42} & Local &
Qwen/\allowbreak{}Qwen2.5-1.5B-Instruct as mapped (abstract says Qwen3-8B) ---
conflict & GRPO / Dr.GRPO & --- & GSM8K CoT & 3 per arm (42, 123, 456) /
30 & GRPO 0.2017→0.2633 (Δ +0.0617); Dr.GRPO 0.2050→0.2550 (Δ +0.0500);
n = 200/seed & X \\
P5-C1 & 2 --- {[}operational identifier withheld{]}, {[}identifier
withheld{]} & Tinker-managed versus TRL backend & Qwen/\allowbreak{}Qwen3-8B-Base
versus Qwen/\allowbreak{}Qwen3-8B --- confounded & GRPO / TRL-GRPO & 8 & GSM8K & 1
per backend (seed 42) / 30 & Final/last-10 training reward 0.050→0.856;
no held-out metric & X \\
P5-C2 & 2 --- {[}identifier withheld{]}, {[}identifier withheld{]} &
Open trainer versus managed stack & Qwen2.5-0.5B-Instruct versus
Qwen3.5-4B (declared non-consistent exhibit) & DAPO label versus
asymmetric-clip surrogate & --- & --- & 2 open versus 3 managed / not
logged (open), 15 (managed) & Mean ZVF 0.00 versus 0.58; no comparable
held-out metric & C \\
P5-C3 & 1 --- {[}identifier withheld{]} & Managed summary &
Qwen/\allowbreak{}Qwen3.5-4B & Four algorithm labels & --- & --- & 3 per arm (42,
123, 456) / 15 & Last-10 reward band 0.710--0.744 (width 0.034); no
held-out metric & C \\
P6-C1 & 48 registry entries & Registry (schema validation) & N/A & N/A &
N/A & N/A & N/A & 44/48 entries pass schema validation in the later
claim-to-run snapshot (\longtok{claim\_\allowbreak{}to\_\allowbreak{}run\_\allowbreak{}table.tsv}, which adds
amendment entries); the 31/31 of §7.2 is the 31 records shipped with P6
& R \\
P6-C2 (source-identity audit) & Private source reconciliation omitted &
Cross-source metadata & N/A & N/A & N/A & N/A & Not a training or
evaluation study & Unresolved attribution limits retained; no public
inventory claim & R \\
P7-C1 & 368 --- {[}operational identifier withheld{]} & W\&B only (0
exact Tinker IDs) & Qwen3-4B-Instruct-2507 (113), Llama-3.2-3B (112),
Llama-3.2-1B (107), 3 further models (36) & --- & 4, 8, 16 (also 6) &
--- (prompt and tokenizer families \texttt{llama3},
\longtok{qwen3\_instruct} recorded, not datasets) & 10 logged labels
(0--9) / 2--12 & ZVF U-shape reproduced; no held-out metric in the W\&B
summaries & C \\
P7-C2 & 1 --- {[}identifier withheld{]} & W\&B (Colab) &
Qwen/\allowbreak{}Qwen2.5-0.5B-Instruct & --- & --- & --- & 3 (0, 1, 2) in one
aggregate / 30 & Gradient magnitude versus p(1−p): Pearson r = +0.71 &
C \\
P7-C3 & 1 --- {[}identifier withheld{]} & W\&B (Colab) &
Qwen/\allowbreak{}Qwen2.5-0.5B-Instruct & Adaptive-G versus fixed recipes & --- & ---
& 1 per arm / not logged & Adaptive-G held-out Δ +0.575; mean ZVF 0.23;
186 rollouts & C \\
P8-C1 & 1 --- \longtok{local:qp8\_fraud:xgboost:seed42} & Local (XGBoost)
& Gradient-boosted trees (200 estimators, depth 6) & --- & N/A &
Synthetic card-fraud: 50,000 rows; 10,000-row held-out split &
random\_state 42 / 200 estimators & AUC 0.7955 on the 10,000-row
held-out split & C \\
P8-C2 & 1 --- {[}operational identifier withheld{]} & Tinker (managed
SFT) & Qwen/\allowbreak{}Qwen3.5-4B & SFT row serialisation (no KL, no reference
policy) & N/A & Row-to-text serialisations; 500-row positive-enriched
held-out subset & 1 (seed 0) / 63 minibatches & AUC 0.482675, accuracy
0.792 (n = 500) --- near chance & C \\
P8-C3 & 0 --- UNLINKED & --- & Proposed sensor / scribe / triage
architecture & --- & --- & --- & 0 / 0 & No execution; operational
benefit over XGBoost alone not measured & X \\
\end{longtable}
\endgroup

Three facts here are load-bearing. First, \textbf{no empirical headline
row in the P1--P8 ledger reaches Tier A or B} under the
statistical-rigor thresholds; the rectified P1--P8 tally is C = 12, R =
2, X = 4. P11 adds an eight-seed matched block but does not meet the
stated Tier A step threshold, as clarified below. Second, the four Tier
X rows are X for specific, checkable reasons: \textbf{P3-C2} rests on G
= 32 \emph{sampling} cells with zero optimizer steps, so no matched G =
4/G = 32 training replication exists; \textbf{P4-C1}'s abstract names
Qwen3-8B while every mapped run is \texttt{Qwen/Qwen2.5-1.5B-Instruct};
\textbf{P5-C1} compares \texttt{Qwen/Qwen3-8B-Base} against
\texttt{Qwen/Qwen3-8B}, so the base checkpoint changes together with the
backend and the pair must not be cited as a backend-only causal effect;
and \textbf{P8-C3} is a proposed architecture with no execution, hence
zero linked runs. Third, some cross-provider model-identity conflicts
remain unresolved. Their private identity reconciliation is outside this
public package; the affected scientific model-attribution claims remain
contested.

Two column caveats travel with the table. The dataset cell for P7-C1 is
\texttt{—} because the ledger's only dataset-adjacent field for those
368 runs is a prompt/tokenizer family (\texttt{llama3},
\texttt{qwen3\_instruct}, \longtok{gpt\_\allowbreak{}oss\_\allowbreak{}low\_\allowbreak{}reasoning},
\longtok{nemotron3\_\allowbreak{}disable\_\allowbreak{}thinking}), not a dataset name. And the P3
sweep here is \textbf{not} the Semester-3 sweep of Table A.4: it is a
different project (\texttt{study\ logging\ group}), a different model
(Qwen3.5-4B), and it reports deltas of 0.0000 rather than held-out
accuracies near 0.99. The two share a group-size grid and nothing else.
Chapters 4 and 6 cite the Table A.4 sweep (Qwen2.5-0.5B) whenever they
give the 0.838 → 0.631 ZVF decline or the 0.98--0.99 held-out figures.
They do not cite P3-C1.

The extension ledger added in August 2026 carries five rows besides P11:
P10-C1 is Tier R (an algebraic identity, not an empirical result),
P10-C2 and P12-C2 are Tier X (placeholder figures and unlinked
prospective claims, zero runs), and P12-C1 is Tier R as a lemma and Tier
X as an empirical headline.

\textbf{P11} is an eight-seed, matched-stack block of 40 arm--seed units
at
\longtok{local:zvf-program/audit/results/full/\{grpo,dapo,gspo,drgrpo,aero\}-seed-\{11,23,37,53,71,89,107,131\}.json},
five arms × eight seeds on one stack (Qwen3-8B, GSM8K, 30 steps, Colab
A100, held-out n = 500). Its headline contrast, DAPO versus GRPO, is Δ =
+0.001, percentile-bootstrap CI {[}−0.0045, +0.00675{]} (paired-t 95\%
CI {[}−0.0063, +0.0083{]}, n = 8 seeds; the bootstrap undercovers at
this n), MDE₈₀ = 0.01012, all raw p \textgreater{} 0.21, with
Benjamini--Hochberg rejecting nothing. The mechanism is visible in the
diagnostic: DAPO's mean ZVF is 0.000 against GRPO's 0.693, bought with
3.61× the rollouts (1,734 versus 480) and 1.44× the wall clock. This
eight-paired-seed contrast supports a null. Every retained P11 final
manifest specifies \texttt{max\_steps=30}; the 40 units are 30-step
runs. Accordingly this block does \textbf{not} satisfy the stated Tier A
requirement of at least 100 steps per run and remains Tier C under the
vocabulary above. Pooling steps across arms or seeds does not satisfy
that per-run requirement. This corrects the evidence tier, not the
recorded rewards, seed counts, interval, or null result.

\subsection{Opening and null-result
runs}\label{opening-and-null-result-runs}

Chapter 6 reports four null families --- a mixed-variance curriculum, a
matched token budget, hard-prompt filtering, and step-1 layer selection
(§6.4.3) --- and each rests on source-specific scientific result records
rather than an account-wide inventory. The study The rows are counted
from the per-run \texttt{name} (or mode) fields of each artefact's JSON
--- \texttt{results.json} for the curriculum, token-budget, hard-prompt
and loss-formulation blocks, \texttt{result.json} and its siblings for
P1 --- not from an inventory API, which is why this cohort is reported
as its own total. Three neighbouring blocks add nothing to it: the P3
group-size artefacts are already indexed ---
\texttt{groupsize\_zvf/results.json} reproduces the per-group deltas of
Table A.5's P3-C1 row exactly, and the second sweep file is discussed
below --- and the P2, P7 and P8 ``openings''
(\longtok{p2\_openings/collapse\_\allowbreak{}baseline\_\allowbreak{}analysis.json},
\texttt{p7\_controller/sim.json}, \texttt{p8\_openings/*.json}) are
zero-compute re-analyses that launch no runs.

\textbf{Semester-4 opening and null-result runs (46 runs).} Model is
\texttt{Qwen/Qwen3.5-4B} and dataset GSM8K unless a row says otherwise.
All rows are Tinker-managed except the P1 layer-freeze block, which ran
on Colab for per-layer gradient access. ``Curriculum'' means training
only on non-collapsed (mixed-variance) groups. Group sizes marked
\emph{default} are the launch defaults of the scripts in
\longtok{platform\_hybrid/experiments/openings/}; only the campaign
artefact records a group size per run. Δ is the change in held-out
accuracy.

\begingroup\scriptsize\setlength{\tabcolsep}{3pt}
\begin{longtable}[]{@{}
  >{\raggedright\arraybackslash\hspace{0pt}}p{(\linewidth - 18\tabcolsep) * \real{0.1400}}
  >{\raggedright\arraybackslash\hspace{0pt}}p{(\linewidth - 18\tabcolsep) * \real{0.1300}}
  >{\raggedright\arraybackslash\hspace{0pt}}p{(\linewidth - 18\tabcolsep) * \real{0.1500}}
  >{\raggedright\arraybackslash\hspace{0pt}}p{(\linewidth - 18\tabcolsep) * \real{0.1100}}
  >{\raggedright\arraybackslash\hspace{0pt}}p{(\linewidth - 18\tabcolsep) * \real{0.0600}}
  >{\raggedright\arraybackslash\hspace{0pt}}p{(\linewidth - 18\tabcolsep) * \real{0.1000}}
  >{\raggedright\arraybackslash\hspace{0pt}}p{(\linewidth - 18\tabcolsep) * \real{0.0500}}
  >{\raggedright\arraybackslash\hspace{0pt}}p{(\linewidth - 18\tabcolsep) * \real{0.0700}}
  >{\raggedright\arraybackslash\hspace{0pt}}p{(\linewidth - 18\tabcolsep) * \real{0.1900}}@{}}
\caption{Semester-4 opening and null-result runs: the curriculum,
token-budget, hard-prompt, loss-formulation and layer-freeze
blocks.}\tabularnewline
\toprule\noalign{}
\begin{minipage}[b]{\linewidth}\raggedright
Run identifier
\end{minipage} & \begin{minipage}[b]{\linewidth}\raggedright
Stack / platform
\end{minipage} & \begin{minipage}[b]{\linewidth}\raggedright
Model
\end{minipage} & \begin{minipage}[b]{\linewidth}\raggedright
Algo
\end{minipage} & \begin{minipage}[b]{\linewidth}\raggedright
G
\end{minipage} & \begin{minipage}[b]{\linewidth}\raggedright
Dataset
\end{minipage} & \begin{minipage}[b]{\linewidth}\raggedright
Seed
\end{minipage} & \begin{minipage}[b]{\linewidth}\raggedright
Steps
\end{minipage} & \begin{minipage}[b]{\linewidth}\raggedright
Outcome / status
\end{minipage} \\
\midrule\noalign{}
\endfirsthead
\toprule\noalign{}
\begin{minipage}[b]{\linewidth}\raggedright
Run identifier
\end{minipage} & \begin{minipage}[b]{\linewidth}\raggedright
Stack / platform
\end{minipage} & \begin{minipage}[b]{\linewidth}\raggedright
Model
\end{minipage} & \begin{minipage}[b]{\linewidth}\raggedright
Algo
\end{minipage} & \begin{minipage}[b]{\linewidth}\raggedright
G
\end{minipage} & \begin{minipage}[b]{\linewidth}\raggedright
Dataset
\end{minipage} & \begin{minipage}[b]{\linewidth}\raggedright
Seed
\end{minipage} & \begin{minipage}[b]{\linewidth}\raggedright
Steps
\end{minipage} & \begin{minipage}[b]{\linewidth}\raggedright
Outcome / status
\end{minipage} \\
\midrule\noalign{}
\endhead
\bottomrule\noalign{}
\endlastfoot
\longtok{campaign/baseline-G4-s0…s5} & Tinker → W\&B
\texttt{study\ logging\ group} (\texttt{campaign}) & Qwen/\allowbreak{}Qwen3.5-4B &
GRPO & 4 & GSM8K & 0--5 & 8 & 6 runs; mean Δ +0.0111 (0.0000--0.0333) \\
\longtok{campaign/curriculum-G4-s0…s5} & Tinker → W\&B
\texttt{study\ logging\ group} (\texttt{campaign}) & Qwen/\allowbreak{}Qwen3.5-4B &
GRPO + curriculum filter & 4 & GSM8K & 0--5 & 8 & 6 runs; mean Δ +0.0167
(−0.0667 to +0.0667); oversample 3.3--7.2×; zero-loss fraction 0.0 \\
\longtok{campaign/baseline-G2-s0} & Tinker → W\&B
\texttt{study\ logging\ group} (\texttt{campaign}) & Qwen/\allowbreak{}Qwen3.5-4B &
GRPO & 2 & GSM8K & 0 & 8 & 1 run; Δ 0.0000; zero-loss fraction 0.7 \\
\longtok{campaign/baseline-G8-s0} & Tinker → W\&B
\texttt{study\ logging\ group} (\texttt{campaign}) & Qwen/\allowbreak{}Qwen3.5-4B &
GRPO & 8 & GSM8K & 0 & 8 & 1 run; Δ −0.1333 \\
\longtok{campaign/baseline-G16-s0} & Tinker → W\&B
\texttt{study\ logging\ group} (\texttt{campaign}) & Qwen/\allowbreak{}Qwen3.5-4B &
GRPO & 16 & GSM8K & 0 & 8 & 1 run; Δ +0.0333 \\
\longtok{curriculum\_opening/baseline} & Tinker → W\&B
\texttt{study\ logging\ group} (ungrouped) & Qwen/\allowbreak{}Qwen3.5-4B & GRPO & 4
\emph{default} & GSM8K & 0 & 8 & 1 run; Δ +0.0500; held-out 20;
zero-loss fraction 0.50 \\
\longtok{curriculum\_opening/curriculum} & Tinker → W\&B
\texttt{study\ logging\ group} (ungrouped) & Qwen/\allowbreak{}Qwen3.5-4B & GRPO +
curriculum filter & 4 \emph{default} & GSM8K & 0 & 8 & 1 run; Δ +0.0500;
held-out 20; oversample 4.81×; zero-loss fraction 0.00; superseded, see
below \\
\longtok{token\_budget/baseline-b30000-s0…s2} & Tinker → W\&B
\texttt{study\ logging\ group} (\longtok{token-budget}) & Qwen/\allowbreak{}Qwen3.5-4B
& GRPO & 4 \emph{default} & GSM8K & 0--2 & 30k-token budget & 3 runs;
mean Δ +0.0278; no groups skipped \\
\longtok{token\_budget/curriculum-b30000-s0…s2} & Tinker → W\&B
\texttt{study\ logging\ group} (\longtok{token-budget}) & Qwen/\allowbreak{}Qwen3.5-4B
& GRPO + curriculum filter & 4 \emph{default} & GSM8K & 0--2 & 30k-token
budget & 3 runs; mean Δ +0.0278; 11--20 groups skipped per seed \\
\longtok{hard\_curriculum/baseline-hard-s0…s2} & Tinker → W\&B
\texttt{study\ logging\ group} (\longtok{hard-curriculum}) &
Qwen/\allowbreak{}Qwen3.5-4B & GRPO & 6 \emph{default} & GSM8K hard-learnable pool
(40 of 300 probed) & 0--2 & 10 & 3 runs; mean Δ +0.0278; zero-loss
fraction 0.1--0.2 \\
\longtok{hard\_curriculum/curriculum-hard-s0…s2} & Tinker → W\&B
\texttt{study\ logging\ group} (\longtok{hard-curriculum}) &
Qwen/\allowbreak{}Qwen3.5-4B & GRPO + curriculum filter & 6 \emph{default} & GSM8K
hard-learnable pool (40 of 300 probed) & 0--2 & 10 & 3 runs; mean Δ
0.0000; zero-loss fraction 0.0 \\
\longtok{p4\_surprise/\{sum,mean,surprise\}-s0…s1} & Tinker → W\&B
\texttt{study\ logging\ group} (\longtok{p4-surprise}) & Qwen/\allowbreak{}Qwen3.5-4B
& GRPO, loss variant & 4 \emph{default} & GSM8K & 0--1 & 8 & 6 runs;
mean Δ +0.1250 (sum) / 0.0000 (mean) / +0.0417 (surprise) \\
\longtok{p1\_layerfreeze/result.json} & Colab L4 &
Qwen/\allowbreak{}Qwen2.5-1.5B-Instruct & --- (gradient probe, no optimizer steps) &
--- & --- & --- & 5 & 1 run; step-1 predicts final top-layer set at
overlap 1.0; top-25\% concentration 0.476 \\
\longtok{p1\_layerfreeze/scaled\_\allowbreak{}4seed\_\allowbreak{}result.json} & Colab L4 &
Qwen/\allowbreak{}Qwen2.5-3B-Instruct & GRPO & 4 & GSM8K (24 problems) & 0--3 & 10 &
4 runs; overlap 0.0833 ± 0.0481 --- approximately chance; concentration
0.3908 \\
\longtok{p1\_layerfreeze/freeze\_\allowbreak{}flop\_\allowbreak{}result.json} & Colab L4 &
Qwen/\allowbreak{}Qwen2.5-1.5B-Instruct & GRPO (step-1-frozen arm) & --- & --- & 0--1
& --- & 4 runs, full and frozen arms; frozen arm gains 0.0000 against
0.0625 at 60.7\% of the parameter count \\
\longtok{p1\_layerfreeze/emergence\_result.json} & Colab L4 &
Qwen/\allowbreak{}Qwen2.5-1.5B-Instruct & GRPO (step-1 layer selection) & --- & --- &
0--1 & --- & 2 runs; mean step-1 overlap 0.643 \\
\end{longtable}
\endgroup

Three provenance notes belong with this table, because two of its blocks
carry superseded siblings and one carries a refuted claim. First,
\texttt{campaign/FINDINGS.md} describes a nine-run, three-seed version
of the campaign on a twelve-prompt held-out set with arm means of +0.028
and −0.028; \texttt{campaign/results.json} holds fifteen runs over six
seeds on a thirty-prompt held-out set (0.8667 = 26/30), and the table
follows the per-run file. The FINDINGS warning that
\longtok{zero\_\allowbreak{}loss\_\allowbreak{}frac} was invalid (an empty-metrics-dict bug,
``since fixed'') belongs to the nine-run batch, not to that file.
Second, the P1 layer-freeze block has the same shape: the two-seed 3B
pass (overlap 0.1111 ± 0.0000) is superseded by the four-seed pass
(0.0833 ± 0.0481), and both sit near chance, so the toy-scale overlap of
1.0 must not be read as a transferable layer-selection signal. Third,
the two single-seed training blocks flagged SUSPECT by independent
verification (kimi, 2026-07-06) are superseded by the multi-seed blocks
above: the curriculum opening's +0.05 (campaign curriculum and baseline
arms separate by 0.006 over six seeds, at 3.3--7.2× oversample cost),
and a single-seed G = 4 gain of +0.125 in
\longtok{p3\_groupsize/sweep\_results.json} --- the campaign's G4 seed-0
baseline gains 0.0000, not 0.125, and that sweep file retains only two
of its twelve entries, so the W\&B group is its record.

\subsection{Same-stack rerun on the unsaturated
configuration}\label{m1b-run-registry}

The rerun of §6.5.5 was executed on Modal on 2 October 2026 by
\longtok{platform\_hybrid/experiments/modal/modal\_\allowbreak{}samestack\_\allowbreak{}gsm8k\_\allowbreak{}cot.py}:
15 runs, one A10G each, about 14 minutes of training per run. Per-run
results, including per-item pre- and post-training correctness and step
logs, are in
\longtok{platform\_hybrid/experiments/results/samestack\_\allowbreak{}gsm8k\_\allowbreak{}cot\_\allowbreak{}full.json};
adapters are on the a retained result volume Modal volume under the
study adapter directory. No W\&B run was logged.

\begingroup\scriptsize\setlength{\tabcolsep}{3pt}
\begin{longtable}[]{@{}
  >{\raggedright\arraybackslash\hspace{0pt}}p{(\linewidth - 18\tabcolsep) * \real{0.1111}}
  >{\raggedright\arraybackslash\hspace{0pt}}p{(\linewidth - 18\tabcolsep) * \real{0.1111}}
  >{\raggedright\arraybackslash\hspace{0pt}}p{(\linewidth - 18\tabcolsep) * \real{0.1111}}
  >{\raggedright\arraybackslash\hspace{0pt}}p{(\linewidth - 18\tabcolsep) * \real{0.1111}}
  >{\raggedright\arraybackslash\hspace{0pt}}p{(\linewidth - 18\tabcolsep) * \real{0.1111}}
  >{\raggedright\arraybackslash\hspace{0pt}}p{(\linewidth - 18\tabcolsep) * \real{0.1111}}
  >{\raggedright\arraybackslash\hspace{0pt}}p{(\linewidth - 18\tabcolsep) * \real{0.1111}}
  >{\raggedright\arraybackslash\hspace{0pt}}p{(\linewidth - 18\tabcolsep) * \real{0.1111}}
  >{\raggedright\arraybackslash\hspace{0pt}}p{(\linewidth - 18\tabcolsep) * \real{0.1111}}@{}}
\caption{Same-stack rerun on the unsaturated configuration
(§6.5.5).}\tabularnewline
\toprule\noalign{}
\begin{minipage}[b]{\linewidth}\raggedright
Run identifier
\end{minipage} & \begin{minipage}[b]{\linewidth}\raggedright
Stack / platform
\end{minipage} & \begin{minipage}[b]{\linewidth}\raggedright
Model
\end{minipage} & \begin{minipage}[b]{\linewidth}\raggedright
Algo
\end{minipage} & \begin{minipage}[b]{\linewidth}\raggedright
G
\end{minipage} & \begin{minipage}[b]{\linewidth}\raggedright
Dataset
\end{minipage} & \begin{minipage}[b]{\linewidth}\raggedright
Seed
\end{minipage} & \begin{minipage}[b]{\linewidth}\raggedright
Steps
\end{minipage} & \begin{minipage}[b]{\linewidth}\raggedright
Outcome / status
\end{minipage} \\
\midrule\noalign{}
\endfirsthead
\toprule\noalign{}
\begin{minipage}[b]{\linewidth}\raggedright
Run identifier
\end{minipage} & \begin{minipage}[b]{\linewidth}\raggedright
Stack / platform
\end{minipage} & \begin{minipage}[b]{\linewidth}\raggedright
Model
\end{minipage} & \begin{minipage}[b]{\linewidth}\raggedright
Algo
\end{minipage} & \begin{minipage}[b]{\linewidth}\raggedright
G
\end{minipage} & \begin{minipage}[b]{\linewidth}\raggedright
Dataset
\end{minipage} & \begin{minipage}[b]{\linewidth}\raggedright
Seed
\end{minipage} & \begin{minipage}[b]{\linewidth}\raggedright
Steps
\end{minipage} & \begin{minipage}[b]{\linewidth}\raggedright
Outcome / status
\end{minipage} \\
\midrule\noalign{}
\endhead
\bottomrule\noalign{}
\endlastfoot
\longtok{samestack\_gsm8k/grpo\_\allowbreak{}g8\_\allowbreak{}s\{42,123,456,789,1024\}} & Modal
A10G, local trainer & Qwen/\allowbreak{}Qwen2.5-1.5B-Instruct & GRPO & 8 & GSM8K
(test 200) & 5 seeds & 30 & post 0.230--0.255, mean 0.245 \\
\longtok{samestack\_gsm8k/grpo\_\allowbreak{}g2\_\allowbreak{}s\{…\}} & Modal A10G, local trainer &
Qwen/\allowbreak{}Qwen2.5-1.5B-Instruct & GRPO & 2 & GSM8K (test 200) & 5 seeds & 30
& post 0.230--0.280, mean 0.250 \\
\longtok{samestack\_gsm8k/ppo\_s\{…\}} & Modal A10G, local trainer &
Qwen/\allowbreak{}Qwen2.5-1.5B-Instruct & PPO, value head & --- & GSM8K (test 200) &
5 seeds & 30 & post 0.160--0.215, mean 0.195 \\
\end{longtable}
\endgroup

\section{E1--E14 evaluation campaign
lanes}\label{e1e14-evaluation-campaign-lanes}

The fourteen evaluation lanes are not training runs; they are
evaluations of a single frozen actor, and they carry the registry's
strictest receipt discipline. The actor is base
\texttt{Qwen/Qwen3.6-35B-A3B} (commit \texttt{995ad96e…}) with adapter
\texttt{campaign-adapter} (commit \texttt{withheld-adapter-revision…}),
bfloat16, trained on Tinker, run seed 809 taken from the adapter name.
The adapter was trained for 40 steps at G = 4 and learning rate 1 × 10⁻⁵
by REINFORCE with a group-mean baseline (no KL, no clipping), as
recorded in the local W\&B run directory for run {[}identifier
withheld{]} and the run receipt
\texttt{withheld\ actor-training\ receipt} (Chapter 8 §8.1.1). Those are
properties of the single actor rather than of any lane, and every lane
carries the same run seed (809), so algorithm, group size, seed and step
columns are omitted from Table A.8. Terminal states are from the ledger
of record, \texttt{withheld\ historical\ campaign\ chronology}; all
scores are from \longtok{outputs/E1\_\allowbreak{}E14\_\allowbreak{}FINAL\_\allowbreak{}RESULTS\_\allowbreak{}2026-09-19.md},
except the six rows marked ``finish 2026-09-27'', which are from
\longtok{outputs/finish\_\allowbreak{}pending\_\allowbreak{}2026-09-27/\textless{}lane\textgreater{}/result.json}.

\textbf{E1--E14 campaign lanes.}

\begingroup\scriptsize\setlength{\tabcolsep}{3pt}
\begin{longtable}[]{@{}
  >{\raggedright\arraybackslash\hspace{0pt}}p{(\linewidth - 10\tabcolsep) * \real{0.2000}}
  >{\raggedright\arraybackslash\hspace{0pt}}p{(\linewidth - 10\tabcolsep) * \real{0.2000}}
  >{\raggedright\arraybackslash\hspace{0pt}}p{(\linewidth - 10\tabcolsep) * \real{0.2000}}
  >{\raggedright\arraybackslash\hspace{0pt}}p{(\linewidth - 10\tabcolsep) * \real{0.2000}}
  >{\raggedright\arraybackslash\hspace{0pt}}p{(\linewidth - 10\tabcolsep) * \real{0.2000}}@{}}
\caption{E1-E14 evaluation campaign lanes with their run identifiers,
stacks and platforms, models, suites and outcomes.}\tabularnewline
\toprule\noalign{}
\begin{minipage}[b]{\linewidth}\raggedright
Lane / run identifier
\end{minipage} & \begin{minipage}[b]{\linewidth}\raggedright
Stack / platform
\end{minipage} & \begin{minipage}[b]{\linewidth}\raggedright
Model / actor
\end{minipage} & \begin{minipage}[b]{\linewidth}\raggedright
Dataset / suite
\end{minipage} & \begin{minipage}[b]{\linewidth}\raggedright
Outcome / terminal state
\end{minipage} \\
\midrule\noalign{}
\endfirsthead
\toprule\noalign{}
\begin{minipage}[b]{\linewidth}\raggedright
Lane / run identifier
\end{minipage} & \begin{minipage}[b]{\linewidth}\raggedright
Stack / platform
\end{minipage} & \begin{minipage}[b]{\linewidth}\raggedright
Model / actor
\end{minipage} & \begin{minipage}[b]{\linewidth}\raggedright
Dataset / suite
\end{minipage} & \begin{minipage}[b]{\linewidth}\raggedright
Outcome / terminal state
\end{minipage} \\
\midrule\noalign{}
\endhead
\bottomrule\noalign{}
\endlastfoot
E1 SWE-bench Pro & native harness, Tinker actor & Qwen3.6-35B-A3B +
adapter & SWE-bench Pro & 2/731 = 0.274\% pass@1 (97.54\% coverage,
713/731 native evals); originals retained \\
E1 replacement (finish 2026-09-27) & native swebench 5.0.2, Modal
sandboxes & as above & SWE-bench Multilingual & 1/190 resolved (57
graded); earlier waves 4/110 kept separate --- \texttt{COMPLETE} \\
E2 FrontierSWE & native harness & as above & FrontierSWE & 1/17 tasks;
replay normalized 0.8628 \\
E2 replacement (finish 2026-09-27) & native grader, GCP VM per capsule &
as above & CORE-Bench hard & 27/45 = 0.600 task accuracy ---
\texttt{COMPLETE} \\
E3 SDAB & --- & as above & SDAB (private bundle) & no local or paid path
--- \longtok{CLOSED\_EXTERNAL} \\
E4 BankerToolBench & Tinker→Modal bridge & as above & BankerToolBench &
1/100 tasks, recovery 0.3115 --- \longtok{CLOSED\_PARTIAL} \\
E4 base-model rerun & Modal \longtok{/v1/responses} & base model,
zero-init LoRA (0 training steps) & BankerToolBench & 100/100 trials,
mean reward 0.0; bridge stop-sequence defect is a plausible mechanism,
with unresolved causal contribution (§8.5) \\
E5 APEX-Agents & native harness & as above & APEX-Agents & 7/480
native-scored; prefix mean 0.050505 \\
E5 replacement (finish 2026-09-27) & native tau2, Modal 65k actor & as
above & Tau3 banking\_knowledge & 8/97 = 0.082 pass\^{}1 (25 infra
errors counted as failures) --- \texttt{COMPLETE} \\
E6 WebArena (finish 2026-09-27) & native evaluator, GCP VM & as above &
WebArena & 90/812 = 0.111, lower bound; 108 ungraded (judge unavailable)
--- \texttt{COMPLETE} \\
E7 BinaryAudit & native harness & as above & BinaryAudit (private
payload) & 1/46 attempted, verifier reward 0.0, no grade ---
\longtok{CLOSED\_EXTERNAL} \\
E8 LAB-Bench (public) & native harness & as above & LAB-Bench public, 8
categories & \textbf{1967/1967 COMPLETE} \\
E8 LifeSciBench (original) & --- & as above & LifeSciBench private &
\longtok{CLOSED\_EXTERNAL} \\
E9 MLE-bench & Modal streaming arm & as above & MLE-bench & 40/75
natively graded (53.33\%); suite score null \\
E9 replacement (finish 2026-09-27) & OpenHands CodeActAgent, GCP VM & as
above & ML-Dev-Bench & 10/34 = 0.294 task success ---
\texttt{COMPLETE} \\
E10 AgentDojo (benign) & native harness & as above & AgentDojo benign
utility & 97/97 tasks completed, utility 88/97 = 0.907 ---
\texttt{COMPLETE} \\
E10 AgentHarm (original) & --- & as above & AgentHarm private &
\longtok{CLOSED\_EXTERNAL} \\
E11 VerilogEval & native harness & as above & VerilogEval (both
framings) & \textbf{312/312 = 129/312 pass@1 (41.35\%)}; 67/156
completion + 62/156 spec-to-RTL \\
E12 AppBench & --- & as above & AppBench deployment &
\longtok{CLOSED\_EXTERNAL} \\
E13 BALROG (finish 2026-09-27) & native BALROG evaluator, local & as
above & BALROG & 255/255 episodes; 26.1\% progression ---
\texttt{COMPLETE} \\
E13 OpenReward (original) & --- & as above & OpenReward held-out games &
externally blocked, no closure record \\
E14 Omni-MATH & native scorer & as above & Omni-MATH & 4426/4428
accepted (99.95\%); accuracy 2271/4428 = 51.29\% (official scorer
51.31\% on 4426 judged rows) --- \longtok{COMPLETE\_\allowbreak{}TERMINAL\_\allowbreak{}NOTE} \\
E14 FrontierMath (original) & --- & as above & FrontierMath hosted &
\longtok{CLOSED\_EXTERNAL} \\
\end{longtable}
\endgroup

Two source documents disagree on two lanes, and both observations are
retained rather than resolved into one, because the disagreement is
itself evidence: the historical snapshot records E7 as ``1/28 errored,
no grade'' against the terminal-state table's ``1/46 attempted'' ---
both real observations at different points in the lane's life. Likewise
E11's canonical figure is 129/312 = 41.35\%; a 129/311 sensitivity
figure survives in non-canonical receipts and is not used as the
headline.

\section{Completeness limitations}\label{completeness-limitations}

The historical registry is an index of its retained snapshot and is
incomplete as a scientific dataset. Many records lack model, dataset,
algorithm, group size, steps or a usable outcome. Missing fields remain
unknown; names and provider state are not used to infer them. The public
version does not include the account-wide field-coverage table or
private per-identity export.

\textbf{Scientific metadata limitations.}

\begin{longtable}[]{@{}
  >{\raggedright\arraybackslash\hspace{0pt}}p{(\linewidth - 4\tabcolsep) * \real{0.5000}}
  >{\raggedright\arraybackslash\hspace{0pt}}p{(\linewidth - 4\tabcolsep) * \real{0.5000}}@{}}
\caption{Metadata limitations in the historical retained registry;
private account-level field counts are omitted.}\tabularnewline
\toprule\noalign{}
\begin{minipage}[b]{\linewidth}\raggedright
Required scientific field
\end{minipage} & \begin{minipage}[b]{\linewidth}\raggedright
Interpretation of missing data
\end{minipage} \\
\midrule\noalign{}
\endfirsthead
\toprule\noalign{}
\begin{minipage}[b]{\linewidth}\raggedright
Required scientific field
\end{minipage} & \begin{minipage}[b]{\linewidth}\raggedright
Interpretation of missing data
\end{minipage} \\
\midrule\noalign{}
\endhead
\bottomrule\noalign{}
\endlastfoot
Model and dataset & Do not infer identity from a display name or
tokenizer family \\
Algorithm and group size & An unrecorded configuration remains
unknown \\
Training steps & Provider state does not establish completed
optimization \\
Outcome & A logged record without a metric is not a scored scientific
result \\
\end{longtable}

Four further limits apply.

\begin{enumerate}
\def\labelenumi{\arabic{enumi}.}
\item
  \textbf{Aggregate presentation.} This appendix names 105 runs
  individually: the 4 cross-framework runs of Table A.2, the 20 Modal
  campaign runs of A.3, the 12 sweep runs (in 4 rows) of A.4, the 46
  opening runs (in 16 rows) of A.6, and the 23 E-lane rows of A.8 ---
  plus the 18 claim rows of Table A.5, which index claims rather than
  runs. Other historical inventory entries are not published in this
  derivative.
\item
  \textbf{Non-runs inside the count.} Three entries are not training
  results: \texttt{verl-dryrun} and \texttt{openrlhf-dryrun} are seeded
  deterministic fallbacks, and the E4 base-model rerun performs zero
  training steps; to these add §A.3.1's P1 gradient probe, which takes
  no optimizer steps at all. They are counted because they are runs in
  the registry's sense, and flagged so no reader treats them as
  measurements.
\item
  \textbf{Identity conflicts.} Historical source-identity disagreements
  remain a limitation; private provider reconciliation and its counts
  are withheld.
\item
  \textbf{Snapshot scope.} Historical source inventories differ in
  scope. They are neither current live accounts nor counts of
  independent experiments, and this public package does not provide a
  complete inventory replay.
\end{enumerate}

None of this changes a reported result, because every result in Chapters
6--8 is bound to named artefacts rather than to the registry as a whole.
It does bound what the registry can be used for: it establishes
provenance and identity, shows that the run population is far larger
than the reported claims, and shows that those claims rest on a small,
explicitly tiered subset with the P1--P8 empirical rows at Tier C or X
and P11 at Tier C under the stated per-run step threshold. The four Tier
X P1--P8 rows must not be cited as supporting evidence.

The run totals that appear in different chapters count different
corpora. Table A.10 lists the scientific study corpora retained
publicly.

\textbf{Scientific study corpora.} Each total is a different corpus, not
a different count of the same one.

\begingroup\scriptsize\setlength{\tabcolsep}{3pt}
\begin{longtable}[]{@{}
  >{\raggedright\arraybackslash\hspace{0pt}}p{(\linewidth - 8\tabcolsep) * \real{0.2500}}
  >{\raggedright\arraybackslash\hspace{0pt}}p{(\linewidth - 8\tabcolsep) * \real{0.2500}}
  >{\raggedright\arraybackslash\hspace{0pt}}p{(\linewidth - 8\tabcolsep) * \real{0.2500}}
  >{\raggedright\arraybackslash\hspace{0pt}}p{(\linewidth - 8\tabcolsep) * \real{0.2500}}@{}}
\caption{Run totals quoted in this report and the corpus each one
counts.}\tabularnewline
\toprule\noalign{}
\begin{minipage}[b]{\linewidth}\raggedright
Corpus
\end{minipage} & \begin{minipage}[b]{\linewidth}\raggedright
Count
\end{minipage} & \begin{minipage}[b]{\linewidth}\raggedright
What it is
\end{minipage} & \begin{minipage}[b]{\linewidth}\raggedright
Used in
\end{minipage} \\
\midrule\noalign{}
\endfirsthead
\toprule\noalign{}
\begin{minipage}[b]{\linewidth}\raggedright
Corpus
\end{minipage} & \begin{minipage}[b]{\linewidth}\raggedright
Count
\end{minipage} & \begin{minipage}[b]{\linewidth}\raggedright
What it is
\end{minipage} & \begin{minipage}[b]{\linewidth}\raggedright
Used in
\end{minipage} \\
\midrule\noalign{}
\endhead
\bottomrule\noalign{}
\endlastfoot
P1 cross-scale roster & 70+ & Runs across seven RL libraries and five
model families; 12 of them are the fitted anchors & Chapter 6 §6.2 \\
Instrumented roster (P8 workshop source) & 79 & The instrumented roster
of \longtok{paper\_\allowbreak{}P8\_\allowbreak{}workshop.tex}, 0.6B to \textasciitilde671B &
Chapter 4 §4.1 \\
P5 audit corpus & 790 & The corpus from which the 368-run audit is drawn
(\longtok{p5\_stack.tex}); its relation to the ledger is not recorded &
Chapter 7 §7.1 \\
Mega-campaign cells & 98 & Sampling-only cells of
\longtok{mega\_20260704} & Chapter 7 §7.1 \\
E1--E14 actor training & 1 & The 40-step \longtok{pavlov\_portfolio} run
of August 2026, after the ledger cut-off & Chapter 8 §8.1.1 \\
Unsaturated same-stack rerun & 15 & Three arms × five seeds on Modal, 2
October 2026 (§A.3.2) & Chapter 6 §6.5.5 \\
\end{longtable}
\endgroup



\chapter{Reproducibility and
Provenance}\label{reproducibility-and-provenance}

\section{What this appendix
establishes}\label{what-this-appendix-establishes}

This appendix records three things that should not be confused with one
another. The first is what a reviewer can independently re-run: a pinned
container, a pinned dependency set, a headline training command, and the
tolerance within which a rerun is expected to land. The second is what
the repository verifies about itself without re-running any model: an
audit suite driven by a single entry point, and a deterministic ledger
check that recomputes the reported arithmetic directly from surviving
evidence files. The third is provenance: which parts of this work were
produced by the Semester 3 group, which by the individual contribution
reported here, and what was lost along the way.

\section{The pinned runtime}\label{the-pinned-runtime}

The container definition is \texttt{Dockerfile}, built from
\longtok{nvidia/cuda:12.4.1-cudnn-devel-ubuntu22.04}. It installs Python
3.12 through the deadsnakes archive rather than the distribution
default, with the in-file comment that Python 3.12 ``matches the frozen
runtime validated 106/106''. It sets \texttt{PYTHONHASHSEED=42} and a
fixed \texttt{LANG}/\texttt{LC\_ALL} of \texttt{C.UTF-8}; stamps the
build commit into an OCI \longtok{org.opencontainers.image.revision}
label and a \texttt{.build\_info} file; fixes \texttt{SEED=42},
\texttt{WANDB\_MODE=offline}, \longtok{HF\_\allowbreak{}HUB\_\allowbreak{}DISABLE\_\allowbreak{}TELEMETRY=1} and
\longtok{OMP\_\allowbreak{}NUM\_\allowbreak{}THREADS=8}; self-checks by importing the core stack at
build time; and runs under \texttt{tini}.

The authoritative pin set is \texttt{requirements.txt}, headed as
``pinned to the frozen E1 runtime lock \ldots{} validated 106/106 on
this stack''. Its exact pins include \texttt{torch==2.7.1},
\texttt{numpy==2.2.6}, \texttt{transformers==5.5.4},
\texttt{trl==1.2.0}, \texttt{datasets==4.8.4},
\texttt{accelerate==1.13.0}, \texttt{peft==0.19.1},
\texttt{huggingface-hub==1.11.0} and \texttt{wandb==0.21.0}, with
bounded ranges for scipy, safetensors, tokenizers, rliable, math-verify,
latex2sympy2-extended and pytest, alongside the framework clients
\texttt{tinker}, \texttt{tinker-cookbook} and \texttt{atroposlib}.

The 106/106 figure is the recorded outcome of the frozen-runtime
validation, which \texttt{execution-notes.md} states as passing the
complete pinned pilot/preregistration suite on Python 3.12.13 with the
pinned versions listed above, and the focused amendment, protocol,
launcher and verifier suite passing 55/55. The same record binds five
artefacts by SHA-256 --- protocol document, screening manifest, launch
manifest, launch fingerprint and source binding --- and states that a
mismatch in package, Python, A100, FLOP, token or gradient conditions
fails closed rather than proceeding.

\section{The reviewer's entry point}\label{the-reviewers-entry-point}

\texttt{REPRODUCE.md} is the reviewer-facing document. Its fastest path
is the structural smoke test at
\longtok{platform\_modal/scripts/smoke\_test.sh}, which runs in about two
minutes without a Tinker key and about eight when a live key is supplied
so that the wire protocol is exercised; the keyless path needs no GPU
and no account, and is offline only once the GSM8K dataset is cached,
because one step loads it.

The headline result is GRPO on Qwen3-8B over GSM8K for 30 steps at LoRA
rank 32, group size 8 and seed 42, recorded as a peak per-step training
accuracy (the binary reward) of 62.5 \% with a last-10 average of 34.4
\%; these summarise the training curve of the seed-42 run, not a
held-out evaluation, and Table 1.1 lists them beside the other Qwen3-8B
GSM8K figures they are sometimes confused with. Reviewer-side reruns are
accepted within ±5 percentage points on the last-10 figure and ±10
points on peak, justified in the same document by the residual
non-determinism discussed in §B.10. Per-seed reference values and W\&B
run identifiers are tabulated for seeds 42, 123, 456 and 789, with seed
1024 marked as a reviewer-only rerun; the canonical multi-seed set is
\texttt{\{42,\ 123,\ 456,\ 789,\ 1024,\ 2048,\ 4096,\ 8192,\ 16384,\ 32768\}}.
The group-size ablation is expected to reproduce a peak of 50.0 \% at
G=2, 75.0 \% at G=4, 100 \% at G=8 and 71.9 \% at G=16, with last-10
values of 37.5 \%, 52.1 \%, 84.4 \% and 38.0 \%. Five-seed aggregation
runs through \texttt{utils/stats.py}, which now accepts
\texttt{-\/-bootstrap-samples} (default 10,000) and prints the mean,
standard error and a percentile-bootstrap 95\% interval; its
\texttt{-\/-format} and \texttt{-\/-rliable} flags are accepted but not
yet acted on, and it reads a
\longtok{results/\textless{}experiment\textgreater{}/seed\_\textless{}N\textgreater{}/*.jsonl}
layout that the trainer does not write, so a reviewer must write one
metrics file per seed first, as \texttt{REPRODUCE.md} §3.4 states.
\texttt{utils/verify\_results.py} checks a rerun against
\longtok{platform\_hybrid/paper/expected\_results.json} at those
tolerances. It parses the per-seed \texttt{avg\_last10} /
\texttt{peak\_reward} block that \texttt{grpo\_cli} actually prints, and
it exits non-zero when a value is outside tolerance or missing, when no
result file matches an expected experiment, or (exit 2) when the
expectations file or results directory is missing. Before the 3 October
fix it parsed a line format the trainer never prints and exited 0 when
nothing matched, so a failed rerun would have been reported as a pass;
\longtok{tests/test\_\allowbreak{}verify\_\allowbreak{}results.py} now covers both formats and the
failure cases. One limitation remains: every seed is checked against the
single seed-42 headline (0.344 / 0.625), so a seed whose recorded values
differ by more than the tolerance, such as seed 456 with a 100\% peak,
is reported as outside tolerance.

The cost is stated plainly: about five GPU-hours on a single A100-80GB
for the headline run, against roughly 446 GPU-hours for a full paper
reproduction. Dataset provenance is tabulated with pinned revisions ---
GSM8K at \texttt{740312a} with 7,473 training and 1,319 test rows, plus
MATH-500, HumanEval, ArenaHard Auto, NoRobots and OpenThoughts3 each at
a named snapshot --- with exact shas in \texttt{ARTIFACT.md} §3.4.

\section{Where the documentation has drifted from the
code}\label{where-the-documentation-has-drifted-from-the-code}

Four discrepancies existed between the repository's prose and its live
tree; \texttt{REPRODUCE.md} was corrected on 3 October 2026 and
\texttt{ARTIFACT.md} was not. The training driver has moved to
\longtok{platform\_tinker/grpo\_\allowbreak{}gsm8k\_\allowbreak{}base.py}. \texttt{REPRODUCE.md}
now invokes it at that path, while \texttt{ARTIFACT.md} still invokes
\longtok{grpo\_\allowbreak{}gsm8k\_\allowbreak{}base.py} at the repository root; the flags are
unchanged, because the CLI in
\longtok{platform\_tinker/tinkerrl/grpo\_cli.py} carries a compatibility
entry point, \texttt{legacy\_main}, that translates the documented flag
spellings onto the current ones. The same correction pointed
\texttt{REPRODUCE.md}'s Modal, figure and Atropos commands at their live
paths (\longtok{platform\_modal/scripts/modal\_\allowbreak{}run\_\allowbreak{}experiments.py},
\longtok{platform\_modal/scripts/make\_\allowbreak{}paper\_\allowbreak{}figures.py},
\texttt{platform\_tinker/atropos}), removed an Atropos \texttt{-\/-seed}
flag the launcher does not accept, replaced the expected-output blocks
with the format the trainer prints, and added a verify step to the quick
start. Every documented command file now exists and the Python entry
points answer \texttt{-\/-help}; the Atropos scripts need
\texttt{tenacity}, \texttt{latex2sympy2\_extended} and
\texttt{atroposlib}, which \texttt{requirements.txt} installs in the
Docker image but the local development environment lacks. The analysis
scripts \longtok{platform\_modal/scripts/p5p8/p5\_\allowbreak{}headline\_\allowbreak{}cis\_\allowbreak{}full.py}
and \longtok{platform\_modal/scripts/length\_\allowbreak{}bias\_\allowbreak{}iter24.py} previously
resolved their repository root and results directory incorrectly and
created directories on import; both now resolve paths from the
repository root, and their regenerated outputs are the source of the
corrected §6.5.3 and §7.1 values. \texttt{ARTIFACT.md} names the smoke
test at \texttt{scripts/smoke\_test.sh} rather than at
\longtok{platform\_modal/scripts/smoke\_test.sh}, though
\texttt{REPRODUCE.md} gives the correct path. The pins quoted in prose
are stale: \texttt{REPRODUCE.md} §9 states that
\texttt{requirements.txt} pins
\longtok{transformers\textgreater{}=4.56.2,\textless{}4.60.0} and a
\texttt{datasets\textless{}3.2} bound, whereas the live file pins
\texttt{transformers==5.5.4} and \texttt{datasets==4.8.4}, and
\texttt{ARTIFACT.md} §3.3 quotes a third, older set; similarly,
\texttt{ARTIFACT.md} describes the environment as Python 3.10 from the
Ubuntu 22.04 archive, whereas the \texttt{Dockerfile} installs Python
3.12 via deadsnakes and the frozen-runtime record validates on Python
3.12.13. The \texttt{ARTIFACT.md} items are explained by that document's
frozen status --- it records a Semester 3 artefact at commit
\longtok{9019b5c946f17ffe4761b5f4eac13b86d1e51cc0}, dated 19 April 2026
--- but \texttt{REPRODUCE.md} §9 is live guidance, not a frozen record,
and carries the same defect.

\section{The audit apparatus}\label{the-audit-apparatus}

The integrity suite is driven by a single entry point,
\longtok{platform\_local/run\_\allowbreak{}all\_\allowbreak{}audits.py}, which imports a shared
result interface from \texttt{utils/audit\_utils.py} and registers its
audits as a tuple of name-and-function pairs. As it stands in this tree,
that tuple contains \textbf{nine} audits: \texttt{claim\_issues},
\texttt{sync\_issues}, \texttt{anon\_issues}, \texttt{strength\_issues},
\texttt{package\_issues}, \texttt{workflow\_issues},
\longtok{export\_\allowbreak{}guard\_\allowbreak{}issues}, \texttt{caveat\_issues} and
\texttt{scientific\_issues}. Each is implemented by a corresponding
\texttt{platform\_local/} module carrying a \texttt{\_\_main\_\_} block,
so each can also run on its own; \longtok{tests/test\_\allowbreak{}audit\_\allowbreak{}runner.py}
asserts that the registered modules and the
\texttt{platform\_local/*\_audit.py} files are the same set.

Some repository documents give different counts, and the difference is
historical. At the blind-review sweep the suite was thirteen audits,
registered as a list of script filenames at the repository root and
executed as subprocesses; the committed log of that run records
\texttt{METRIC\ audits\_total=13} with ten passing and three failing,
and the same sweep states that all seventeen \texttt{*\_audit.py}
scripts then at the root were executed, with the suite covering thirteen
of them. The registry now holds nine audits, all passing: a pruning
commit removed obsolete audits, and a later restructuring replaced the
subprocess runner with the in-process interface, so the seventeen-script
description has narrowed by the same route. No script now prints a
metric line directly; the line is constructed in one place,
\texttt{AuditResult.metric\_line()} in \texttt{utils/audit\_utils.py},
and emitted by the nine audit modules plus the runner. A repository-wide
scan finds several dozen files matching \texttt{*\_audit.py}, but that
population is not a suite: it includes unrelated analyses, test modules
and a third-party file in the virtual environment.

The output contract is uniform. \texttt{AuditResult.metric\_line()}
emits
\texttt{METRIC\ \textless{}name\textgreater{}=\allowbreak{}\textless{}issue\_\allowbreak{}count\textgreater{}},
so every audit reports one machine-readable line whose right-hand side
is the number of issues it found, and an audit passes exactly when that
count is zero. \texttt{render\_suite} prints each audit's metric line
and issue codes, then three suite-level lines ---
\texttt{METRIC\ suite\_issues}, \texttt{METRIC\ audits\_total} and
\texttt{METRIC\ audits\_passing} --- before either listing the failing
audits by name or stating that all audits are passing. The entry point
exits \texttt{0} when the suite passes and \texttt{1} otherwise, which
makes it usable as a gate; the packaging step runs the suite before
exporting and refuses to package if it fails. The runner's own test
asserts this rendering against a synthetic two-audit suite, checking for
\texttt{METRIC\ audits\_total=2} and \texttt{METRIC\ suite\_issues=1}.

\section{The deterministic ledger
check}\label{the-deterministic-ledger-check}

Where the suite above audits the submission package's internal
consistency, a second instrument audits the campaign's reported figures.
\longtok{outputs/verification/LEDGER\_\allowbreak{}CODE\_\allowbreak{}CHECK\_\allowbreak{}2026-09-19.json},
under schema \texttt{ledger-code-check-v1}, records eleven checks of
which eleven passed and none failed. They cover E1, E8 coverage, E9, E11
and E14 counts and hashes and the lost-source probes; they do not cover
the E10 score, E8 accuracy, the E2, E4, E5 and E7 lanes or anything
recorded after 19 September, and the E14 check confirmed the official
scorer's 4,426-row denominator rather than the 4,428-row denominator
this report uses. Each recomputes a claim from the underlying evidence
rather than trusting a summary of it, and the four classes of check
prove four different things.

\begin{longtable}[]{@{}
  >{\raggedright\arraybackslash\hspace{0pt}}p{(\linewidth - 6\tabcolsep) * \real{0.3333}}
  >{\raggedright\arraybackslash\hspace{0pt}}p{(\linewidth - 6\tabcolsep) * \real{0.3333}}
  >{\raggedright\arraybackslash\hspace{0pt}}p{(\linewidth - 6\tabcolsep) * \real{0.3333}}@{}}
\caption{The deterministic ledger check by class, listing each specific
check and what it establishes.}\tabularnewline
\toprule\noalign{}
\begin{minipage}[b]{\linewidth}\raggedright
Class of check
\end{minipage} & \begin{minipage}[b]{\linewidth}\raggedright
Specific checks
\end{minipage} & \begin{minipage}[b]{\linewidth}\raggedright
What it proves
\end{minipage} \\
\midrule\noalign{}
\endfirsthead
\toprule\noalign{}
\begin{minipage}[b]{\linewidth}\raggedright
Class of check
\end{minipage} & \begin{minipage}[b]{\linewidth}\raggedright
Specific checks
\end{minipage} & \begin{minipage}[b]{\linewidth}\raggedright
What it proves
\end{minipage} \\
\midrule\noalign{}
\endhead
\bottomrule\noalign{}
\endlastfoot
Exact divisions & E1 \texttt{2/731}; E11 \texttt{129/312}; E11
\longtok{67+62=129}; E14 \longtok{4426/4428} and accuracy
\longtok{0.5131043831902395} (which is 2271/4426, the scorer's judged-row
denominator; the reported headline counts the 2 unjudged rows as
failures, 2271/4428 = 0.5129); E9 \texttt{40\ of\ 75} & That each
reported ratio is the quotient of the two integers the report claims,
and that a partition sums to its whole \\
Receipt presence & E8 eight categories summing to 1967 with every
\texttt{evaluated\ =\allowbreak{}=\allowbreak{}\ expected\_\allowbreak{}total}; E9 graded-row counts; the
strict-completion set & That every claimed unit of coverage has a
surviving graded row behind it, and that the set of strictly complete
suites is exactly \texttt{{[}E8,\ E10,\ E11{]}} \\
Sealed-artefact hash equality & E1 v6 sealed request and fixture receipt
& That the sealed request and the validation receipt are byte-identical
to the artefacts recorded at sealing time \\
Lost-source absence probes & Source availability check & That four
historical temporary execution workspace paths are genuinely absent, so
a later re-creation cannot be mistaken for the lost original \\
\end{longtable}

The distinction matters under questioning. An exact division proves the
arithmetic is not a rounding convenience or a transcription error; it
cannot prove the underlying evidence is sound. A receipt-presence check
proves coverage claims are backed by artefacts; it cannot prove those
artefacts were produced correctly. A hash-equality check proves identity
between two artefacts --- here, between what was sealed and what was
later validated --- and is the strongest of the four because it does not
depend on re-running anything. An absence probe proves a negative: it
prevents a later reconstruction from being passed off as the original.
Two instances show the classes in use. The historical E1 records
describe matching request and validation bindings without establishing
that execution occurred. Private request identifiers are withheld. The
E9 records show 41 graded rows resolving to 40 unique competitions
against a denominator of 75 in the streaming arm and exactly one valid
grade in the merged-VLLM arm, which is the check that confirms the two
serving arms were never pooled.

\section{The receipt model}\label{the-receipt-model}

A receipt in this campaign is a durable JSON artefact binding an outcome
to the inputs that produced it, and the binding carries the evidential
weight rather than the file's existence. The clearest statement of the
principle is the receipt-bundle auditor's docstring, which records that
the campaign and the receipt builders are kept deliberately separate
from the auditor, and that the auditor ``never imports a provider SDK,
contacts a service, launches a job, or treats
\texttt{status}/\texttt{online}/a URL/ID alone as provenance''. The
auditor reconstructs contract coverage, receipt predicates, canonical
hash rules, budget caps and cross-artefact bindings, carrying its
expectations as constants rather than reading them from what it checks
--- suite counts of 12 training, 14 primary-evaluation and 26 total ---
so a receipt set that quietly loses a suite fails the audit rather than
redefining it.

Receipts bind four dimensions of a run: protocol, runtime, source and
event order. The frozen-runtime record applies that discipline to the
runtime dimension alone, binding five artefacts --- protocol document,
screening manifest, launch manifest, launch fingerprint and source
binding --- by SHA-256 and failing closed on any mismatch. The judgment
layer for subjectively scored questions follows the same pattern: one
durable receipt per judgment binds the state bytes, the question text
and the raw answer, and battery execution is gated on a health check
rather than the service's availability, so a degraded answer cannot be
laundered into a valid one. Where a lane could not be scored, its
terminal state is recorded in the ledger of record with a named reopen
condition, so the absence of a result is itself receipted.

\section{Evidence loss in the finish
era}\label{evidence-loss-in-the-finish-era}

The historical record reports that four finish-era execution-source
components are unavailable: an E1 recovery driver, restored E9 runtime
dependencies, an E6 receipt validator and an E1 multilingual metadata
driver. This limits end-to-end reproduction. Operational
scratch-workspace names, executed request identifiers and private source
paths are omitted from this public derivative.

The four lost execution-source paths do not themselves contain the score
artefacts, but their loss limits end-to-end reproducibility. The
surviving evidence supports three narrower observations. First, none of
the four paths is the source of any reported score. The E1 and E9
replacement lanes were eventually run on 27 September with new thin
runners rather than the lost code (§8.4). The loss still bears on E1:
the source-context collector had to be re-implemented for that run and
was checked byte-for-byte against the original on only three recorded
tasks, which is one reason the 1/190 and 4/110 figures are not pooled.
Second, the sealed requests, validation receipts, regression fixtures
and the 10-of-10 targeted-test log from the lost E1 code all survive, so
a faithful re-implementation is constrained by surviving fixtures rather
than free to invent behaviour; that re-implementation has since landed
under a tracked path and passes the lost suite's recorded regressions,
four process-group and six cleanup. Third, the loss is confined to
finish-era iteration code: the replacement-lane native runners survive
in tracked locations, including a 515-line multilingual runner carrying
no reference to the deleted tree. The policy adopted in response ---
that execution sources live in tracked repository paths and never in
untracked scratch trees --- is the durable outcome of the incident.

\section{The Semester 3 and Semester 4 provenance
boundary}\label{the-semester-3-and-semester-4-provenance-boundary}

The boundary between inherited and individual work is a git tag.
\texttt{capstone-final-2026-04-25} resolves to commit
\longtok{21a99ef766543a1ca86bfbad188445d3596ea73c}, dated 23 April 2026,
whose subject line records the final metric-separation rewrite that
segregated training rewards, held-out capability and proxy diagnostics.
Two hundred and fifty commits precede that tag. A reviewer can reproduce
the state at the boundary by adding it as a worktree,
\texttt{git\ worktree\ add\ ../tinker-rl-lab-sem3\ capstone-final-2026-04-25},
which materialises the Semester 3 tree without disturbing the current
checkout.

The Semester 3 group comprised six students --- Arvind C R, Sandhya
Jeyaraj, Madhu Kumara L, Mohammad Rafi, Dhruva N Murthy and Arumugam
Chetty K --- under guides Anwesh Reddy Paduri and Narayana Darapaneni.
Their work is the multi-framework benchmark infrastructure and the
initial GRPO baseline experiments. Semester 4 was carried out
individually under Ramesh Prakash Guledgudd, and the declaration facing
this report assigns to that contribution the research direction, the
experimental expansion, the statistical analyses, the P1--P8 paper
series and the E1--E14 campaign. The solo-authorship commit
\longtok{282e6b4f0f158de11b3f5bb4656e2d871f68f5c0}, dated 7 July 2026,
records that division at paper level. The provenance record states that
584 commits and 5,050 changed files separated the boundary from the tree
as it then stood, and that the eight paper PDFs were compiled at commit
\longtok{006e90a38e4dabe2d0b4778be01dc7b27f2d3462} on 9 July 2026, one
recompiled the following day after a findings-ledger rename. Those
counts carry an as-of date rather than describing the current tree, so a
reviewer recomputing them will obtain a larger number.

One consequence of this history shapes the repository's structure: the
implementation remains shared, with experiment drivers, result tables,
figures, LaTeX sources and audit scripts referencing one another across
the boundary, and the provenance record is explicit that this is why the
code is not split into separate trees. What the boundary separates is
authorship of the research contribution, not access to the machinery.
Two reporting rules follow and are enforced throughout this report:
original-contract and replacement-scope results are never pooled, with
no cross-suite aggregate, which the ledger check independently verifies
through the arm-separation and component-split checks; and no
replacement benchmark's score is presented as the score of the benchmark
it replaces, with every substitution named as such at its point of use.

\section{Residual non-determinism and the limits of
reproduction}\label{residual-non-determinism-and-the-limits-of-reproduction}

Not everything here is reproducible in the strict sense, and the
tolerances exist because of it. Three sources of residual
non-determinism survive even with fixed seeds and deterministic cuDNN
settings: remote LoRA training is not bit-deterministic across scheduler
decisions, with runs replayed months apart drifting by two to four
percentage points on GSM8K last-10, documented through three reruns of
the same seed-123 configuration; some matmul operations on H100 and A100
hardware introduce non-determinism through CUDA and cuBLAS; and sampling
at temperature 0.8 with top-p 0.95 executes in a remote virtual machine,
so a fixed seed does not fix the sampled trajectory. The record states
the consequence itself: run-to-run variance on this setup is of the same
order as the effect reported in the paper's PPO-versus-GRPO comparison,
which limits the historical saturated comparison. The newer unsaturated
comparison (§6.5.5) is baseline-specific and test-sensitive, with
paired-t p = 0.016 but exact sign-flip p = 0.0625; neither supports a
general algorithm ranking.

Three limits remain. Lanes blocked on external access cannot be
reproduced from this repository, because the blocking resource is not in
it. The six lanes completed on 27 September can be rerun from the code
in their lane directories under
\longtok{outputs/finish\_\allowbreak{}pending\_\allowbreak{}2026-09-27/}, but only with cloud
accounts, a served copy of the actor and, for E6, a judge account with
credit; CORE-Bench's answer keys are deliberately kept out of the
repository. The evidence-loss paths cannot be reproduced in their
original form, only re-implemented against surviving fixtures, and that
distinction should be preserved in any future citation. And figures
originating at the Semester 3 boundary describe the tree at that tag, so
reproducing them requires the worktree checkout rather than the current
HEAD.

\section{Primary-artefact recovery and
availability}\label{primary-artefact-recovery-and-availability}

The 2 October checkout audit correctly recorded five named primary paths
as absent from that checkout. A bounded read-only search of the author's
Mac on 3 October recovered all five originals in a separate live
research checkout, plus an earlier pending-experiments snapshot. The
files were ignored by the local \texttt{outputs/*} rule. Their absence
from the frozen cloud checkout and inspected Git/archive corpus was
therefore a scope boundary, not evidence of deletion. The frozen
repository is unchanged; exact recovered copies are retained privately
and are not included in this public derivative.

\begin{longtable}[]{@{}
  >{\raggedright\arraybackslash\hspace{0pt}}p{(\linewidth - 4\tabcolsep) * \real{0.5000}}
  >{\raggedright\arraybackslash\hspace{0pt}}p{(\linewidth - 4\tabcolsep) * \real{0.5000}}@{}}
\caption{Recovered primary campaign files. Original-source preservation
metadata is retained privately; the public package does not include the
recovery export or raw operational chronology.}\tabularnewline
\toprule\noalign{}
\begin{minipage}[b]{\linewidth}\raggedright
Recovered primary source
\end{minipage} & \begin{minipage}[b]{\linewidth}\raggedright
Verified file-level evidence and remaining limit
\end{minipage} \\
\midrule\noalign{}
\endfirsthead
\toprule\noalign{}
\begin{minipage}[b]{\linewidth}\raggedright
Recovered primary source
\end{minipage} & \begin{minipage}[b]{\linewidth}\raggedright
Verified file-level evidence and remaining limit
\end{minipage} \\
\midrule\noalign{}
\endhead
\bottomrule\noalign{}
\endlastfoot
\longtok{agentdojo\_\allowbreak{}native\_\allowbreak{}receipt.json} & 2,228 bytes; 97 completed
benign episodes, 88 utility passes; native v1.2.2, security evaluation
false, no held-out claim \\
\longtok{labbench\_\allowbreak{}final\_\allowbreak{}diagnostics.json} & 3,317 bytes; 1,967
evaluated, 450 correct, 1,517 incorrect; all eight category totals and
diagnostics retained \\
\longtok{training\_\allowbreak{}identity\_\allowbreak{}verified.json} & 815 bytes; run
\longtok{withheld-id}, recorded dataset and file hashes;
identity/configuration verification, not content decontamination \\
\longtok{finish/Pending\_Experiments.md} & 14,027 bytes; September 19
terminal ledger updated through September 21; historical readiness, not
later execution proof \\
\longtok{finish/e14\_\allowbreak{}terminal\_\allowbreak{}note\_\allowbreak{}2026-09-19.json} & 1,564 bytes;
records two truncated-judge exclusions and terminal closure; contains an
internal denominator error discussed below \\
\end{longtable}

The additional September 12 \texttt{Pending\_Experiments.md} is a
distinct 1,535-byte snapshot, not a duplicate of the finish ledger. All
six UTF-8 bodies total 23,486 bytes; independently recomputed byte
lengths and SHA-256 values match the authenticated source export and the
materialized copies. Source filesystem timestamps and existence come
from the Mac export; no historical Git commit or pre-execution
attestation is invented. These are recovered files, not newly generated
experiment receipts.

Two source interpretation corrections follow. The E14 terminal note
writes \texttt{2271/4428\ =\ 0.5131043831902395}, but that decimal
equals 2,271/4,426. The unchanged thesis keeps the official judged-row
rate of 51.31\% separate from the all-disposition rate of 51.29\%; the
original note is preserved with its inconsistency. The finish ledger
explicitly says the E4 rerun had not launched by its September 21
update, so it cannot substantiate the later detailed zero-reward rerun
account. That account retains its separate source and raw-payload
limitations in Appendix D.

A separate read-only payload check on the Mac verified all 29 explicit
reference records covering 27 distinct file paths. For E8, 1,979
declared archive members match their retained local copies, including
1,967 unique native-grade records; independent retallying of the
exported per-task records gives 450 correct, 1,517 incorrect, 1,086
length finishes, 881 stop finishes, 1,960 sure and seven unsure
responses, and zero parser-error flags. The 1,259 unparsed-answer total
remains supported by the recovered diagnostic rather than re-counted
from the compact grade export. For E10, all 525 declared archive-member
bindings match and the selected 97 distinct native traces each match
their receipt hash; the exported trace records independently sum to 88
utility passes and zero handled errors. Historical partial traces are
not pooled into that selection.

The Mac also inspected the retained E14 disposition and judge files,
reporting 4,428 distinct dispositions, 4,426 accepted and 2,271 correct.
The two excluded rows, 1195 and 2044, have recorded skip reasons and no
Equivalence Judgement heading in the retained judge text. The compact
export carries those counts, the two exclusion records and source
hashes; it does not include all judge text or establish independent
mathematical correctness. The source-native score distinguishes the
4,426 judged denominator from the 4,428-disposition diagnostic,
resolving the terminal note's arithmetic conflict without changing its
bytes. No scorer, judge or model was rerun. The training-identity file
also matches its archived copy, but its original configuration, metadata
and output-log payload paths and dirty source diff remain unresolved.

Recovery closes the five named file-availability gaps and the bounded
stored-grade/hash checks support the indicated counts. These checks do
not reproduce model execution, regrade answers or establish
decontamination. The original checkout audit remains a dated record in
\longtok{campaign\_\allowbreak{}evidence\_\allowbreak{}availability.json}; the new source export is
\texttt{withheld\ primary-source\ recovery\ record}, the companion
\texttt{withheld\ original-source\ preservation\ record} binds every
recovered copy, and
\texttt{withheld\ stored-payload\ validation\ record} records the
bounded payload checks. The separately lost execution-source incident in
§B.8 remains unresolved by this recovery.

\section{Fresh-Sampling Diagnostic: Frozen Contract and
Audit}\label{fresh-sampling-diagnostic-frozen-contract-and-audit}

The new evidence in §7.4 is run
\texttt{V2\ frozen-policy\ sampling\ study}, a sampling-only experiment
with no weight updates. Its protocol, code, tests, expected cohort and
parser were hash-bound before allocation or generation at 18:15:03 UTC
on 2 October 2026. The freeze covers 26 files. Prompt ordering was
deterministic and independent of scores, using seed 20261002; the V2
smoke and main cohorts are disjoint from each other and from the failed
V1 smoke, and the request-seed namespace is disjoint from V1. The main
cohort is 512 prompts and the independent smoke is 32 prompts, each with
eight initial plus 24 fresh draws.

\textbf{Model and decoding.} Model and tokenizer are
\texttt{Qwen/Qwen3-8B}, both at revision
\longtok{b968826d9c46dd6066d109eabc6255188de91218}. The dataset is
\texttt{openai/gsm8k}, configuration \texttt{main}, \textbf{train}
split, revision \longtok{740312add88f781978c0658806c59bc2815b9866} (7,473
rows). The model used BF16 with thinking disabled, temperature 0.7,
top-p 0.8, top-k 20, min-p 0, a 1,024-generated-token cap and
4,096-token context. Each request had n = 1 and its own seed; batching
grouped four prompts. The recorded stack used vLLM 0.11.0, PyTorch
2.8.0, Transformers 4.57.1, eager execution and disabled prefix caching
on one NVIDIA A100-SXM4-80GB. The supervised process completed in
approximately 26.6 minutes, including loading and both cohorts.
Hardware/kernel changes need not reproduce samples bit-for-bit.

\textbf{Parser and review contract.} The V2 parser accepts specified
explicit final-answer forms, normalises exact rational numbers and fails
closed on malformed, conflicting or invalid claims; it does not search
arbitrary prose for a number. Gold parsing passed offline on all 7,473
training rows before generation. This is a stricter contract than the
existing GSM8K test evaluator and must not be retrospectively
substituted for it. Automated validation checked the frozen source and
manifest bindings, cohort identity, completeness, ordered draw indices,
unique seeds, reward recomputation and quality counts. A separate
set-based computation reproduced 461 eligible groups and 33 events. The
audit covered the first reward-changing response plus all eight initial
responses for every event. Two assistant reviewers handled disjoint
event partitions; this does not mean that each event received two
independent reviews. The coordinator's additional arithmetic checks
covered both recoveries and the ambiguous bread-rise case. The
implementation field \longtok{manual\_\allowbreak{}review\_\allowbreak{}complete} records
completion of this assistant review, not human validation; residual
model-review error is possible.

\textbf{Descriptive checks, not controller validation.} Among the
original 461 homogeneous groups, the combined-16 prefix became mixed for
15 (3.25\%); the fresh-16-only comparison was mixed for 24 (5.21\%). The
all-32-parseable-and-uncapped subgroup was 12/439, selected on later
outcomes and therefore not an unbiased replacement primary. Fixed
Beta(1,1) and Beta(1/2,1/2) posterior-predictive assumptions assigned
rescue probabilities of 72.73\% and 49.41\%, respectively, versus
observed strict-reward mixedness of 7.16\%. Their descriptive Brier
scores were 0.4964 and 0.2450, against 0.0716 for plug-in zero. These
uncalibrated fixed-prior assumptions substantially overpredict this
endpoint; the exercise does not fit or validate a controller and does
not vindicate certainty of zero rescue. The post-hoc later-change check
in the 20 parser-first groups is exploratory and remains separate from
the predeclared first-event numeric subset.

\textbf{Receipt locations and hash boundary.} Raw ordered completions
are in the run directory's \texttt{study.jsonl} and
\texttt{smoke.jsonl}; the per-prompt analysis is
\texttt{analysis\_final/per\_prompt.csv}. The final summary, validation
and completed audit preserve the unchanged raw primary endpoint
alongside the conservative numeric subset. The main \texttt{study.jsonl}
has SHA-256
\longtok{859ee12d9c146a3e9db65882b0a305d123bd385c573fa6ac4692b25220ac422f}.
The summary and validation list the freeze, manifest and source hashes,
distinguishing the manifest's canonical-content binding from its
file-byte checksum. Hashes bind the receipts to the trusted
pre-execution freeze; they are not cryptographic attestation of GPU
execution. Token counts and identities were cross-checked, not
independently decoded. The failed V1 smoke is retained as engineering
evidence and is never pooled with V2. These new sampling receipts are
separate from the original historical files subsequently recovered in
§B.11.

\section{C1 scientific source
availability}\label{c1-scientific-source-availability}

The C1 analysis preserves the frozen 11/64 endpoint, the separate three
conservative numerical witnesses, all sixty-four case classifications
and all eleven event records. The included case ledgers and public
scientific reference bindings retain disagreements and hash-level
scientific source descriptions. This public derivative excludes runtime
allocation, session, deletion and raw-transfer records, account resource
limits and operational provenance. No live runtime or GPU execution is
attested by publication of this document.

The historical thesis and evidence archives remain unchanged outside
this public package. The public derivative has its own editable
Markdown, LaTeX, builder and selected scientific ledgers. Appendix G is
the complete scientific case index; Appendix I and the public
availability manifest explain omitted source categories.



\chapter{Notation, Symbols and
Abbreviations}\label{notation-symbols-and-abbreviations}

\section{Scope and convention}\label{scope-and-convention}

This appendix records the symbols and abbreviations used across this
report and the papers behind it, as the sources actually use them rather
than as a house style would prefer. Where a symbol carries more than one
meaning across documents, both are listed and the collision flagged,
because a reader carrying one definition into another chapter will
misread a figure.

\section{Symbols}\label{symbols}

\begin{longtable}[]{@{}
  >{\raggedright\arraybackslash\hspace{0pt}}p{(\linewidth - 6\tabcolsep) * \real{0.3333}}
  >{\raggedright\arraybackslash\hspace{0pt}}p{(\linewidth - 6\tabcolsep) * \real{0.3333}}
  >{\raggedright\arraybackslash\hspace{0pt}}p{(\linewidth - 6\tabcolsep) * \real{0.3333}}@{}}
\caption{Notation: symbols used in this report, their meaning, and any
divergence in how the underlying sources use them.}\tabularnewline
\toprule\noalign{}
\begin{minipage}[b]{\linewidth}\raggedright
Symbol
\end{minipage} & \begin{minipage}[b]{\linewidth}\raggedright
Meaning
\end{minipage} & \begin{minipage}[b]{\linewidth}\raggedright
As used in the sources
\end{minipage} \\
\midrule\noalign{}
\endfirsthead
\toprule\noalign{}
\begin{minipage}[b]{\linewidth}\raggedright
Symbol
\end{minipage} & \begin{minipage}[b]{\linewidth}\raggedright
Meaning
\end{minipage} & \begin{minipage}[b]{\linewidth}\raggedright
As used in the sources
\end{minipage} \\
\midrule\noalign{}
\endhead
\bottomrule\noalign{}
\endlastfoot
\(G\) & Group size: completions sampled per prompt & Headline runs
\(G=8\); sweeps \(G \in \{2,4,8,16\}\) \\
\(K\) & Rollouts per group in the ZVF formalisation;
\textbf{separately}, successes within a group & Elsewhere
\(K = \sum_i R_i \sim \mathrm{Binomial}(G,p)\) \\
\(B_t\), \(\mathcal{P}_t\) & Batch at optimizer step \(t\); prompts at
step \(t\) & Competing denominators of ZVF \\
\(p\), \(\hat{p}\) & Latent success probability; sample rate \(K/G\) &
\(\mathrm{Bernoulli}(p)\) model; \(\hat{p}\) centres the advantage \\
\(r_i\), \(R_i\) & Scalar reward of completion \(i\); binary reward &
Verifiable reward; derivations appendix \\
\(h_G(p)\) & \(p^G + (1-p)^G\) & Bernoulli null for ZVF \\
\(\mathrm{ZVF}_t\), \(\mathrm{GU}_t\) & Thresholded Zero-Variance
Fraction; Gradient Utilization \(1-\mathrm{ZVF}_t\) & Sample variance ≤
10⁻⁶; \(\mathrm{ZVF}_0\) explicitly denotes exact ties (§4.3) \\
\(A_i\) & Advantage of completion \(i\) & Unnormalised \(R_i-\hat{p}\);
normalised \((r_i-\mu_g)/(\sigma_g+\varepsilon)\); clipped Dr.GRPO \\
\(\mu_g\), \(\sigma_g\) & Group mean and standard deviation of reward &
\(\mu_g = \frac{1}{G}\sum_j r_j\); \(\sigma_g\) in population form
(\(1/G\)) throughout, as every GRPO-style runner computes it (§4.2) \\
\(\pi_\theta\), \(\pi_{\theta_{\mathrm{old}}}\), \(\pi_{\mathrm{ref}}\)
& Trained, sampling and frozen reference policy; per-token importance
ratio \(\rho_{i,t}\) & Ratio against the sampling policy; KL penalty
against the reference \\
\(\mathcal{L}(\theta)\) & Policy loss & Original GRPO: clipped surrogate
plus KL penalty. Inspected custom-loss Tinker loops use REINFORCE with a
group-mean baseline; KL/reference status is row-specific, with TRL
defaults/config evidence and managed unknowns retained (Table 4.2) \\
\(\beta\) & KL coefficient; \textbf{separately}, a regression slope &
TRL and veRL 0.04; OpenRLHF 0.02; Tinker managed; also a slope in
M1/M2 \\
\(\varepsilon\) & Three distinct meanings --- see §C.3 & Clip range 0.2;
ZVF tolerance \(10^{-6}\); Adam term \(10^{-8}\) \\
\(\gamma\), \(\lambda\), \(T\) & Discount factor; GAE parameter;
sampling temperature & \(0.99\); \(0.95\); canonical \(0.8\) \\
\(\mathrm{SI}\), \(\mathrm{PTD}\) & Stability index; peak-to-tail drift
& Reported in place of direct KL \\
PCD & Pairwise-contrast density &
\(\mathbb{E}[\hat{p}(1-\hat{p})] = \frac{G-1}{G}\,p(1-p)\) \\
pass@1 & Pass rate under the suite's native scorer & E11
\(129/312 = 41.35\,\%\); E1 \(2/731 = 0.274\,\%\) \\
\(r\), \(\rho\) & Correlation coefficients, Pearson and rank &
\(r = -0.769\), \(p = 0.0008\); \(\rho \approx 0.65\) \\
\end{longtable}

\section{Four collisions worth
flagging}\label{four-collisions-worth-flagging}

The symbol \(\varepsilon\) carries three unrelated meanings. In the
hyperparameter mapping it is the PPO/GRPO clip range, default 0.2, swept
over \(\{0.05, 0.1, 0.2, 0.3, 0.5\}\). In the ZVF definition it is a
numerical tolerance, \(10^{-6}\), deciding whether a group's sample
variance counts as degenerate. In the optimiser configuration it is the
Adam epsilon, \(10^{-8}\). It also appears additively in the advantage
denominator, \(\sigma_g + \varepsilon\).

\(K\) is used two ways: in the formal treatment of ZVF it denotes the
rollouts in a group, a synonym for the group size elsewhere called
\(G\), whereas in the derivations appendix \(G\) is the group size and
\(K\) is the number of successes --- a design parameter in the one case,
an outcome in the other. The letter \(r\) is similarly the reward, the
LoRA rank (swept over \(\{4,8,16,32,64\}\)) and the Pearson correlation.

Finally, ZVF is expanded two ways in the sources. The formal definition
and the abstract both read ``Zero-Variance Fraction'', the expansion
used here; one figure caption reads ``Zero-Value-Filtered''. The caption
variant is noted only so that a reader meeting it is not misled.

\section{Abbreviations}\label{abbreviations}

\begin{longtable}[]{@{}
  >{\raggedright\arraybackslash\hspace{0pt}}p{(\linewidth - 4\tabcolsep) * \real{0.5000}}
  >{\raggedright\arraybackslash\hspace{0pt}}p{(\linewidth - 4\tabcolsep) * \real{0.5000}}@{}}
\caption{Abbreviations and their expansions, with names that no source
expands recorded as identities rather than glossed.}\tabularnewline
\toprule\noalign{}
\begin{minipage}[b]{\linewidth}\raggedright
Abbreviation
\end{minipage} & \begin{minipage}[b]{\linewidth}\raggedright
Expansion or identity
\end{minipage} \\
\midrule\noalign{}
\endfirsthead
\toprule\noalign{}
\begin{minipage}[b]{\linewidth}\raggedright
Abbreviation
\end{minipage} & \begin{minipage}[b]{\linewidth}\raggedright
Expansion or identity
\end{minipage} \\
\midrule\noalign{}
\endhead
\bottomrule\noalign{}
\endlastfoot
GRPO & Group Relative Policy Optimization --- critic-free,
group-relative \\
PPO & Proximal Policy Optimization \\
DPO & Direct Preference Optimization \\
RLHF & Reinforcement learning from human feedback \\
IS Loss & Importance-sampling loss, the distillation objective \\
SFT & Supervised fine-tuning \\
LoRA & Low-rank adaptation of large language models \\
KL & Kullback--Leibler divergence; never spelled out in the sources \\
CoT & Chain-of-thought; spelled out in prose, written ``GSM8K CoT'' for
the variant \\
ZVF, GU & Zero-Variance Fraction; Gradient Utilization \\
PCD, SI, PTD & Pairwise-contrast density; stability index; peak-to-tail
drift \\
Dr.GRPO & GRPO with the per-response length normalisation and the
division by the group standard deviation removed \\
GAE & Generalised advantage estimation; \(\lambda = 0.95\) \\
MoE & Mixture-of-experts \\
BF16 & bfloat16, the serving precision for the LAB-Bench, AgentDojo and
Omni-MATH lanes \\
IQM & Interquartile mean, reported via \texttt{rliable} \\
CI & Confidence interval; 95 \%, 10,000 bootstrap resamples \\
W\&B & Weights \& Biases, the experiment-tracking service \\
MLE-bench & The suite forming lane E9's original contract; MLE never
appears standalone \\
MATH-500 & The mathematics evaluation suite \\
GSM8K & The grade-school arithmetic suite; 7,473 training and 1,319 test
rows \\
HumanEval & The code-generation suite used for the code arm \\
ArenaHard, NoRobots, OpenThoughts3 & Dataset names used as-is, each at a
pinned revision \\
MiniWoB & Browser micro-task suite used for the proxy smoke test \\
SWE-bench (Pro, Multilingual) & The repository-repair suite forming lane
E1 \\
LAB-Bench, AgentDojo, VerilogEval, Omni-MATH, BALROG, WebArena,
APEX-Agents, CORE-Bench, Tau3, MLDevBench, SDAB, BankerToolBench,
BinaryAudit, LifeSciBench, AgentHarm, WebBench, AppBench, OpenReward
Games, FrontierSWE, FrontierMath & Suite names used as names, with no
expansion in any source; see §C.5 \\
Tinker & The remote training and sampling API used to train the
adapter \\
TRL & HuggingFace's LLM post-training library, in both framework
rosters \\
SkyRL, verl (veRL), OpenRLHF, Atropos & Launchers in the cross-launcher
roster; only Tinker and TRL produced completed runs, the veRL and
OpenRLHF entries being dry-run placeholders \\
\end{longtable}

\section{The E1--E14 lane identifiers, scopes and
statuses}\label{the-e1e14-lane-identifiers-scopes-and-statuses}

Each lane's short code names one original benchmark contract and is not
a rank or a score. Receipt paths in Table C.3 are historical
identifiers, not guarantees that every referenced payload is present in
the frozen checkout. The six originals covering five previously missing
basenames are now preserved in the withheld source archive; Appendix B
§B.11 maps their historical paths and hashes and records supplementary
stored-payload checks. This recovery does not recreate missing training
logs or E4 raw rerun evidence. Several lanes hold two states at once ---
an original contract and a replacement scope --- so the code alone does
not say which scope a figure comes from. Table C.3 is the single record
of lane scope, headline figure, terminal state and receipt; the chapters
point here rather than restating it. ``orig'' is the original contract
and ``repl'' a declared replacement scope. Terminal states are from the
recovered \texttt{withheld\ historical\ campaign\ chronology}, whose
latest annotations run through 21 September; the E1, E2, E5, E6, E9 and
E13 replacement scopes were completed on 2026-09-27, after that ledger,
so their states are assigned in this report from each lane's
\texttt{result.json} (§8.4); the result files themselves carry no state
field. Receipt paths are relative to \texttt{outputs/}.

\begingroup\scriptsize\setlength{\tabcolsep}{3pt}
\begin{longtable}[]{@{}
  >{\raggedright\arraybackslash\hspace{0pt}}p{(\linewidth - 12\tabcolsep) * \real{0.1667}}
  >{\raggedright\arraybackslash\hspace{0pt}}p{(\linewidth - 12\tabcolsep) * \real{0.1667}}
  >{\raggedright\arraybackslash\hspace{0pt}}p{(\linewidth - 12\tabcolsep) * \real{0.1667}}
  >{\raggedright\arraybackslash\hspace{0pt}}p{(\linewidth - 12\tabcolsep) * \real{0.1667}}
  >{\raggedright\arraybackslash\hspace{0pt}}p{(\linewidth - 12\tabcolsep) * \real{0.1667}}
  >{\raggedright\arraybackslash\hspace{0pt}}p{(\linewidth - 12\tabcolsep) * \real{0.1667}}@{}}
\caption{The E1--E14 lanes: original contract, replacement scope,
headline figure, terminal state and receipt. E14 retains 51.29\% over
all 4,428 dispositions and 51.31\% over 4,426 judged rows; the recovered
terminal note's mixed denominator label is corrected in Table E.1. The
recovered E4 pending-state ledger predates its later reported 100-trial
rerun (§8.5).}\tabularnewline
\toprule\noalign{}
\begin{minipage}[b]{\linewidth}\raggedright
Lane
\end{minipage} & \begin{minipage}[b]{\linewidth}\raggedright
Original contract
\end{minipage} & \begin{minipage}[b]{\linewidth}\raggedright
Replacement scope
\end{minipage} & \begin{minipage}[b]{\linewidth}\raggedright
Headline figure (scope)
\end{minipage} & \begin{minipage}[b]{\linewidth}\raggedright
Terminal state
\end{minipage} & \begin{minipage}[b]{\linewidth}\raggedright
Receipt
\end{minipage} \\
\midrule\noalign{}
\endfirsthead
\toprule\noalign{}
\begin{minipage}[b]{\linewidth}\raggedright
Lane
\end{minipage} & \begin{minipage}[b]{\linewidth}\raggedright
Original contract
\end{minipage} & \begin{minipage}[b]{\linewidth}\raggedright
Replacement scope
\end{minipage} & \begin{minipage}[b]{\linewidth}\raggedright
Headline figure (scope)
\end{minipage} & \begin{minipage}[b]{\linewidth}\raggedright
Terminal state
\end{minipage} & \begin{minipage}[b]{\linewidth}\raggedright
Receipt
\end{minipage} \\
\midrule\noalign{}
\endhead
\bottomrule\noalign{}
\endlastfoot
E1 & SWE-bench Pro & SWE-bench Multilingual & 2/731 = 0.274\% pass@1
(orig, full suite); 1/190 resolved on the tasks not attempted before,
4/110 earlier, not pooled (repl) & orig: complete, no ledger label;
repl: \texttt{COMPLETE} on coverage (all 300 attempted, two contracts) &
\longtok{modal\_\allowbreak{}e1\_\allowbreak{}e14/2026-08-16/e1\_\allowbreak{}swe\_\allowbreak{}bench\_\allowbreak{}pro\_\allowbreak{}full/seed1818/receipt.json};
\longtok{finish\_\allowbreak{}pending\_\allowbreak{}2026-09-27/E1/result.json} \\
E2 & FrontierSWE & CORE-Bench & 1/17 tasks, replay 0.8628 (orig,
partial); 27/45 = 0.600 task accuracy (repl, complete) & orig: partial;
repl: \texttt{COMPLETE} &
\longtok{e2\_\allowbreak{}frontier\_\allowbreak{}swe/e2\_\allowbreak{}terminal\_\allowbreak{}attempt\_\allowbreak{}receipt\_\allowbreak{}2026-08-22.json};
\longtok{finish\_\allowbreak{}pending\_\allowbreak{}2026-09-27/E2/result.json} \\
E3 & SDAB (private bundle) & --- & none & \longtok{CLOSED\_EXTERNAL} &
\texttt{withheld\ operational\ closure\ record} \\
E4 & BankerToolBench & --- & 1/100 tasks, recovery metric 0.3115 (orig,
partial) & \longtok{CLOSED\_PARTIAL} &
\longtok{modal\_\allowbreak{}e1\_\allowbreak{}e14/2026-08-16/e4\_\allowbreak{}recovery\_\allowbreak{}pass16\_\allowbreak{}receipt.json} \\
E5 & APEX-Agents & Tau3 & 7/480 natively scored (orig, partial); 8/97 =
0.082 pass\^{}1 (repl, complete) & orig: partial; repl:
\texttt{COMPLETE} &
\longtok{e5\_\allowbreak{}apex\_\allowbreak{}agents/e5\_\allowbreak{}exact\_\allowbreak{}sequential\_\allowbreak{}prefix\_\allowbreak{}aggregate\_\allowbreak{}2026-08-22.json};
\longtok{finish\_\allowbreak{}pending\_\allowbreak{}2026-09-27/E5/result.json} \\
E6 & WebBench & WebArena & 90/812 = 0.111 task success, lower bound; 108
judge-dependent tasks ungraded (repl, complete in coverage) & orig: no
original-contract result (ledger N/A); repl: \texttt{COMPLETE} on
coverage; 108 ungraded &
\longtok{finish\_\allowbreak{}pending\_\allowbreak{}2026-09-27/E6/result.json} \\
E7 & BinaryAudit (private payload) & --- & 1/46 attempted, verifier
reward 0.0, no grade (orig) & \longtok{CLOSED\_EXTERNAL} &
\longtok{e7\_binaryaudit/2026-08-22\_\allowbreak{}e7\_\allowbreak{}paid\_\allowbreak{}attempt\_\allowbreak{}receipt.json} \\
E8 & LifeSciBench (private package) & LAB-Bench public split & 450/1,967
= 0.2288 (repl, complete; end-to-end harness score) & repl:
\texttt{COMPLETE}; orig: \longtok{CLOSED\_EXTERNAL} &
\longtok{public\_\allowbreak{}portfolio\_\allowbreak{}2026-09-05/labbench\_\allowbreak{}final\_\allowbreak{}diagnostics.json}
(recovered);
\longtok{public\_\allowbreak{}portfolio\_\allowbreak{}2026-09-05/labbench\_\allowbreak{}native\_\allowbreak{}receipt.json}
(hash-checked on Mac) \\
E9 & MLE-bench & MLDevBench & 40/75 competitions graded, suite score
null (orig, partial); 10/34 = 0.294 task success (repl, complete) &
orig: partial; repl: \texttt{COMPLETE} &
\longtok{e9\_\allowbreak{}mle\_\allowbreak{}bench/modal\_streaming/};
\longtok{finish\_\allowbreak{}pending\_\allowbreak{}2026-09-27/E9/result.json} \\
E10 & AgentHarm (private tasks) & AgentDojo benign utility & 88/97 =
0.9072 utility, 97/97 completed (repl, complete) & repl:
\texttt{COMPLETE}; orig: \longtok{CLOSED\_EXTERNAL} &
\longtok{public\_\allowbreak{}portfolio\_\allowbreak{}2026-09-05/agentdojo\_\allowbreak{}native\_\allowbreak{}receipt.json} \\
E11 & VerilogEval & none (two native framings) & 129/312 = 41.35\%
pass@1 (orig, complete; end-to-end harness score) & \texttt{COMPLETE} &
\longtok{modal\_\allowbreak{}e1\_\allowbreak{}e14/2026-08-16/e11\_\allowbreak{}full\_\allowbreak{}receipt.json} \\
E12 & AppBench (deployment) & --- & none & \longtok{CLOSED\_EXTERNAL} &
\texttt{withheld\ operational\ closure\ record} \\
E13 & OpenReward Games (held-out games) & BALROG & 26.1\% BALROG
progression, 255/255 episodes (repl, complete) & repl:
\texttt{COMPLETE}; orig: externally blocked, no closure record &
\longtok{finish\_\allowbreak{}pending\_\allowbreak{}2026-09-27/E13/result.json} \\
E14 & FrontierMath (hosted) & Omni-MATH & 2,271 correct of 4,428
dispositions (repl, terminal-complete; §8.2) & repl:
\longtok{COMPLETE\_\allowbreak{}TERMINAL\_\allowbreak{}NOTE}; orig: \longtok{CLOSED\_EXTERNAL} &
\longtok{PES\_\allowbreak{}Phase2\_\allowbreak{}Review\_\allowbreak{}2026-09-12/finish/e14\_\allowbreak{}terminal\_\allowbreak{}note\_\allowbreak{}2026-09-19.json}
(recovered);
\longtok{public\_\allowbreak{}portfolio\_\allowbreak{}2026-09-05/native\_\allowbreak{}finish\_\allowbreak{}v5/e14\_\allowbreak{}official\_\allowbreak{}collector/official\_result01/native\_score.json}
(Mac validation) \\
\end{longtable}
\endgroup

Counted from this table: three scopes are strictly complete (E8 and E10
replacement scopes, E11 original contract), E14 is terminal-complete,
and six original contracts are \texttt{CLOSED\_EXTERNAL} with written
closure records (E3, E7, E8-original, E10-original, E12, E14-original).
E13's original contract is also blocked on external access but carries
no closure record, because the lane continues through its replacement
scope. Original-contract full-coverage scores exist for E1 and E11.



\chapter{Campaign Lane Detail}\label{campaign-lane-detail}

This appendix holds the lane-by-lane record behind Chapter 8: provenance
of the training run, per-lane serving and scoping notes, the engineering
history behind individual figures, the per-lane narratives of the six
replacement scopes completed on 27 September 2026, the full E4 rerun
audit, and the per-lane caveats of the small-scale reruns. The
descriptive figures agree with Chapter 8, and the two reporting rules of
Chapter 8 (no pooling of original-contract and replacement-scope
results, no unsupported precision) apply throughout. The E11 uncertainty
calculation below is a new, reproducible analysis of surviving paired
verdicts. Six original files covering five previously missing basenames
are now preserved in the withheld source archive (Appendix B §B.11).
Receipt-level fields and the supplementary stored-payload checks are
distinguished below from historical claims whose raw trajectories,
training logs or supporting artifacts remain unvalidated; recovery and
count/hash checks are not new execution or regrading of the experiments.

\section{Actor identity, inference conditions and scoping
notes}\label{actor-identity-inference-conditions-and-scoping-notes}

The design consequence of evaluating one actor for seed variance is
specific: every suite number reported in Chapter 8 comes from one actor,
and where a lane sampled more than one completion per task it did so to
satisfy the suite's own protocol, not to estimate variance.

Inference conditions were reported per lane rather than imposed
globally, and the lanes differ widely. The recovered E10 native receipt
declares a temperature of 0, a 32,768-token context, 4,096 output tokens
per request, 3,600 seconds of wall time. It also declares
\longtok{max\_\allowbreak{}generation\_\allowbreak{}requests} = 1,967, a ceiling on generation
requests rather than a task count: the lane has 97 episodes. The earlier
report separately gives 414 native requests (405 in the main run and 9
in the continuation) and says the collector does not admit shorter
native conditions. Those execution counts and collector claims are not
fields of the recovered receipt, and the cited terminal-readiness source
remains unrecovered. The E11 run was sampled at temperature 0.2, top-p
0.95, 1,024 maximum output tokens and seed 1816, one sample per problem
across all 312 prompts. The point of recording them per lane is that the
sampling regime is part of the result, not a detail beneath it.

Two scoping notes complete the picture of what the actor has actually
been measured on. The E1 original-contract evaluation was a
one-checkpoint evaluation whose 731 generations were produced by a mixed
backend within a single run --- 476 Tinker generations and 255 Modal
vLLM generations --- and that backend mix is recorded in the original
receipt. Separately, two small probes exist at other parameter settings
but are not suite results and are not reported as such: a base-model
smoke run of the E11 harness measured zero scored tasks, and a four-task
pilot on the step-40 checkpoint (the last training step, saved
separately from \texttt{stepfinal}) returned four passes out of four
with zero extraction failures while carrying an unresolved outstanding
blocker --- the decontamination receipt identifier for that pilot does
not exist, and the split-manifest validator still fails on it.

\section{Training-run provenance}\label{training-run-provenance}

The earlier report attributed the adapter configuration to the local
Weights \& Biases directory for run \texttt{withheld-id}
(\longtok{withheld-training-log-directory}), the run receipt
\texttt{withheld\ actor-training\ receipt}, and Hugging Face checkpoint
branches. The original
\longtok{outputs/public\_\allowbreak{}portfolio\_\allowbreak{}2026-09-05/training\_\allowbreak{}identity\_\allowbreak{}verified.json}
is now recovered, with SHA-256
\longtok{e340d5fe41b7b028dd78ea9f01f6da787c7a8cb968f5a0ac1b103a92b546ee86}.
It records run \texttt{withheld-id}, dataset revision
\longtok{bd1f0db6056091e99d5e28e3a7b3d05d77b21876ad65e5b31bf6f1b470e73e5c},
the two training-suite IDs and hashes for \texttt{config.yaml},
\texttt{wandb-metadata.json} and \texttt{output.log}. It explicitly
labels this
\longtok{TRAINING\_\allowbreak{}CONFIG\_\allowbreak{}VERIFIED\_\allowbreak{}NOT\_\allowbreak{}CONTENT\_\allowbreak{}DECONTAMINATION} and
states that W\&B captured the dirty checkout's base commit but did not
upload the dirty source diff. The underlying W\&B files and local run
receipt have not been recovered in this package, so a new byte-level
match against those three recorded hashes is not established here. The
surviving E11 receipt verifies the evaluated base/adapter identity, but
does not reconstruct the missing training logs. Table 8.1 retains
detailed training fields as earlier reported except where the recovered
identity record directly corroborates them.

Four things remain unavailable or were never created, and they are
listed so that no reader fills them in:

\begin{itemize}
\tightlist
\item
  \textbf{The exact training code.} The run records git commit
  \texttt{8fdade2f} with uncommitted changes that were never saved. The
  \texttt{pavlov\_portfolio} preset and its dataset builder do not exist
  at that commit; they first appear in commit \texttt{72b9336d3} on
  2026-08-14. The reward and data-pool code in Table 8.1 is therefore
  taken from the current tree, whose dataset hash does not match the
  recorded \texttt{bd1f0db6…}. It is the best reconstruction, not proof.
  The loss function and training loop are identical at
  \texttt{8fdade2f}.
\item
  \textbf{Which 80 training prompts were drawn.} The draw used
  \texttt{random.Random(809)} over the 512-row pool, but the pool
  construction that ran is the unrecoverable code above.
\item
  \textbf{The original sampler execution.} The public package does not
  include execution-specific operational access paths and cannot replay
  that original route.
\item
  \textbf{Any held-out training-time evaluation.} None was run
  (\texttt{evaluate\_heldout} false).
\end{itemize}

The run also departs from
\longtok{zvf-program/flagship/PAVLOV\_\allowbreak{}EXPERIMENT\_\allowbreak{}PROTOCOL\_\allowbreak{}2026-08-09.md}.
The protocol plans a learning-rate search and a twelve-suite curriculum
with 1,200/128 prompt and response tokens. This run used two suites, 40
steps, 1,536/384 tokens and a learning rate of 1 × 10⁻⁵, and the
protocol does not mention it; only the budget matches.

\section{Serving conditions: format limits and the scope of the
stop-sequence
mechanism}\label{serving-conditions-format-limits-and-the-scope-of-the-stop-sequence-mechanism}

Table 8.2 reports substantial extraction and length failures; E11's
aggregate extraction count is verifiable in the surviving receipt and
E8's aggregate diagnostics are now corroborated by its recovered
original. The supplementary validation independently corroborates E8's
stored grade and finish-reason counts, but the 1,259 parse-none count
remains diagnostic-only. These checks do not validate every serving
field, response body or cited finish-era source. The reported absence of
role-token run-on in E1's single-turn Tinker generations does not
establish a necessary condition for the mechanism. The E4 bridge code
and recorded examples support a concrete failure mode, while the
corrected shim's nonzero outcome comes from one task against 100 bridge
trials. Task selection, serving path and caps differ, so the fraction of
the original aggregate zero caused by the defect remains unidentified.
For other bridge-served lanes exposure is possible, not demonstrated.
This matches the scope qualification in
\longtok{outputs/e1\_\allowbreak{}e14\_\allowbreak{}small\_\allowbreak{}scale\_\allowbreak{}2026-09-26/E4/BRIDGE\_\allowbreak{}STOP\_\allowbreak{}TOKEN\_\allowbreak{}FINDING.md}.

Three non-bridge failure attributions are corrected (Table E.1). E11's
150 extraction failures are attributed to the 1,024-token cap with
thinking on, not to answer format; E14's capped responses are correct at
34.5\% against 85.9\% for uncapped ones; and 593 of E1's 713 submitted
patches are corrupt for \texttt{git\ apply}. Lane detail is in §D.4 and
§D.5. The runners now record per-sample finish reasons and truncation
(E11, E14) and repair and record patch-apply status (E1), but no stored
score was regenerated.

\section{Strictly complete scopes: lane
detail}\label{strictly-complete-scopes-lane-detail}

The ledger lists E11 both among its ``replacement scopes'' and among its
original-contract results; Table 8.A lists it as original-public-subset.
It is one run and one result, not two.

\textbf{E8 --- LAB-Bench, public split.} The recovered
\longtok{outputs/public\_\allowbreak{}portfolio\_\allowbreak{}2026-09-05/labbench\_\allowbreak{}final\_\allowbreak{}diagnostics.json}
(SHA-256
\longtok{7d7d0bd47162534212f7526db548e1f5d42bbe16554bc2e29115f0d9cf5af941})
directly records the following per-category accuracy, correct counts and
coverage: TableQA 0.7869 (192/244); FigQA 0.5138 (93/181); ProtocolQA
0.2500 (27/108); SuppQA 0.1829 (15/82); LitQA2 0.1759 (35/199); DbQA
0.1308 (68/520); SeqQA 0.0333 (20/600); and CloningScenarios 0 (0/33).
No category is incomplete. Coverage per category is at or above 0.9878
everywhere, and the CloningScenarios coverage is exactly 1.0 with zero
correct. The receipt also records zero parser errors, 1,960 responses
marked ``sure'' against 7 marked ``unsure'', and a completion profile of
1,086 length-truncated finishes against 881 stop finishes. The
diagnostic's correct and incorrect fields partition all 1,967 rows, so
the headline is an end-to-end harness score over the full split. The
earlier account treats the 1,259 unparsed responses as incorrect and
supplies the conditional 450/708 sensitivity; that response-level
membership has not been independently regraded. The diagnostic binds a
native-receipt hash (\texttt{b08995d5…}) and a graded-responses hash
(\texttt{8a61def9…}). The supplementary Mac validation locates the
native receipt at the declared hash and checks 1,979 declared archive
members against both archive and local bytes. Independently tallying its
1,967 exported grade records reproduces 450 correct, 1,517 incorrect,
all category denominators and correct counts, 1,960 sure, 7 unsure,
1,086 length finishes, 881 stop finishes and zero parser errors. Those
exported records have no parsed-answer field, so they do not
independently establish the 1,259 parse-none count or its response-level
membership. No native scorer was rerun.

\textbf{E10 --- AgentDojo, benign utility scope.} The recovered native
receipt (SHA-256
\longtok{a40c90015200a389ce0767b14a78f9409330a0ae8615cb4822ef9b54656ebaa0})
records \longtok{missing\_\allowbreak{}evaluation\_\allowbreak{}ids} empty and
\longtok{native\_\allowbreak{}handled\_\allowbreak{}error\_\allowbreak{}count} zero. It carries run-manifest
hash
\longtok{00e4228f4b757dc32b69b74f05a88348151f13185b470c5f00e31fdaa830ef30}
and fixes the metric as \longtok{benign\_\allowbreak{}task\_\allowbreak{}utility}, benchmark
version v1.2.2, attack and defence null, \texttt{security\_evaluation}
false and \texttt{heldout\_claim} false. The earlier statement that all
three protocol fields were null is corrected (Table E.1). There is no
upstream benchmark commit field in the recovered receipt, so the earlier
attribution of commit \longtok{089ed468cf3ed0322acc66b0211f26d9d90dbf60}
to it is unsupported. The episode decomposition, banking 16, slack 21,
travel 20 and workspace 40, is absent from the native receipt but is
independently corroborated by the supplementary export's 97 selected
trace records. Their utility/error fields tally to 88 passes and zero
handled errors. All 97 evaluation IDs map uniquely to exact-hash
archived traces, selecting 94 from the original collection and 3 from
the explicit continuation; extra historical partial traces are not
pooled. The Mac validation also reports all 525 archive-member pins
matching. These are stored-trace checks, not re-execution of the native
evaluator or independent semantic grading. Its decontamination block
records \longtok{TRAINING\_\allowbreak{}INVENTORY\_\allowbreak{}ABSENT}; empty match lists without
the training inventory do not establish absence of overlap.

The following collector history is retained from the earlier report. The
supplementary validation locates the continuation terminal-evidence file
on the Mac and matches its declared hash, and the exported 94-plus-3
trace mapping supports the completed-episode accounting. It does not
reproduce the collector's 18 offline tests or export the full snapshot,
readiness and continuation-result contents, so the detailed repair and
test claims below remain historical. These sources are still absent from
the frozen repository. An earlier attempt reportedly terminated with 94
of 97 episodes and score null, missing workspace tasks 37 through 39.
The cause was a confirmed collector fault rather than a model failure:
verified task receipts recorded absolute \texttt{/cache/...} trace
paths, and the frozen summary recomputed those paths on the local host,
so every completed receipt failed comparison after extraction. The
repair was deliberately minimal --- the original archive and every
extracted receipt and trace byte are unchanged, and a narrowly scoped
adapter rewrites only the freshly computed
\longtok{task\_result.native\_\allowbreak{}trace\_\allowbreak{}path} to the verified canonical
remote path before the frozen summary compares it, leaving native task
code, environment, prompt, result parsing, utility values and
handled-error policy untouched. The revised collector passed 18 offline
tests, independently re-run in 0.507 seconds, and the continuation run
that closed episodes 37 through 39 was executed in the reported
continuation with the native wrapper hash
\longtok{11a8b61631390d9b20decd4263008e5a4c6f648a033d314384bac9e5d2456d77}.

\textbf{E11 --- VerilogEval.} The upstream NVlabs/verilog-eval harness
ran at revision \longtok{c498220d0a52248f8e3fdffe279075215bde2da6}, using
Icarus Verilog 12.0 at revision
\longtok{4fd5291632232fbe1ba49b2c26bb6b2bf1c6c9cf} and Verilator 5.050.
Pass verdicts were determined deterministically: for each pinned problem
the single \texttt{sample01} \texttt{sv-iv-test} log was read and marked
PASS only when it contained a full-line ``Mismatches: 0 in samples''
verdict, with missing or non-matching logs marked FAIL. The upstream
\texttt{sv-iv-analyze} summariser could not be used because it fails on
a missing optional langchain dependency; that failure is recorded in the
receipt and does not affect the aggregation, which reads the simulator
logs directly. An extraction failure writes an empty candidate file that
always fails.

\textbf{E11 design-cluster uncertainty.} The
\longtok{verifier.receipts.code-complete-iccad2023.log\_receipts} and
\longtok{verifier.receipts.spec-to-rtl.log\_receipts} mappings in the
surviving full receipt contain exactly the same 156 design IDs. Matching
by ID gives 54 designs passing both framings, 81 failing both, 13
passing code-completion only and 8 passing specification-to-RTL only.
The 312 prompt outcomes are therefore not 312 independent design
observations. The descriptive pass fraction is unchanged at (2 × 54 + 13
+ 8)/312 = 129/312. For uncertainty over design composition, resample
156 whole design pairs uniformly with replacement per replicate, retain
both observed outcomes, and calculate their summed passes divided by
312. With 200,000 replicates from Python
\texttt{random.Random(20261002)}, the 2.5th and 97.5th percentiles are
0.3429487 and 0.4871795, reported as {[}0.343, 0.487{]}. Quantiles use
linear interpolation at index (B − 1)p.~The standard-library script
\longtok{platform\_hybrid/experiments/scripts/e11\_\allowbreak{}design\_\allowbreak{}cluster\_\allowbreak{}bootstrap.py}
reproduces the calculation; the adjacent
\longtok{e11\_\allowbreak{}design\_\allowbreak{}cluster\_\allowbreak{}bootstrap.json} preserves method, seed,
receipt SHA-256, pair counts and all 156 paired verdicts, and the script
falls back to it when the ignored receipt is absent. Run it from the
repository root with
\texttt{python\ platform\_\allowbreak{}hybrid/\allowbreak{}experiments/\allowbreak{}scripts/\allowbreak{}e11\_\allowbreak{}\allowbreak{}design\_\allowbreak{}\allowbreak{}cluster\_\allowbreak{}\allowbreak{}bootstrap.py}.
The exact upper bound depends on drawing with \texttt{random.choices},
because the 97.5th percentile falls on the boundary between 151/312 and
152/312; drawing with \texttt{randrange} under the same seed gives
0.4839744. The script previously cited under
\longtok{outputs/PES\_\allowbreak{}Phase2\_\allowbreak{}Third\_\allowbreak{}Review\_\allowbreak{}2026-09-24/analysis/} was
never committed.

This interval assumes exchangeable underlying designs and conditions on
one checkpoint, one generation per framing and the observed protocol. It
does not estimate variation across training seeds, decoding seeds or
protocols, and does not establish generalization beyond comparable
designs. Conditional on the complete fixed benchmark and these recorded
outputs, 129/312 is simply the observed fraction. The parseable-only
129/162 remains descriptive: the receipt records 150 extraction failures
only in aggregate, not per-design extraction indicators, so it cannot
support a cluster-aware interval for that conditional rate.

\textbf{E11 failure attribution.} The 150 extraction failures are not
primarily answer-format failures (corrected; see Table E.1).
\longtok{run\_\allowbreak{}e11\_\allowbreak{}full} in
\longtok{zvf-program/flagship/modal\_\allowbreak{}e1\_\allowbreak{}e14.py} used
\texttt{max\_tokens} = 1,024 and did not pass \texttt{enable\_thinking},
so the model's template default, thinking on, applied. No per-sample
finish reason was recorded, but the receipt's 312,622 response tokens
over 312 samples give a mean of 1,002 of 1,024, so nearly every
generation reached the cap. On the 50 items of the 26 September
small-scale rerun, the receipt's verdicts give 23/50 and the same
trained adapter, served by vLLM with thinking off, temperature 0 and a
4,096-token cap, scores 33/50 with no extraction failures and every
response ending on a stop; the discordant pairs are 11 against 1 (exact
McNemar p ≈ 0.006). Engine and temperature changed along with the
budget, and 50 of 312 prompts is a subset, so 129/312 stays the
canonical score and no budget-corrected score is reported. The extractor
also stripped reasoning only when a closing
\longtok{\textless{}/think\textgreater{}} was present; for a response
truncated inside its reasoning it graded any
\texttt{module\ …\ endmodule} found in the draft, and only response
hashes were kept, so it cannot be determined how many of the 162
extracted candidates or the 129 passes came from reasoning drafts. The
runner now refuses extraction from an unterminated reasoning block,
makes the thinking mode explicit and records response text, finish
reason and truncation per sample; the stored result was produced before
those changes. The earlier unclustered Wilson intervals on pooled and
parseable-only outcomes are not used.

The E11 lane also carries the campaign's clearest example of a scoring
decision taken for evidential rather than numerical reasons. The
all-problems reference sweep --- running the shipped reference
implementation of every problem through both simulators --- returned 311
of 312. One problem,
\longtok{verilog\_eval/spec-to-rtl/Prob099\_\allowbreak{}m2014\_\allowbreak{}q6c}, fails to
elaborate against its own test bench under both simulators, with Icarus
reporting that ports \texttt{Y2} and \texttt{Y4} are not ports of
\texttt{good1}; no candidate implementation can pass that task, and the
upstream defect biases pass@k downwards by at most 0.64 percentage
points on the specification-to-RTL split (1/156) and by 0.13 percentage
points overall (129/311 − 129/312). Excluding it yields 129/311 =
0.4148. That corrected figure is \emph{not} canonical. The recorded
amendment of 2026-08-22 states the reason plainly: no immutable
benchmark-level justification was captured for the exclusion, so the raw
129/312 result is canonical and the 129/311 value is retained only as a
noncanonical sensitivity; no samples and no verifier artifacts were
changed by the amendment. The ledger's integrity check confirms that the
component split sums correctly to the canonical total.

\textbf{E14 --- Omni-MATH replacement scope.} The surviving consolidated
ledger and
\longtok{outputs/verification/LEDGER\_\allowbreak{}CODE\_\allowbreak{}CHECK\_\allowbreak{}2026-09-19.json}
record 4,426/4,428 accepted rows and accuracy 0.5131043831902395,
consistent with 2,271 passes among 4,426 judged rows. Counting the
remaining two as failures gives 2,271/4,428 = 51.29 per cent. The
recovered terminal note (SHA-256
\longtok{14b8308df159f9c86a3a58d724b8152a5f44f098c57a824637fcbe5ae6f3602a})
identifies \texttt{omni-01195-a6cbdc4ee3c8694b} and
\texttt{omni-02044-612ddacaa39d91fe} as the two unaccepted rows and
records that Omni-Judge outputs were truncated before an Equivalence
Judgement section. It states that both had recorded skip reasons, that
omission followed the native scorer, that no re-judge was performed and
that the lane closed with all 4,428 dispositions. This now supports the
terminal decision and explains the earlier parser-failure terminology at
receipt level. The supplementary Mac-side validation now inspects the
stored rows and exports the two exclusion records, each lacking the
Equivalence Judgement heading and carrying the recorded omission reason;
no mathematical answer was independently rejudged. The note's
\texttt{score\_impact} nevertheless writes
\texttt{2271/4428\ =\ 0.5131043831902395}; the decimal is instead
2,271/4,426. The recovered finish ledger repeats the inconsistent
pairing of 2,271/4,428 with 51.31\%. Both sources remain unchanged, and
the thesis preserves the corrected distinction between 51.29\% over all
dispositions and 51.31\% over judged rows (Table E.1). The supplementary
export identifies the disposition and judged-row payloads and the
scoring-recheck file on the Mac, with hashes. Its traversal records
4,428 valid rows in each of the disposition, judged-row and
new-judge-response JSONL files, with 4,426 accepted, 2,271 correct and
two skipped dispositions. The recorded native-score fields explicitly
distinguish denominator 4,426 from the all-4,428 diagnostic. The full
E14 rows are not exported here; the cloud check can independently verify
the exported arithmetic and two exclusion records, but cannot
independently recount all 4,428 rows from this export.

\textbf{E14 generation budget.} The actor cap was 2,048 tokens
(\longtok{zvf-program/flagship/public\_\allowbreak{}omni\_\allowbreak{}math\_\allowbreak{}native.py}), and 2,982
of 4,428 responses reached it. Recounting
\texttt{native\_dispositions.jsonl}, capped responses are correct in
1,029 of 2,982 (34.5\%) and uncapped ones in 1,242 of 1,446 (85.9\%)
(corrected reading; see Table E.1). Difficulty confounds the split,
because harder problems both run longer and fail more often, so it does
not estimate what a larger budget would score. It does mean the 51.29\%
headline is substantially a property of the budget and not a budget-free
capability estimate. The judge's 300-token cap is what truncated the two
unjudged reports. The runner did not record capped status; it now writes
a per-task \texttt{cap\_status.json} and an
\longtok{accuracy\_\allowbreak{}by\_\allowbreak{}actor\_\allowbreak{}cap} receipt field, without altering the
stored grades.

\section{Original-contract results: lane
detail}\label{original-contract-results-lane-detail}

\textbf{E1 --- SWE-bench Pro.} The denominator is the complete 731-row
test split of \texttt{ScaleAI/SWE-bench\_Pro} at revision
\longtok{7ab5114912baf22bb098818e604c02fe7ad2c11f}. Within that
denominator the run produced 713 native evaluations, fourteen generation
failures and four lost generation artifacts; those eighteen non-native
outcomes are retained in the denominator rather than trimmed, so the
figure is a rate over all 731 tasks and not over the subset that graded
cleanly. Those 713 native evaluations are mostly not evaluations of an
applied fix (corrected reading; see Table E.1). The prompt asked for
unified diffs with concrete line ranges but showed source without line
numbers and with elided regions, and the runner's diff validator checked
syntax but not that hunk-header line counts match the hunk bodies.
Parsing the stored \texttt{candidates.json} with
\texttt{git\ apply\ -\/-stat} rejects 593 of 713 patches as corrupt
(``corrupt patch at line \ldots{}''), and 133 carry a placeholder
\texttt{index\ 1234567..abcdefg} line. The official entry script runs a
plain \texttt{git\ apply\ -v} without \texttt{-\/-recount} and without
stopping on failure, so for those 593 instances the tests ran against
the unpatched base. With \texttt{-\/-recount}, 711 of 713 parse. Both
resolved instances are among the 120 that parse as submitted. Apply
status was not captured by the run, so whether every parseable patch
applied to its repository is not known. The 2/731 is retained as the
measured score, but it mostly reflects diff-header arithmetic rather
than issue resolution. The runner now recounts hunk headers before
submission, rejects count mismatches and records parse status for the
raw and submitted patch; the stored run predates that repair. The
declared licence for the dataset is null in the lane's own receipt. The
evaluation harness is \texttt{scaleapi/SWE-bench\_Pro-os} at commit
\longtok{ca10a60a5fcae51e6948ffe1485d4153d421e6c5}, and the 731-task
scoring run used the 731 official run scripts unmodified on native
\texttt{linux/amd64} Modal workers. A 1,200-second timeout, against the
official 120, was used only in an earlier single-task harness check that
ran the gold patch under arm64 QEMU emulation. That check is recorded as
not a model score, so the timeout deviation does not apply to the 2/731
figure.

\textbf{E2 --- FrontierSWE.} The scientific result remains partial at
1/17 tasks. The public package does not supply the original execution
route or privately held operational records.

\textbf{E4 --- BankerToolBench.} The surviving historical summaries and
recovered finish ledger report restored public artifacts and 1,175
manifest checks passing. The recovered ledger specifically reports 0/6
pass3/pass16 archive artifacts available, with only
\texttt{TINKER\_BRIDGE.json} verifying. This supports the historical
archive-loss account, not a fresh verification of the archive or a
replacement of the receipt-only 0.3115 metric.

\textbf{E7 --- BinaryAudit.} The recovered finish ledger preserves both
the terminal table's ``1/46 attempted'' and its superseded historical
snapshot's ``1/28 errored''. The terminal figure is retained in Chapter
8. Recovery makes both source statements directly inspectable, but the
ledger does not explain their denominator difference and the raw attempt
inventory has not been revalidated here.

\textbf{E9 --- MLE-bench.} Coverage without complete grading is not a
score, and the lane records it as such. Behind the coverage number sit
97 legacy receipts across 50 competition identifiers, of which 41 graded
receipts cover 40 unique competitions; the ledger's integrity check
verifies precisely that decomposition. The separate merged-vLLM arm
holds seven receipts, of which exactly one is a valid grade (H\&M,
0.02132) and six are invalid submissions, and the integrity check
confirms that arm is recorded separately with one graded row. One old
Spooky Author submission is absent locally although the later run's
complete submission is retained.

\section{Replacement-scope lane narratives, 27 September
2026}\label{replacement-scope-lane-narratives-27-september-2026}

Three things changed in how the lanes were executed, and each is
recorded as a deviation in the lane's own result file. First, where a
fail-closed controller blocked only on paperwork that the directive had
superseded, the lane used a thin runner that calls the benchmark's
native loop and native grader directly. Second, compute moved off AWS:
E2, E6 and E9 ran on Google Cloud virtual machines, E1's grading ran in
Modal sandboxes, and E13 ran locally. Third, every lane was served the
same actor, \texttt{campaign-actor-bf16} (the base model with the
seed809 adapter merged, §8.1), through vLLM with the non-thinking chat
template. No lane went through the Tinker bridge implicated in the E4
failure mechanism (§8.1.2, §8.5), so that defect is not expected here,
although the lanes' result files do not record their stop settings. The
shared endpoint served at a 32,768-token context. E5 used a dedicated
65,536-token copy of the same weights because its banking prompts exceed
32k.

\textbf{E1 --- SWE-bench Multilingual.} This run covers the 16 wave-10
recovery tasks and the 174 tasks that the 12 September waves never
reached, 190 tasks in all. The actor resolved one
(\texttt{phpoffice\_\_phpspreadsheet-3463}), which is 1/190 = 0.53 per
cent. Only 57 of the 190 produced a patch that applied and could be
graded: 120 failed at patch application or evaluation and 13 produced an
empty patch. Over graded tasks alone the rate is 1/57 = 1.75 per cent
({[}0.3, 9.3{]}). Thirteen tasks produced no generation and were scored
as failures: eight exceeded the 32,768-token serving context and five
hit provider errors. The earlier waves covered the other 110 tasks, with
4 resolved (3.64 per cent, {[}1.42, 8.98{]}). They used the same task
universe, adapter and grader family, but served at a 65,536-token
context, and their source-context collector was lost and had to be
re-implemented for this run. Because the contract differs, the two runs
are not combined, and the lane does not report a 300-task figure.

\textbf{E2 --- CORE-Bench, hard difficulty.} The actor answered every
question correctly on 27 of the 45 \texttt{codeocean\_hard} capsules, a
task accuracy of 0.600. At the question level it answered 53 of 79
correctly: 28 of 33 written questions and 25 of 46 vision questions.
Thirty-two tasks ended with the agent calling \texttt{finish}, ten at
the 150-turn cap and three at the 8,100-second time limit. The agent was
a thin tool loop (\texttt{bash}, \texttt{query\_image}, \texttt{finish})
rather than the benchmark's reference CORE-Agent scaffold, and the same
actor answered the vision queries. Thirty tasks never reached a machine
because of a local DNS outage, and a driver crash then interrupted ten
of their reruns. All thirty were rerun from scratch, and every aborted
attempt is archived. Because a leak of the answers would inflate this
figure, the lane checked for one. No agent command referenced any host
that serves the answers. Each capsule's \texttt{results/} directory is
stripped before the agent starts, and the answer keys are held off the
agent's machine. Two qualifications remain. Twelve of the 27 correct
tasks left no files in the capsule's results directory, which the native
grader does not require. Some vision questions are categorical and could
be guessed. With n = 45 the interval spans 27 percentage points.

\textbf{E5 --- Tau3.} Eight of 97 tasks succeeded, a pass\^{}1 of 0.082.
Twenty-five tasks ended in infrastructure errors: fourteen exceeded the
65,536-token context and eleven returned an empty assistant message.
Their cause is not established, and they count as failures. An
accidental resume of the run overwrote 24 of those error records in the
results file; their identifiers and causes were rebuilt from the run
log. Over the 72 tasks without such an error the native pass\^{}1 is
8/72 = 0.111, reported as a sensitivity figure only. The 20 tasks
cleaned before 19 September were rerun under the same protocol rather
than carried over, and scored 1/20. The user simulator was
\texttt{gpt-4.1-2025-04-14} at temperature 0.

\textbf{E6 --- WebArena.} All 812 tasks were run, and 90 succeeded,
giving 0.111. The figure is a lower bound. Of the 114 tasks whose native
evaluator calls an LLM judge, 108 could not be graded, because the judge
was unavailable for the run; the other six were resolved to zero by a
non-judge check. Forty-one tasks failed with a harness-level browser
timeout. Both groups count as failures. The graded-only rate is 90/704 =
0.128 ({[}0.105, 0.155{]}). If every ungraded task had passed, the rate
would be 198/812 = 0.244. With the 41 harness errors counted as
failures, the true figure lies between 0.111 and 0.244, and nearer the
lower end if the ungraded tasks succeed at the graded rate. The native
judge model (\texttt{gpt-4-1106-preview}) is retired, so a later regrade
would substitute a different judge. Most failures are early stops by the
native loop: the same action three times (237), the same typing three
times (103), the step limit (101), and three consecutive parse failures
(98). The agent was limited to 384 output tokens per step, the few-shot
examples were passed as user and assistant turns because the chat
template rejects mid-conversation system messages, and sites were not
reset between tasks, as in the native runner. The last 103
shopping-admin tasks were split across three containers of the same
official site image to finish in time.

\textbf{E9 --- ML-Dev-Bench.} Ten of the 34 released tasks passed their
native validators, 0.294. The agent was the benchmark's OpenHands
CodeActAgent, limited to 50 iterations. No task errored or timed out,
but only after one harness repair: OpenHands 0.21.0 hangs when it
truncates its context window, and a one-line upstream fix was backported
without touching prompts, validators or tasks. Thirty-nine of 1,045
model calls exceeded the 32,768-token context and were absorbed by that
truncation path. A relay between the agent and the actor forced the
non-thinking template and restored tool-call text that the server's
parser had dropped. An earlier four-task batch that hung before the
repair is excluded.

\textbf{E13 --- BALROG.} All 255 native episodes were played and graded,
and the BALROG score is 26.1 per cent progression (standard error 1.77;
a stratified bootstrap gives {[}22.7, 29.5{]}). BALROG's score is the
mean of six per-environment means, not a mean over episodes. By
environment the actor reached BabyAI 82.0 per cent (n = 50), Baba Is AI
25.0 (n = 120), Crafter 22.7 (n = 10), TextWorld 16.9 (n = 30), MiniHack
10.0 (n = 40) and NetHack 0.0 (n = 5). Sampling followed the recorded
protocol at temperature 1.0, but with thinking off where the original
run had it on. A 2,000-step cap was placed on NetHack for cost; all five
NetHack episodes ended on their own before it, so the cap never took
effect. Episodes stopped by an operational stop or a hung worker were
restarted from scratch with fresh seeds. The result file records 262
episode attempt records for the 255 scored episodes, of which 7 are
marked killed for restart (\longtok{n\_killed\_\allowbreak{}for\_\allowbreak{}restart\_\allowbreak{}records});
its caveat text instead counts 11 restarts (one NetHack and five Baba Is
AI episodes at the 01:40 UTC stop, plus five hung workers), and the file
does not reconcile the two counts. If hangs correlate with episode
difficulty, that is a mild selection risk. The run used macOS arm64
builds of the pinned sources rather than the Linux image. The 26.1
figure is on a different scope from the 26-episode small-scale E13 rows
of Tables 8.A and 8.B and is not compared with them.

The launch-pending and quota labels of 19 September no longer describe
any lane.

\section{E4 rerun: reported audit and causal
limits}\label{e4-rerun-reported-audit-and-causal-limits}

The recovered finish ledger records E4 as \texttt{CLOSED\_PARTIAL},
preserving the 19 September close-out. It reports the pass3/pass16
archive as missing (0/6 archive artifacts; only
\texttt{TINKER\_BRIDGE.json} verifies), and keeps the 0.3115 recovery
metric as receipt-only evidence. Its updates through 21 September state
that the full 100-trial rerun had not launched. That snapshot cannot
substantiate the later completed rerun. The surviving 26 September E4
result's \longtok{prior\_\allowbreak{}base\_\allowbreak{}run} field instead reports a 21 September
run with 100 trials, one error and mean reward 0.0. Its source job
receipt remains unrecovered in this package, so the detailed account
below remains historical rather than a fresh trajectory-level
verification.

The earlier report describes the rerun as the base model with a
zero-initialised LoRA snapshot, equivalent to base sampling with no
training steps, through the Modal bridge \texttt{/v1/responses} path at
\texttt{-\/-n-concurrent\ 2} It reports 100 of 100 trials, one
\texttt{NonZeroAgentExitCodeError}, a mean reward of exactly 0.0, all
one hundred trials in the 0.0 bucket and a runtime of 3 hours 29
minutes. The surviving
\longtok{outputs/e1\_\allowbreak{}e14\_\allowbreak{}small\_\allowbreak{}scale\_\allowbreak{}2026-09-26/E4/result.json}
(\longtok{prior\_\allowbreak{}base\_\allowbreak{}run}) corroborates the 100 trials, one error, zero
mean and reward buckets, and identifies the serving path. The exact
launch configuration, named error type, cap and runtime are not
established by that secondary field and cannot be rechecked against the
unrecovered
\longtok{outputs/e4\_\allowbreak{}banker\_\allowbreak{}toolbench/official\_\allowbreak{}repo\_\allowbreak{}ff6db552/jobs/btb-banking-tasks-tinker-bridge/result.json}.
None of those later-run details is attributed to the recovered
pending-state snapshot.

The historical audit reports a failure mechanism across the one hundred
trajectories. Agents reportedly terminate after 2 to 14 steps, with a
median of about 6, and produce no \texttt{deliverables/} directory; the
verifier consequently short-circuits to 0.0 on the message ``No
deliverables found''. Twenty-seven or more of the trajectories end in
role-token degeneration of the form \texttt{user\ user\ assistant\ ...}
repetition. The record at the time read this as the base model being
unable to sustain the opencode tool-use loop. A later bridge check
exposes an alternative explanation, so that reading is not identified by
this uncontrolled run. The bridge's \texttt{/v1/responses} path passed
no stop sequences and decoded with special tokens removed, so generation
ran past the end-of-turn token into invented \texttt{user},
\texttt{tool\_response} and \texttt{assistant} turns, which is exactly
the role-token degeneration observed. The corrected shim scored 0.4087
on \texttt{btb-19b3361c}, but that single-task success is compared with
100 bridge trials under different tasks, serving path and caps. The
contrast supports a plausible mechanism that can cause zero-reward
trajectories; it does not quantify the defect's contribution to the full
zero or rule out model limitations. The reported no-deliverables check
is consistent with fail-closed grading, but the absent raw trajectories
prevent a fresh verification of the cause of every zero. No unconfounded
full-suite BankerToolBench capability estimate exists for this actor
family; the other measurement is the six-task subset of Tables 8.A and
8.B. The same bridge served E2, E5, E7 and the E9 streaming arm, so
Table 8.2 marks those lanes as exposed to the same defect. In
particular, the rerun's 0.0 must not be set against the original
single-task outcome as though the two were a controlled pairing: they
are different configurations under different runtime conditions, and the
campaign makes no comparison between them.

The historical rerun report records 16.35 million prompt tokens and 2.18
million completion tokens. These token details cannot be rechecked
against the absent job receipt. The surviving
\longtok{outputs/e1\_\allowbreak{}e14\_\allowbreak{}small\_\allowbreak{}scale\_\allowbreak{}2026-09-26/E4/result.json}
(\longtok{prior\_\allowbreak{}base\_\allowbreak{}run}) repeats the 100 trials, one error and zero
mean; its \longtok{verifiable\_\allowbreak{}from\_\allowbreak{}raw} flag is a historical assertion
and does not establish raw-artifact availability in this repository.

\section{Externally blocked lanes and verification
limits}\label{externally-blocked-lanes-and-verification-limits}

Seven original contracts remain externally blocked: SDAB, BinaryAudit,
LifeSciBench, AgentHarm, AppBench, OpenReward and FrontierMath. These
are coverage limitations, not scored failures or evidence about model
ability. Replacement scopes retain their separate names and
denominators. Operational access correspondence, contact destinations,
account capacity, request identifiers and account costs are omitted from
this public derivative.

The historical record reports eleven deterministic checks covering
selected arithmetic, coverage and source bindings. This is not an
independent re-execution of the benchmarks, and it does not establish
universal receipt availability. The original campaign recovery and
stored-payload checks have the bounded scope stated in Appendix B.

\section{Small-scale reruns: detail and per-lane
caveats}\label{small-scale-reruns-detail-and-per-lane-caveats}

Each lane's items, raw outputs and grader logs are kept so that every
figure in Tables 8.A and 8.B can be recomputed.

For E12 the 151 rubric items come from only six generated applications,
so they are not independent. The cluster-level test is an exact
sign-flip test over the six per-app pass-rate differences (−0.042,
+0.030, 0.000, −0.160, 0.000, −0.042), and it gives p = 0.25. The E12
grader is an LLM judge that is the same \texttt{Qwen/Qwen3.6-35B-A3B}
base model as the base arm, so self-preference could favour the base
arm. Two judge versions disagree widely (0.536 against 0.808 on the
Tinker base run; \texttt{E12/result.json}). The E12 rates and difference
are therefore judge-dependent, and no independent judge was run. For E4,
all six tasks score 0 in both arms except \texttt{btb-19b3361c}; a
Wilson interval is not valid for a mean of fractional {[}0, 1{]}
rewards, and a bootstrap over one nonzero difference in six is
degenerate. For E13, in 22 of the 26 episodes the progression is
identical in the two arms. No lane-level difference is interpreted as
evidence that the adapter improves or degrades performance unless its
test excludes zero, and none is extrapolated to the full suite.

The principal caveat for each lane, as recorded in its result file, is
listed below; the full list is in each lane's \texttt{result.json}.

\begin{itemize}
\tightlist
\item
  \textbf{E1.} New base-model arm; not comparable to or poolable with
  the original 2/731. Original inference temperature was 0.2; this
  public subset cannot replay its private execution route.
\item
  \textbf{E2.} Substitute benchmark; never pool with or present as the
  FrontierSWE lane score (original: 1/17 tasks, replay 0.8628). Original
  execution access is outside this public package.
\item
  \textbf{E3.} New base-model arm; never pool with or present as the
  original SDAB/seed809 lane score.
\item
  \textbf{E4.} Small n=6 of 100; one nonzero task, so no interval is
  reported (the result file's Wilson95 on summed fractional reward is
  not valid for a {[}0,1{]} mean).
\item
  \textbf{E5.} n=10, not the requested 20: the E4+E5 6M Tinker-token cap
  was exhausted (E4 1.56M, abandoned APEX attempt 2.33M, tau2 2.05M).
\item
  \textbf{E6.} Tasks drawn by random.Random(20260926).sample from all
  128 sorted miniwob env ids; episode seeds 0-2.
\item
  \textbf{E7.} New base-model arm; never pool with or present as the
  original lane score.
\item
  \textbf{E8.} New base-model arm; not comparable to or poolable with
  the seed809-adapter 1967/1967 run (22.88\%); the public subset cannot
  replay the original execution route.
\item
  \textbf{E9.} Competition pool: 6 low-split competitions with small
  non-image-heavy data (detecting-insults, jigsaw-toxic,
  leaf-classification, nomad2018, random-acts-of-pizza, spooky); 5 drawn
  by random.Random(20260926).sample. n=5 so the CI is very wide.
\item
  \textbf{E10.} New base-model arm on AgentHarm itself; the
  seed809-adapter lane figure (AgentDojo benign utility 88/97 = 90.72\%;
  original execution route not supplied publicly) is a different
  benchmark and not comparable.
\item
  \textbf{E11.} New base-model arm; not comparable to the retained
  adapter receipt 129/312 = 41.35\% (that run used thinking-default
  template, temperature 0.2, all 312).
\item
  \textbf{E12.} Public tasks, not held-out. Label mandatory. New
  base-model arm; never pool with the original lane.
\item
  \textbf{E13.} New base-model arm; not comparable to or pooled with the
  original seed809-LoRA E13 plan (21-episode BabyAI admission set).
\item
  \textbf{E14.} New base-model arm; the prior 2271/4428 = 51.29\% figure
  (51.31\% on the official scorer's 4426 judged rows) is not comparable
  (different actor, full split).
\end{itemize}



\chapter{Corrections to Internal Source
Documents}\label{corrections-to-internal-source-documents}

The numbers in this thesis are drawn from the project's internal source
documents: the unpublished P1--P8 reports and their section files, audit
summaries and result files. Some of those documents contain errors --- a
wrong sign, a test applied to the wrong unit of analysis, an interval
that does not describe its own estimate, or a label that does not match
the rows it covers. The chapters report only the corrected value and
mark it ``(corrected; see Table E.1)''. This appendix records, for each
correction, what the source states, what the thesis uses instead, and
the basis for the change, so that a reader who opens the source can
reconcile it with the thesis. Caveats about the evidence itself, as
opposed to errors in a source document, stay in the chapters. Paths are
relative to the repository root; ``sections/'' is
\longtok{platform\_hybrid/paper/sections/} and ``results/'' is
\longtok{platform\_hybrid/experiments/results/}.

\textbf{Corrections applied to internal source documents.}

\begingroup\scriptsize\setlength{\tabcolsep}{3pt}
\begin{longtable}[]{@{}
  >{\raggedright\arraybackslash\hspace{0pt}}p{(\linewidth - 10\tabcolsep) * \real{0.1200}}
  >{\raggedright\arraybackslash\hspace{0pt}}p{(\linewidth - 10\tabcolsep) * \real{0.1800}}
  >{\raggedright\arraybackslash\hspace{0pt}}p{(\linewidth - 10\tabcolsep) * \real{0.2200}}
  >{\raggedright\arraybackslash\hspace{0pt}}p{(\linewidth - 10\tabcolsep) * \real{0.2600}}
  >{\raggedright\arraybackslash\hspace{0pt}}p{(\linewidth - 10\tabcolsep) * \real{0.2200}}@{}}
\caption{Corrections applied to internal source documents. Each row's
corrected value is the one used in the chapter named in the first
column.}\tabularnewline
\toprule\noalign{}
\begin{minipage}[b]{\linewidth}\raggedright
Location in thesis
\end{minipage} & \begin{minipage}[b]{\linewidth}\raggedright
Source document
\end{minipage} & \begin{minipage}[b]{\linewidth}\raggedright
What the source states
\end{minipage} & \begin{minipage}[b]{\linewidth}\raggedright
Corrected value / reading
\end{minipage} & \begin{minipage}[b]{\linewidth}\raggedright
Basis
\end{minipage} \\
\midrule\noalign{}
\endfirsthead
\toprule\noalign{}
\begin{minipage}[b]{\linewidth}\raggedright
Location in thesis
\end{minipage} & \begin{minipage}[b]{\linewidth}\raggedright
Source document
\end{minipage} & \begin{minipage}[b]{\linewidth}\raggedright
What the source states
\end{minipage} & \begin{minipage}[b]{\linewidth}\raggedright
Corrected value / reading
\end{minipage} & \begin{minipage}[b]{\linewidth}\raggedright
Basis
\end{minipage} \\
\midrule\noalign{}
\endhead
\bottomrule\noalign{}
\endlastfoot
§1.2 (variant audit) &
\longtok{zvf-program/audit/STATISTICAL\_REANALYSIS.md} &
Percentile-bootstrap 95\% intervals over the eight seeds, e.g.~Dr.~GRPO
−0.00200, {[}−0.00950, +0.00725{]} & Paired-\(t\) 95\% intervals,
e.g.~Dr.~GRPO \(-0.00200\), \([-0.01272, +0.00872]\), \(p = 0.673\);
DAPO \(+0.00100\), \([-0.00632, +0.00832]\) (also in Table 1.1) & A
percentile bootstrap undercovers at \(n = 8\); recomputed from
\longtok{zvf-program/audit/results/full/*.json} \\
§4.1, §6.2.3 (same-stack PPO vs GRPO) &
\longtok{sections/frontier\_\allowbreak{}synthesis\_\allowbreak{}scaling.tex} & Reports paired Δ =
−0.002, p = 0.374 without naming the model or task &
\longtok{Qwen/Qwen2.5-0.5B}, two-digit addition, 40 steps, five paired
seeds, 200 fresh draws per seed &
\longtok{results/samestack\_\allowbreak{}ppo\_\allowbreak{}grpo.json};
\longtok{platform\_hybrid/experiments/modal/modal\_\allowbreak{}samestack\_\allowbreak{}ppo\_\allowbreak{}grpo.py} \\
§1.2, §4.1, §6.2.3, §6.4.2, §6.6 (same-stack PPO vs GRPO pairing) &
\longtok{sections/frontier\_\allowbreak{}synthesis\_\allowbreak{}scaling.tex};
\longtok{results/samestack\_\allowbreak{}ppo\_\allowbreak{}grpo.json} & A paired comparison over
five seeds, each seed's held-out accuracy on 200 items & Paired by seed
only: the GRPO and PPO evaluation sets of a seed share 6, 6, 4, 7 and 3
of 200 items; evaluation items also drawn in training are 12/16/13/15/13
of 200 for GRPO against 90/86/80/83/74 for PPO; the stored −0.002, p =
0.374 is unchanged & Exact replay of the Python RNG stream in
\longtok{modal\_\allowbreak{}samestack\_\allowbreak{}ppo\_\allowbreak{}grpo.py}: the evaluation set is drawn
from the shared stream after training, and GRPO consumes 16 draws per
step against PPO's 128 \\
§6.3.2 (60-run ZVF ranges) & \longtok{sections/p2\_\allowbreak{}results\_\allowbreak{}intro.tex} &
The 60 runs behind the late-phase and early-phase ZVF ranges and the
lag-1 autocorrelation of ≈ 0.94 are ``measured'' & 45 of the 60 are
five-seed variance-mitigation projection rows; only the remainder, which
includes the 12-run group-size sweep, is measured & The nine-row panel
comes from the dry-run file
\longtok{experiments/results/variance\_mitigation.tsv}
(\longtok{results/zvf\_\allowbreak{}by\_\allowbreak{}library.tsv}) \\
§6.3.3 (ZVF against the Bernoulli null) &
\longtok{sections/group\_size.tex} & Closed form evaluated at each arm's
held-out accuracy (0.964, 0.954, 0.923 and 0.704 at G = 2, 4, 8, 16),
giving negative residuals & Per-step-average null 0.828, 0.709, 0.597,
0.522 against measured 0.838, 0.764, 0.691, 0.631; residuals +0.010,
+0.055, +0.094, +0.109, positive at every G & ZVF is a training-batch
quantity and held-out accuracy is not; recomputed from
\longtok{results/groupsize\_\allowbreak{}zvf\_\allowbreak{}sweep.json}; the per-step-average
baseline is the one named in
\longtok{sections/appendix\_\allowbreak{}zvf\_\allowbreak{}formalization.tex} \\
§6.3.4 (naive ZVF--accuracy correlation) & \longtok{sections/zvf.tex} & r
= 0.22 over n = 23 pooled rows, bootstrap 95\% CI {[}−0.29, 0.94{]} & n
= 17 per-experiment rows; Fisher-z 95\% CI {[}−0.29, +0.63{]} & The
residual table on disk, \longtok{results/zvf\_\allowbreak{}iter26\_\allowbreak{}residual.tsv},
holds 17 rows; a bootstrap upper bound near +0.95 at this n indicates an
unstable bootstrap \\
§6.3.4 (residualised ZVF) & \longtok{sections/zvf.tex} & After regressing
out mean reward, residualised ZVF remains positively associated with
last-10 accuracy, r = +0.80, 95\% CI {[}+0.54, +1.00{]} & That slice is
the n = 9 variance-mitigation rows, a simulation projection; the
measured pooled residual is r = +0.205, Fisher 95\% CI {[}−0.31,
+0.62{]}, p = 0.43, n = 17 & \longtok{results/zvf\_\allowbreak{}iter26\_\allowbreak{}residual.tsv};
projection status in §6.3.2 \\
§6.4.1 (group-size TOST) & \longtok{sections/group\_\allowbreak{}size\_\allowbreak{}iter27.tex} &
TOST p-values against G = 2 of p \textless{} 10⁻⁹ at G = 4 and G = 8 and
p \textless{} 10⁻⁴ at G = 16, with accompanying intervals & Seed-level
paired TOST (margin ±0.02, n = 3): p = 0.012 (G = 4), 0.043 (G = 8),
0.056 (G = 16); equivalence not established at G = 16; the source
intervals are withdrawn & p-values that small cannot come from three
paired seeds, so the source does not use the seed as the unit;
recomputed per seed \\
§6.4.1 (retention asymptote) & \longtok{sections/group\_\allowbreak{}size\_\allowbreak{}iter27.tex}
& Asymptote interval {[}0.712, 0.745{]} for
\(R(T) = 0.729 + 0.247\,e^{-T/5.76\mathrm{M}}\) & No confidence
interval; the fit is illustrative only & A three-parameter curve through
four reconstructed points leaves one residual degree of freedom \\
§6.4.2 (algorithm axis of the factorization audit) &
\longtok{results/berkeley/unpacking\_\allowbreak{}dpo\_\allowbreak{}ppo\_\allowbreak{}factorization.json} &
``DECISIVE'' equivalence of GRPO and PPO & Equivalence not established;
paired \(t\) \(p = 0.374\) is the primary test, exact sign-flip
\(p = 0.625\) & A failure to reject is not equivalence; paired TOST at
±0.005 gives \(p = 0.104\) (§6.2.3) \\
§6.4.2 (H1, gradient-utilisation ratio) &
\longtok{results/berkeley/adding\_\allowbreak{}error\_\allowbreak{}bars\_\allowbreak{}summary.json};
\longtok{sections/p3\_conclusion.tex} & Graded DECISIVE, with interval
{[}4.16, 4.82{]} & Held at hypothesis-generator status; point estimate
5.03, no interval reported & The interval excludes its own point
estimate; it is taken over n = 4 budgets around their mean of 4.456, and
the audit's own note says ``n=4 budgets \ldots{} cannot support CI on
the ratio'' \\
§6.4.2 (H2, retention-versus-token slope) &
\longtok{results/berkeley/adding\_\allowbreak{}error\_\allowbreak{}bars\_\allowbreak{}summary.json};
\longtok{sections/p3\_conclusion.tex} & Graded DECISIVE, with interval
{[}−0.237, −0.038{]} & Held at hypothesis-generator status & The slope
is fitted to the four reconstructed budgets of the token-budget grid,
not to trained arms \\
§6.5.3 (sliding-window trap rate) & P4 source,
\longtok{sections/length\_bias.tex} & The 6/16 (37.5\%) sliding-window
rate is an upper bound on trap prevalence in the controlled data & A
point estimate; Wilson 95\% CIs {[}5.7\%, 51\%{]} (2/10), {[}30\%,
90\%{]} (4/6), {[}18.5\%, 61\%{]} (6/16) & A point estimate is not a
bound \\
§6.5.3 (iteration-20 mechanism regression) &
\longtok{sections/length\_bias.tex} & Arithmetic OLS slope ``reduced by
0.43'' with 95\% CI {[}−0.58, −0.28{]}, paired p = 0.036; Theil--Sen
``reduced by 1.11'' with CI {[}−1.59, −0.63{]}, p = 0.047 & Dr.~GRPO
minus GRPO, paired by seed, n = 5: OLS +0.43 {[}+0.18, +0.68{]}, paired
\(t\) \(p = 0.009\); Theil--Sen +1.11 {[}+0.33, +1.88{]}, \(p = 0.017\);
permutation \(p = 0.0625\) for both & With mean OLS slopes of −3.65
(GRPO) and −3.22 (Dr.~GRPO), Dr.~GRPO minus GRPO is positive, so the
printed intervals have the wrong sign, and the source's CIs and p-values
do not come from the same test; recomputed from
\longtok{results/length\_\allowbreak{}bias\_\allowbreak{}mechanism\_\allowbreak{}per\_\allowbreak{}run.tsv} \\
§6.5.3 (iteration-24 late-window sign-flip) &
\longtok{sections/length\_bias.tex};
\longtok{results/length\_\allowbreak{}bias\_\allowbreak{}iter24\_\allowbreak{}signflip.tsv} & Sign-flip p =
0.468, 0.562, 0.011 and p \textless{} 0.001 across the four cells; unit
the within-run window & Exact two-sided sign-flip permutation over
seeds: p = 0.688 (arithmetic GRPO), 0.625 (arithmetic Dr.~GRPO), 0.250
(GSM8K GRPO), 0.250 (GSM8K Dr.~GRPO); unit the seed, floor 0.25 at n = 3
& The source labelled a one-sample t-test over per-seed fractions as a
sign-flip test; on three identical values it returns t = −∞, p = 0;
recomputed by
\longtok{platform\_modal/scripts/length\_\allowbreak{}bias\_\allowbreak{}iter24.py} \\
§6.5.3 (Benjamini--Hochberg over fourteen exploratory tests) &
\longtok{sections/length\_bias.tex} & GSM8K Dr.~GRPO sign-flip the sole
survivor at q = 0.014; arithmetic OLS and Theil--Sen at q = 0.051 and
0.058 & No test survives; smallest q = 0.116 (arithmetic OLS and
Theil--Sen), and 0.116 for every fourteen-of-sixteen family selection &
The source's q-values are reproduced exactly only by entering the
degenerate p = 0 as 0.001; recomputed with the corrected sign-flip
p-values \\
§6.5.4 (GSM8K-CoT pre-to-post gains) &
\longtok{sections/p4\_\allowbreak{}results\_\allowbreak{}intro.tex} & McNemar exact
\(p = 1.1 \times 10^{-4}\) (GRPO) and \(7.6 \times 10^{-4}\) (Dr.~GRPO),
pooled over seeds & Seed-level one-sample \(t\) on the three per-seed
deltas: \(p = 0.0051\) (GRPO), \(p = 0.0099\) (Dr.~GRPO) & The pooled
test treats the 600 seed-item pairs as independent, but the same 200
items recur across seeds; recomputed from
\longtok{results/drgrpo\_\allowbreak{}gsm8k\_\allowbreak{}cot.json} \\
§6.5.4 (iteration-136 late-efficiency contrast) &
\longtok{results/berkeley/adding\_\allowbreak{}error\_\allowbreak{}bars\_\allowbreak{}summary.json};
\longtok{sections/p4\_conclusion.tex} & Graded DECISIVE, with ``one-sided
parametric p = 0.031'' & Suggestive: \(p = 0.031\) is an exact one-sided
Wilcoxon signed-rank p at its \(n = 5\) floor (1/32); paired
\(t(4) = 5.96\), two-sided \(p = 0.004\); two-sided sign-flip
\(p = 0.0625\); Holm across the eight-hypothesis family gives 0.25 &
Recomputed from \longtok{length\_\allowbreak{}bias\_\allowbreak{}iter136\_\allowbreak{}step\_\allowbreak{}coupling.tsv} \\
§7.1 (variance decomposition, algorithm axis) &
\longtok{sections/p5\_evidence.tex};
\longtok{results/p5p8/p5\_\allowbreak{}headline\_\allowbreak{}cis\_\allowbreak{}full.tsv} & η² ZVF bootstrap CI
{[}0.0014, 0.0585{]}, graded DECISIVE (algorithm explains at most about
6\% of ZVF variance); reward CI {[}0.0014, 0.0577{]} & Within-method
step-bootstrap CIs: ZVF {[}0.0094, 0.1476{]}, graded SUGGESTIVE; reward
{[}0.0019, 0.0790{]}; point estimates 0.0454 and 0.0075 unchanged; still
too narrow because steps are autocorrelated & The source bootstrap
resampled the pooled steps, which gives a null band rather than a CI; on
the \texttt{loss} column it gave {[}0.0014, 0.0582{]} around an estimate
of 0.9867; recomputed by
\longtok{platform\_modal/scripts/p5p8/p5\_\allowbreak{}headline\_\allowbreak{}cis\_\allowbreak{}full.py} \\
§7.1 (variance decomposition, group-size axis) &
\longtok{sections/p5\_evidence.tex} & η² = 1.0000 for the group-size
axis, and group-size-to-algorithm ratios of 22× (ZVF) and 133× (reward)
& No group-size η² and no ratio reported & Computed on
\longtok{groupsize\_\allowbreak{}zvf\_\allowbreak{}sweep.tsv}, which holds one row per G level (4
cells, df = 0 within), so η² = 1 by construction \\
§7.1 (98-cell decomposition) & \longtok{sections/p5\_evidence.tex} &
Stack-to-seed variance ratios of 503× to 96,128× & Not quoted as effect
sizes & The seed axis rounds to 0.000, so the ratios divide by a
near-zero quantity \\
§7.1 (manifest--outcome coupling) & \longtok{sections/p5\_stack.tex},
§Live-Corpus Coupling & Hamming distance predicts
\textbar ΔZVF\textbar{} with Spearman ρ = +0.529 and
\textbar ΔPCD\textbar{} with ρ = +0.541, p ≤ 10⁻³ & ρ = +0.44 (ZVF),
+0.39 (PCD), +0.28 (reward) over all 4,753 pairs of the 98 cells; Mantel
permutation p \textless{} 10⁻³ for each & The source script
\longtok{platform\_modal/scripts/p5p8/p5\_\allowbreak{}manifest\_\allowbreak{}outcome\_\allowbreak{}coupling.py}
and its committed output give +0.407 and +0.336 on a 2,000-pair sample;
no variant reproduces +0.529 or +0.541; pair-level p-values are invalid
because pairs share cells; recomputed in
\longtok{thesis\_\allowbreak{}review\_\allowbreak{}2026-09-27/stats\_followup/m18\_\allowbreak{}mantel\_\allowbreak{}summary.json} \\
§7.2 (same-stack method deltas) &
\longtok{sections/p6\_\allowbreak{}measured\_\allowbreak{}evidence.tex} & AERO and AREAL reward
decreases with intervals excluding zero, read as resampling
prompt-steps, uncorrected & Intervals resample the last 10 of 40 steps;
exact sign-flip Holm p = 0.047 (AERO), 0.047 (AREAL), 0.23 (GIFT); at
prompt level the AERO interval includes zero & Recomputed from
\longtok{n2\_\allowbreak{}reward\_\allowbreak{}tensor\_\allowbreak{}resume/*\_s0\_tensors.jsonl};
\longtok{thesis\_\allowbreak{}review\_\allowbreak{}2026-09-27/stats\_followup/m16\_summary.json} \\
§7.3 (hindsight-optimal shrink rule) &
\longtok{sections/p7\_controller.tex} & The 0.797 cost ratio of the
per-prompt hindsight optimum is a ceiling on rollout spend & The best
cost ratio under a no-shrinking-of-saturated-prompts constraint, not a
lower bound on cost for any adaptive rule & Dualformer-Auto's observed
0.66 to 0.68 on the four-bin rule is cheaper \\
§8.1, §8.3, Appendix D §D.5, §§9.1--9.3 (E1 interpretation) &
\longtok{seed1818/receipt.json}; \longtok{E1\_\allowbreak{}E14\_\allowbreak{}Results.md} & 713
native evaluations of model patches behind 2/731 & 593 of 713 submitted
patches are corrupt for \texttt{git\ apply} (hunk-header counts) and
were tested on the unpatched base; 711/713 parse with
\longtok{-\/-recount}; both passes are among the 120 that parse; 2/731
unchanged but mostly measures patch formatting &
\texttt{git\ apply\ -\/-stat} over the stored \longtok{candidates.json};
official entry script runs plain \texttt{git\ apply} and continues on
failure \\
§8.1, §8.2, Appendix D §D.3--D.4, §9.3 (E11 failure attribution) &
\longtok{e11\_\allowbreak{}full\_\allowbreak{}receipt.json}; earlier campaign transcription & 150
of 183 failures are format failures; parseable-only 129/162 a format
sensitivity & The 150 extraction failures follow mainly from a
1,024-token cap with thinking on (mean 1,002 of 1,024 tokens); same
adapter 33/50 with thinking off at 4,096 tokens against 23/50 for this
run on a 50-item subset (McNemar p ≈ 0.006); some extracted code may
come from unterminated reasoning drafts; 129/312 unchanged &
\longtok{outputs/e1\_\allowbreak{}e14\_\allowbreak{}small\_\allowbreak{}scale\_\allowbreak{}2026-09-26/E11/paired.json}
against the receipt's per-problem verdicts; extractor in
\longtok{zvf-program/flagship/modal\_\allowbreak{}e1\_\allowbreak{}e14.py} \\
§8.1, §8.2, Appendix D §D.4 (E14 generation budget) &
\longtok{native\_dispositions.jsonl}; earlier transcription & 51.29\%
read as the replacement-scope accuracy with a 67.3\% cap rate disclosed
& Capped responses 1,029/2,982 = 34.5\% correct, uncapped 1,242/1,446 =
85.9\%; confounded by difficulty; 2,271/4,428 unchanged & Recount of
stored dispositions by \longtok{completion\_tokens} = 2,048 \\
§8.2, Appendix D §D.4 (E11 bootstrap script) & Earlier thesis citation &
Script at
\longtok{outputs/PES\_\allowbreak{}Phase2\_\allowbreak{}Third\_\allowbreak{}Review\_\allowbreak{}2026-09-24/analysis/e11\_\allowbreak{}design\_\allowbreak{}cluster\_\allowbreak{}bootstrap.py}
&
\longtok{platform\_hybrid/experiments/scripts/e11\_\allowbreak{}design\_\allowbreak{}cluster\_\allowbreak{}bootstrap.py}
and \texttt{.json}; interval {[}0.3429487, 0.4871795{]} reproduced
exactly & The cited path was never committed and is gitignored \\
§8.2, Appendix D §D.4 (E14 denominator label) & Recovered
\longtok{outputs/PES\_\allowbreak{}Phase2\_\allowbreak{}Review\_\allowbreak{}2026-09-12/finish/e14\_\allowbreak{}terminal\_\allowbreak{}note\_\allowbreak{}2026-09-19.json};
recovered finish \longtok{Pending\_Experiments.md} & Terminal note writes
\texttt{2271/\allowbreak{}4428\ =\allowbreak{}\ 0.5131043831902395}; finish table pairs
\longtok{2271/4428} with 51.31\% & Existing thesis arithmetic is
retained: 2,271/4,428 = 51.29\% over all dispositions; 2,271/4,426 =
51.31\% over judged rows & Exact arithmetic and supplementary stored
native-score fields distinguish denominator 4,426 from the all-4,428
diagnostic; originals preserved unchanged; no re-judge \\
§8.2, Appendix D §D.4 (E10 protocol attribution) & Earlier campaign
transcription; recovered \longtok{agentdojo\_\allowbreak{}native\_\allowbreak{}receipt.json} &
Attack, defence and security-evaluation fields all null; upstream commit
attributed to native receipt & \texttt{attack} and \texttt{defense} are
null; \longtok{security\_evaluation} is false; receipt binds v1.2.2 and
contains no upstream benchmark commit field & Direct inspection of the
byte-preserved original; no new security or held-out claim \\
§8.5, Appendix D §D.7 (E4 later-rerun attribution) & Earlier campaign
transcription citing finish \longtok{Pending\_Experiments.md} for the
completed 100-trial rerun & Pending-state ledger cited as evidence for
later zero-reward rerun details & Recovered ledger says NOT launched at
its 21 September snapshot; 100 trials, one error and zero mean are
supported here by surviving E4 \longtok{result.json} →
\longtok{prior\_\allowbreak{}base\_\allowbreak{}run} & Temporal source correction only; raw rerun
receipt remains unrecovered and detailed launch/trajectory/spend fields
remain historical \\
\end{longtable}
\endgroup

\section{Evidence-scope repairs after the 2 October
review}\label{evidence-scope-repairs-after-the-2-october-review}

The following corrections were made without new model training or
regenerated historical receipts.

\begin{itemize}
\tightlist
\item
  \textbf{ZVF definition (§§4.3, 6.3.1, 7.1, 7.3).} Exact ties and
  sample variance ≤ 10⁻⁶ are now separate metrics. On the 600 stored
  groups, seeded micro-jitter changes exact ties from 95 to zero, but
  leaves 95 below the declared threshold. The epsilon sweep and source
  hashes are recorded in
  \longtok{pcd\_\allowbreak{}vs\_\allowbreak{}zvf\_\allowbreak{}tolerance\_\allowbreak{}analysis.json}.
\item
  \textbf{Controller interpretation (§7.3).} Zero rescue after plugging
  \ensuremath{\hat{p}} = 0 or 1 into a Bernoulli model, or resampling eight identical
  rewards, is estimator-imposed. It does not show latent certainty or
  impossibility of recovery with independent continuations.
\item
  \textbf{Loss provenance (§§1.1, 4.2--4.3, 6.3.1).} Zero-KL claims are
  row-specific. TRL nonzero-KL configuration evidence, adaptive-KL
  defaults, managed unknowns and actor source-binding gaps remain
  visible; zero reward-relative advantage does not imply a zero full
  optimiser update.
\item
  \textbf{PPO and group-size scope (§6.5.5; Chapter 9).} The unsaturated
  rerun is present in the synthesis. Its five-seed PPO comparison is
  baseline-specific and test-sensitive (paired-t p = 0.016; exact p =
  0.0625), with unequal prompt exposure and near-cap lengths. G8 minus
  G2 remains inconclusive, without equivalence.
\item
  \textbf{E11 uncertainty (§8.2; Appendix D).} The point score stays
  129/312. Resampling the 156 paired underlying designs gives a 95\%
  percentile interval {[}34.3\%, 48.7\%{]}, replacing the iid-312 Wilson
  interval. The conditional extracted-only score remains descriptive.
\item
  \textbf{E4 mechanism (§8.5; Chapter 9).} Missing stops are a verified
  defect, but the one corrected-path task changes several conditions
  together and cannot attribute the full 100-trial zero to stops alone.
\item
  \textbf{Source access (Appendix B §B.11; recovery update, 3 October).}
  Six exact originals covering all five previously missing basenames
  were recovered from the Mac, including both pending-state snapshots,
  and preserved with hashes in a separate private package. All remain
  absent from the frozen checkout. Supplementary read-only checks
  corroborate stored E8/E10 grades and E14 disposition/denominator
  reporting; they do not rerun scorers or establish content
  decontamination. Training raw logs and E4 raw rerun evidence remain
  unresolved. The E14 denominator and E10/E4 attribution repairs above
  change no scientific counts or existing thesis denominator arithmetic.
\item
  \textbf{Fraud side-study intervals (§F.4).} Sensor-minus-full AUC and
  accuracy intervals are {[}−0.0320, −0.0174{]} and {[}−0.0160,
  −0.0109{]}, reversing and reordering the source's opposite-direction
  intervals.
\end{itemize}



\chapter{Applied Side Study: Credit-Card Fraud Detection
(P8)}\label{applied-side-study-credit-card-fraud-detection-p8}

\textbf{Scope statement.} This appendix reports a side study. It uses
synthetic data and sits outside the RL contribution of this thesis. The
dataset is generated with scikit-learn's \texttt{make\_classification};
the language-model arm is trained by supervised fine-tuning, not by
reinforcement learning; and no result in this appendix is used to
support any GRPO, PPO, ZVF, group-size or campaign claim elsewhere in
the report. The study is retained, in condensed form, as the P8 slot of
the internal report series and as one worked example of measuring a
claimed capability rather than accepting its label.

The study compares a gradient-boosted tree with a fine-tuned language
model on credit-card fraud detection. Its conclusion is that the tree
keeps the real-time scoring seat, and that the language model's
plausible roles lie in tasks the tree cannot perform at all.

\section{Dataset, splits, and the two
arms}\label{dataset-splits-and-the-two-arms}

The dataset is synthetic: \texttt{make\_classification} (seed 42),
50,000 transactions, 20 anonymised numeric features V1--V20 (10
informative, 2 redundant), two clusters per class, class separation 0.8
and \texttt{flip\_y=0.01}, with a realised fraud rate of 1.44\%. Four
per-row aggregates (mean, standard deviation, maximum, minimum) are
appended, giving 24 features. The data reproduce class imbalance and
anonymised feature shape, and none of the temporal drift, verification
latency or adversarial adaptation of real card-fraud data.

\begin{longtable}[]{@{}
  >{\raggedright\arraybackslash\hspace{0pt}}p{(\linewidth - 6\tabcolsep) * \real{0.3333}}
  >{\raggedright\arraybackslash\hspace{0pt}}p{(\linewidth - 6\tabcolsep) * \real{0.3333}}
  >{\raggedright\arraybackslash\hspace{0pt}}p{(\linewidth - 6\tabcolsep) * \real{0.3333}}@{}}
\caption{The two arms of P8.}\tabularnewline
\toprule\noalign{}
\begin{minipage}[b]{\linewidth}\raggedright
Arm
\end{minipage} & \begin{minipage}[b]{\linewidth}\raggedright
Configuration
\end{minipage} & \begin{minipage}[b]{\linewidth}\raggedright
Evaluation split
\end{minipage} \\
\midrule\noalign{}
\endfirsthead
\toprule\noalign{}
\begin{minipage}[b]{\linewidth}\raggedright
Arm
\end{minipage} & \begin{minipage}[b]{\linewidth}\raggedright
Configuration
\end{minipage} & \begin{minipage}[b]{\linewidth}\raggedright
Evaluation split
\end{minipage} \\
\midrule\noalign{}
\endhead
\bottomrule\noalign{}
\endlastfoot
XGBoost (scorer) & 200 estimators, depth 6, lr 0.05, subsample/colsample
0.8, \longtok{scale\_\allowbreak{}pos\_\allowbreak{}weight=7}; no search & 80/20 stratified split,
10,000 test rows, 144 positives \\
Qwen3.5-4B SFT & TabLLM-style row serialisation; answer-token
cross-entropy, no reference policy, no KL; 63 steps at batch 32 &
500-row positive-enriched held-out subset at 20\% fraud, disjoint from
training \\
\end{longtable}

\emph{Sources:
platform\_hybrid/paper/archive/absorbed/P08\_fraud/sections/p8\_setup.tex;
platform\_hybrid/experiments/results/quick\_20260704/qp8-fraud-sft\_manifest.json;
platform\_hybrid/experiments/results/quick\_20260704/qp8-fraud-sft.tsv.}

\section{The head-to-head result}\label{the-head-to-head-result}

\begin{longtable}[]{@{}lllll@{}}
\caption{Head-to-head result on the synthetic card-fraud set, by arm and
evaluation split.}\tabularnewline
\toprule\noalign{}
Arm & Eval split & \emph{n} & Accuracy & ROC-AUC \\
\midrule\noalign{}
\endfirsthead
\toprule\noalign{}
Arm & Eval split & \emph{n} & Accuracy & ROC-AUC \\
\midrule\noalign{}
\endhead
\bottomrule\noalign{}
\endlastfoot
XGBoost (scorer) & full test & 10,000 & n/a & 0.7955 \\
Qwen3.5-4B SFT & positive-enriched test & 500 & 0.7920 & 0.4827 \\
\end{longtable}

\emph{Source:
platform\_hybrid/experiments/results/quick\_20260704/qp8\_fraud.tsv.}

The two AUCs sit on different held-out splits, so this is not a
like-for-like ranking. The language model's AUC of 0.4827 is
chance-level ranking: its Hanley--McNeil 95\% CI, with 100 positives and
400 negatives, is {[}0.42, 0.55{]} and includes 0.5. Its accuracy of
0.792 is slightly \emph{below} the 0.800 that a constant ``legitimate''
answer would score at 20\% prevalence, so it is not evidence of
discrimination. The tree's AUC interval is {[}0.75, 0.84{]}, with 144
positives among 10,000 test rows. The tree's precision 0.723 (34/47,
Wilson 95\% CI {[}0.58, 0.83{]}), recall 0.236 (34/144, {[}0.17,
0.31{]}) and F1 0.356 at the default threshold come from the released
script's own run (AUC 0.7942, 0.49 s training, about 6 ms inference per
10,000 rows), which uses a slightly different configuration from the
0.7955 artefact. An older record of 21 June 2026 (XGBoost AUC 0.975,
language model 0.948) could not be reconstructed and is not used.

\section{Why the scoring seat stays with the
tree}\label{why-the-scoring-seat-stays-with-the-tree}

Four operational arguments hold even at AUC parity:

\begin{itemize}
\tightlist
\item
  \textbf{Latency.} Card authorisation needs a score within tens of
  milliseconds; the tree takes about 6 ms per 10,000 rows, against
  roughly 300 ms for a language-model pass.
\item
  \textbf{Cost.} The released cost table puts the synchronous
  language-model path at about 35 times the per-transaction cost of the
  tree path.
\item
  \textbf{Calibration.} Downstream threshold and expected-loss rules
  consume the score as a probability; verbalised language-model
  confidences are poorly calibrated.
\item
  \textbf{Prompt injection.} A model that reads attacker-controlled text
  fields in a synchronous decision loop opens an injection channel that
  a tree scoring numbers does not have.
\end{itemize}

\clearpage

\section{Supporting measurements}\label{supporting-measurements}

\begingroup\small

\begin{longtable}[]{@{}
  >{\raggedright\arraybackslash\hspace{0pt}}p{(\linewidth - 6\tabcolsep) * \real{0.3333}}
  >{\raggedright\arraybackslash\hspace{0pt}}p{(\linewidth - 6\tabcolsep) * \real{0.3333}}
  >{\raggedright\arraybackslash\hspace{0pt}}p{(\linewidth - 6\tabcolsep) * \real{0.3333}}@{}}
\caption{Supporting P8 measurements. Bracketed ranges are
paired-bootstrap 95\% intervals. The first row reports sensor minus
full-feature performance; the AUC and accuracy intervals are reversed
and reordered from the source's full-feature-minus-sensor convention.
Source files are under
\protect\longtok{platform\_hybrid/experiments/results/} and
\protect\longtok{platform\_hybrid/paper/archive/absorbed/P08\_fraud/sections/}.}\tabularnewline
\toprule\noalign{}
\begin{minipage}[b]{\linewidth}\raggedright
Check
\end{minipage} & \begin{minipage}[b]{\linewidth}\raggedright
Result
\end{minipage} & \begin{minipage}[b]{\linewidth}\raggedright
Source
\end{minipage} \\
\midrule\noalign{}
\endfirsthead
\toprule\noalign{}
\begin{minipage}[b]{\linewidth}\raggedright
Check
\end{minipage} & \begin{minipage}[b]{\linewidth}\raggedright
Result
\end{minipage} & \begin{minipage}[b]{\linewidth}\raggedright
Source
\end{minipage} \\
\midrule\noalign{}
\endhead
\bottomrule\noalign{}
\endlastfoot
Aggregate-only tree (stand-in for a language-model ``sensor'') vs
24-feature tree & Worse on every axis: AUC −0.0245 {[}−0.0320,
−0.0174{]}, accuracy −0.0135 {[}−0.0160, −0.0109{]}, Brier +0.0109
{[}0.0099, 0.0120{]} & p5p8/p8\_\allowbreak{}headline\_\allowbreak{}cis.tsv \\
Adding the four aggregates to the 20 raw features & AUC +0.0002
{[}−0.0002, 0.0007{]}, within noise & p5p8/p8\_\allowbreak{}headline\_\allowbreak{}cis.tsv \\
Leakage caveat & Tables reporting AUC 0.998--0.999 fit on
\longtok{fraud\_data.csv} and score a subset of it; they are
within-protocol relative comparisons only. The leakage-free split gives
0.7955 & p8\_evidence.tex \\
Cost per fraud caught (\$0.50/alert, \$100 missed-fraud loss) & Raw
features \$2.68; sensor-augmented \$3.28; sensor-only \$14.92; the
sensor never pays for itself in the 25-cell (σ, L) grid &
p8\_evidence.tex \\
Sensor-noise break-even & \$5.29 (oracle), \$26.44 (σ ≤ 0.02), \$36.40
(σ = 0.05) & p8\_\allowbreak{}noisy\_\allowbreak{}sensor.tex \\
\end{longtable}

\endgroup

\section{Proposed roles for the language
model}\label{proposed-roles-for-the-language-model}

The study argues, without measuring, that the language model's useful
roles are tasks the tree cannot enter: document and image fraud through
a vision-language model, drafting compliance narratives (Suspicious
Activity Reports under 31 C.F.R. § 1020.320; adverse-action reasons
under 12 C.F.R. § 1002.9), cold-start triage of new fraud typologies
before labels exist, and post-score investigation of the alert queue.
The proposed architecture keeps XGBoost as the scorer and uses the
language model as a sensor (typed features, no authority) and a scribe
(narratives conditioned on the tree's own attributions, with a human
signature). The ``85 times cheaper than an analyst'' triage figure is an
internal estimate from June-2026 token prices and assumed investigation
times, and has not been validated.

\section{Limitations}\label{limitations}

The dataset is synthetic and single-generator, so no number here is a
production expectation, and there is no cross-institution or
public-benchmark replication. The challenger is one model family, one
serialisation and one fine-tuning budget. Of the proposed roles, only
the scorer comparison carries a measurement of the study's own; the
others rest on external literature and regulatory text, and nothing in
the compliance discussion is legal advice.



\chapter{C1 Evidence and Case Ledger}\label{c1-evidence-and-case-ledger}

Benchmark reference errors, ambiguous questions and equivalent units
limit what this diagnostic can establish. Its frozen parser-defined
result is \textbf{11/64 questions (17.19\%; Wilson 95\% interval
9.88-28.21\%)}. Assistant-assisted reconciliation corroborates
\textbf{three numeric recovery witnesses on the same 64-question
denominator}. This 3/64 count (4.69\%) describes the support found under
a conservative review rule; it does not replace the frozen endpoint or
estimate a population recovery rate. The other eight parser events
include three convention-dependent interpretations, four materially
ambiguous cases and one match to an incorrect reference answer.
{[}S1-S4{]}

The study contains no training update. It does not measure useful
reinforcement-learning gradients, reasoning improvement, controller
superiority, held-out accuracy or general model capability. The prompts
are public GSM8K training questions and are fresh only to this campaign.
Historical V2 evidence remains separate and unpooled.

\section{Event ledger}\label{event-ledger}

The following ledger accounts for every frozen primary event. C labels
use the zero-based selected index. Fresh indices identify all
parser-success outputs among the 24 fresh draws. Complete question IDs,
reference hashes, response hashes and review notes accompany the
editable case ledger. {[}S3{]}

\begingroup\scriptsize\setlength{\tabcolsep}{3pt}
\begin{longtable}[]{@{}
  >{\raggedright\arraybackslash\hspace{0pt}}p{(\linewidth - 8\tabcolsep) * \real{0.2308}}
  >{\raggedleft\arraybackslash}p{(\linewidth - 8\tabcolsep) * \real{0.3077}}
  >{\raggedright\arraybackslash\hspace{0pt}}p{(\linewidth - 8\tabcolsep) * \real{0.2308}}
  >{\raggedright\arraybackslash\hspace{0pt}}p{(\linewidth - 8\tabcolsep) * \real{0.2308}}@{}}
\caption{Complete accounting of all eleven parser-defined
events.}\tabularnewline
\toprule\noalign{}
\begin{minipage}[b]{\linewidth}\raggedright
Case and source rank
\end{minipage} & \begin{minipage}[b]{\linewidth}\raggedleft
Reference
\end{minipage} & \begin{minipage}[b]{\linewidth}\raggedright
Fresh success indices
\end{minipage} & \begin{minipage}[b]{\linewidth}\raggedright
Interpretation
\end{minipage} \\
\midrule\noalign{}
\endfirsthead
\toprule\noalign{}
\begin{minipage}[b]{\linewidth}\raggedright
Case and source rank
\end{minipage} & \begin{minipage}[b]{\linewidth}\raggedleft
Reference
\end{minipage} & \begin{minipage}[b]{\linewidth}\raggedright
Fresh success indices
\end{minipage} & \begin{minipage}[b]{\linewidth}\raggedright
Interpretation
\end{minipage} \\
\midrule\noalign{}
\endhead
\bottomrule\noalign{}
\endlastfoot
C00 Santana presents, 676 & 8 & 12, 15, 23 & Corroborated \\
C07 Suzanne bread, 912 & 26 & 2, 5, 10, 15, 16, 21 & Corroborated \\
C08 Citations, 923 & 4 & 0, 1, 3, 15, 23 & Convention / disagreement \\
C19 Pet owners, 1482 & 100 & 4, 16 & Ambiguous \\
C20 Raft capacity, 1544 & 17 & 23 & Convention / disagreement \\
C29 Fundraiser, 2044 & 475 & 2, 4, 7, 13, 21, 22 & Convention /
disagreement \\
C35 Checkpoints, 2593 & 6 & 6, 7, 10, 22 & Wrong reference \\
C37 Walking distance, 2660 & 36 & 0, 22 & Ambiguous \\
C38 Hank reading, 2699 & 810 & 2 & Corroborated \\
C47 Cake and ham, 3203 & 12 & 0, 4, 5, 9, 22 & Ambiguous \\
C50 Reggie marbles, 3255 & 1 & 2, 15, 22 & Ambiguous \\
\end{longtable}
\endgroup

\textbf{C00 Santana presents.} Gold 8 is (4 second-half birthdays+7
Christmas gifts)-3 March gifts. Both reviews support the quantity. Fresh
12 matches 8 through a false 8-0 argument; fresh 15 and 23 give the
correct 11-3 calculation.

\textbf{C07 Suzanne bread.} 3 slices per meal × 2 weekend meals × 52
weeks/12 slices per loaf = 26. Sharing one slice does not require
knowing the number of daughters. All 8 initial 90/91 claims use seven
days per week; every fresh 26 uses the correct weekend schedule.

\textbf{C08 Citations.} Other citations most naturally means the
combined non-parking count: L+L+2(L+L) = 24 gives 4. All initial 6
claims substitute twice one category. First treats that as material
ambiguity; second supports 4. Ordinary-reading support is stronger than
the per-category alternative, but disagreement remains and this is not
added to the corroborated recovery count.

\textbf{C19 Pet owners.} Half of pet owners plus disjoint exhaustive
dog/cat ownership gives 100; half of all citizens gives 300 under the
same partition. Overlap allows other values. The denominator ambiguity
is material even if ordinary classroom partition assumptions are
granted.

\textbf{C20 Raft capacity.} The ordinary exact additive-space model
makes 14 jackets occupy 7 person-spaces; 8 jackets reduce 21 by 4,
yielding 17. Arbitrary real packing/safety possibilities do not by
themselves constitute a defect. First conservatively withheld support;
second accepts the conventional model. Retain that disagreement and
conditional event support only in the ledger.

\textbf{C29 Fundraiser.} Next 4 days naturally names a single comparison
period, making 950 the four-day total and shortfall 475 per student.
Initial 0 answers insert per day before multiplying by 4; the text does
not say each day. Retain the first alternative-scope objection and
second conventional support, not a revised recovery estimate.

\textbf{C35 Checkpoints.} Interior route checkpoints at mile 1 and mile
25 enclose 24 miles and 3 intervening gaps, hence 8 miles. The closing
2-mile arc is outside those gaps. All 8 initials give 8; all 4 gold
matches 6 divide by 4 and are wrong. Screen 1 has an internally
contradictory parenthetical despite its correct final 8.

\textbf{C37 Walking distance.} Summed person-distance is 5 × 3 × 6+12+6
= 108; shared outings counted once give 36. The aggregate quantity is
genuinely different. Initial 0-6 are correct under person-mile reading;
initial 7 incorrectly repeats weekly extras and gets 198. Neither
gold-matching response proves the 108 answers wrong.

\textbf{C38 Hank reading.} 5 × (30+60)+2 × 2 × (30+60) = 810 minutes.
Reference final sentence mislabels subtotals but final 810 is correct.
Initial 630/690/1350 claims are wrong; fresh 2 visibly supplies the
correct 450+360 calculation.

\textbf{C47 Cake and ham.} Cake-day overlap with any ham day during the
week has probability 3/5 = 60\%. A specified/randomly sampled test day
has joint probability 3/25 = 12\%. Both reviewer rounds retain
sampling-unit ambiguity; all 5 fresh 12 matches use the latter event.

\textbf{C50 Reggie marbles.} With starting Reggie count R and losses
L,90 = R+10 × (9-2 L), so R = 20 L. His friend's 100 does not set R.
Multiple feasible starts/losses exist. Gold's net-loss shortcut is
invalid; all 3 fresh 1 matches have defective justifications. There is
no uniquely correct loss count.

\section{Complete reference classification
index}\label{complete-reference-classification-index}

This index preserves all 64 selected questions, including non-events. W
= second-review wrong final value; A = material ambiguity; I =
inconsistent premises; R = correct final value with flawed reference
explanation; S = supported; C = supported with conventional assumptions.
The category column uses the refined mutually exclusive mapping. An
asterisk marks a retained first/second category disagreement. E marks a
frozen parser event; N marks a conservative numeric witness. Neither
marker changes the reference classification. {[}S2, S4{]}

\begingroup\scriptsize\setlength{\tabcolsep}{3pt}
\begin{longtable}[]{@{}
  >{\raggedright\arraybackslash\hspace{0pt}}p{(\linewidth - 10\tabcolsep) * \real{0.1429}}
  >{\raggedleft\arraybackslash}p{(\linewidth - 10\tabcolsep) * \real{0.1905}}
  >{\raggedleft\arraybackslash}p{(\linewidth - 10\tabcolsep) * \real{0.1905}}
  >{\centering\arraybackslash}p{(\linewidth - 10\tabcolsep) * \real{0.2381}}
  >{\centering\arraybackslash}p{(\linewidth - 10\tabcolsep) * \real{0.2381}}@{}}
\caption{Complete bounded-cohort index. The JSON and CSV ledgers retain
full identifiers, both review labels and case-specific
explanations.}\tabularnewline
\toprule\noalign{}
\begin{minipage}[b]{\linewidth}\raggedright
Selected case
\end{minipage} & \begin{minipage}[b]{\linewidth}\raggedleft
Source rank
\end{minipage} & \begin{minipage}[b]{\linewidth}\raggedleft
Recorded gold
\end{minipage} & \begin{minipage}[b]{\linewidth}\centering
Category
\end{minipage} & \begin{minipage}[b]{\linewidth}\centering
Event and witness
\end{minipage} \\
\midrule\noalign{}
\endfirsthead
\toprule\noalign{}
\begin{minipage}[b]{\linewidth}\raggedright
Selected case
\end{minipage} & \begin{minipage}[b]{\linewidth}\raggedleft
Source rank
\end{minipage} & \begin{minipage}[b]{\linewidth}\raggedleft
Recorded gold
\end{minipage} & \begin{minipage}[b]{\linewidth}\centering
Category
\end{minipage} & \begin{minipage}[b]{\linewidth}\centering
Event and witness
\end{minipage} \\
\midrule\noalign{}
\endhead
\bottomrule\noalign{}
\endlastfoot
C00 & 676 & 8 & S & E N \\
C01 & 706 & 42 & C & - \\
C02 & 754 & 200 & W* & - \\
C03 & 803 & 200 & A & - \\
C04 & 846 & 40 & A & - \\
C05 & 855 & 625 & S* & - \\
C06 & 878 & 480 & W & - \\
C07 & 912 & 26 & S & E N \\
C08 & 923 & 4 & S* & E \\
C09 & 925 & 50 & I & - \\
C10 & 945 & 1886 & S & - \\
C11 & 950 & 41 & W & - \\
C12 & 1078 & 2460 & W & - \\
C13 & 1139 & 25 & W & - \\
C14 & 1227 & 95 & C* & - \\
C15 & 1266 & 49 & W & - \\
C16 & 1307 & 301 & S* & - \\
C17 & 1357 & 560 & W & - \\
C18 & 1377 & 96 & W & - \\
C19 & 1482 & 100 & A & E \\
C20 & 1544 & 17 & C* & E \\
C21 & 1739 & 750 & A & - \\
C22 & 1743 & 50 & W & - \\
C23 & 1780 & 3 & A & - \\
C24 & 1872 & 35 & S & - \\
C25 & 1897 & 2100 & C* & - \\
C26 & 1921 & 23 & W & - \\
C27 & 1932 & 140 & A & - \\
C28 & 1996 & 99 & R & - \\
C29 & 2044 & 475 & C* & E \\
C30 & 2095 & 4 & W & - \\
C31 & 2142 & 540 & A & - \\
C32 & 2211 & 42 & A & - \\
C33 & 2446 & 280 & W & - \\
C34 & 2514 & 3 & A* & - \\
C35 & 2593 & 6 & W & E \\
C36 & 2617 & 24 & C & - \\
C37 & 2660 & 36 & A & E \\
C38 & 2699 & 810 & R & E N \\
C39 & 2748 & 56 & A & - \\
C40 & 2769 & 40 & C & - \\
C41 & 2899 & 400 & A & - \\
C42 & 2989 & 1100 & A & - \\
C43 & 3021 & 9 & C* & - \\
C44 & 3089 & 8000 & A & - \\
C45 & 3118 & 350 & R & - \\
C46 & 3119 & 163 & S & - \\
C47 & 3203 & 12 & A & E \\
C48 & 3209 & 1000 & A & - \\
C49 & 3216 & 5 & A* & - \\
C50 & 3255 & 1 & A & E \\
C51 & 3259 & 983 & I & - \\
C52 & 3277 & 84 & C & - \\
C53 & 3299 & 128 & W & - \\
C54 & 3335 & 30 & W & - \\
C55 & 3409 & 7 & R* & - \\
C56 & 3421 & 57 & S & - \\
C57 & 3464 & 90 & W & - \\
C58 & 3533 & 25 & C* & - \\
C59 & 3536 & 6 & S & - \\
C60 & 3579 & 40 & S & - \\
C61 & 3581 & 4 & A & - \\
C62 & 3593 & 400 & S & - \\
C63 & 3597 & 52 & C & - \\
\end{longtable}
\endgroup

\section{Public evidence
availability}\label{public-evidence-availability}

The included case ledger preserves all sixty-four classifications, all
eleven parser-defined events, scientific hashes and review
disagreements. Source labels S1-S9 refer to frozen scientific summaries
and reviews; the public reference-binding file records their
descriptions and checksums without private operational paths. These
checksum references do not assert that every underlying payload is
publicly included. Operational closeout records, runtime/session
identifiers, raw archive transfer bindings and allocation details are
excluded. S12 and S13 identify the retained generation configuration and
scientific manifest; this public derivative reports their scientific
settings in the methodology without publishing their private transfer
wrappers.



\chapter{Source Reference Map}\label{source-reference-map}

This map collects source references retained in the sanitized
manuscript. A reference is not a promise that its payload is bundled or
publicly available. Repository-relative scientific paths are retained
where safe; private operational identifiers and source-export paths are
omitted. The public availability manifest distinguishes included
document sources from withheld evidence categories. Original records
have not been regenerated to fill gaps, and no complete public
provenance or universal experiment-coverage claim is made. Table sources
remain under their tables.

\section{Chapter 1: Introduction}\label{chapter-1-introduction}

\textbf{1.1 Background: critic-free RL post-training and the
group-relative estimator}

\begin{itemize}
\tightlist
\item
  \longtok{platform\_tinker/reports/esa\_phase1/Phase1\_\allowbreak{}Project\_\allowbreak{}Report\_\allowbreak{}ZVF.tex}
\item
  \longtok{platform\_hybrid/paper/sections/group\_size.tex}
\item
  \longtok{platform\_hybrid/paper/sections/zvf\_iter122.tex}
\item
  \longtok{platform\_hybrid/paper/main\_\allowbreak{}eai\_\allowbreak{}body.tex}
\end{itemize}

\textbf{1.2 A Motivating Example}

\begin{itemize}
\tightlist
\item
  \longtok{platform\_tinker/reports/esa\_phase1/Phase1\_\allowbreak{}Project\_\allowbreak{}Report\_\allowbreak{}ZVF.tex}
\item
  \longtok{platform\_hybrid/paper/main\_\allowbreak{}eai\_\allowbreak{}body.tex}
\item
  \longtok{platform\_hybrid/paper/sections/\_shared\_methods.tex}
\item
  \longtok{autoresearch/deli-neurips-tmlr-260802/audits/statistical\_refutation.md}
\end{itemize}

\textbf{1.3 Problem Statement}

\begin{itemize}
\tightlist
\item
  \longtok{platform\_hybrid/paper/main\_\allowbreak{}eai\_\allowbreak{}body.tex}
\item
  \longtok{platform\_hybrid/paper/sections/\_shared\_methods.tex}
\item
  \longtok{platform\_hybrid/experiments/results/samestack\_\allowbreak{}ppo\_\allowbreak{}grpo.json}
\item
  \longtok{platform\_tinker/reports/esa\_phase1/Phase1\_\allowbreak{}Project\_\allowbreak{}Report\_\allowbreak{}ZVF.tex}
\end{itemize}

\textbf{1.5 Scope and Non-Goals}

\begin{itemize}
\tightlist
\item
  \longtok{outputs/wrong\_\allowbreak{}cohort\_\allowbreak{}2026-10-02/continuation\_c1/final\_analysis/summary.json}
\item
  \longtok{outputs/wrong\_\allowbreak{}cohort\_\allowbreak{}2026-10-02/continuation\_c1/assistant\_audit/final\_reconciliation/final\_\allowbreak{}evidence\_\allowbreak{}summary.json}
\end{itemize}

\textbf{1.6 Contribution of the Candidate}

\begin{itemize}
\tightlist
\item
  \texttt{platform\_hybrid/sem} 4 \texttt{work/EXPERIMENTS.md}
\item
  \texttt{platform\_hybrid/sem} 4 \texttt{work/README.md}
\item
  \texttt{platform\_hybrid/sem} 4 \texttt{work/PROVENANCE.md}
\item
  \longtok{outputs/E1\_\allowbreak{}E14\_\allowbreak{}FINAL\_\allowbreak{}RESULTS\_\allowbreak{}2026-09-19.md}
\item
  \longtok{autoresearch/deli-neurips-tmlr-260802/audits/statistical\_refutation.md}
\end{itemize}

\section{Chapter 2: Literature
Survey}\label{chapter-2-literature-survey}

\begin{itemize}
\tightlist
\item
  \longtok{platform\_tinker/reports/esa\_phase1/Phase1\_\allowbreak{}Project\_\allowbreak{}Report\_\allowbreak{}ZVF.tex}
\end{itemize}

\textbf{2.1 From RLHF and PPO to Group-Relative Methods}

\begin{itemize}
\tightlist
\item
  \longtok{capstone-literature-survey/section\_\allowbreak{}2\_\allowbreak{}2\_\allowbreak{}rlhf.tex}
\item
  \longtok{capstone-literature-survey/section\_\allowbreak{}2\_\allowbreak{}3\_\allowbreak{}ppo.tex}
\item
  \longtok{capstone-literature-survey/section\_\allowbreak{}2\_\allowbreak{}5\_\allowbreak{}grpo.tex}
\item
  \longtok{platform\_tinker/reports/esa\_phase1/Phase1\_\allowbreak{}Project\_\allowbreak{}Report\_\allowbreak{}ZVF.tex}
\item
  \longtok{capstone-literature-survey/section\_\allowbreak{}2\_\allowbreak{}6\_\allowbreak{}deepseek\_\allowbreak{}r1.tex}
\end{itemize}

\textbf{2.2 Direct Preference Optimisation as a Contrast}

\begin{itemize}
\tightlist
\item
  \longtok{capstone-literature-survey/section\_\allowbreak{}2\_\allowbreak{}4\_\allowbreak{}dpo.tex}
\item
  \longtok{platform\_hybrid/paper/sections/related\_\allowbreak{}work\_\allowbreak{}v2.tex}
\item
  \longtok{platform\_hybrid/paper/sections/p7\_intro.tex}
\end{itemize}

\textbf{2.3 Fixing GRPO's Biases and Collapse}

\begin{itemize}
\tightlist
\item
  \longtok{platform\_hybrid/paper/sections/p4\_intro.tex}
\item
  \longtok{platform\_hybrid/paper/sections/p6\_intro.tex}
\item
  \longtok{platform\_hybrid/paper/sections/related\_\allowbreak{}work\_\allowbreak{}v2.tex}
\item
  \longtok{platform\_hybrid/paper/sections/p7\_intro.tex}
\item
  \longtok{platform\_tinker/reports/esa\_phase1/Phase1\_\allowbreak{}Project\_\allowbreak{}Report\_\allowbreak{}ZVF.tex}
\item
  \longtok{platform\_hybrid/paper/sections/variance\_\allowbreak{}mitigation\_\allowbreak{}comparison.tex}
\item
  \longtok{platform\_hybrid/experiments/results/drgrpo\_\allowbreak{}vs\_\allowbreak{}grpo.json}
\item
  \longtok{platform\_hybrid/experiments/results/drgrpo\_\allowbreak{}gsm8k\_\allowbreak{}cot.json}
\end{itemize}

\textbf{2.4 Verifiable Rewards, Frameworks and Reporting Practice}

\begin{itemize}
\tightlist
\item
  \longtok{platform\_tinker/reports/esa\_phase1/Phase1\_\allowbreak{}Project\_\allowbreak{}Report\_\allowbreak{}ZVF.tex}
\item
  \longtok{platform\_hybrid/paper/sections/related\_\allowbreak{}work\_\allowbreak{}v2.tex}
\item
  \longtok{capstone-literature-survey/section\_\allowbreak{}2\_\allowbreak{}8\_\allowbreak{}mathematical\_\allowbreak{}reasoning\_\allowbreak{}benchmarks.tex}
\item
  \longtok{platform\_hybrid/paper/sections/p5\_related.tex}
\item
  \longtok{platform\_hybrid/experiments/results/p5p8/p5\_\allowbreak{}iter109\_\allowbreak{}minreport\_\allowbreak{}coupling.tsv}
\item
  \longtok{platform\_hybrid/experiments/results/p5p8/p5\_\allowbreak{}iter109\_\allowbreak{}crossref\_\allowbreak{}verify.tsv}
\item
  \longtok{platform\_hybrid/experiments/results/p5p8/p5\_\allowbreak{}iter109\_\allowbreak{}bib\_\allowbreak{}audit.tsv}
\item
  \longtok{platform\_hybrid/paper/sections/p6\_intro.tex}
\item
  \longtok{platform\_hybrid/paper/sections/p6\_\allowbreak{}related\_\allowbreak{}resources.tex}
\end{itemize}

\textbf{2.5 Scaling and Layer-Adaptation Context}

\begin{itemize}
\tightlist
\item
  \longtok{platform\_hybrid/paper/sections/related\_\allowbreak{}work\_\allowbreak{}v2.tex}
\item
  \longtok{platform\_hybrid/paper/sections/p1\_intro.tex}
\item
  \longtok{platform\_hybrid/paper/sections/p2\_intro.tex}
\item
  \longtok{platform\_hybrid/paper/sections/p3\_intro.tex}
\item
  \longtok{platform\_hybrid/paper/sections/p4\_intro.tex}
\item
  \longtok{platform\_hybrid/paper/sections/frontier\_\allowbreak{}synthesis\_\allowbreak{}scaling.tex}
\item
  \longtok{platform\_hybrid/paper/sections/p1\_conclusion.tex}
\item
  \longtok{capstone-literature-survey/section2\_\allowbreak{}7\_\allowbreak{}peft.tex}
\item
  \longtok{platform\_tinker/reports/esa\_phase1/Phase1\_\allowbreak{}Project\_\allowbreak{}Report\_\allowbreak{}ZVF.tex}
\end{itemize}

\textbf{2.6 Process-, Tree- and Stability-Aware Variants}

\begin{itemize}
\tightlist
\item
  \longtok{platform\_hybrid/paper/sections/related\_\allowbreak{}work\_\allowbreak{}v2.tex}
\item
  \longtok{platform\_tinker/reports/esa\_phase1/Phase1\_\allowbreak{}Project\_\allowbreak{}Report\_\allowbreak{}ZVF.tex}
\item
  \longtok{platform\_hybrid/paper/sections/p7\_intro.tex}
\end{itemize}

\textbf{2.7 Positioning the Portfolio}

\begin{itemize}
\tightlist
\item
  \longtok{platform\_hybrid/paper/sections/p6\_\allowbreak{}related\_\allowbreak{}resources.tex}
\item
  \longtok{platform\_tinker/reports/esa\_phase1/Phase1\_\allowbreak{}Project\_\allowbreak{}Report\_\allowbreak{}ZVF.tex}
\item
  \longtok{outputs/E1\_\allowbreak{}E14\_\allowbreak{}FINAL\_\allowbreak{}RESULTS\_\allowbreak{}2026-09-19.md}
\item
  \longtok{outputs/e1\_\allowbreak{}e14\_\allowbreak{}small\_\allowbreak{}scale\_\allowbreak{}2026-09-26/E4/BRIDGE\_\allowbreak{}STOP\_\allowbreak{}TOKEN\_\allowbreak{}FINDING.md}
\item
  \longtok{platform\_hybrid/paper/sections/related\_\allowbreak{}work\_\allowbreak{}v2.tex}
\end{itemize}

\section{Chapter 3: System Requirements
Specification}\label{chapter-3-system-requirements-specification}

\textbf{3.1 Overview and Actors}

\begin{itemize}
\tightlist
\item
  \longtok{platform\_tinker/reports/esa\_phase1/Phase1\_\allowbreak{}Project\_\allowbreak{}Report\_\allowbreak{}ZVF.tex}
\item
  \longtok{zvf-program/flagship/e1\_\allowbreak{}e14\_\allowbreak{}completion\_\allowbreak{}gate.py}
\end{itemize}

\textbf{3.4 Test-Case Specification}

\begin{itemize}
\tightlist
\item
  \longtok{platform\_tinker/reports/esa\_phase1/Phase1\_\allowbreak{}Project\_\allowbreak{}Report\_\allowbreak{}ZVF.tex}
\item
  \longtok{platform\_local/run\_\allowbreak{}all\_\allowbreak{}audits.py}
\item
  \texttt{utils/audit\_utils.py}
\item
  \longtok{platform\_local/claim\_\allowbreak{}strength\_\allowbreak{}audit.py}
\item
  \longtok{platform\_local/paper\_\allowbreak{}sync\_\allowbreak{}audit.py}
\item
  \longtok{platform\_local/export\_\allowbreak{}guard\_\allowbreak{}audit.py}
\item
  \longtok{platform\_local/scientific\_audit.py}
\item
  \longtok{zvf-program/flagship/e1\_\allowbreak{}e14\_\allowbreak{}completion\_\allowbreak{}gate.py}
\item
  \longtok{platform\_hybrid/paper/paper\_\allowbreak{}P6\_\allowbreak{}registry.tex}
\item
  \longtok{platform\_hybrid/registry/query.py}
\end{itemize}

\textbf{3.7 Software Requirements}

\begin{itemize}
\tightlist
\item
  \longtok{platform\_tinker/reports/esa\_phase1/Phase1\_\allowbreak{}Project\_\allowbreak{}Report\_\allowbreak{}ZVF.tex}
\item
  \longtok{platform\_local/run\_\allowbreak{}all\_\allowbreak{}audits.py}
\item
  \texttt{utils/audit\_utils.py}
\item
  \longtok{zvf-program/flagship/e1\_\allowbreak{}e14\_\allowbreak{}completion\_\allowbreak{}gate.py}
\item
  \longtok{platform\_hybrid/registry/query.py}
\item
  \longtok{platform\_local/scientific\_audit.py}
\item
  \longtok{outputs/PES\_\allowbreak{}Phase2\_\allowbreak{}Third\_\allowbreak{}Review\_\allowbreak{}2026-09-24/thesis/assemble\_thesis.py}
\item
  \longtok{outputs/E1\_\allowbreak{}E14\_\allowbreak{}FINAL\_\allowbreak{}RESULTS\_\allowbreak{}2026-09-19.md}
\end{itemize}

\textbf{3.8 Constraints and Assumptions}

\begin{itemize}
\tightlist
\item
  \longtok{platform\_tinker/reports/esa\_phase1/Phase1\_\allowbreak{}Project\_\allowbreak{}Report\_\allowbreak{}ZVF.tex}
\item
  \longtok{platform\_hybrid/paper/paper\_\allowbreak{}P7\_\allowbreak{}zvf\_\allowbreak{}controller.tex}
\item
  \longtok{platform\_hybrid/paper/sections/\_shared\_methods.tex}
\item
  \longtok{platform\_hybrid/experiments/results/framework\_comparison.json}
\item
  \longtok{platform\_hybrid/paper/paper\_\allowbreak{}P1\_\allowbreak{}scaling.tex}
\item
  \longtok{outputs/E1\_\allowbreak{}E14\_\allowbreak{}FINAL\_\allowbreak{}RESULTS\_\allowbreak{}2026-09-19.md}
\item
  \longtok{platform\_hybrid/paper/paper\_\allowbreak{}P6\_\allowbreak{}registry.tex}
\end{itemize}

\section{Chapter 4: Proposed
Methodology}\label{chapter-4-proposed-methodology}

\textbf{4.1 The Design Commitment to Attribution}

\begin{itemize}
\tightlist
\item
  \longtok{platform\_tinker/reports/esa\_phase1/Phase1\_\allowbreak{}Project\_\allowbreak{}Report\_\allowbreak{}ZVF.tex}
\item
  \texttt{platform\_hybrid/sem} 4
  \longtok{work/submissions/mtech-final-review/build\_\allowbreak{}review\_\allowbreak{}pptx.py},
  slide 7
\item
  \longtok{platform\_hybrid/experiments/results/samestack\_\allowbreak{}ppo\_\allowbreak{}grpo.json}
\item
  \longtok{platform\_hybrid/paper/sections/\_shared\_methods.tex}, its
  status is recorded in §6.2.3
\item
  \longtok{platform\_hybrid/paper/neurips\_\allowbreak{}2026\_\allowbreak{}variants/paper\_\allowbreak{}P8\_\allowbreak{}workshop.tex}
\item
  \longtok{platform\_hybrid/paper/sections/\_shared\_methods.tex}
\item
  \longtok{platform\_hybrid/paper/sections/\_shared\_methods.tex}, §6.1
\end{itemize}

\textbf{4.2 Per-Step Telemetry}

\begin{itemize}
\tightlist
\item
  \longtok{platform\_tinker/reports/esa\_phase1/Phase1\_\allowbreak{}Project\_\allowbreak{}Report\_\allowbreak{}ZVF.tex}
\item
  \longtok{platform\_hybrid/paper/sections/p2\_abstract.tex}
\end{itemize}

\textbf{4.3 Formalising the Zero-Variance Fraction}

\begin{itemize}
\tightlist
\item
  \longtok{platform\_hybrid/paper/sections/appendix\_\allowbreak{}zvf\_\allowbreak{}formalization.tex}
\item
  \longtok{platform\_hybrid/paper/neurips\_\allowbreak{}2026\_\allowbreak{}variants/paper\_\allowbreak{}P8\_\allowbreak{}workshop.tex}
\item
  \longtok{platform\_hybrid/paper/sections/zvf\_\allowbreak{}pipeline\_\allowbreak{}spec.tex}
\item
  \longtok{platform\_hybrid/paper/sections/p2\_abstract.tex}
\item
  \longtok{platform\_hybrid/paper/sections/appendix\_\allowbreak{}zvf\_\allowbreak{}formalization.tex},
  results in Chapter 6 §6.3.3
\item
  \longtok{platform\_hybrid/paper/sections/p7\_abstract.tex}
\item
  \longtok{platform\_hybrid/paper/sections/p2\_abstract.tex}, results in
  Chapter 6 §6.3.4
\item
  \longtok{platform\_hybrid/paper/sections/zvf.tex}, §6.3.2
\end{itemize}

\textbf{4.4 The Eight-Study Portfolio}

\begin{itemize}
\tightlist
\item
  \texttt{platform\_hybrid/sem} 4 \texttt{work/README.md}
\item
  \longtok{platform\_hybrid/paper/sections/p1\_abstract.tex}
\item
  \longtok{platform\_hybrid/paper/sections/p2\_abstract.tex}
\item
  \longtok{platform\_hybrid/paper/sections/p3\_abstract.tex}
\item
  \longtok{platform\_hybrid/paper/sections/p4\_abstract.tex}
\item
  \longtok{platform\_hybrid/paper/sections/p5\_abstract.tex}
\item
  \longtok{platform\_hybrid/paper/sections/p6\_abstract.tex}
\item
  \longtok{platform\_hybrid/paper/sections/p7\_abstract.tex}
\item
  \texttt{platform\_hybrid/sem} 4
  \longtok{work/submissions/mtech-final-review/README.md}
\item
  \longtok{platform\_hybrid/paper/archive/absorbed/P08\_fraud/sections/p8\_abstract.tex}
\item
  \longtok{platform\_hybrid/paper/PAPERS\_README.md}
\end{itemize}

\textbf{4.5 The Design of the E1--E14 Evaluation Campaign}

\begin{itemize}
\tightlist
\item
  \longtok{outputs/E1\_\allowbreak{}E14\_\allowbreak{}FINAL\_\allowbreak{}RESULTS\_\allowbreak{}2026-09-19.md}
\item
  \longtok{outputs/public\_\allowbreak{}portfolio\_\allowbreak{}2026-09-05/agentdojo\_\allowbreak{}native\_\allowbreak{}receipt.json}
\item
  \longtok{outputs/modal\_\allowbreak{}e1\_\allowbreak{}e14/2026-08-16/e11\_\allowbreak{}full\_\allowbreak{}receipt.json}
\item
  \texttt{zvf-program/flagship/}
\item
  \longtok{outputs/public\_\allowbreak{}portfolio\_\allowbreak{}2026-09-05/decontamination/README.md}
\item
  \longtok{zvf-program/flagship/PAVLOV\_\allowbreak{}EXPERIMENT\_\allowbreak{}PROTOCOL\_\allowbreak{}2026-08-09.md}
\end{itemize}

\section{Chapter 5: Implementation
Details}\label{chapter-5-implementation-details}

\textbf{5.1 Two Rosters, One Dispatch Surface}

\begin{itemize}
\tightlist
\item
  \longtok{platform\_hybrid/paper/sections/\_shared\_methods.tex}
\item
  \texttt{platform\_tinker/atropos/}
\item
  \longtok{platform\_tinker/atropos/README.md}
\end{itemize}

\textbf{5.2 The Adapter Architecture: What versus Where}

\begin{itemize}
\tightlist
\item
  \longtok{platform\_local/unified/backends/base.py}
\item
  \longtok{platform\_local/unified/matrix.json}
\item
  \longtok{platform\_local/unified/launcher.py}
\item
  \longtok{platform\_hybrid/experiments/framework\_\allowbreak{}config\_\allowbreak{}dumps/README.md}
\end{itemize}

\textbf{5.3 The Tinker GRPO Loop}

\begin{itemize}
\tightlist
\item
  \longtok{platform\_tinker/reports/esa\_phase1/Phase1\_\allowbreak{}Project\_\allowbreak{}Report\_\allowbreak{}ZVF.tex}
\item
  \longtok{platform\_tinker/tinkerrl/grpo.py},
  \longtok{make\_\allowbreak{}grpo\_\allowbreak{}loss\_\allowbreak{}fn}
\item
  \longtok{platform\_tinker/tinkerrl/grpo\_cli.py}
\end{itemize}

\textbf{5.4 The Parallel Runner and a Real Logging Defect}

\begin{itemize}
\tightlist
\item
  \longtok{platform\_hybrid/experiments/results/p3\_groupsize/FINDINGS.md}
\end{itemize}

\textbf{5.5 The Colab White-Box Path (Per-Layer Gradient Profile)}

\begin{itemize}
\tightlist
\item
  \longtok{platform\_hybrid/experiments/results/p1\_layerfreeze/FINDINGS.md}
\item
  \longtok{platform\_tinker/reports/esa\_phase1/Phase1\_\allowbreak{}Project\_\allowbreak{}Report\_\allowbreak{}ZVF.tex}
\item
  \longtok{platform\_hybrid/experiments/results/p1\_layerfreeze/scaled\_\allowbreak{}4seed\_\allowbreak{}result.json}
\item
  \longtok{platform\_hybrid/experiments/results/p1\_layerfreeze/FREEZE\_FINDINGS.md}
\end{itemize}

\textbf{5.6 Software Modules}

\begin{itemize}
\tightlist
\item
  \longtok{platform\_tinker/reports/esa\_phase1/Phase1\_\allowbreak{}Project\_\allowbreak{}Report\_\allowbreak{}ZVF.tex}
\item
  \texttt{platform\_local/unified/}
\end{itemize}

\textbf{5.7 The E1--E14 Evaluation Harness}

\begin{itemize}
\tightlist
\item
  \longtok{zvf-program/flagship/e1\_\allowbreak{}e14\_\allowbreak{}completion\_\allowbreak{}gate.py}
\item
  \longtok{outputs/E1\_\allowbreak{}E14\_\allowbreak{}FINAL\_\allowbreak{}RESULTS\_\allowbreak{}2026-09-19.md}
\end{itemize}

\textbf{5.8 The Bridge and Sampler Path}

\begin{itemize}
\tightlist
\item
  \longtok{zvf-program/flagship/modal\_\allowbreak{}tinker\_\allowbreak{}openai\_\allowbreak{}bridge.py}
\item
  \longtok{zvf-program/e1\_wave10/README.md}
\end{itemize}

\textbf{5.9 Recovery Re-implementations under \texttt{zvf-program/}}

\begin{itemize}
\tightlist
\item
  \longtok{outputs/E1\_\allowbreak{}E14\_\allowbreak{}FINAL\_\allowbreak{}RESULTS\_\allowbreak{}2026-09-19.md}
\item
  \longtok{zvf-program/e2\_core/launch\_gate.py}
\end{itemize}

\section{Chapter 6: Results I: Core GRPO Studies
(P1--P4)}\label{chapter-6-results-i-core-grpo-studies-p1p4}

\textbf{6.1 Scope, Sources and Evidence Boundary}

\begin{itemize}
\tightlist
\item
  \longtok{platform\_hybrid/paper/sections/p1\_abstract.tex}
\item
  \longtok{platform\_hybrid/paper/sections/zvf.tex}
\item
  \longtok{platform\_hybrid/paper/sections/p3\_intro.tex}
\item
  \longtok{platform\_hybrid/paper/sections/length\_bias.tex}
\item
  \longtok{platform\_hybrid/paper/sections/\_shared\_methods.tex}
\end{itemize}

\textbf{6.2 P1 --- Cross-Library and Cross-Scale GRPO Behaviour}

\begin{itemize}
\tightlist
\item
  \longtok{platform\_hybrid/paper/sections/p1\_abstract.tex}
\end{itemize}

\textbf{6.2.1 The anchor pool and the fitted saturation forms}

\begin{itemize}
\tightlist
\item
  \longtok{platform\_hybrid/paper/sections/p1\_abstract.tex}
\item
  \longtok{platform\_hybrid/paper/sections/scaling\_laws.tex}
\end{itemize}

\textbf{6.2.2 The cross-scale test and its nulls}

\begin{itemize}
\tightlist
\item
  \longtok{platform\_hybrid/paper/sections/scaling\_laws.tex}
\end{itemize}

\textbf{6.2.3 Where the geometric template fails}

\begin{itemize}
\tightlist
\item
  \longtok{platform\_hybrid/paper/sections/p1\_abstract.tex}
\item
  \longtok{platform\_hybrid/paper/sections/\_shared\_methods.tex},
  Chapter 1 §1.2
\end{itemize}

\textbf{6.2.4 A withdrawn architectural contrast}

\begin{itemize}
\tightlist
\item
  \longtok{platform\_hybrid/paper/sections/scaling\_laws.tex}
\item
  provenance in
  \longtok{autoresearch/deli-neurips-tmlr-260802/state/nemotron\_\allowbreak{}architecture\_\allowbreak{}resolution.md}
\item
  \longtok{platform\_hybrid/paper/sections/scaling\_laws.tex}, Table
  \texttt{tab:scaling-iter17-kappa}
\end{itemize}

\textbf{6.3.1 Definition and measurement}

\begin{itemize}
\tightlist
\item
  \longtok{platform\_hybrid/paper/sections/zvf.tex}
\item
  \longtok{platform\_hybrid/paper/sections/p2\_abstract.tex}
\item
  \longtok{platform\_modal/scripts/pcd\_\allowbreak{}vs\_\allowbreak{}zvf.py}
\item
  \longtok{platform\_hybrid/experiments/results/pcd\_\allowbreak{}vs\_\allowbreak{}zvf\_\allowbreak{}tolerance\_\allowbreak{}analysis.json}
\item
  \longtok{platform\_hybrid/experiments/results/pcd\_\allowbreak{}vs\_\allowbreak{}zvf\_\allowbreak{}epsilon\_\allowbreak{}sensitivity.tsv}
\end{itemize}

\textbf{6.3.2 Measured magnitudes across libraries}

\begin{itemize}
\tightlist
\item
  \longtok{platform\_hybrid/paper/sections/zvf.tex}
\item
  \longtok{platform\_hybrid/experiments/results/zvf\_\allowbreak{}by\_\allowbreak{}library.tsv}
\item
  \longtok{platform\_hybrid/paper/sections/p2\_\allowbreak{}results\_\allowbreak{}intro.tex}
\end{itemize}

\textbf{6.3.3 The all-correct wall and the mastery--incapacity aliasing}

\begin{itemize}
\tightlist
\item
  \longtok{platform\_hybrid/paper/sections/zvf.tex}
\item
  \longtok{platform\_hybrid/paper/sections/p2\_conclusion.tex}
\item
  \longtok{platform\_hybrid/paper/sections/frontier\_\allowbreak{}synthesis\_\allowbreak{}zvf.tex}
\item
  \longtok{platform\_hybrid/paper/sections/appendix\_\allowbreak{}zvf\_\allowbreak{}formalization.tex}
\end{itemize}

\textbf{6.3.4 Prognostic value}

\begin{itemize}
\tightlist
\item
  \longtok{platform\_hybrid/experiments/results/zvf\_\allowbreak{}iter26\_\allowbreak{}residual.tsv},
  corrected, see Table E.1
\item
  \longtok{platform\_hybrid/paper/sections/zvf.tex}
\item
  \longtok{platform\_hybrid/paper/neurips\_\allowbreak{}2026\_\allowbreak{}variants/paper\_\allowbreak{}P8\_\allowbreak{}workshop.tex}
\item
  \longtok{platform\_hybrid/paper/sections/appendix\_\allowbreak{}zvf\_\allowbreak{}formalization.tex}
\item
  \longtok{platform\_hybrid/paper/sections/p2\_\allowbreak{}results\_\allowbreak{}intro.tex},
  Appendix A, claim P2-C2
\item
  \longtok{platform\_modal/scripts/zvf\_\allowbreak{}diagnostic\_\allowbreak{}iter130.py}
\item
  \longtok{platform\_hybrid/paper/sections/p2\_conclusion.tex}
\item
  \longtok{platform\_hybrid/paper/sections/frontier\_\allowbreak{}synthesis\_\allowbreak{}zvf.tex}
\end{itemize}

\textbf{6.4.1 The measured sweep}

\begin{itemize}
\tightlist
\item
  \longtok{platform\_hybrid/paper/sections/group\_size.tex}
\item
  \longtok{platform\_hybrid/experiments/results/group\_\allowbreak{}size\_\allowbreak{}effect\_\allowbreak{}dpo\_\allowbreak{}check.tsv}
\item
  \longtok{platform\_hybrid/paper/sections/group\_\allowbreak{}size\_\allowbreak{}iter27.tex}
\item
  \longtok{platform\_hybrid/paper/sections/p3\_abstract.tex}
\end{itemize}

\textbf{6.4.2 The same-stack four-axis audit and the reconstructed
budgets}

\begin{itemize}
\tightlist
\item
  \longtok{platform\_hybrid/experiments/results/berkeley/unpacking\_\allowbreak{}dpo\_\allowbreak{}ppo\_\allowbreak{}factorization.json}
\item
  as summarised in
  \longtok{platform\_hybrid/paper/sections/p3\_conclusion.tex}
\item
  \longtok{platform\_hybrid/experiments/results/berkeley/adding\_\allowbreak{}error\_\allowbreak{}bars\_\allowbreak{}summary.json}
\item
  \longtok{platform\_hybrid/experiments/results/group\_\allowbreak{}size\_\allowbreak{}effect.tsv}
\item
  \longtok{platform\_hybrid/paper/sections/group\_size.tex}
\item
  \longtok{platform\_hybrid/paper/sections/group\_\allowbreak{}size\_\allowbreak{}iter23.tex}
\item
  \longtok{platform\_hybrid/paper/sections/p3\_conclusion.tex}
\item
  \longtok{platform\_hybrid/paper/sections/p3\_abstract.tex}
\end{itemize}

\textbf{6.4.3 The curriculum and matched-token null results}

\begin{itemize}
\tightlist
\item
  \longtok{platform\_hybrid/experiments/results/curriculum\_opening/FINDINGS.md}
\item
  as summarised in
  \longtok{platform\_hybrid/paper/sections/robustness\_\allowbreak{}smallscale\_\allowbreak{}nulls.tex}
\item
  \longtok{platform\_hybrid/paper/sections/robustness\_\allowbreak{}smallscale\_\allowbreak{}nulls.tex}
\item
  \longtok{platform\_hybrid/experiments/results/p1\_layerfreeze/scaled\_\allowbreak{}4seed\_\allowbreak{}result.json}
\end{itemize}

\textbf{6.5.1 The controlled null}

\begin{itemize}
\tightlist
\item
  \longtok{platform\_hybrid/experiments/results/drgrpo\_\allowbreak{}vs\_\allowbreak{}grpo.json}
\item
  \longtok{platform\_hybrid/paper/sections/length\_bias.tex}
\item
  \longtok{platform\_hybrid/experiments/results/drgrpo\_\allowbreak{}gsm8k\_\allowbreak{}cot\_\allowbreak{}full.json}
\item
  \longtok{platform\_hybrid/paper/sections/p4\_abstract.tex}
\item
  \longtok{platform\_hybrid/paper/sections/p4\_\allowbreak{}results\_\allowbreak{}intro.tex}
\end{itemize}

\textbf{6.5.2 The roster-level verbosity-trap audit}

\begin{itemize}
\tightlist
\item
  \longtok{platform\_hybrid/paper/archive/absorbed/U01\_\allowbreak{}main\_\allowbreak{}compendium/main.tex}
\item
  the same rates appear in
  \longtok{platform\_tinker/reports/combined/combined\_capstone.tex},
  where the trap rates are given as GRPO 36\% and PPO 71\%
\item
  \longtok{platform\_tinker/reports/combined/combined\_capstone.tex}
\item
  \longtok{platform\_hybrid/paper/sections/length\_\allowbreak{}bias\_\allowbreak{}iter48.tex}
\end{itemize}

\textbf{6.5.3 Elevated analyses}

\begin{itemize}
\tightlist
\item
  \longtok{platform\_hybrid/paper/sections/length\_bias.tex}
\end{itemize}

\textbf{6.5.4 Held-out generalisation and the scale extension}

\begin{itemize}
\tightlist
\item
  \longtok{platform\_hybrid/paper/archive/absorbed/U01\_\allowbreak{}main\_\allowbreak{}compendium/main.tex}
\item
  \longtok{platform\_hybrid/paper/sections/p4\_\allowbreak{}results\_\allowbreak{}intro.tex}
\item
  \longtok{platform\_hybrid/paper/sections/p4\_intro.tex}
\item
  \longtok{platform\_hybrid/paper/sections/p4\_abstract.tex}
\item
  \longtok{platform\_hybrid/paper/sections/length\_bias.tex}
\item
  \longtok{platform\_hybrid/experiments/results/berkeley/adding\_\allowbreak{}error\_\allowbreak{}bars\_\allowbreak{}summary.json}
\item
  \longtok{platform\_hybrid/paper/sections/p4\_conclusion.tex}
\end{itemize}

\textbf{6.5.5 Same-stack rerun on the unsaturated configuration}

\begin{itemize}
\tightlist
\item
  \longtok{platform\_hybrid/experiments/modal/modal\_\allowbreak{}samestack\_\allowbreak{}gsm8k\_\allowbreak{}cot.py}
\item
  \longtok{platform\_hybrid/experiments/results/samestack\_\allowbreak{}gsm8k\_\allowbreak{}cot.json}
\item
  \longtok{platform\_hybrid/experiments/results/samestack\_\allowbreak{}scope\_\allowbreak{}check.json}
\item
  regenerated by
  \longtok{platform\_hybrid/experiments/scripts/samestack\_\allowbreak{}scope\_\allowbreak{}check.py}
  from samestack\_\allowbreak{}gsm8k\_\allowbreak{}cot\_\allowbreak{}full.json
\end{itemize}

\section{Chapter 7: Results II: Reporting, Registry and Controller
(P5--P7)}\label{chapter-7-results-ii-reporting-registry-and-controller-p5p7}

\textbf{7.1 A Minimum Reporting Standard for Stack-Conditioned Results
(P5)}

\begin{itemize}
\tightlist
\item
  \longtok{platform\_hybrid/paper/sections/p5\_intro.tex}
\item
  \longtok{platform\_hybrid/paper/sections/p5\_evidence.tex}, Table
  \texttt{tab:p5-h2h}
\item
  \longtok{platform\_hybrid/paper/sections/p5\_stack.tex}, §Item 3
\item
  \longtok{platform\_hybrid/paper/sections/p5\_abstract.tex}
\item
  \longtok{platform\_hybrid/paper/sections/p5\_evidence.tex}, Exhibits 5
  and 7
\item
  \longtok{platform\_hybrid/experiments/results/groupsize\_\allowbreak{}zvf\_\allowbreak{}sweep.tsv}
\item
  \longtok{platform\_hybrid/paper/sections/p5\_stack.tex}, Table
  \texttt{tab:p5-minreport}
\item
  \longtok{platform\_hybrid/paper/sections/p5\_stack.tex}, Table
  \texttt{tab:p5-ushape}
\item
  \longtok{platform\_hybrid/paper/sections/p5\_stack.tex}, Table
  \texttt{tab:p5-gsize}
\item
  \longtok{platform\_hybrid/paper/sections/p5\_stack.tex}, §Item 7
\item
  \longtok{platform\_hybrid/paper/sections/p5\_stack.tex}, §Item 8
\item
  \longtok{platform\_hybrid/paper/sections/p5\_evidence.tex}, Exhibit 8
\item
  \longtok{platform\_hybrid/paper/sections/p5\_evidence.tex}, Exhibits 6
  and 18
\item
  \longtok{platform\_hybrid/paper/sections/p6\_\allowbreak{}registry\_\allowbreak{}health.tex},
  Table \texttt{tab:p6-minreport-subfield}, reported for the same corpus
\item
  \longtok{platform\_hybrid/paper/sections/p5\_evidence.tex}, Table
  \longtok{tab:p5-discriminative-entropy}
\item
  \longtok{platform\_hybrid/paper/sections/p5\_evidence.tex}, Table
  \texttt{tab:p5-field-suff}
\item
  \longtok{platform\_hybrid/paper/sections/p5\_evidence.tex}, Table
  \texttt{tab:p5-manifest-self-suff}
\item
  \longtok{platform\_hybrid/paper/sections/p5\_stack.tex},
  §Placebo-Replacement Bound
\item
  \longtok{platform\_hybrid/paper/sections/p5\_stack.tex}, §Live-Corpus
  Coupling
\item
  recomputed in
  \longtok{outputs/PES\_\allowbreak{}Phase2\_\allowbreak{}Third\_\allowbreak{}Review\_\allowbreak{}2026-09-24/thesis\_\allowbreak{}review\_\allowbreak{}2026-09-27/stats\_followup/m18\_\allowbreak{}mantel\_\allowbreak{}summary.json}
\end{itemize}

\textbf{7.2 The GRPO-Registry (P6)}

\begin{itemize}
\tightlist
\item
  \longtok{platform\_hybrid/paper/sections/p6\_abstract.tex}
\item
  \longtok{platform\_hybrid/paper/paper\_\allowbreak{}P6\_\allowbreak{}registry.tex}
\item
  \longtok{platform\_hybrid/paper/sections/p6\_schema.tex}
\item
  \longtok{platform\_hybrid/paper/sections/p6\_population.tex}
\item
  \longtok{platform\_hybrid/paper/sections/p6\_\allowbreak{}registry\_\allowbreak{}health.tex}
\item
  \longtok{platform\_hybrid/experiments/results/p5p8/registry\_\allowbreak{}stress\_\allowbreak{}per\_\allowbreak{}run.tsv}
  and registry\_\allowbreak{}stress\_\allowbreak{}summary.json
\item
  \longtok{platform\_hybrid/paper/sections/p6\_\allowbreak{}registry\_\allowbreak{}health.tex},
  §Stress Test
\item
  \longtok{outputs/PES\_\allowbreak{}Phase2\_\allowbreak{}Third\_\allowbreak{}Review\_\allowbreak{}2026-09-24/thesis\_\allowbreak{}review\_\allowbreak{}2026-09-27/stats\_followup/m16\_summary.json}
\item
  \longtok{platform\_hybrid/paper/sections/p6\_\allowbreak{}measured\_\allowbreak{}evidence.tex},
  Table \texttt{tab:p6-measured-deltas}
\item
  \longtok{platform\_hybrid/paper/sections/p6\_\allowbreak{}registry\_\allowbreak{}health.tex},
  §Contrastive Yield and Anti-Herding
\item
  \longtok{platform\_hybrid/paper/sections/p6\_\allowbreak{}query\_\allowbreak{}auditor.tex}
\end{itemize}

\textbf{7.3 From Diagnostic to Controller (P7)}

\begin{itemize}
\tightlist
\item
  \longtok{platform\_hybrid/paper/sections/p7\_abstract.tex}
\item
  \longtok{platform\_hybrid/paper/sections/p7\_theory.tex}
\item
  \longtok{platform\_hybrid/paper/sections/p7\_theory.tex}, Tables
  \texttt{tab:p7-modelG} and \texttt{tab:p7-gsweep}
\item
  \longtok{platform\_hybrid/paper/sections/p7\_controller.tex}, Table
  \texttt{tab:p7-e3}
\item
  \longtok{platform\_hybrid/paper/sections/p7\_controller.tex}, Tables
  \texttt{tab:p7-cf-sweep} and \texttt{tab:p7-cf-granularity}
\item
  \longtok{platform\_hybrid/paper/sections/p7\_controller.tex}, §Bayesian
  Refinement
\item
  \longtok{platform\_hybrid/paper/sections/p7\_controller.tex},
  §Posterior-Predictive Contrast Restoration and §Statistically
  Validated Pareto Frontier
\item
  \longtok{platform\_hybrid/paper/sections/p7\_controller.tex}, §Exact
  Finite-Pool Correction
\item
  \longtok{platform\_hybrid/paper/sections/p7\_\allowbreak{}design\_\allowbreak{}rules.tex}
\end{itemize}

\textbf{7.4 Historical V2 Frozen-Policy Fresh-Rollout Diagnostic}

\begin{itemize}
\tightlist
\item
  \longtok{outputs/rescue\_\allowbreak{}colab\_\allowbreak{}v2\_\allowbreak{}2026-10-02/analysis\_final/summary.json}
\item
  \longtok{outputs/rescue\_\allowbreak{}colab\_\allowbreak{}v2\_\allowbreak{}2026-10-02/audit/completed\_\allowbreak{}primary\_\allowbreak{}event\_\allowbreak{}audit.json}
\item
  \longtok{outputs/rescue\_\allowbreak{}colab\_\allowbreak{}v2\_\allowbreak{}2026-10-02/analysis\_final/validation.json}
\item
  \longtok{outputs/rescue\_\allowbreak{}colab\_\allowbreak{}v2\_\allowbreak{}2026-10-02/audit\_\allowbreak{}method\_\allowbreak{}disclosure.json}
\item
  \longtok{outputs/rescue\_\allowbreak{}colab\_\allowbreak{}v2\_\allowbreak{}2026-10-02/audit/independent\_\allowbreak{}arithmetic\_\allowbreak{}corroboration.json}
\item
  \longtok{outputs/rescue\_\allowbreak{}colab\_\allowbreak{}v2\_\allowbreak{}2026-10-02/independent\_\allowbreak{}count\_\allowbreak{}and\_\allowbreak{}later\_\allowbreak{}change\_\allowbreak{}check.json}
\item
  \longtok{platform\_tinker/reports/final/evaluate\_\allowbreak{}gsm8k\_\allowbreak{}test.py}
\end{itemize}

\section{Chapter 8: The E1--E14 Evaluation
Campaign}\label{chapter-8-the-e1e14-evaluation-campaign}

\textbf{8.1 The evaluated actor and the one-actor design}

\begin{itemize}
\tightlist
\item
  \longtok{outputs/E1\_\allowbreak{}E14\_\allowbreak{}FINAL\_\allowbreak{}RESULTS\_\allowbreak{}2026-09-19.md}, the same
  identity pair is verifiable in the surviving E11 full receipt
\item
  which also preserves the training sampler path {[}operational
  identifier withheld{]}: source
  \longtok{outputs/modal\_\allowbreak{}e1\_\allowbreak{}e14/2026-08-16/e11\_\allowbreak{}full\_\allowbreak{}receipt.json},
  the corroborating E10 receipt
\item
  \longtok{outputs/public\_\allowbreak{}portfolio\_\allowbreak{}2026-09-05/agentdojo\_\allowbreak{}native\_\allowbreak{}receipt.json}
\item
  is now preserved in the withheld source archive and binds the same
  \texttt{base/adapter} identity
\end{itemize}

\textbf{8.2 Suites with strictly complete scopes}

\begin{itemize}
\tightlist
\item
  \longtok{outputs/verification/LEDGER\_\allowbreak{}CODE\_\allowbreak{}CHECK\_\allowbreak{}2026-09-19.json},
  check \longtok{strict\_\allowbreak{}completed\_\allowbreak{}suites\_\allowbreak{}are\_\allowbreak{}E8\_\allowbreak{}E10\_\allowbreak{}E11}
\item
  \longtok{outputs/modal\_\allowbreak{}e1\_\allowbreak{}e14/2026-08-16/e11\_\allowbreak{}full\_\allowbreak{}receipt.json}
\item
  \longtok{outputs/e1\_\allowbreak{}e14\_\allowbreak{}small\_\allowbreak{}scale\_\allowbreak{}2026-09-26/E11/paired.json}
\item
  bootstrap script and paired-outcome artifact
  \longtok{platform\_hybrid/experiments/scripts/e11\_\allowbreak{}design\_\allowbreak{}cluster\_\allowbreak{}bootstrap.py}
  and \texttt{.json}
\item
  \longtok{outputs/e11\_\allowbreak{}verilog\_\allowbreak{}eval/e11\_\allowbreak{}verilog\_\allowbreak{}eval\_\allowbreak{}rerun\_\allowbreak{}receipt.json}
\item
  \longtok{outputs/e11\_\allowbreak{}verilog\_\allowbreak{}eval/lane\_\allowbreak{}status\_\allowbreak{}2026-08-09.md}
\item
  \longtok{outputs/public\_\allowbreak{}portfolio\_\allowbreak{}2026-09-05/native\_\allowbreak{}finish\_\allowbreak{}v5/e14\_\allowbreak{}official\_\allowbreak{}collector/official\_result01/native\_dispositions.jsonl}
\end{itemize}

\textbf{8.3 Original-contract results}

\begin{itemize}
\tightlist
\item
  \longtok{outputs/modal\_\allowbreak{}e1\_\allowbreak{}e14/2026-08-16/e1\_\allowbreak{}swe\_\allowbreak{}bench\_\allowbreak{}pro\_\allowbreak{}full/seed1818/receipt.json}
\item
  \longtok{outputs/e1\_\allowbreak{}e14\_\allowbreak{}results\_\allowbreak{}2026-09-05/E1\_\allowbreak{}E14\_\allowbreak{}Results.md}
\item
  parse of the stored \texttt{candidates.json} beside the receipt
\item
  the runner
  \longtok{zvf-program/flagship/modal\_\allowbreak{}e1\_\allowbreak{}swe\_\allowbreak{}bench\_\allowbreak{}pro\_\allowbreak{}full.py} now
  repairs hunk headers before submission and records apply status, but
  the stored score was produced without that repair
\item
  \longtok{outputs/E1\_\allowbreak{}E14\_\allowbreak{}FINAL\_\allowbreak{}RESULTS\_\allowbreak{}2026-09-19.md}
\item
  \longtok{outputs/verification/LEDGER\_\allowbreak{}CODE\_\allowbreak{}CHECK\_\allowbreak{}2026-09-19.json}
\end{itemize}

\Needspace{8\baselineskip}

\textbf{8.4 Replacement-scope lanes completed on 27 September 2026}

\textbf{8.5 The E4 rerun as a diagnostic}

\begin{itemize}
\tightlist
\item
  \longtok{outputs/e1\_\allowbreak{}e14\_\allowbreak{}small\_\allowbreak{}scale\_\allowbreak{}2026-09-26/E4/BRIDGE\_\allowbreak{}STOP\_\allowbreak{}TOKEN\_\allowbreak{}FINDING.md}
\end{itemize}

\textbf{8.7 Small-scale reruns of all fourteen lanes}

\begin{itemize}
\tightlist
\item
  \longtok{outputs/e1\_\allowbreak{}e14\_\allowbreak{}small\_\allowbreak{}scale\_\allowbreak{}2026-09-26/PROTOCOL.md}
\item
  \longtok{outputs/e1\_\allowbreak{}e14\_\allowbreak{}small\_\allowbreak{}scale\_\allowbreak{}2026-09-26/TRAINED\_\allowbreak{}ACTOR\_\allowbreak{}ENDPOINT.md}
\end{itemize}

\section{Chapter 9: Cross-Study Synthesis, Threats to Validity,
Conclusions and Future
Work}\label{chapter-9-cross-study-synthesis-threats-to-validity-conclusions-and-future-work}

\textbf{9.1 Cross-Study Synthesis}

\begin{itemize}
\tightlist
\item
  \longtok{platform\_tinker/reports/esa\_phase1/Phase1\_\allowbreak{}Project\_\allowbreak{}Report\_\allowbreak{}ZVF.tex}
\item
  \longtok{platform\_hybrid/paper/sections/p4\_conclusion.tex}
\item
  \longtok{platform\_hybrid/experiments/results/berkeley/adding\_\allowbreak{}error\_\allowbreak{}bars\_\allowbreak{}summary.json}
\item
  \longtok{platform\_hybrid/paper/sections/p3\_abstract.tex}
\item
  \longtok{platform\_hybrid/paper/sections/p7\_abstract.tex}
\item
  \longtok{platform\_hybrid/paper/sections/p5\_abstract.tex}
\item
  \longtok{platform\_hybrid/paper/sections/p6\_abstract.tex}
\item
  \longtok{platform\_hybrid/paper/archive/absorbed/P08\_fraud/sections/p8\_abstract.tex}
\item
  \longtok{outputs/E1\_\allowbreak{}E14\_\allowbreak{}FINAL\_\allowbreak{}RESULTS\_\allowbreak{}2026-09-19.md}, withheld
  historical campaign chronology
\end{itemize}

\textbf{9.2 Rigour and Threats to Validity}

\begin{itemize}
\tightlist
\item
  \longtok{outputs/verification/LEDGER\_\allowbreak{}CODE\_\allowbreak{}CHECK\_\allowbreak{}2026-09-19.json}
\item
  LIMITATIONS\_\allowbreak{}AND\_\allowbreak{}IMPACT.md §7.1, Table 1.1
\item
  LIMITATIONS\_\allowbreak{}AND\_\allowbreak{}IMPACT.md §6.2
\item
  LIMITATIONS\_\allowbreak{}AND\_\allowbreak{}IMPACT.md §7.1
\item
  LIMITATIONS\_\allowbreak{}AND\_\allowbreak{}IMPACT.md §6.1
\item
  \longtok{platform\_tinker/reports/combined/combined\_capstone.tex}
\item
  \longtok{outputs/E1\_\allowbreak{}E14\_\allowbreak{}FINAL\_\allowbreak{}RESULTS\_\allowbreak{}2026-09-19.md}
\item
  LIMITATIONS\_\allowbreak{}AND\_\allowbreak{}IMPACT.md §2.2, §2.3
\item
  \longtok{platform\_hybrid/paper/archive/absorbed/U01\_\allowbreak{}main\_\allowbreak{}compendium/main.tex}
\item
  \longtok{outputs/E1\_\allowbreak{}E14\_\allowbreak{}FINAL\_\allowbreak{}RESULTS\_\allowbreak{}2026-09-19.md} §6
\item
  \longtok{outputs/E1\_\allowbreak{}E14\_\allowbreak{}FINAL\_\allowbreak{}RESULTS\_\allowbreak{}2026-09-19.md} §9.5
\item
  \longtok{outputs/E1\_\allowbreak{}E14\_\allowbreak{}FINAL\_\allowbreak{}RESULTS\_\allowbreak{}2026-09-19.md} §8
\item
  LIMITATIONS\_\allowbreak{}AND\_\allowbreak{}IMPACT.md §7.2, §7.3
\item
  LIMITATIONS\_\allowbreak{}AND\_\allowbreak{}IMPACT.md §3, §4
\end{itemize}

\textbf{9.3 Conclusions}

\begin{itemize}
\tightlist
\item
  \longtok{platform\_hybrid/paper/sections/p7\_abstract.tex}
\item
  LIMITATIONS\_\allowbreak{}AND\_\allowbreak{}IMPACT.md §2.2
\item
  \longtok{platform\_tinker/reports/esa\_phase1/Phase1\_\allowbreak{}Project\_\allowbreak{}Report\_\allowbreak{}ZVF.tex}
\item
  \longtok{platform\_hybrid/experiments/results/p1\_layerfreeze/scaled\_\allowbreak{}4seed\_\allowbreak{}result.json}
\end{itemize}

\textbf{9.4 Future Work}

\begin{itemize}
\tightlist
\item
  \longtok{outputs/finish\_\allowbreak{}pending\_\allowbreak{}2026-09-27/E6/result.json}
\item
  \longtok{outputs/E1\_\allowbreak{}E14\_\allowbreak{}FINAL\_\allowbreak{}RESULTS\_\allowbreak{}2026-09-19.md} §5, withheld
  historical campaign chronology
\item
  \longtok{platform\_tinker/reports/esa\_phase1/Phase1\_\allowbreak{}Project\_\allowbreak{}Report\_\allowbreak{}ZVF.tex}
\end{itemize}

\section{Appendix B: Reproducibility and
Provenance}\label{appendix-b-reproducibility-and-provenance}

\textbf{B.12 Fresh-Sampling Diagnostic: Frozen Contract and Audit}

\begin{itemize}
\tightlist
\item
  \longtok{research/rescue\_\allowbreak{}sampling\_\allowbreak{}v2\_\allowbreak{}protocol\_\allowbreak{}2026\_\allowbreak{}10\_\allowbreak{}02.md}
\item
  \longtok{research/rescue\_\allowbreak{}sampling\_\allowbreak{}v2\_\allowbreak{}freeze\_\allowbreak{}2026\_\allowbreak{}10\_\allowbreak{}02.json}
\item
  \longtok{research/rescue\_\allowbreak{}v2\_\allowbreak{}expected\_\allowbreak{}cohort\_\allowbreak{}2026\_\allowbreak{}10\_\allowbreak{}02.json}
\item
  \longtok{research/rescue\_\allowbreak{}answer\_\allowbreak{}parser\_\allowbreak{}v2.py}
\item
  \longtok{outputs/rescue\_\allowbreak{}colab\_\allowbreak{}v2\_\allowbreak{}2026-10-02/preflight/gold\_\allowbreak{}parser\_\allowbreak{}preflight.json}
\item
  \longtok{outputs/rescue\_\allowbreak{}colab\_\allowbreak{}v2\_\allowbreak{}2026-10-02/analysis\_final/validation.json}
\item
  \longtok{outputs/rescue\_\allowbreak{}colab\_\allowbreak{}v2\_\allowbreak{}2026-10-02/audit\_\allowbreak{}method\_\allowbreak{}disclosure.json}
\item
  \longtok{outputs/rescue\_\allowbreak{}colab\_\allowbreak{}v2\_\allowbreak{}2026-10-02/audit/completed\_\allowbreak{}primary\_\allowbreak{}event\_\allowbreak{}audit.json}
\item
  \longtok{outputs/rescue\_\allowbreak{}colab\_\allowbreak{}v2\_\allowbreak{}2026-10-02/audit/independent\_\allowbreak{}arithmetic\_\allowbreak{}corroboration.json}
\item
  \longtok{outputs/rescue\_\allowbreak{}colab\_\allowbreak{}v2\_\allowbreak{}2026-10-02/independent\_\allowbreak{}count\_\allowbreak{}and\_\allowbreak{}later\_\allowbreak{}change\_\allowbreak{}check.json}
\item
  \longtok{outputs/rescue\_\allowbreak{}colab\_\allowbreak{}v2\_\allowbreak{}2026-10-02/analysis\_final/summary.json}
\item
  \longtok{outputs/rescue\_\allowbreak{}colab\_\allowbreak{}v2\_\allowbreak{}2026-10-02/analysis\_final/per\_prompt.csv}
\item
  \longtok{research/rescue\_\allowbreak{}v1\_\allowbreak{}manifest\_\allowbreak{}reference\_\allowbreak{}2026\_\allowbreak{}10\_\allowbreak{}02.json}
\end{itemize}

\section{Appendix D: Campaign Lane
Detail}\label{appendix-d-campaign-lane-detail}

\textbf{D.1 Actor identity, inference conditions and scoping notes}

\begin{itemize}
\tightlist
\item
  recovered
  \longtok{outputs/public\_\allowbreak{}portfolio\_\allowbreak{}2026-09-05/agentdojo\_\allowbreak{}native\_\allowbreak{}receipt.json},
  unrecovered \texttt{withheld\ AgentDojo\ collector\ source}
\item
  \longtok{outputs/modal\_\allowbreak{}e1\_\allowbreak{}e14/2026-08-16/e11\_\allowbreak{}full\_\allowbreak{}receipt.json}
\item
  \longtok{outputs/e1\_\allowbreak{}e14\_\allowbreak{}results\_\allowbreak{}2026-09-05/E1\_\allowbreak{}E14\_\allowbreak{}Results.md}
\item
  \longtok{outputs/e11\_\allowbreak{}verilog\_\allowbreak{}eval/e11\_\allowbreak{}smoke\_\allowbreak{}receipt.json}
\item
  \longtok{outputs/e11\_\allowbreak{}verilog\_\allowbreak{}eval/e11\_\allowbreak{}trained\_\allowbreak{}step40\_\allowbreak{}receipt.json}
\end{itemize}

\textbf{D.4 Strictly complete scopes: lane detail}

\begin{itemize}
\tightlist
\item
  recovered
  \longtok{outputs/public\_\allowbreak{}portfolio\_\allowbreak{}2026-09-05/agentdojo\_\allowbreak{}native\_\allowbreak{}receipt.json}
\item
  \longtok{outputs/modal\_\allowbreak{}e1\_\allowbreak{}e14/2026-08-16/e11\_\allowbreak{}full\_\allowbreak{}receipt.json}
\item
  \longtok{outputs/e1\_\allowbreak{}e14\_\allowbreak{}small\_\allowbreak{}scale\_\allowbreak{}2026-09-26/E11/paired.json}
\item
  \longtok{outputs/e11\_\allowbreak{}verilog\_\allowbreak{}eval/e11\_\allowbreak{}verilog\_\allowbreak{}eval\_\allowbreak{}rerun\_\allowbreak{}receipt.json}
\item
  \longtok{outputs/e11\_\allowbreak{}verilog\_\allowbreak{}eval/lane\_\allowbreak{}status\_\allowbreak{}2026-08-09.md}
\item
  \longtok{outputs/verification/LEDGER\_\allowbreak{}CODE\_\allowbreak{}CHECK\_\allowbreak{}2026-09-19.json},
  check \longtok{E11\_\allowbreak{}component\_\allowbreak{}split\_\allowbreak{}sums\_\allowbreak{}67\_\allowbreak{}plus\_\allowbreak{}62\_\allowbreak{}equals\_\allowbreak{}129}
\end{itemize}

\textbf{D.5 Original-contract results: lane detail}

\begin{itemize}
\tightlist
\item
  \longtok{outputs/modal\_\allowbreak{}e1\_\allowbreak{}e14/2026-08-16/e1\_\allowbreak{}swe\_\allowbreak{}bench\_\allowbreak{}pro\_\allowbreak{}full/seed1818/receipt.json}
\item
  \longtok{outputs/e1\_\allowbreak{}swe\_\allowbreak{}bench\_\allowbreak{}pro/lane\_\allowbreak{}receipt\_\allowbreak{}2026-08-09.json}
\item
  \longtok{outputs/e1\_\allowbreak{}swe\_\allowbreak{}bench\_\allowbreak{}pro/preflight\_2026-08-09.json}
\item
  \longtok{outputs/e1\_\allowbreak{}e14\_\allowbreak{}results\_\allowbreak{}2026-09-05/E1\_\allowbreak{}E14\_\allowbreak{}Results.md}
\item
  \longtok{outputs/e1\_\allowbreak{}e14\_\allowbreak{}results\_\allowbreak{}2026-09-05/E1\_\allowbreak{}E14\_\allowbreak{}Results.md},
  recovered \texttt{withheld\ historical\ campaign\ chronology}
\item
  \longtok{outputs/E1\_\allowbreak{}E14\_\allowbreak{}FINAL\_\allowbreak{}RESULTS\_\allowbreak{}2026-09-19.md}
\item
  \longtok{outputs/verification/LEDGER\_\allowbreak{}CODE\_\allowbreak{}CHECK\_\allowbreak{}2026-09-19.json}
\end{itemize}

\textbf{D.6 Replacement-scope lane narratives, 27 September 2026}

\begin{itemize}
\tightlist
\item
  \longtok{outputs/finish\_\allowbreak{}pending\_\allowbreak{}2026-09-27/E1/result.json}
\item
  \longtok{outputs/finish\_\allowbreak{}pending\_\allowbreak{}2026-09-27/E2/result.json}
\item
  \longtok{outputs/finish\_\allowbreak{}pending\_\allowbreak{}2026-09-27/E5/result.json}
\item
  \longtok{outputs/finish\_\allowbreak{}pending\_\allowbreak{}2026-09-27/E6/result.json}
\item
  \longtok{outputs/finish\_\allowbreak{}pending\_\allowbreak{}2026-09-27/E9/result.json}
\item
  \longtok{outputs/finish\_\allowbreak{}pending\_\allowbreak{}2026-09-27/E13/result.json}
\end{itemize}

\textbf{D.7 E4 rerun: reported audit and causal limits}

\begin{itemize}
\tightlist
\item
  \longtok{outputs/e1\_\allowbreak{}e14\_\allowbreak{}small\_\allowbreak{}scale\_\allowbreak{}2026-09-26/E4/BRIDGE\_\allowbreak{}STOP\_\allowbreak{}TOKEN\_\allowbreak{}FINDING.md}
\end{itemize}

\textbf{D.9 Small-scale reruns: detail and per-lane caveats}

\begin{itemize}
\tightlist
\item
  \longtok{outputs/e1\_\allowbreak{}e14\_\allowbreak{}small\_\allowbreak{}scale\_\allowbreak{}2026-09-26/PROTOCOL.md}
\end{itemize}

\section{Appendix F: Applied Side Study: Credit-Card Fraud Detection
(P8)}\label{appendix-f-applied-side-study-credit-card-fraud-detection-p8}

\begin{itemize}
\tightlist
\item
  \longtok{platform\_hybrid/paper/archive/absorbed/P08\_fraud/sections/p8\_intro.tex}
\end{itemize}

\textbf{F.1 Dataset, splits, and the two arms}

\begin{itemize}
\tightlist
\item
  \longtok{platform\_hybrid/paper/archive/absorbed/P08\_fraud/sections/p8\_setup.tex}
\item
  \texttt{fraud\_data.csv}
\item
  \longtok{platform\_hybrid/paper/archive/absorbed/P08\_fraud/sections/p8\_limitations.tex}
\end{itemize}

\textbf{F.2 The head-to-head result}

\begin{itemize}
\tightlist
\item
  \longtok{platform\_hybrid/paper/archive/absorbed/P08\_fraud/sections/p8\_abstract.tex}
\item
  \texttt{xgboost\_results.json}
\item
  \longtok{platform\_hybrid/paper/archive/absorbed/P08\_fraud/sections/p8\_setup.tex}
\item
  \longtok{platform\_hybrid/paper/archive/absorbed/P08\_fraud/sections/p8\_limitations.tex}
\end{itemize}

\textbf{F.3 Why the scoring seat stays with the tree}

\begin{itemize}
\tightlist
\item
  \longtok{platform\_hybrid/paper/archive/absorbed/P08\_fraud/sections/p8\_scorer.tex}
\item
  \longtok{platform\_hybrid/paper/archive/absorbed/P08\_fraud/sections/p8\_evidence.tex}
\end{itemize}

\textbf{F.5 Proposed roles for the language model}

\begin{itemize}
\tightlist
\item
  \longtok{platform\_hybrid/paper/archive/absorbed/P08\_fraud/sections/p8\_taxonomy.tex}
\item
  \longtok{platform\_hybrid/paper/archive/absorbed/P08\_fraud/sections/p8\_compliance.tex}
\item
  \longtok{platform\_hybrid/paper/archive/absorbed/P08\_fraud/sections/p8\_architecture.tex}
\end{itemize}

\textbf{F.6 Limitations}

\begin{itemize}
\tightlist
\item
  \longtok{platform\_hybrid/paper/archive/absorbed/P08\_fraud/sections/p8\_limitations.tex}
\end{itemize}



\chapter{Public Coverage and Source
Availability}\label{public-coverage-and-source-availability}

This appendix is the publication-safe derivative of a bounded source
audit completed on 3 October 2026. It reports scientific source
availability and limitations. It does not publish authenticated provider
inventories, private project names, account-wide activity totals,
collaborator identities or operational access correspondence. The
underlying private audit records are withheld; these paragraphs do not
turn their contents into independently reproducible public evidence.

\section{Counting units and publication
boundary}\label{counting-units-and-publication-boundary}

A scientific execution, a logging record, a training identity, a
checkpoint revision and a claim-to-source link are different units. One
execution can generate multiple records, retries and checkpoint
versions. These units must not be added as independent trials. Appendix
A retains scientific study scope and qualitative registry limitations;
no account-wide inventory is included.

Provider inventory work used bounded pagination, immutable revisions
where available, explicit source joins and missing-field handling. An
inventory match establishes recorded identity within the inspected
scope, not successful training, benchmark correctness or availability of
all historical weights. Unavailable or inaccessible records, unresolved
joins, unproven consumer-list exhaustion and missing history remain
limitations. Account-wide counts and exact operational identities are
omitted.

The editable sources and selected scientific case ledgers in this
package can rebuild the public document. They cannot reproduce every
experiment. The source-reference map is a map of reported scientific
sources, including references whose original payloads are retained
outside this public package. The public availability manifest is the
authority on which files are actually included.

\section{Bounded checkpoint and provider
evidence}\label{bounded-checkpoint-and-provider-evidence}

Training metadata, checkpoint bytes and an executable serving route are
distinct evidence. This public derivative does not publish private
provider state, checkpoint-retention findings or exact execution
identities. The public package cannot reproduce the original managed
execution route, and no current account or private repository
availability is asserted.

Public upstream model and benchmark revision references remain where
scientifically relevant. Published Git or LFS identifiers are metadata,
not a fresh byte-integrity test of a large tensor payload. Small public
card, configuration and manifest checks do not establish numerical
equivalence across serving stacks. No training, sampling or benchmark
evaluation was rerun by this source audit. Per-step metric histories
were not reanalysed.

\section{Public P11 confirmatory
bindings}\label{public-p11-confirmatory-bindings}

The forty public manifest bindings below belong to P11, the
five-algorithm by eight-seed confirmatory study. The repository naming
convention includes e1, but these are not the E1 SWE-bench lane of the
later fourteen-suite campaign. At the recorded audit date, all forty
public pinned commits matched the public heads, and all forty downloaded
run manifests matched both their expected SHA-256 and retained local
bytes. The companion public binding file retains scientific units,
immutable commit and checksum fields. Execution-specific repository
access paths are omitted to avoid publishing embedded provider
identifiers; this limits independent public replay from this package
alone.

Each accepted unit records checkpoints at steps 5, 10, 15, 20, 25 and 30
plus a final adapter. Optional evaluation-resume sidecars were recorded
for 36 of 40 units; that difference does not remove the matching final
manifests. These are forty accepted scientific units, not forty
additional trials inferred from provider logging.

\begin{longtable}[]{@{}lll@{}}
\caption{Public P11 manifest checks. Full immutable revisions and
manifest SHA-256 values are in
public\_\allowbreak{}p11\_\allowbreak{}manifest\_\allowbreak{}bindings.json.}\tabularnewline
\toprule\noalign{}
Scientific unit & Public commit prefix & Manifest verification \\
\midrule\noalign{}
\endfirsthead
\toprule\noalign{}
Scientific unit & Public commit prefix & Manifest verification \\
\midrule\noalign{}
\endhead
\bottomrule\noalign{}
\endlastfoot
aero/seed-107 & \longtok{42980267a519} & Exact SHA-256 match \\
aero/seed-11 & \longtok{2affa816a795} & Exact SHA-256 match \\
aero/seed-131 & \longtok{ec9931f7fe85} & Exact SHA-256 match \\
aero/seed-23 & \longtok{92123d35e401} & Exact SHA-256 match \\
aero/seed-37 & \longtok{613d3f70cc0f} & Exact SHA-256 match \\
aero/seed-53 & \longtok{448707a7a1dc} & Exact SHA-256 match \\
aero/seed-71 & \longtok{5168da9f98c3} & Exact SHA-256 match \\
aero/seed-89 & \longtok{43421d47afa9} & Exact SHA-256 match \\
dapo/seed-107 & \longtok{6e60484cfd40} & Exact SHA-256 match \\
dapo/seed-11 & \longtok{33f158dfabf2} & Exact SHA-256 match \\
dapo/seed-131 & \longtok{39f916902470} & Exact SHA-256 match \\
dapo/seed-23 & \longtok{08d3a34da450} & Exact SHA-256 match \\
dapo/seed-37 & \longtok{bb7fb354be9f} & Exact SHA-256 match \\
dapo/seed-53 & \longtok{194c9d2ea0d7} & Exact SHA-256 match \\
dapo/seed-71 & \longtok{59406ac609a1} & Exact SHA-256 match \\
dapo/seed-89 & \longtok{c290319779e4} & Exact SHA-256 match \\
drgrpo/seed-107 & \longtok{8978e087b78f} & Exact SHA-256 match \\
drgrpo/seed-11 & \longtok{dbcd0740bc95} & Exact SHA-256 match \\
drgrpo/seed-131 & \longtok{4b05547ae98a} & Exact SHA-256 match \\
drgrpo/seed-23 & \longtok{4318ffd9026b} & Exact SHA-256 match \\
drgrpo/seed-37 & \longtok{26e36e5e915f} & Exact SHA-256 match \\
drgrpo/seed-53 & \longtok{5bef04e01ba9} & Exact SHA-256 match \\
drgrpo/seed-71 & \longtok{a8cd65ea26f0} & Exact SHA-256 match \\
drgrpo/seed-89 & \longtok{1977a08e30b3} & Exact SHA-256 match \\
grpo/seed-107 & \longtok{ffcbfffd322a} & Exact SHA-256 match \\
grpo/seed-11 & \longtok{b39028ea32b0} & Exact SHA-256 match \\
grpo/seed-131 & \longtok{7786d5e94527} & Exact SHA-256 match \\
grpo/seed-23 & \longtok{de7dc778c412} & Exact SHA-256 match \\
grpo/seed-37 & \longtok{2755e14f26f2} & Exact SHA-256 match \\
grpo/seed-53 & \longtok{0a4145834057} & Exact SHA-256 match \\
grpo/seed-71 & \longtok{821fa6b4504a} & Exact SHA-256 match \\
grpo/seed-89 & \longtok{f88da35f8939} & Exact SHA-256 match \\
gspo/seed-107 & \longtok{2b521a728689} & Exact SHA-256 match \\
gspo/seed-11 & \longtok{f064274b7c79} & Exact SHA-256 match \\
gspo/seed-131 & \longtok{5753775e1976} & Exact SHA-256 match \\
gspo/seed-23 & \longtok{d1a6f1879bde} & Exact SHA-256 match \\
gspo/seed-37 & \longtok{a9c9ee88756d} & Exact SHA-256 match \\
gspo/seed-53 & \longtok{6b58a60fb73d} & Exact SHA-256 match \\
gspo/seed-71 & \longtok{c0c0b968a61a} & Exact SHA-256 match \\
gspo/seed-89 & \longtok{40695b92925a} & Exact SHA-256 match \\
\end{longtable}

The manifests specify thirty optimizer steps. P11 therefore does not
satisfy the report's Tier A threshold of at least one hundred steps,
despite its eight-seed design. The reported contrast, confidence
interval, power estimate and null result remain unchanged. Failed and
superseded attempts are not pooled with the accepted forty units.

\section{Recovered campaign evidence}\label{recovered-campaign-evidence}

Six unchanged original files resolved five named primary-file gaps in
the retained-source audit. Their private filesystem and export metadata
are not included here. The scientific availability result is narrower
than validation of every claim: LAB-Bench records 450 correct of 1,967;
AgentDojo records 88 utility passes across 97 benign episodes, with
security evaluation false and no held-out claim. The recovered
training-identity note explicitly disclaims content decontamination and
does not provide the dirty source diff.

Bounded stored-payload checks cover 29 explicit file-reference records
across 27 targets, 1,979 E8 archive/local bindings and 525 E10
archive-member bindings. The selected 1,967 E8 grades and 97 E10 traces
reproduce the stored success counts of 450 and 88. E8's 1,259
unparsed-answer total remains diagnostic-supported rather than
independently recalculated from the reduced grade rows. These are
structural, hash and tally checks, not model execution, scorer reruns or
independent semantic grading.

For E14, the source-side record contains 4,428 dispositions, 4,426
accepted judged rows and 2,271 successes. The official judged-row score
is 2,271/4,426 = 51.31\%; the all-disposition diagnostic is 2,271/4,428
= 51.29\%, counting two missing-verdict rows as failures. The recovered
terminal note incorrectly associates the official decimal with the full
denominator. The erratum remains explicit, the original source is
unchanged, and strict-contract score is not inferred from terminal
closure.

The recovered finish-era chronology precedes the later E4 rerun and
cannot verify that later execution. The later one-hundred-trial
zero-reward account remains a separate reported result with
raw-trajectory limits. Original training-log payloads, the dirty source
diff and content-level decontamination remain unresolved.

\section{Preserved C1 result and residual
limits}\label{preserved-c1-result-and-residual-limits}

The C1 parser-defined endpoint remains 11/64. The conservative review
supports three numerical recovery witnesses on the same denominator;
these counts are separate and neither establishes training improvement.
The sixty-four reference classifications, eleven event records,
scientific hashes and disagreements remain in the included case ledger.
Historical V2 findings remain separate and unpooled.

This package excludes private provider exports, authenticated inventory
and reconciliation ledgers, raw access/budget chronologies, private
notebook links, original recovery-export wrappers, personal signature
data and account identifiers. Omission is a publication boundary, not
evidence that a scientific result failed. Conversely, a retained source
reference or checksum is not proof that the cited bytes are publicly
available or that all claims have been independently verified. The
public PDF has a distinct title and filename and does not replace the
preserved historical manuscript or certify institutional approval.



\end{document}
